
\documentclass[a4paper]{tufte-book}%symmetric
\hypersetup{colorlinks}% uncomment this line if you prefer colored hyperlinks (e.g., for onscreen viewing)

%%
% Book metadata
\title{EMS \\Physics Notes \\ $\beta$ Release}
\author[Joe Rowing]{Joe Rowing}
\publisher{Exeter Mathematics School}

%%
% If they're installed, use Bergamo and Chantilly from www.fontsite.com.
% They're clones of Bembo and Gill Sans, respectively.
%\IfFileExists{bergamo.sty}{\usepackage[osf]{bergamo}}{}% Bembo
%\IfFileExists{chantill.sty}{\usepackage{chantill}}{}% Gill Sans

%\usepackage{microtype}

%%
%%
% For nicely typeset tabular material
\usepackage{booktabs}
\setcounter{secnumdepth}{0}
%%
% For graphics / images
\usepackage{graphicx}
\setkeys{Gin}{width=\linewidth,totalheight=\textheight,keepaspectratio}
\graphicspath{{graphics/}}

% The fancyvrb package lets us customize the formatting of verbatim
% environments.  We use a slightly smaller font.
\usepackage{fancyvrb}

\usepackage{fancyvrb,cprotect}
%\usepackage{fancyhdr}
%\pagestyle{fancy}
%\lhead{ \fancyplain{}{EMS Notes} }
%\rhead{ \fancyplain{}{\leftmark} }
%\cfoot{ \fancyplain{}{\thepage} }
\usepackage[unicode]{hyperref}
\usepackage{amsmath}
\usepackage{amsfonts}
\usepackage{latexsym}
\usepackage{pifont}
\usepackage{fourier}
\usepackage{wasysym}
\usepackage{stmaryrd}
\usepackage[dvipsnames,svgnames]{xcolor}
\usepackage{extarrows}
\usepackage{chemarrow}
\usepackage{graphicx}
\usepackage{tikz}
\usepackage{pgfplots}
\usetikzlibrary{arrows, positioning, calc, fadings, shapes, decorations.markings, decorations.pathmorphing}
\usepgfmodule{shapes}
\usepgfmodule{plot}
\usepackage{pgffor}
\usepackage{array}
\usepackage{anyfontsize}
\usepackage{circuitikz}
\usepackage{multirow}
%\usepackage{dingbat}   % special symbols
\usepackage{indentfirst}
\usepackage{paralist}
\usepackage{enumitem}
\usepackage{framed}
\usepackage{physics}
\usepackage{easyphys}
\usepackage{wrapfig}
\usepackage{caption}
\usepackage{titletoc}
%\usepackage[explicit]{titlesec}
\usepackage{makeidx,xfrac,xspace}
\usepackage{ifthen}
\usepackage[makeroom]{cancel}
\fvset{fontsize=\normalsize}
\usepackage[]{geometry}
\usepackage{textalpha}
\usepackage{float}
%\usepackage[alsoload=abbr]{siunitx}
\usepackage{empheq}
\usepackage[most]{tcolorbox}
\usepackage{todonotes}
\usepackage{pdfpages}
\usepackage[export]{adjustbox}
% \makeatother


% \geometry{
%   left=24.8mm, % left margin
%   textwidth=100mm, % main text block
%   marginparsep=8.2mm, % gutter between main text block and margin notes
%   marginparwidth=49.4mm % width of margin notes
% }


\colorlet{myred}{red!65!black}
\tikzstyle{ground}=[preaction={fill,top color=black!10,bottom color=black!5,shading angle=20},
                    fill,pattern=north east lines,draw=none,minimum width=0.3,minimum height=0.6]
\tikzstyle{mass}=[line width=0.6,red!30!black,fill=red!40!black!10,rounded corners=1,
                  top color=red!40!black!20,bottom color=red!40!black!10,shading angle=20]
\tikzstyle{rope}=[brown!70!black,line width=1.2,line cap=round] %very thick

% \newcommand\rightAngle[4]{
%   \pgfmathanglebetweenpoints{\pgfpointanchor{#2}{center}}{\pgfpointanchor{#3}{center}}
%   \coordinate (tmpRA) at ($(#2)+(\pgfmathresult+45:#4)$);
%   %\draw[white,line width=0.6] ($(#2)!(tmpRA)!(#1)$) -- (tmpRA) -- ($(#2)!(tmpRA)!(#3)$);
%   \draw[blue!40!black] ($(#2)!(tmpRA)!(#1)$) -- (tmpRA) -- ($(#2)!(tmpRA)!(#3)$);
% }
%%%%%%%%%%%%%%%%%%%quote


\pgfkeys{%
    /calloutquote/.cd,
    width/.code                   =  {\def\calloutquotewidth{#1}},
    position/.code                =  {\def\calloutquotepos{#1}}, 
    author/.code                  =  {\def\calloutquoteauthor{#1}},
    /calloutquote/.unknown/.code   =  {\let\searchname=\pgfkeyscurrentname
                                 \pgfkeysalso{\searchname/.try=#1,                                
    /tikz/\searchname/.retry=#1},\pgfkeysalso{\searchname/.try=#1,
                                  /pgf/\searchname/.retry=#1}}
                            }  


\newcommand\calloutquote[2][]{%
       \pgfkeys{/calloutquote/.cd,
         width               = 5cm,
         position            = {(0,-1)},
         author              = {}}
  \pgfqkeys{/calloutquote}{#1}                   
  \node [rectangle callout,callout relative pointer={\calloutquotepos},text width=\calloutquotewidth,/calloutquote/.cd,
     #1] (tmpcall) at (0,0) {#2};
  \node at (tmpcall.pointer){\calloutquoteauthor};    
}  





%%%%%%%%%%%%%%%%%%
% PULLEY
\def\r{0.05} % pulley small radius
\tikzset{
  pics/pulley/.style={
    code={
      \draw[pulcol,line width=0.6] (0,0) circle (#1);
      \draw[pulcol,thick] (0,0) circle (\r);
  }},
  pics/mount/.style args={#1:#2}{ % angle, length
    code={
      \draw[mount] (0,0)++(#1-90:0.9*\r) arc (#1-90:#1-270:0.9*\r) --++ (#1:#2) --++ (#1-90:1.8*\r) -- cycle;
  }},
  pics/weight/.style args={#1,#2,#3}{ % bottom width, top width, height
    code={
      \draw[mass] (0,0) -- (#2/2,0) -- (#1/2,-0.7*#3)
        |- (-#1/2,-#3) -- (-#1/2,-0.7*#3) -- (-#2/2,0) -- cycle;
      \path[mass] (0,0) -- (0,-#3) node[pos=0.52] {$m$};
  }},
  pics/pulley/.default=0.3,
}

% FORCES SWITCH
\tikzstyle{force}=[->,myred,thick,line cap=round]
\tikzstyle{Fproj}=[force,myred!40]
\newcommand{\vbF}{\vb{F}}
\usetikzlibrary{calc} % for pic
\usetikzlibrary{angles,quotes} % for pic
\usetikzlibrary{patterns}
\tikzstyle{metal}=[fill,top color=black!40,bottom color=black!20,shading angle=10]
\tikzstyle{pulcol}=[draw=blue!30!black,%fill=blue!40!black!10
                    top color=blue!40!black!20,bottom color=blue!40!black!10,shading angle=20]
\def\rope#1{ \draw[black,line width=1.5] #1; \draw[rope] #1; }
\tikzstyle{mount}=[blue!20!black,fill,top color=blue!20!black!70,bottom color=blue!20!black!40,shading angle=10] %,line width=1.8,line cap=round
%\tikzstyle{mount}=[color=black!60,line width=1.8,line cap=round]
\tikzstyle{spring}=[line width=0.8,black!80,snake=coil,segment amplitude=5,segment length=5,line cap=round]
\pgfdeclarelayer{back} % to draw on background
\pgfsetlayers{back,main}

\tikzstyle{myarr}=[{Latex[length=3,width=3]}-,xcol]
\tikzstyle{myarr2}=[{Latex[length=2,width=3]}-{Latex[length=2,width=3]}]

\colorlet{xcol}{blue!70!black}
\colorlet{vcol}{green!60!black}
\colorlet{myred}{red!65!black}
\colorlet{mypurple}{blue!60!red!80}
\colorlet{acol}{red!50!blue!80!black!80}
\tikzstyle{CM}=[red!40!black,fill=red!80!black!80]
\tikzstyle{rvec}=[->,xcol,very thick,line cap=round]
\tikzstyle{force}=[->,myred,very thick,line cap=round]
\tikzstyle{mass}=[line width=0.6,red!30!black,fill=red!40!black!10,rounded corners=1,
                  top color=red!40!black!20,bottom color=red!40!black!10,shading angle=20]
\tikzstyle{rope}=[brown!30!black,fill=brown!70!black,line width=0.4]
\tikzstyle{limb}=[thick,line cap=round]
\tikzstyle{force}=[->,myred,very thick,line cap=round]
\tikzstyle{Fproj}=[force,myred!40]
\tikzstyle{myarr}=[{Latex[length=3,width=3]}-]
\tikzstyle{CM}=[red!40!black,fill=red!80!black!80]
\tikzstyle{mydashed}=[dash pattern=on 2 off 2]
\tikzstyle{mass}=[line width=0.6,draw=red!30!black,rounded corners=1,
                  top color=red!40!black!30,bottom color=red!40!black!10,shading angle=30]
\tikzstyle{ground}=[preaction={fill,top color=black!10,bottom color=black!5,shading angle=20},
                    fill,pattern=north east lines,draw=none,minimum width=0.3,minimum height=0.6]
\tikzstyle{metal}=[fill,top color=black!40,bottom color=black!20,shading angle=10]
\def\tick#1#2{\draw[thick] (#1) ++ (#2:0.1) --++ (#2-180:0.2)}


\usetikzlibrary{angles,quotes} % for pic
\usetikzlibrary{arrows.meta}
\usetikzlibrary{patterns}
\tikzset{>=latex} % for LaTeX arrow head

\tikzset{
  pics/Tin/.style={
    code={
      \def\R{0.12}
      \draw[pic actions,line width=0.6,#1,fill=white] % ,thick
        (0,0) circle (\R) (-135:.75*\R) -- (45:.75*\R) (-45:.75*\R) -- (135:.75*\R);
  }},
  pics/Tout/.style={
    code={
      \def\R{0.12}
      \draw[pic actions,line width=0.6,#1,fill=white] (0,0) circle (\R);
      \fill[pic actions,#1] (0,0) circle (0.3*\R);
  }},
  pics/Tin/.default=mypurple,
  pics/Tout/.default=mypurple,
}












\titleformat{\chapter}
  [block]% shape
  {\relax\ifthenelse{\NOT\boolean{@tufte@symmetric}}{\begin{fullwidth}}{}}% format applied to label+text
  {\itshape\huge\thechapter}% label
  {1em}% horizontal separation between label and title body
  {\huge\rmfamily\itshape}% before the title body
  [\ifthenelse{\NOT\boolean{@tufte@symmetric}}{\end{fullwidth}}{}]% after the title body


\tikzset{>=stealth', pil/.style={ ->, thick, shorten <=2pt, shorten >=2pt,} }
% tikz decription node style
\tikzset{note/.style={rectangle, rounded corners, minimum size=6mm, draw=black, fill=white, align=center,execute at begin node=\setlength{\baselineskip}{1.2em}}}
\tikzset{twoline/.style={align=center,execute at begin node=\setlength{\baselineskip}{1.2em}}}
\tikzset{twolinecap/.style={align=center,execute at begin node=\setlength{\baselineskip}{1.6em}}}
\tikzset{graphlabel/.style={postaction={decorate,transform shape, decoration={markings, mark=at position .5 with \node #1;}}}}
%%%%%%%%%%%%%%%%%%%%%

%%%%%%%%%%%%%%%%%%%%%%
\usetikzlibrary{decorations}
\usetikzlibrary{arrows}
\usetikzlibrary{snakes}
\usetikzlibrary{decorations.pathmorphing}

% % % % % % % % % % % % % % % % % % % % % % % % % % %
\makeatletter
% counter for margin figures
\newcounter{marfigure}[chapter]
\renewcommand\themarfigure
     {\ifnum \c@chapter>\z@ \thechapter.\fi \@arabic\c@marfigure}
% definitions for the new float marfigure
\def\fps@marfigure{tbp}
\def\ftype@marfigure{1}
\def\ext@marfigure{lmf}
\def\fnum@marfigure{\figurename\nobreakspace\themarfigure}
\renewenvironment{marginfigure}[1][-1.2ex]%
  {\begin{@tufte@margin@float}[#1]{marfigure}}
  {\end{@tufte@margin@float}}
% definitions for the new list of margin figures
\newcommand\listmarfigurename{List of Margin Figures}
\let\l@marfigure\l@figure
\newcommand\listofmarfigures{%
  \ifthenelse{\equal{\@tufte@class}{book}}%
    {\chapter*{\listmarfigurename}}%
    {\section*{\listmarfigurename}}%
%  \begin{fullwidth}%
    \@starttoc{lmf}%
%  \end{fullwidth}%
}
\makeatother



% % % % % % % % % % % % % % % % % % % % % % % % % % %

\newcommand{\syl}[1]{\footnote{Specification statement: #1}}
\newcommand{\sq}{\begin{Large} $\Square$ \end{Large}}
%%%%%%%%%%%%%%%%%%%%%%%%

%%
% Prints argument within hanging parentheses (i.e., parentheses that take
% up no horizontal space).  Useful in tabular environments.
\newcommand{\hangp}[1]{\makebox[0pt][r]{(}#1\makebox[0pt][l]{)}}


%%
% Prints an asterisk that takes up no horizontal space.
% Useful in tabular environments.
\newcommand{\hangstar}{\makebox[0pt][l]{*}}
%%
% Prints a trailing space in a smart way.
\usepackage{xspace}
% example + question + exercise environment
\newcounter{example}[chapter]
\newcounter{exercise}[chapter]
\newcounter{question}[chapter]
\renewcommand{\theexample}{\arabic{chapter}.\arabic{example}}
\renewcommand{\theexercise}{\arabic{chapter}.\arabic{exercise}}
\renewcommand{\thequestion}{\arabic{chapter}.\arabic{question}}

\newcommand{\example}[1]{%
	\refstepcounter{example}
	\noindent {\textcolor{cyan}{\textsf{\textbf{Example \theexample }}} \hspace*{1pt} #1 %
	} \vspace{5pt}}
 
\newcommand{\question}[1]{%
	\refstepcounter{question}
	\noindent{\textcolor{magenta}{\textsf{\textbf{Question \thequestion }}} \hspace*{1pt} #1 %
}}
\newcommand{\exercise}[1]{%
	\refstepcounter{exercise}
	\noindent{\textsf{\textbf{Exercise \theexercise }} \hspace*{1pt} #1 %
}}

\newcommand{\cmt}{}
\newcommand{\piste}{{\fontencoding{U}\fontfamily{futs}\selectfont\char 49\relax}}

\newcommand{\sol}{\hspace{2cm}\textcolor{Green}{\ding{52}}}
\newcommand{\solc}{\noindent\hspace{0.12em}\textcolor{Green}{\ding{52}} \hspace{0.2em} \vspace*{-\baselineskip}} 

% highlight environment
\newenvironment{ilight}
  {\centering
  	\vspace*{6pt}
  	\begin{tcolorbox}[colframe=Gray,colback=LightGrey!15]
  \setlength{\baselineskip}{\baselineskip}%
  }
  {\end{tcolorbox}\vspace*{-4pt}}

\newtcbox{\mymath}[1][]{%
    nobeforeafter, math upper, tcbox raise base,
    enhanced, colframe=blue!30!black,
    colback=blue!30, boxrule=1pt,
    #1}


\newenvironment{soln}
  {\centering
  	\vspace*{6pt}
  	\begin{tcolorbox}[colframe=Gray,colback=LightGreen!15]
  \setlength{\baselineskip}{\baselineskip}%
  }
  {\end{tcolorbox}\vspace*{-4pt}}

% columns with customised width in tabulars
\newcolumntype{C}[1]{>{\centering\arraybackslash}p{#1}}
\newcolumntype{D}[1]{>{\centering\arraybackslash}m{#1}}

\pgfplotsset{compat=1.18}

\newcommand{\keypoint}[1]{\textbf{\textcolor{red}{#1}}}


\newcommand{\TL}{Tufte-\LaTeX\xspace}

% Prints the month name (e.g., January) and the year (e.g., 2008)
\newcommand{\monthyear}{%
  \ifcase\month\or January\or February\or March\or April\or May\or June\or
  July\or August\or September\or October\or November\or
  December\fi\space\number\year
}


% Prints an epigraph and speaker in sans serif, all-caps type.
\newcommand{\openepigraph}[2]{%
  %\sffamily\fontsize{14}{16}\selectfont
  \begin{fullwidth}
  \sffamily\large
  \begin{doublespace}
  \noindent\allcaps{#1}\\% epigraph
  \noindent\allcaps{#2}% author
  \end{doublespace}
  \end{fullwidth}
}

% Inserts a blank page
\newcommand{\blankpage}{\newpage\hbox{}\thispagestyle{empty}\newpage}

\usepackage{units}

% Typesets the font size, leading, and measure in the form of 10/12x26 pc.
\newcommand{\measure}[3]{#1/#2$\times$\unit[#3]{pc}}

% Macros for typesetting the documentation
\newcommand{\hlred}[1]{\textcolor{Maroon}{#1}}% prints in red
\newcommand{\hangleft}[1]{\makebox[0pt][r]{#1}}
\newcommand{\hairsp}{\hspace{1pt}}% hair space
\newcommand{\hquad}{\hskip0.5em\relax}% half quad space
\newcommand{\TODO}{\textcolor{red}{\bf TODO!}\xspace}
\newcommand{\ie}{\textit{i.\hairsp{}e.}\xspace}
\newcommand{\eg}{\textit{e.\hairsp{}g.}\xspace}
\newcommand{\na}{\quad--}% used in tables for N/A cells
\providecommand{\XeLaTeX}{X\lower.5ex\hbox{\kern-0.15em\reflectbox{E}}\kern-0.1em\LaTeX}
\newcommand{\tXeLaTeX}{\XeLaTeX\index{XeLaTeX@\protect\XeLaTeX}}
% \index{\texttt{\textbackslash xyz}@\hangleft{\texttt{\textbackslash}}\texttt{xyz}}
\newcommand{\tuftebs}{\symbol{'134}}% a backslash in tt type in OT1/T1
\newcommand{\doccmdnoindex}[2][]{\texttt{\tuftebs#2}}% command name -- adds backslash automatically (and doesn't add cmd to the index)
\newcommand{\doccmddef}[2][]{%
  \hlred{\texttt{\tuftebs#2}}\label{cmd:#2}%
  \ifthenelse{\isempty{#1}}%
    {% add the command to the index
      \index{#2 command@\protect\hangleft{\texttt{\tuftebs}}\texttt{#2}}% command name
    }%
    {% add the command and package to the index
      \index{#2 command@\protect\hangleft{\texttt{\tuftebs}}\texttt{#2} (\texttt{#1} package)}% command name
      \index{#1 package@\texttt{#1} package}\index{packages!#1@\texttt{#1}}% package name
    }%
}% command name -- adds backslash automatically
\newcommand{\doccmd}[2][]{%
  \texttt{\tuftebs#2}%
  \ifthenelse{\isempty{#1}}%
    {% add the command to the index
      \index{#2 command@\protect\hangleft{\texttt{\tuftebs}}\texttt{#2}}% command name
    }%
    {% add the command and package to the index
      \index{#2 command@\protect\hangleft{\texttt{\tuftebs}}\texttt{#2} (\texttt{#1} package)}% command name
      \index{#1 package@\texttt{#1} package}\index{packages!#1@\texttt{#1}}% package name
    }%
}% command name -- adds backslash automatically
\newcommand{\docopt}[1]{\ensuremath{\langle}\textrm{\textit{#1}}\ensuremath{\rangle}}% optional command argument
\newcommand{\docarg}[1]{\textrm{\textit{#1}}}% (required) command argument
\newenvironment{docspec}{\begin{quotation}\ttfamily\parskip0pt\parindent0pt\ignorespaces}{\end{quotation}}% command specification environment
\newcommand{\docenv}[1]{\texttt{#1}\index{#1 environment@\texttt{#1} environment}\index{environments!#1@\texttt{#1}}}% environment name
\newcommand{\docenvdef}[1]{\hlred{\texttt{#1}}\label{env:#1}\index{#1 environment@\texttt{#1} environment}\index{environments!#1@\texttt{#1}}}% environment name
\newcommand{\docpkg}[1]{\texttt{#1}\index{#1 package@\texttt{#1} package}\index{packages!#1@\texttt{#1}}}% package name
\newcommand{\doccls}[1]{\texttt{#1}}% document class name
\newcommand{\docclsopt}[1]{\texttt{#1}\index{#1 class option@\texttt{#1} class option}\index{class options!#1@\texttt{#1}}}% document class option name
\newcommand{\docclsoptdef}[1]{\hlred{\texttt{#1}}\label{clsopt:#1}\index{#1 class option@\texttt{#1} class option}\index{class options!#1@\texttt{#1}}}% document class option name defined
\newcommand{\docmsg}[2]{\bigskip\begin{fullwidth}\noindent\ttfamily#1\end{fullwidth}\medskip\par\noindent#2}
\newcommand{\docfilehook}[2]{\texttt{#1}\index{file hooks!#2}\index{#1@\texttt{#1}}}
\newcommand{\doccounter}[1]{\texttt{#1}\index{#1 counter@\texttt{#1} counter}}
%%%%% code boxes
\tcbuselibrary{minted,breakable,xparse,skins}

\definecolor{bg}{gray}{0.95}
\DeclareTCBListing{mintedbox}{O{}m!O{}}{%
  breakable=true,
  listing engine=minted,
  listing only,
  minted language=#2,
  minted style=default,
  minted options={%
    linenos,
    gobble=0,
    breaklines=true,
    breakafter=,,
    fontsize=\small,
    numbersep=8pt,
    #1},
  boxsep=0pt,
  left skip=0pt,
  right skip=0pt,
  left=25pt,
  right=0pt,
  top=3pt,
  bottom=3pt,
  arc=5pt,
  leftrule=0pt,
  rightrule=0pt,
  bottomrule=2pt,
  toprule=2pt,
  colback=bg,
  colframe=orange!70,
  enhanced,
  overlay={%
    \begin{tcbclipinterior}
    \fill[orange!20!white] (frame.south west) rectangle ([xshift=20pt]frame.north west);
    \end{tcbclipinterior}},
  #3}

\tikzstyle{wave}=[myblue,thick]
\tikzstyle{mydashed}=[black!70,dashed,thin]
\tikzstyle{mymeas}=[{Latex[length=3,width=2]}-{Latex[length=3,width=2]},thin]
\tikzstyle{mysmallarr}=[-{Latex[length=3,width=2]}]

\tikzset{
  declare function={
    int_arg(\y,\lam,\a,\L) = \a*\y/sqrt(\L*\L+\y*\y)/\lam; %sin(\x);
    int_one(\y,\lam,\a,\L) = (sin(180*int_arg(\y,\lam,\a,\L))/(pi*int_arg(\y,\lam,\a,\L)))^2;
    int_two(\y,\lam,\a,\L) = cos(180*int_arg(\y,\lam,\a,\L))^2;
    int_arg_ang(\t,\lam,\a) = \a*sin(\t)/\lam;
    int_one_ang(\t,\lam,\a) = (sin(180*int_arg_ang(\t,\lam,\a))/(pi*int_arg_ang(\t,\lam,\a)))^2;
    int_two_ang(\t,\lam,\a) = cos(180*int_arg_ang(\t,\lam,\a))^2;
  }
}

\usetikzlibrary{fadings}
\tikzset{>=latex} % for LaTeX arrow head

\colorlet{wall}{blue!30!black}
\colorlet{myblue}{blue!70!black}
\colorlet{myred}{red!65!black}
\colorlet{mypurple}{red!50!blue!95!black!75}
\colorlet{myshadow}{blue!30!black!90}
\colorlet{mydarkred}{red!50!black}
\colorlet{mylightgreen}{green!60!black!70}
\colorlet{mygreen}{green!60!black}
\colorlet{myredgrey}{red!50!black!80}

% Generates the index
\usepackage{makeidx}
\makeindex

\makeatletter
% Paragraph indentation and separation for normal text
\renewcommand{\@tufte@reset@par}{%
  \setlength{\RaggedRightParindent}{0pc}%
  \setlength{\JustifyingParindent}{0pc}%
  \setlength{\parindent}{0pt}%
  \setlength{\parskip}{0pt}%
}
\@tufte@reset@par

% Paragraph indentation and separation for marginal text
\renewcommand{\@tufte@margin@par}{%
  \setlength{\RaggedRightParindent}{0pc}%
  \setlength{\JustifyingParindent}{0pc}%
  \setlength{\parindent}{0pc}%
  \setlength{\parskip}{0pt}%
}
\makeatother

\begin{document}

% Front matter
%\frontmatter

% r.1 blank page
%\blankpage

% v.2 epigraphs
%\newpage\thispagestyle{empty}

% r.3 full title page
\maketitle


% v.4 copyright page
\newpage
\begin{fullwidth}
~\vfill
\thispagestyle{empty}
\setlength{\parindent}{0pt}
\setlength{\parskip}{\baselineskip}
Copyright \copyright\ \the\year\ Joe Rowing

\par\smallcaps{Published by Exeter Maths School}

\par\smallcaps{exetermathematicsschool.com}

\par Creative Commons – Attribution Non Commercial Share Alike – CC BY-NC-SA

CC BY-NC-SA 4.0

The CC BY-NC-SA licence allows the work to be re-mixed, re-used, and re-shared so long as  attribution is provided to the creator, the work is not used for commercial purposes, and any new work is re-shared freely under the same licence..\index{license}

\par\textit{First printing, \monthyear}
\end{fullwidth}

% r.5 contents
\tableofcontents

%\listoffigures

%\listoftables

% r.7 dedication
%\cleardoublepage
%~\vfill
%\begin{doublespace}
%\noindent\fontsize{18}{22}\selectfont\itshape
%\nohyphenation
%\vfill
%\vfill
%\mainmatter
%%
% The back matter contains appendices, bibliographies, indices, glossaries, etc.
%%%%%%%%%%%       
\chapter{Openings}

This book/these notes free for use and modification by students and teachers at Exeter Maths School. Source code available on request


Mostly this book is self-explanatory but beware - "\piste" symbol denotes material off-spec.



\include{physical-quantities}
\chapter{Kinematics}

Kinematics is the study of motion. We will introduce three kinematic quantities, displacement, velocity and acceleration, and see how these terms are used to describe an object's motion.

\section{Kinematic quantities}

\subsection{Displacement \& Distance}

In everyday language, we talk about the \keypoint{distance} travelled by an object, which usually refers to the length travelled by an object without considering in which direction it moves, to fully describe position of an object however, we also need specify where it is heading.

\begin{ilight}
	\centering \keypoint{displacement} is defined as the distance moved by an object in a specific direction\index{displacement}
\end{ilight}

 Displacement is the \emph{straight-line} distance pointing from starting point towards end point -even if actual path taken is curved, ie it is a vector quantity, while distance is a scalar. Displacement is always the straight-line distance. Displacement and distance are measured in metres, or any reasonable length units.


 \begin{figure*}[!ht]
	\centering
	\begin{tikzpicture}[scale=1]
%	\draw[step=0.5] (0,0) grid (5,5);
	 \draw[thick, dashed,->] (0,0) node[below]{start} [out=150, in=140] to (0.5,2) [out=-40, in=210] to (2.5,4) [out=30, in=185] to (4,3) [out=-5, in=160] to (5.5,4) [out=-20, in=55] to (5,2) node[right]{end};
	 \draw[thick, blue, ->] (0,0) -- (5,2) ;
	 \node[note] at (0.6,2.6) {distance};
	 \node[note] at (3.2,0.5) {displacement};
	 \draw (5.7,3) --++ (1,0.6) node[right] {actual path taken};
	 \draw (4,1.55) --++ (2.2,-0.4) node[right] {straight-line connector};
	\end{tikzpicture}
	
	\caption{difference between displacement and distance}
\end{figure*}


\example{An athlete is running around a circular track of radius 60 m. When he completes one lap, what is the distance moved out? What about his displacement?}

\begin{soln}
    
Distance moved is the perimeter of the circle: $s = 2\pi r = 2\pi\times 60 \approx 380 \text{ m}$

athlete returns to same starting point after one lap, so displacement is zero
\end{soln}

\begin{figure}
%	\vspace*{-12pt}
	\centering
	\begin{tikzpicture}[scale=1.1]
	\draw[thick,->] (0,0) -- (0,3) node[midway, left]{30 km};
	\draw[thick,->] (0,3) -- (4,3) node[midway, above]{40 km};
	\draw[thick,blue,->] (0,0) -- (4,3) node[midway, below, rotate=36.89]{50 km};
	\draw (0,0.5) arc(90:36.89:0.5);
	\node at (63:0.8) {$\theta$};
	\end{tikzpicture}
	\vspace*{-16pt}
\end{figure}

\example{A ship travels 30 km north, takes a right, and then travels 40 km east to reach its destination. Compare the distance and the displacement travelled.}

\begin{soln}
Summing of all lengths gives distance: $30+40 = 70 \text{ km}$

The displacement vector is shown on the diagram.

$\text{The magnitude of displacement} = \sqrt{30^2 + 40^2} = 50 \text{ km}$

at an angle of $\theta = \tan^{-1}\frac{40}{30} \approx 53^\circ$ east of north
\end{soln} 



\subsection{Velocity \& speed}

The displacement of a moving object may change with respect to time, the rate of this change in displacement gives us the velocity with which it's moving.

\begin{ilight}
	The change in displacement per unit time is called the \keypoint{velocity} of the object velocity: \begin{empheq}[box=\tcbhighmath]{equation*}{v = {\Delta s}/{\Delta t}}\end{empheq}
\end{ilight}

Velocity is a vector quantity which follows from the fact that displacement is a vector quantity. The SI unit of measurement for velocity: $[v]= \mps$ 

In the case of vector quantities in one dimension, \emph{linear} motion, we always assign a direction and indicate this with positive and negative signs. A negative velocity implies motion in the opposite direction to positive.

Defining equation $v=\frac{\Delta s}{\Delta t}$ gives the \emph{average} value for velocity or speed over an interval $\Delta t$ more precisely: \begin{empheq}[box=\tcbhighmath]{equation*}{\text{average velocity} = \frac{\text{total displacement}}{\text{time taken}}}\end{empheq} and... \begin{empheq}[box=\tcbhighmath]{equation*}{\text{average speed} = \frac{\text{total distance}}{\text{time taken}}}\end{empheq}

This should be distinguished from the notion of \emph{instantaneous velocity}instantaneous velocity; defined as the rate of change in displacement at a particular instant.

If we take a very short interval $\Delta t$, as $\Delta t \to 0$, average velocity tends to instantaneous velocity, this is expressed in a compact differential form: $v = \lim_{\Delta t \to 0} \frac{\Delta s}{\Delta t} \ra$ \begin{empheq}[box=\tcbhighmath]{equation*}{v=\frac{\dd s}{\dd t}} \end{empheq}


\newpage

It is also common to use \keypoint{speed} to describe how fast an object moves, the speed is defined as the change of the distance travelled per unit time. Velocity can be thought as speed in a particular direction.

\example{A cyclist travels a distance of 3.0 km in 20 minutes. She rests for 15 minutes. She then covers a further distance of 5.1 km in a time of 40 minutes. Calculate the average speed of the cyclist in m s$^{-1}$: (a) during the first 20 minutes of the journey, (b) for the whole journey.}

\begin{soln} For the first 20 minutes: $ v = \frac{3.0\times10^3}{20 \times 60} = 2.5 \mps$

\eqyskip For whole journey: $v = \frac{(3.0+0+5.1)\times10^3}{(20+15+40)\times60} = 1.8 \mps $ 
\end{soln}

\example{A man walks along a straight road for a distance of 800 m in 5.0 minutes. He then turns around, and walks along the same road for a distance of 280 m in 3.0 minutes. What is the average speed and the average velocity of this man during the 8.0 minutes?}

\begin{soln} $\text{Total distance travelled} = 800 + 280 = 1080 \text{m}$, so average speed: $v = \frac{1080}{8.0\times60} = 2.25 \mps$

\eqyskip  $\text{change of displacement} = 800 + (-280) = 520 \text{m}$, so average velocity: $v = \frac{520}{8.0\times60} \approx 1.08 \mps$ \end{soln}

\example{A maglev train travels at an average speed of $60\mps$ from the city centre to the airport, and at $40\mps$ on its return journey over the same distance. What is the average speed of the round-trip? What about the average velocity?}

\begin{soln} Suppose the distance between airport and city centre is $S$.

Average speed: $v = \frac{2S}{t_1 + t_2} = \frac{2S}{\frac{S}{60} + \frac{S}{40}} = 48 \mps$

For a round-trip, train returns to same staring position so change in displacement is zero. Average velocity is therefore \emph{Zero} \end{soln} 


\subsection{Acceleration}

The velocity of a moving object may change, i.e., objects can speed up, slow down \textbf{or make turns}

\begin{ilight}
	\centering Change in velocity per unit time is defined as the \keypoint{acceleration}\index{acceleration}: \begin{empheq}[box=\tcbhighmath]{equation*}{a = \frac{\Delta v}{\Delta t}}\end{empheq}
\end{ilight}

\titem{unit of measurement for acceleration: $[a]= \mpss$}

Acceleration\footnote{Joe often writes this as "$Accl^{n}$" on the board. Useful shorthand that!} is a vector quantity, it has both magnitude and direction. Because of vector nature of velocity, change in velocity must also have direction. When we're considering \emph{linear} motion, we usually pick direction of initial velocity as positive direction, often left to right.

\titem $a>0$ would imply acceleration in the normal sense, i.e., motion with an increasing speed.

\titem $a<0$ would imply deceleration, i.e., motion with a decreasing speed.

When velocity changes, it could be change in magnitude or/and change in direction\footnote{Acceleration of an object can be considered as the combination of two components. One component is known as the \emph{normal} acceleration or the \emph{centripetal} acceleration, which is at right angle to the velocity and is responsible for the change in direction of motion. The other component is called the \emph{tangential} acceleration, which is parallel to the direction of motion and causes change in magnitude of object's velocity. You will learn more about these in further mechanics.} for example, for an object moving along a curved path, its velocity is constantly changing direction, so it must have a non-zero acceleration. An acceleration of Zero on the other hand would imply no change in speed and no change in direction of motion.

The defining equation, which gives average acceleration over time interval $\Delta t$ is: $$a = \frac{\Delta v}{\Delta t}$$ 

We can likewise introduce \emph{instantaneous acceleration} as the rate of change in velocity. Taking the limit where the time interval $\Delta t \to 0$. we have: $a = \lim_{\Delta t \to 0} \frac{\Delta v}{\Delta t} \ra$ \begin{empheq}[box=\tcbhighmath]{equation*}{a=\frac{\dd v}{\dd t}} \end{empheq}

\example{A ball hits a barrier at right angles with a speed of $15 \mps$. It makes contact with the barrier for 30 ms and then rebounds with a speed of $12 \mps$. What is the average acceleration during the time of contact?}

\begin{soln} 
note that direction of velocity changed during rebound, so $\Delta v = 15 - (-12) = 27 \mps$

average acceleration: $a = \frac{\Delta v}{\Delta t} = \frac{27}{30 \times 10^{-3}} = 900 \mpss$ \end{soln}



\subsection{Motion graphs}

How one quantity changes with another quantity can be visually shown on a \emph{graph}, changes in displacement, velocity or acceleration over time can be shown on \emph{motion graphs}\index{motion graphs}. As we will see, $s$-$t$ graph, $v$-$t$ graph and $a$-$t$ graphs are closely interrelated to one another.

\subsection{Displacement-time graphs}

A displacement-time graph shows an object's position at any given time:

\colorlet{myblue}{blue!70!black}
\colorlet{mydarkblue}{blue!40!black}
\colorlet{mygreen}{green!60!black}
\colorlet{myred}{red!65!black}
\colorlet{mypurple}{red!50!blue!95!black!75}
\tikzstyle{wave}=[myblue,thick]
\tikzstyle{xline}=[very thick,myblue]
\tikzstyle{vline}=[very thick,mygreen]
\tikzstyle{aline}=[very thick,mypurple]
\tikzstyle{mydashed}=[mydarkblue,dashed]

\def\tick#1#2{\draw[thick] (#1) ++ (#2:0.1) --++ (#2-180:0.2)}
\def\tlabel{$t\,\left[\text{s}\right]$}
\def\vlabel{$d\,\left[\text{m}\right]$}
\def\alabel{$a\,\left[\si{m/s^2}\right]$}

\begin{figure}
    \centering
    \def\xmax{5.5} % maximum x axis (time)
\def\ymax{3.0} % maximum y axis (velocity)
\def\tmax{16} % maximum time on x axis
\begin{tikzpicture}
  \def\vmax{10} % maximum velocity on x axis
  \def\xscale{\xmax/\tmax}
  \def\yscale{\ymax/\vmax}
  
  % GRID
  \foreach \t in {0,2,...,\tmax}{
    \draw[black!80,dotted] (\t*\xscale,0) --++ (0,\vmax*\yscale+0.3);
    \tick{\t*\xscale,0}{90} node[below=0,scale=0.77] {\t};
  }
  \foreach \v in {0,2,...,\vmax}{
    \draw[black!80,dotted] (0,\v*\yscale) --++ (\tmax*\xscale+0.3,0);
    \tick{0,\v*\yscale}{0} node[left=0,scale=0.77] {\v};
  }
  
  % AXES
  \draw[->,thick]
    (0,0) -- (\xmax+0.5,0) node[midway,below=1] {\tlabel};
  \draw[->,thick]
    (0,0) -- (0,\ymax+0.5) node[midway,rotate=90,above=0.8] {\vlabel};
  
  % GRAPH
  \draw[vline,line cap=round,xscale=\xscale,yscale=\yscale]
    (0,0) -- (2,8) -- (8,8) -- (12,4) -- (16,4);
  
\end{tikzpicture}
    \caption{A generic displacement-time graph}
    \label{fig:enter-label}
\end{figure}
The gradient of the graph (or the tangent) gives rate of change in displacement, this is instantaneous velocity, which can be given by $v = \frac{\dd s}{\dd t}$, so we have:

{
	\centering
	
	\eqyskip \begin{empheq}[box=\tcbhighmath]{equation*}{\text{velocity} = \text{gradient of $s$-$t$ graph}} \end{empheq}
	
}

\newpage


\example{Describing the motion from the displacement-time graph shown.}

\begin{marginfigure}
	\vspace{-24pt}
	\begin{center}
		\begin{tikzpicture}[xscale=1.1,yscale=0.95]
		\draw[<->] (0,3.2)node[left]{s} -- (0,0) -- (4.2,0)node[below]{$t$};
		\draw[blue, thick] plot [smooth] coordinates {(0,0) (.5,.2) (1,.8) (1.5,2.0) (2,2.5) (2.5,2.7) (3.0,2.8) (4,2.8)};
		\node at (.8,1) {$A$};
		\node at (2,2.8) {$B$};
		\node at (3.5,3) {$C$};
		\end{tikzpicture}
	\end{center}
	\vspace{-27pt}
\end{marginfigure}

	stage $A$: gradient of the graph is increasing, showing that the object is speeding up
		
	stage $B$: gradient starts to decrease, so the object gradually slows down
		
	stage $C$: curve becomes horizontal, gradient becomes zero, means that the object eventually comes to a stop 
	
\begin{marginfigure}[1cm]
	\centering
	\begin{tikzpicture}[scale=1.1]
		\draw[gray!50,step=0.5] (0,0) grid (4,5);
		\draw[thick, blue,smooth, samples=20, domain=0:4] plot (\x, 0.3*\x*\x);
		\foreach \x in {0,0.5,1,1.5,2,2.5,3,3.5,4} \node[below] at (\x,0){$\x$};
		\foreach \y in {0,2,...,20} \node[left] at (0,\y/4) {$\y$};
		\node at (2,-0.6) {$t$/s};
		\node[left] at (-0.5,2.5) {$s$/m};
		\draw[thick,red] (0.75,0) -- (4,2.925);
	\end{tikzpicture}
	
\end{marginfigure}

\example{The diagram shows the displacement-time graph for a vehicle travelling along a straight road. Use the graph to find, (a) the average velocity during the first 4.0 s of the motion, (b) the velocity of the vehicle \emph{at} time $t=1.5$ s.}

\begin{soln} 
during first 4.0 s, average velocity is

{
	\centering
	
	$ v = \frac{\Delta s}{\Delta t} = \frac{19.2}{4.0} \approx 4.8 \mps$
	
}

to find velocity at $t=1.5$ s, a tangent is drawn

gradient of tangent gives instantaneous velocity:

{
	\centering
	
	$ v = \frac{11.6 - 0}{4.0-0.75} \approx 3.6 \mps $
	
}
\end{soln}

\subsection{Velocity-time graphs}

A velocity-time graph shows the velocity of a moving object at any instant, since the rate of change of velocity gives the acceleration;

\begin{empheq}[box=\tcbhighmath]{equation*}\text{acceleration} = \text{gradient of $v$-$t$ graph}\end{empheq}


\noindent A $v$-$t$ graph also gives information about the change in displacement, that is:
\begin{empheq}[box=\tcbhighmath]{equation*}{\text{change in displacement} = \text{area under $v$-$t$ graph}}\end{empheq}


The reasoning here is that in very short interval $\Delta t_i$, change in velocity is small so $v(t_i)\approx\text{constant}$ during $\Delta t_i$. Following the logic, the displacement changed by $\Delta s_i \approx v(t_i) \Delta t_i$, which corresponds to area of a thin rectangle. The sum of all these small $\Delta s_i$'s gives total change in displacement over a period of time. If we now consider the limit where each of the time interval $\Delta t_i \to 0$ then the total area of these rectangles approximates area bounded by the $v$-$t$ curve and time axis.
\footnote{Mathematically, integration is the inverse operation of taking derivatives. By definition $v=\frac{\dd s}{\dd t}$, then it follows naturally that $\Delta s = \int v\dd t$. While the derivative of a given function gives the gradient of tangent at each point on its graph, integrating a function gives the signed area bounded by the graph. The reader may find the formal treatment of this relationship in any calculus textbook.}



\begin{figure*}[!ht]
	\noindent\centering
	\begin{minipage}{0.32\linewidth}
	\centering
	\begin{tikzpicture}[scale=1.4]
	\draw[thick,fill=blue!25] (1.5,0) rectangle (1.9,0.9375);
	\draw (1.7,0.5) -- (2.1,0.7) node[right]{$v(t)\Delta t$};
	\draw[<->] (1.5,1.2) -- (1.9,1.2) node[midway,above]{$\Delta t$};
	\draw[<->] (0.5,2.7) node[left]{$v$} -- (0.5,0) -- (3.6,0) node[below]{$t$};
	\draw [thick,domain=0.5:3.2,samples=15,smooth,variable=\x] plot (\x,{\x*(\x-1)*(\x-2)/6+1});
	\draw[thick,dashed] (1,0) node[below]{$t_1$} -- (1,1);
	\draw[thick,dashed] (3,0) node[below]{$t_2$} -- (3,2);
	\node[below] at (1.5,0) {$t$};
	\node[below] at (1.9,0.02) {$t+\Delta t$};
	\node at (1.9,2.1){$\Delta s_i \approx v(t_i)\Delta t_i$};
	\end{tikzpicture}
	(a) displacement $\Delta s$
	
	in short interval $\Delta t$
	\end{minipage}\hfill
	\begin{minipage}{0.34\linewidth}
		\centering
		\begin{tikzpicture}[scale=1.4]
		\foreach \x in {1.0,1.2,1.4,...,2.8}
		{ \draw[thick,fill=blue!25] (\x,0) rectangle (\x+.2,{\x*(\x-1)*(\x-2)/6+1});
			}
		\draw[thick,dashed] (1,0) node[below]{$t_1$} -- (1,1);
		\draw[thick,dashed] (3,0) node[below]{$t_2$} -- (3,2);
		\draw[<->] (0.5,2.7) node[left]{$v$} -- (0.5,0) -- (3.6,0) node[below]{$t$};
		\draw [thick,domain=0.5:3.2,samples=15,smooth,variable=\x] plot (\x,{\x*(\x-1)*(\x-2)/6+1});
		\node at (1.9,2.1){$\Delta s \approx \sum_i v(t_i)\Delta t_i$};
		\end{tikzpicture}
		(b) total displacement estimated	
		
		by summing the many $\Delta s_i$'s
	\end{minipage}\hfill
	\begin{minipage}{0.32\linewidth}
		\centering
		\begin{tikzpicture}[scale=1.4]
		\fill[fill=blue!25] (1,0) -- plot[domain=1:3,samples=15,smooth,variable=\x] (\x,{\x*(\x-1)*(\x-2)/6+1}) -- (3,0) -- cycle;
		\draw[thick] (1,0) node[below]{$t_1$} -- (1,1);
		\draw[thick] (3,0) node[below]{$t_2$} -- (3,2);
		\draw[<->] (0.5,2.7) node[left]{$v$} -- (0.5,0) -- (3.6,0) node[below]{$t$};
		\draw [thick,domain=0.5:3.2,samples=15,smooth,variable=\x] plot (\x,{\x*(\x-1)*(\x-2)/6+1});
		\node at (1.9,2.1){$\Delta s = \int_{t_1}^{t_2} v \dd t$};
		\end{tikzpicture}
		(c) total displacement as
		
		area under $v$-$t$ graph
	\end{minipage}
	\caption{calculating change in displacement by finding the area under velocity-time graph}\label{fig:area_of_vt}
\end{figure*}

\begin{marginfigure}[2cm]
	\vspace*{0pt}
	\centering
	\begin{tikzpicture}[scale=1.15]
	\draw[gray!50,step=0.5] (0,-2) grid (4,2);
	\foreach \x in {0,0.5,1,1.5,2,2.5,3,3.5,4} {
	\draw[white,fill] (\x-0.1,-0.35) rectangle (\x+0.1,-0.05);
	\node[below] at (\x,0){$\x$};
	}
	\foreach \y in {-4,-2,0,2,4} \node[left] at (0,\y/2) {$\y$};
	\draw[blue,thick] (0,0) -- (1,2) -- (2,2) -- (3.5,-1) -- (4,0);
	\node at (2,-2.4) {$t$/s};
	\node[left] at (-0.4,0) {$v$/m s$^{-1}$};
	\end{tikzpicture}
	\vspace*{-16pt}
\end{marginfigure}


\example{The velocity of a toy car is shown. For the journey shown on the graph, use the graph to find (a) the total distance travelled, and (b) the total displacement travelled.}

\begin{soln} Distance is estimated using area under $v$-$t$ graph

0$\sim$3.0 s: $s_1 = \frac{1}{2}\times{1.0+3.0}\times4.0 = 8.0 \text{ m}$

3.0$\sim$4.0 s: $s_2 = \frac{1}{2}\times{1.0}\times2.0 = 1.0 \text{ m}$

$\text{total distance} = 8.0 + 1.0 = 9.0 \text{ m}$

$\text{total displacement} = (+8.0) + (-1.0) = 7.0 \text{ m}$ 
\end{soln}

\subsection{Acceleration-time graphs}

In the same way as above, one can similarly plot an acceleration-time graph to give the changes in acceleration, ie. $a$-$t$ graphs can give information about changes in velocity. Similar discussions lead to the following conclusion:\footnote{Using area under $a$-$t$ graph to find changes in velocity is not required in the AS-Level physics syllabus. I am putting this in the notes mainly for the completeness of the discussions on motion graphs.}
\begin{empheq}[box=\tcbhighmath]{equation*}\text{change in velocity} = \text{area under $a$-$t$ graph}
\end{empheq}

\newpage

You'll meet these again \footnote{and properly!} in Mechanicslater on, however the relationships between displacement, velocity and acceleration graphs are summarised below:

\begin{figure*}[!ht]
	\centering
	\begin{tikzpicture}[force/.style={align=center,execute at begin node=\setlength{\baselineskip}{1.2em},draw,thick,rounded corners,inner sep=.3cm},scale=0.85]
	
	\node [force] (displacement) at (0,0) {displacement-time\\ graph};
	\node [force] (velocity) at (6.4,0) {velocity-time\\graph};
	\node [force] (acceleration) at (12.8,0) {acceleration-time\\graph};
	
	\draw[->] (displacement.north east) to [bend left] node[midway,above]{gradient of tangent} node[midway,below]{$v=\frac{\dd s}{\dd t}$} (velocity.north west);
	\draw[<-] (displacement.south east) to [bend right] node[midway,below]{area under curve} node[midway,above]{$\Delta s=\int v\dd t$} (velocity.south west);
	\draw[->] (velocity.north east) to [bend left] node[midway,above]{gradient of tangent} node[midway,below]{$a=\frac{\dd v}{\dd t}$} (acceleration.north west);
	\draw[<-] (velocity.south east) to [bend right] node[midway,below]{area under curve} node[midway,above]{$\Delta v=\int a\dd t$} (acceleration.south west);
	\end{tikzpicture}
%	\caption{relationships between motion graphs}
\end{figure*}

\example{Given the displacement-time graph as shown, check yourself that this $s$-$t$ graph leads to the velocity-time graph and the acceleration-time graph shown.}

\begin{figure*}[ht]
	\centering
	\begin{minipage}{0.3\textwidth}
		\centering
		\begin{tikzpicture}[xscale=0.6,yscale=0.85]
		\draw[<->] (0,2) node[left]{$s$} -- (0,0) -- (7,0) node[below]{$t$};
		\draw (0,0) -- (0,-2);
		\draw [thick,blue,domain=0:6.7,smooth] plot (\x, {1.6*sin(\x r)});
		\end{tikzpicture}
	\end{minipage}\hfil
	\begin{minipage}{0.3\textwidth}
		\centering
		\begin{tikzpicture}[xscale=0.6,yscale=0.85]
		\draw[<->] (0,2) node[left]{$v$} -- (0,0) -- (7,0) node[below]{$t$};
		\draw (0,0) -- (0,-2);
		\draw [thick,red,domain=0:6.7,smooth] plot (\x, {1.6*cos(\x r)});
		\end{tikzpicture}
	\end{minipage}\hfil
	\begin{minipage}{0.3\textwidth}
		\centering
		\begin{tikzpicture}[xscale=0.6,yscale=0.85]
		\draw[<->] (0,2) node[left]{$a$} -- (0,0) -- (7,0) node[below]{$t$};
		\draw (0,0) -- (0,-2);
		\draw [thick,red,domain=0:6.7,smooth] plot (\x, {-1.6*sin(\x r)});
		\end{tikzpicture}
	\end{minipage}
\end{figure*}

\example{Given the velocity-time graph as shown, check yourself that this $v$-$t$ graph leads to the displacement-time graph as shown.}\label{ex-vgraph}

\begin{figure*}[ht]
	\centering
	\begin{minipage}{0.45\textwidth}
		\centering
		\begin{tikzpicture}[xscale=0.8,yscale=1.1]
		\draw[<->] (0,3) node[left]{$v$} -- (0,0) -- (7,0) node[below]{$t$};
		\draw [thick,blue] (0,0) -- (2,2.5) -- (4,2.5) [out=-75,in=180] to (6,0);
		\draw[dashed] (2,0) -- (2,2.5)  (4,0) -- (4,2.5);
		\end{tikzpicture}
	\end{minipage}\hfil
	\begin{minipage}{0.45\textwidth}
		\centering
		\begin{tikzpicture}[xscale=0.8,yscale=1.1]
		\draw[<->] (0,3) node[left]{$s$} -- (0,0) -- (7,0) node[below]{$t$};
		\draw [thick,red,domain=0:2] plot (\x,0.2*\x*\x);
		\draw [thick,red] (2,0.8) -- (4,2.4);
		\draw [thick,red,domain=4:6] plot (\x,{2.4+0.8*(1-exp(-2.5*\x+2.5*4))/2.5});
		\draw[dashed] (2,0) -- (2,0.8)  (4,0) -- (4,2.4);
		\end{tikzpicture}
	\end{minipage}
\end{figure*}





\subsection{Linear motion with constant velocity}

let's look at the simplest kind of motion, an object moving at constant speed in a straight line: $\, v=\text{constant}$.

In this case the equations of motion are the straightforward ones met at GCSE: $$\, \tcbhighmath{a=0}$$ 
and $$ \tcbhighmath{s=vt}$$ \footnote[][-1cm]{It is implicitly assumed that the motion starts from the origin with respect to which displacement is defined. More generically, we should write: $s = s_0 + vt$, where $s_0$ is the initial displacement.}
Where S, represents displacement, V, the velocity and t, the time. You have met these a lot, we'll not say much more about them...


\begin{figure*}[htp]
	\centering
	\begin{minipage}{0.32\textwidth}
		\begin{tikzpicture}[scale=0.84]
		\draw [<->] (0,3.5) node[above] {displacement} -- (0,0) node[below]{0} -- (4,0) node[below] {time};
		\draw[very thick,blue] (0,0) -- (3.6,2.7);
		\end{tikzpicture}
	\end{minipage}\hfill
	\begin{minipage}{0.32\textwidth}
		\begin{tikzpicture}[scale=0.84]
		\draw [<->] (0,3.5) node[above] {velocity} -- (0,2.5) node[left]{$v$} -- (0,0) node[below]{0} -- (4,0) node[below] {time};
		\draw[very thick,blue] (0,2.5) -- (3.6,2.5);
		\draw[dashed] (3,2.5) -- (3,0) node[below]{$t$};
		\end{tikzpicture}
	\end{minipage}\hfill
	\begin{minipage}{0.32\textwidth}
		\begin{tikzpicture}[scale=0.84]
		\draw [<->] (0,3.5) node[above] {acceleration} -- (0,0) node[below]{0} -- (4,0) node[below] {time};
		\draw[very thick,blue] (0,0.01) -- (3.6,0.01);
		\end{tikzpicture}
	\end{minipage}
	\caption{Motion graphs for linear motion at constant velocity}
\end{figure*}



\subsection{Linear motion with constant acceleration}

More complex motion will involve a change in velocity and we limit ourselves at A-level to situations where this change in velocity is constant. ie we have linear motion with constant acceleration $a=\text{constant}$. When we do this we see that our motion graphs look as below:

\begin{figure*}[htp]
	\centering
	\begin{minipage}{0.32\textwidth}
		\begin{tikzpicture}[scale=0.84]
		\draw [<->] (0,3.5) node[above] {displacement} -- (0,0) node[below]{0} -- (4,0) node[below] {time};
		\draw [very thick,blue,domain=0:3.6,samples=12,smooth,variable=\x] plot (\x,{\x*\x/4});
		\end{tikzpicture}
	\end{minipage}\hfill
	\begin{minipage}{0.32\textwidth}
		\begin{tikzpicture}[scale=0.84]
		\draw [<->] (0,3.5) node[above] {velocity} -- (0,1) node[left]{$u$} -- (0,0) node[below]{0} -- (4,0) node[below]{time};
		\draw[very thick,blue] (0,1) -- (4,3);
		\draw[dashed] (0,2.5) node[left]{$v$}  -- (3,2.5);
		\draw[dashed] (3,2.5) -- (3,0) node[below]{$t$};
		\end{tikzpicture}
	\end{minipage}\hfill
	\begin{minipage}{0.32\textwidth}
		\begin{tikzpicture}[scale=0.84]
		\draw [<->] (0,3.5) node[above] {acceleration} -- (0,2.5) node[left]{$a$} -- (0,0) node[below]{0} -- (4,0) node[below] {time};
		\draw[very thick,blue] (0,2.5) -- (3.6,2.5);
		\end{tikzpicture}
	\end{minipage}
	\caption{motion graphs for linear motion at uniform acceleration}\label{fig:const_a_graph}
\end{figure*}

\subsection{Equations of motion}

In order to be able to solve problems relating to motion we need more powerful tools than we have met at GCSE. By combining the equation for acceleration with the definitive equation for velocity, we can come up with a set of motion equations commonly referred to as "the SUVAT equations". 

During a time interval $t$, suppose velocity changes from initial value $u$ to final value $v$. From the defining equation of acceleration $a=\frac{\Delta v}{\Delta t} = \frac{v-u}{t}$, we get\footnote{Proof with calculus: $\dd v = a \dd t \ra \Delta v = \int_{u}^{v} \dd v = \int_{0}^{t} a\dd t \ra v-u=at \ra v=u+at$}:

\begin{empheq}[box=\tcbhighmath]{equation}\label{eqn:const_a_eq1}{v=u+at}
\end{empheq}

This is the first equation in our series. 

To find total displacement travelled, we compute the area under the $v$-$t$ graph
\begin{empheq}[box=\tcbhighmath]{equation}\label{eqn:const_a_eq2}{s=\frac{1}{2}(u+v)t}
\end{empheq}
for which we can interpret $\bar{v}=\frac{1}{2}(u+v)$ as the average velocity during that time

plug \eqref{eqn:const_a_eq1} into \eqref{eqn:const_a_eq2}, we find an expression for the displacement travelled in terms of $u$ and $a$:\footnote{Proof with calculus: $\dd s = v \dd t \ra \Delta s = \int_0^s \dd s = \int_0^t v\dd t \ra s = \int_0^t (u+at)\dd t = \big(ut+\frac{1}{2}at^2\big)\bigg|_0^t \ra s = ut + \frac{1}{2}at^2$}
\footnote{Equation \eqref{eqn:const_a_eq3} assumes a zero initial displacement at $t=0$. If there is a non-zero initial displacement, one should write $s = s_0 + ut + \frac{1}{2}at^2$. Similar discussion applies to equation \eqref{eqn:const_a_eq2}.}

\begin{empheq}[box=\tcbhighmath]{equation}\label{eqn:const_a_eq3}{s = ut + \frac{1}{2}at^2}
\end{empheq}

This shows the displacement $s$ is a quadratic function in time $t$, consistent with the parabolic shape of the $s$-$t$ graph shown. We can also eliminate the time variable $t$ to derive one last equation from \eqref{eqn:const_a_eq1} we have $t=\frac{v-u}{a}$, substitute this into \eqref{eqn:const_a_eq2}, we find
\begin{equation}\label{eqn:const_a_eq4}
	s=\frac{1}{2}(u+v)\times\frac{v-u}{a} \RA \end{equation}
 \begin{empheq}[box=\tcbhighmath]{equation*}{2as = v^2 - u^2}
\end{empheq}

\example{A car starts from rest and accelerates uniformly at 5.0 m s$^{-2}$ for 6.0 s.\\ (a) How fast is the car travelling at $t=8.0$ s? \\(b) What is the distance travelled by the car in this time?}
\begin{soln}
\begin{equation*}
v=u+at \RA v = 0 + 5.0\times 6.0 = 30 \mps
\end{equation*}
\begin{equation*}
s=ut+\frac{1}{2}at^2 \RA s= 0+\frac{1}{2}\times5.0\times6.0^2 = 90 \text{ m} 
\end{equation*}
\end{soln}

\example{A car is travelling at 30\mps. A hazard appears in front of the car, and the driver takes immediate action to stop the car. When brakes are applied, deceleration of the car is 5.0\mpss. What is the braking distance?}
\begin{soln}
\begin{equation*}
2as=v^2-u^2 \RA s = \frac{v^2-u^2}{2a} = \frac{0^2 - 30^2}{2\times(-5.0)} = 90 \text{ m} 
\end{equation*}
\end{soln}

\example{At the instant the traffic light turns green, a motorcycle waiting at the stop line starts with a constant acceleration of $2.0 \mpss$. At the same instant, a truck at a constant speed of $16 \mps$ overtakes and passes the motorcycle. How far beyond the stop line will the motorcycle overtake the truck?}

\begin{soln} Lets suppose overtake occurs at time $t$ after motorcycle starts to accelerate, so the distance travelled by motorcycle: $$s_m = u_mt + \frac{1}{2}at^2  \RA s_m = 0 + \frac{1}{2}\times 2.0 \times t^2$$

In the same time the distance travelled by truck: $$s_t = v_t t \RA s_t = 16t$$

The overtake when: $$s_m = s_t \RA \frac{1}{2}\times 2.0 \times t^2 = 16 t \RA t=16 \text{ s}$$

Substitute $t$ into either $s_m$ or $s_t$, one finds distance travelled: $s = 256 \text{ m}$ 
\end{soln}




\subsection{Free fall}\label{ch_freefall}

A typical example of uniformly accelerated motion is the free fall\index{free fall} everything has the tendency to fall towards ground due to earth's gravity, in this section, we assume that effects of air resistance are negligible, acceleration of free fall is then constant $a=g$. The value of $g$ does not depend on mass of the falling object\footnote{The reason for this constant acceleration of free fall will be elaborated in $\S$\ref{ch_weight}.} Near surface of earth, the value of the acceleration due to gravity is: $g\approx9.81\mpss$, of course value of $g$ could be different on a different planet\footnote[][]{Note that Maths will make you use 9.8\mpss}

For a freely-falling object released from rest, its velocity increases with time as
\begin{equation*}
	v=u+at = 0 + gt \RA v = gt
\end{equation*}

and the distance it has fallen from the point of release is:
\begin{equation*}
s=ut+\frac{1}{2} at^2 = 0\cdot t + \frac{1}{2}gt^2 \RA s=\frac{1}{2}gt^2
\end{equation*}

\example{An object is released from rest from a height of $h=24$ m and falls freely under gravity. Air resistance is negligible.\\ (a) How long does it take to hit the ground? \\(b) What is its speed when hitting the ground?}
\begin{soln}
\begin{equation*}
h = \frac{1}{2}gt^2 \RA t = \sqrt{\frac{2h}{g}} = \sqrt{\frac{2\times24}{9.81}} \approx 2.21 \text{ s}
\end{equation*}
\begin{equation*}
v = gt = 9.81 \times 2.21 \RA v \approx 21.7 \mps
\end{equation*}

to find final velocity, one can also use $v^2 = u^2 + 2as$:
\begin{equation*}
v^2 = \cancelto{0}{u^2} + 2gh = 2\times9.81\times24 \RA v \approx 21.7 \mps 
\end{equation*}
\end{soln}

\example{A photograph is taken for a small particle falling from rest. The photograph is taken at 0.400 s after the object is released. Since the particle is still moving when the photograph is being taken, the image is blurred. It is found the blurred part corresponds to a length of 20.8 cm moved out by the particle, what is time of exposure for the photograph?}

\begin{soln} from $t=0$ to right before photo is taken: 
{	\centering
	
	$s_1 = \frac{1}{2}gt_1^2 = \frac{1}{2}\times9.81\times0.400^2 \approx 0.785 \text{ m}$
	
}

from $t=0$ to right after photo has been taken: 

{
	\centering
	
	$s_2 = s_1 + \Delta s = \frac{1}{2}gt^2_2 \RA t_2 = \sqrt{\frac{2(s_1+\Delta s)}{g}} = \sqrt{\frac{2\times(0.785+0.208)}{9.81}} \approx 0.450 \text{ s}$
	
}

time of exposure: $\Delta t = t_2 - t_1 = 0.450 - 0.400 \approx 0.050 \text{ s}$ 
\end{soln}


\subsection{Upward projection}

Like a freely-falling object, an object thrown upwards experiences the same constant downward acceleration $a=g\approx9.81\mpss$, as long as drag forces can be ignored. Remember that we pretend that we live in a "Perfect Physics World"\texttrademark for most of the A-level course. No air resistance, no friction, no inconvenient complexities. \begin{marginfigure}
        \centering
        \includegraphics[width=0.5\linewidth]{sphericalcow.png}
        \caption{Milk production at a dairy farm was at an all-time low, so the farmer decided to call on the local university for help. A team of the institutions' top scientists was assembled to work on the problem. The team of scientists visited the farm and took extensive data carefully examining every step of the production process. The leader of the team, a theoretical physicist, offered to write the report. A week later the physicist returned to the farm, saying to the farmer, "I have found a solution for your problem, but it only works for spherical cows in vacuum."}
        \label{fig:enter-label}
\end{marginfigure}

Note that initial velocity $u$ is upwards, but acceleration $a$ is downwards
so we will have different signs for $u$ and $a$ in the equations, conventionally, positive direction is taken as same direction as initial velocity. In our case, positive direction is upwards, the acceleration is then negative: $a=-g$ so the velocity-time relation and displacement-time relation are:
\begin{equation*}
v = u - gt \quad \quad s = ut - \frac{1}{2}gt^2
\end{equation*}

\begin{figure*}[ht]
	\centering
	\begin{minipage}{0.45\textwidth}
		\centering
		\begin{tikzpicture}[xscale=1.2]
		\draw[->] (0,-3) -- (0,3) node[left]{$v$};
		\draw[->] (0,0) -- (5,0) node[below]{$t$};
		\draw[thick,blue] (0,1.8) -- (5,-2.7);
		\draw (0.8,1.2) --++ (1,0.8) node[above,note]{object is moving\\ upwards ($v>0$)};
		\draw (2,0.05) --++ (1.2,0.4) node[above,note]{object at maximum\\height ($v=0$)};
		\draw (3.2,-1.2) --++ (-1,-0.6) node[below,note]{object is moving\\downwards ($v<0$)};
		\end{tikzpicture}
	\end{minipage}\hfil
	\centering
	\begin{minipage}{0.45\textwidth}
		\centering
		\begin{tikzpicture}[xscale=1.2]
		\draw[->] (0,-3) -- (0,3) node[left]{$s$};
		\draw[->] (0,0) -- (5,0) node[below]{$t$};
		\draw[thick,blue,domain=0:5,smooth] plot (\x,{-0.5*\x*(\x-4)});
		\draw (2.05,2.05) --++ (1,0.4) node[right,note]{maximum height\\ $h_\tmax$};
		\draw (3.3,1.1) --++ (-0.15,-0.4) node[left,note]{object above point\\of release ($h>0$)};
		\draw (4.4,-1) --++ (-0.8,-0.6) node[left,note]{object below point\\of release ($h<0$)};
		\end{tikzpicture}
	\end{minipage}

\caption{$v$-$t$ graph and $s$-$t$ graph for upward projectile motion}
\end{figure*}

The sign of $v$ now tells direction of motion, $v>0$ means object is moving upwards, $v<0$ means it has reversed direction and starts falling, in particular, an object attains greatest height when $v=0$.

The sign of $s$ gives whether object is at a higher or lower position with respect to point of release.

\titem $s>0$ means the object is above the position from which it is projected

\titem $s<0$ means it is now below the point of release

\example{A ball is projected vertically upwards at $12 \mps$. Air resistance is negligible. (a) Find the time taken for the ball to reach the highest position. (b) Find the greatest height.}
	
\begin{soln} maximum height is reached when $v=0$, so

	\centering
	
	$v = u-gt = 0 \RA t= \frac{u}{g} = \frac{12}{9.81} \approx 1.22 \text{ s}$
	
	$H_\tmax = ut-\frac{1}{2}gt^2 = 12\times1.22 - \frac{1}{2}\times9.81\times1.22^2 \approx 7.34 \text{ m}$
	


It is also possible to use $v^2 - u^2 = 2as$ to find $H_\tmax$:
\begin{equation*}
0^2 - u^2 = 2(-g)H_\tmax \RA H_\tmax = \frac{u^2}{2g} = \frac{12^2}{2\times9.81} \approx 7.34 \text{ m} 
\end{equation*}
\end{soln}

\example{A stone is thrown vertically upwards with an initial velocity of $14.0\mps$	from the edge of a cliff that is 35 m from the sea below.  \\(a) Find the speed at which it hits the sea. \\(b) Find the time taken for the stone to hit the sea.}
\footnote{Note that we have substituted  $a=-g$ since acceleration of free fall always points downwards, and $s=-35\text{ m}$ since sea is below point of release. Also final velocity when hitting water is downwards, which should take a negative sign, so we discarded the positive solution for $v$.}
\begin{soln}
 Take positive direction to point upwards, we use $v^2 - u^2 = 2as$ to find $v^2 = 14.0^2 + 2\times(-9.81)\times(-35) \approx 883 \text{ m}^2 \text{ s}^{-2} \RA v\approx -29.7\mps$.
 To find time, we can use $v = u-gt$, hence: $ t = \frac{v-u}{-g} = \frac{-29.7-14.0}{-9.81} \approx 4.46 \text{ s} $.

One can also attempt $s = ut - \frac{1}{2}gt^2$, this leads to the equation: $-35 = 14.0t - \frac{1}{2}\times9.81t^2$.

This quadratic equation in $t$ gives two roots: $t_1 \approx 4.46 \text{ s}$, and $t_2 \approx -1.60 \text{ s}$.

Note that a negative root should be discarded since it means stones hits the sea before it is thrown so time taken for stone to hit the sea is $t \approx 4.46\text{ s}$ \end{soln}

\subsection{Motion in two dimensions -- projectile motion}\label{ch:projectile}

For projectile motion, we assume no air resistance and no other forces
gravity causes a constant acceleration of free fall that acts vertically downwards. The curved path of a projectile is the combination of its \emph{horizontal} and \emph{vertical} motion.

\titem horizontally: no acceleration, so horizontal component of velocity $v_x = \text{constant}$

\titem vertically: constant acceleration, vertical component of velocity $v_y$ varies over time

As a consequence, a projectile would follows a \emph{parabolic} path as it travels\footnote{{You may be able to prove this statement in Question \ref{Q-projectile-eqn}.}{You may try to prove this statement with a bit algebra.}}

Let's consider a projectile launched at initial velocity $u$ at angle $\theta$ to the horizontal:

\begin{figure*}[htp]
	\centering
	\begin{tikzpicture}[scale=.9]
	\draw[->] (0,-1) -- (0,4.2) node[left]{$y$};
	\draw[->] (0,0) -- (10,0) node[below]{$x$};
	\draw [thick,blue,domain=0:5*sqrt(3),samples=16,smooth,variable=\x] plot (\x,{sqrt(3)*\x-\x*\x/5});
	\draw[->,purple,very thick] (0, 0) -- ++(1,0) node[below]{$u_x$};
	\draw[->,purple,very thick] (0,0) -- ++(0,1.732) node[left]{$u_y$};
	\draw[->,red,very thick] (0,0) -- ++(1,1.732) node[above]{$u$};
	\foreach \idx in {1.2}
	{
		\draw[->,purple,very thick] (1.732*\idx, 3*\idx-.6*\idx*\idx) -- ++(1,0) node[below]{$v_x$};
		\draw[->,purple,very thick] (1.732*\idx, 3*\idx-.6*\idx*\idx) -- ++(0,1.732-0.693*\idx) node[left]{$v_y$};
		\draw[->,red,very thick] (1.732*\idx, 3*\idx-.6*\idx*\idx) -- ++(1,1.732-0.693*\idx) node[above]{$v$};
	}
	\foreach \idx in {3.2,4.2}
	{
		\draw[->,purple,very thick] (1.732*\idx, 3*\idx-.6*\idx*\idx) -- ++(1,0) node[above]{$v_x$};
		\draw[->,purple,very thick] (1.732*\idx, 3*\idx-.6*\idx*\idx) -- ++(0,1.732-0.693*\idx) node[left]{$v_y$};
		\draw[->,red,very thick] (1.732*\idx, 3*\idx-.6*\idx*\idx) -- ++(1,1.732-0.693*\idx) node[right]{$v$};
	}
	\draw (0.5,0) arc(0:60:0.5);
	\node at (30:0.7) {$\theta$};
	\end{tikzpicture}
	\caption{components of the velocity of a projectile at different points along its path}
\end{figure*}

Horizontally, projectile maintains a constant velocity, so:
\begin{equation*}
v_x = u_x \qquad x=u_x t
\end{equation*}

Where $u_x = u\cos\theta$ is horizontal component of initial velocity vertically, if upward direction is taken to be positive, then acceleration $a=-g$, so
\begin{equation*}
v_y = u_y + at = u_y - gt  \qquad y = u_y t + \frac{1}{2}at^2 = u_y t - \frac{1}{2}gt^2
\end{equation*}

where $u_y = u\sin\theta$ is vertical component of initial velocity - components can be combined to give resultant velocity or resultant displacement:
\begin{equation*}
 v = \sqrt{v_x^2 + v_y^2} \qquad s = \sqrt{x^2 + y^2}
\end{equation*}


\subsection*{Maximum height reached by an projectile}

When a projectile reaches the highest position, its instantaneous vertical velocity becomes zero, we can then find the time it takes to attain this maximum height.
\begin{equation*}
	v_y = u\sin\theta - gt = 0 \RA t = \frac{u\sin\theta}{g}
\end{equation*}
	
To find $H_\tmax$, one can use either equation \eqref{eqn:const_a_eq2} or \eqref{eqn:const_a_eq3} 
	\begin{equation*}
	H_\tmax = \frac{1}{2}(u_y+v_y) t = \frac{1}{2}u_y t = \frac{1}{2} \times u \sin\theta \times \frac{u\sin\theta}{g} = \frac{u^2 \sin^2\theta}{2g}
	\end{equation*}
	\begin{equation*}
	H_\tmax = u_y t - \frac{1}{2}gt^2 = u \sin\theta \times \frac{u\sin\theta}{g} - \frac{g}{2} \times \left(\frac{u\sin\theta}{g}\right)^2 = \frac{u^2 \sin^2\theta}{2g}
	\end{equation*}

For any given initial speed $u$, the greater the angle of projection, the higher the object will reach. In the extremal case where $\theta = 90 ^\circ$, it simply becomes an upwardly projected motion.

\subsection*{Airborne time and horizontal range of an projectile}

A ball projected from the ground will first rise in height but it will eventually fall to the ground due to the gravitational pull after a period of time $T$. When it lands, its vertical displacement is zero, so:
	\begin{equation*}
	Y = u_y T - \frac{1}{2}gT^2 = u \sin \theta T - \frac{1}{2}gT^2 = 0 \RA T = \frac{2u\sin\theta}{g}
	\end{equation*}
	
The horizontal range is given by:
	\begin{equation*}
	X = u_x T = u\cos\theta \times \frac{2u\sin\theta}{g} = \frac{2u^2\sin\theta\cos\theta}{g} \RA X = \frac{u^2\sin2\theta}{g}
	\end{equation*}
	
Where in the last step the trigonometric identity $\sin2\alpha = 2\sin\alpha\cos\alpha$ has been used.

For a given initial speed $u$, a projectile launched at a greater angle will stay in the air for longer time, the greatest airborne time is obtained if object is projected straight up, i.e.,  $\theta =90^\circ$. For a given initial speed, horizontal range of projectile depends on angle $\theta$ of projection, so to obtain the greatest horizontal range, two things are required:

\titem sufficiently large horizontal velocity $v_x$

\titem sufficiently long time $T$ staying in the air

however, a larger $v_x$ requires a smaller $\theta$, which means a shorter airborne time $T$, therefore, there is a compromise between the two. The optimal angle should be neither be too large nor too small, and can be shown to be $45^\circ$ \footnote{You could prove this algebraicly if you were so inclined}

\begin{figure*}[ht]
	\centering
	\begin{tikzpicture}[scale=1]
	\draw[<->] (0,5.4) node[left]{$y$} -- (0,0) -- (11,0) node[below]{$x$};
	\draw [thick,red,domain=0:10,samples=16,smooth,variable=\x,graphlabel={[above]{$45^\circ$}}] plot (\x,{\x-0.1*\x*\x});
	\draw [thick,purple,domain=0:5*sqrt(3),samples=16,smooth,variable=\x,graphlabel={[above]{$60^\circ$}}] plot (\x,{sqrt(3)*\x-\x*\x/5});
	\draw [thick,purple,domain=0:5*sqrt(3),samples=16,smooth,variable=\x,graphlabel={[above]{$30^\circ$}}] plot (\x,{\x/sqrt(3)-\x*\x/15});
	\draw [thick,blue,domain=0:5,samples=16,smooth,variable=\x,graphlabel={[above]{$75^\circ$}}] plot (\x,{3.732*\x-\x*\x/1.34});
	\draw [thick,blue,domain=0:5,samples=16,smooth,variable=\x,graphlabel={[above]{$15^\circ$}}] plot (\x,{0.268*\x-\x*\x/18.66});
	\end{tikzpicture}
	\caption{trajectories of projectiles launched at the same speed but different angles}\label{fig:projectile_range}
\end{figure*}


\example{A ball is thrown from a point $O$ at $15 \mps$ at an angle of $40^\circ$ to the horizontal. The ball reaches its highest position at point $P$. Ignore the effects of air resistance. (a) How long does it take to reach $P$? (b) What is the magnitude of the displacement $OP$?}
		
\begin{soln} At the highest point: $v_y = u_y - gt = 0 \RA t = \frac{u\sin\theta}{g} = \frac{15\times\sin40^\circ}{9.81} \approx 0.983 \text{ s}$

Vertical displacement: $y = u_y t - \frac{1}{2}gt^2 = 15\sin40^\circ \times 0.983 - \frac{1}{2}\times 9.81 \times 0.983^2 \approx 4.74 \text{ m}$

Horizontal displacement: $x = u_x t = 15\cos40^\circ \times 0.983 \approx 11.3 \text{ m}$

Resultant displacement: $|OP| = \sqrt{x^2 +y^2} = \sqrt{11.3^2 + 4.74^2} \approx 12.2 \text{ m}$ 
\end{soln}
\newpage

\example{A small object is horizontally projected at 7.20 ms$^{-1}$ from a surface at a height of $h=1.2 \text{ m}$ above the ground. Assume there is no air resistance. (a) What is the time taken for the object to hit the ground? (b) What is the horizontal range? (c) Find the velocity at which the object hits the ground.}

\begin{soln}
    
Vertically, take downward as positive: $h = \cancelto{0}{u_y t} + \frac{1}{2}gt^2 \ra t = \sqrt{\frac{2h}{g}} = \sqrt{\frac{2\times1.2}{9.81}} \approx 0.495 \text{ s} $

Horizontal range: $x = u_x t = 7.20 \times 0.495 \approx 3.56 \text{ m}$

Final vertical velocity: $v_y = \cancelto{0}{u_y} + gt = 9.81 \times 0.495 \approx 4.85 \mps$

Magnitude of resultant velocity: $v = \sqrt{v_x^2 + v_y^2} = \sqrt{7.20^2 + 4.85^2} \approx 8.68 \mps$

Angle to which resultant velocity makes with horizontal: $\phi = \tan^{-1}\frac{v_y}{v_x} = \tan^{-1}\frac{4.85}{7.20} \approx 34^\circ$ 
\end{soln} 

\section{End-of-chapter questions}

\subsection*{Kinematic quantities}

\question{What is the distance covered for a car that travels half a lap along a circular path of radius of 200 m. What about the displacement?}

\question{A ball is released from a height of 2.0 m above the ground. It bounces vertically for quite a number of times before coming to rest. (a) State the change of displacement for the ball. (b) Explain how the distance travelled is different from the change in displacement.}


\question{For an athlete running around a track for many laps, suggest how his average velocity could be zero?}

\question{A car travels 2400 m east in 3.0 minutes, then takes a left turn, and then travels 700 m north in 1.5 minutes. What is the average speed and the average velocity for this journey?}

\question{Is it possible for an object moving at a steady speed to have acceleration?}


\subsection*{Motion graphs}

\question{For the $v$-$t$ graph given in Example \ref{ex-vgraph}, sketch the $a$-$t$ graph for this motion.}

\question{If the tangent of a displacement-time graph at one particular instant is sloping downwards, what does that imply about the velocity at that instant?}

\question{A vehicle initially travels at a steady speed of $15 \mps$. It accelerates uniformly for 10 s to reach a higher speed of $20 \mps$. It maintains at this speed for 20 s, and then decelerates uniformly to a stop in the last 10 s. (a) Sketch the velocity-time graph for this motion. (b) Sketch the acceleration-time graph. (c) Find the distance travelled during the 40 s.}


\subsection*{Linear motion with constant velocity}

\question{Sonar is a technique that uses sound waves to detect objects. It can be used to measure the depth of the seabed. Given that speed of sound in water is $1500 \mps$, and reflected waves sent from a submarine are detected 0.50s after they are transmitted. How deep is the water below the submarine?}

\question{Given that the speed of sound in air is $340 \mps$ and the speed of light in air is $3.0\times10^8 \mps$. If a person hears the sound of a thunder 5.0 seconds after seeing a lightning flash, how far away from this person is did the lightning strike?}



\subsection*{Linear motion with constant acceleration}


\question{A train initially travels at a speed of $40\mps$. It starts to decelerate at 0.50\mpss. (a) What is the distance travelled in 50 s? (b) When it comes to a stop, how far out has it travelled?}

\question{A vehicle moving at $14\mps$ accelerate uniformly to $26\mps$ in 6.0 s. (a) What is the average velocity during this time? (b) What is the acceleration during this time (c) What is the distance travelled by the vehicle? (d) The vehicle then braked with constant deceleration to stop in another 8.0 s. What is the distance travelled during the time when brakes are applied?}


\subsection*{Free fall}


\question{Two balls are dropped from rest from the same height. The second ball is released 0.80 s after the first one. What is their separation 1.5 s after the second ball is dropped?}

\question{A golf ball is dropped from the top of a tower of height 30 m. The ball falls from rest and air resistance is negligible. What time is taken for the ball to fall (a) the first 10 m from rest, (b) the last 10 m to the ground?}

\question{The acceleration of free fall on Pluto is about one-fifteenth of that on Earth. If it takes a time of $T$ for a rock to fall from rest a distance of $S$, what is the time taken, in terms of $T$, for a rock to fall from rest through the same distance $S$ on Pluto?}

\question{In an experiment is carried out to determine the acceleration of free fall $g$ using a falling body. The body is released from rest from a height of $h$, the time taken $t$ for it to hit the floor is measured. (a) Find the expression that can be used to calculate the value of $g$? (b) Suggest what could lead to an overestimation for the value of $g$.}


\subsection*{Upward projection}

\question{A ball is tossed upwards with a speed of $9.0 \mps$. (a) How long does it take to return to the same point if air resistance is negligible? (b) How does the return velocity compare with its initial velocity?}

\question{Someone wants to toss a ball onto a platform that is at a height of 20 m above him. What is the minimum initial velocity needed to launch the ball?}

\question{Someone standing at the top of a high building throws a ball straight up and another ball straight down with the same initial speed. Assume that air drag is negligible, which ball will have a greater speed when it hits the ground?}

\question{In basketball games, hang time refers to the length of time a player stays in the air after jumping from the floor. (a) It is reported that Michael Jordan is able to stay in the air for $T=1.0\text{ s}$ to do his slam-dunk tricks, estimate how high he can jump. (b) The acceleration of fee fall on the Moon is about one-sixth of that on Earth. What is the hang time of Michael Jordan if he takes off from the surface of the Moon?}


\begin{marginfigure}
	\centering
	\vspace*{-10pt}
	\begin{tikzpicture}[scale=1.08]
	\draw[thin, gray!50,step=0.2] (0,-4) grid (4,2);
	\draw[step=1] (0,-4) grid (4,2);
	\foreach \y in {10,5,...,-20} \node[left] at (0,\y/5){$\y$};
	\foreach \x in {0,2,...,8} 
	{
		\draw[white,fill] (\x/2-0.1,-0.4) rectangle (\x/2+0.1,-0.05);
		\node[below] at (\x/2,0) {$\x$};
	}
	\draw[thick,blue,domain=0:3.6] plot (\x, 1.48-1.48*\x); 
	\node at (-0.9,1.5) {$v$/m s$^{-1}$};
	\draw[white,fill] (3.2,-0.7) rectangle (3.8,-0.3);
	\node at (3.5,-0.5) {$t$/s};
	\end{tikzpicture}
	\vspace*{-15pt}
\end{marginfigure}

\question{Mark Watney\footnote{A fictional character in the science fiction movie \emph{The Martian} (2015) based on the novel of the same name written by \emph{Andy Weir}.} stands at the edge of a cliff on the Mars and throws a rock vertically upwards with a speed of $7.4 \mps$. The graph shows the variation with the time $t$ of the rock's velocity $v$. (a) What is the acceleration of free fall on the Mars? (b) When does the rock reach the maximum height? (c) What is the height above the base of the cliff the moment when the rock is thrown? (d) What is the maximum height above the base of the cliff to which the rock rises? (e) What is the total distance travelled by the rock before it strikes the ground?}

\subsection*{Projectile motion}

\question{A ball rolls off a table and lands at a position of a horizontal distance of 1.2 m from the table. The table is 0.95 m high. Find the speed at which the ball leaves the table.}

\question{A ball is kicked from the ground towards a vertical barrier. The barrier is at a horizontal distance of 18 m from the initial position of the ball. The ball strikes the barrier after 1.5 s at a height of 2.5 m above the ground. (a) Find the magnitude and the direction of the initial velocity. (b) Find the magnitude and the direction of the velocity at which the ball hits the barrier.}

\question{When a particle is launched from the origin at an angle $\theta$ with the horizontal at a speed of $u$, show that its trajectory is a parabola given by the equation: $y = x \tan\theta - \frac{gx^2}{2u^2\cos^2\theta}$.}\label{Q-projectile-eqn}

\question{Two golf players each hit a ball at the same speed. One at $30^\circ$ with the horizontal, the other at $60^\circ$. Which ball hits the ground first? Which ball goes farther?}

\question{(a) State the difference between the displacement of a projectile and the distance it travels. (b) Suggest in what situation a projectile's displacement could have the same magnitude as the distance.}

\question{An archer always aims slight higher than the distant target that she wants to hit. Why isn't the bow lined up such that it points exactly at the target?}


\chapter{Force \& Motion}

\section{Force \& motion: an introduction}

In physics, a force appears when two bodies interact with one another - you will encounter various types of forces in this course, some of which are:
\begin{compactitem}
	\item[--] \emph{weight}: the gravitational attraction acting on any object exerted by the earth
	
	\item[--] \emph{tension}: a force in a string, a rope, a chain, etc. when it is being pulled
	
	\item[--] \emph{normal contact}: when a body's surface is compressed, there reacts a normal contact force\footnote{Examples of normal contact force are support force that stops a desk from sinking into the ground, the impact on a football when you kick it, etc.}
	
	\item[--] \emph{friction}: a force that resists relative motion when two surfaces tend to slide over one another
	
	\item[--] \emph{resistance}: also called drag force, this is experienced when a body travels through a medium

	
	\item[--] \emph{upthrust}: an upward force acting on an object immersed in a fluid
	
	\item[--] \emph{electric force}: an attractive or repulsive interaction between electrically charged objects
	\item[--] \emph{magnetic force}: an attractive or repulsive interaction between magnets or electric currents
\end{compactitem}

Detailed features of these forces follow later in the notes.\\

A force can produce various effects to the object, the effect could be:

\begin{compactitem}
	\item[--] an increase/decrease in speed
	
	\item[--] a change in the direction of motion
	
	\item[--] causing the object to rotate
	
	\item[--] a change in shape of the object
\end{compactitem}

In the next few chapters, we will be looking into each of these aspects. When more than one force act on a body, it is useful to find their \emph{resultant}, or the \emph{net force}

\begin{ilight}
	\keypoint{resultant force}, or \keypoint{net force}, is a single force that has the same effect as all forces acting on an object combined. The emph{vector sum} of all of the individual forces gives the resultant force.
\end{ilight}

In this chapter, we will study the dynamics of \emph{point masses} - a \keypoint{point mass} is an idealization that the object has a mass but does not take up any space, furthermore an object treated as a point mass has a position specified with a geometric point in space. This is, of course, another simplification when size, shape, rotation, or structure of object are not important (Spherical cows..).\begin{marginfigure}
        \includegraphics[width=0.5\linewidth]{sphericalcow.png}
        \caption{Perfect physics world strikes again...}
\end{marginfigure}

\subsection{Newton's laws of motion}\index{Newton's laws of motion}

Newton's laws of motion\footnote{These three laws were first addressed by \emph{Isaac Newton} in his famous work \emph{Mathematical Principles of Natural Philosophy}, or simply the \emph{Principia}. The three-volume work was first published in 1687, and was soon recognised as one of the most important works in the history of science. Apart the from the three laws that laid the foundations for classical mechanics, the \emph{Principia} also stated \emph{the law of gravitation}, and accounted for planetary orbits and tides and other phenomena.} are three laws that form the basis of classical mechanics, they describe the relationships between motion of an object and forces acting on it.

\subsection{Newton's first law}

\begin{ilight}
	\keypoint{Newton's first law} states that an object continues in its state of rest or uniform motion at constant velocity if there is no resultant force acting
\end{ilight}


Any object `dislikes' any change to its state of motion, uniform motion tends to persist forever, this tendency to resist changes in motion is called the \keypoint{inertia}, in some textbooks, Newton's first law is also called \emph{the law of inertia}. If there is no change in state of motion, the object is said to be in \keypoint{equilibrium}\index{equilibrium}.

\subsection{Newton's Second law}

If resultant force is non-zero, velocity of the object will change, i.e., force produces acceleration.

\begin{ilight}
	\keypoint{Newton's second law} states that the acceleration is proportional to the resultant force and inversely proportional to the mass of the object. Symbolically, we write $a \propto \frac{\fnet}{m}$
\end{ilight}

With consistent units of measure, this proportionality can be written as an exact equation:
\begin{equation*}
a=\frac{\fnet}{m} \quad \text{or} \end{equation*}\begin{empheq}[box=\tcbhighmath]{equation} \quad {\fnet=ma}
\end{empheq}

\titem The SI unit of measurement for force $F$ is \keypoint{newton} (N)

A force of one newton acting a body of 1 kg produces an acceleration of 1 $\mpss$, acceleration produced is always in same direction of the net force. For the same force, an object of greater mass has a smaller acceleration, hence mass is a measure of the \emph{inertia} of this object in response to a net force\index{inertia}.

The way that the mass of an object is defined from the point of view of Newton's laws can be stated as\footnote{The concept of mass can be defined in many different ways. You might be familiar with the definition for mass as the amount of matter an object possesses. I personally think this definition is a bit vague and does not tell you anything new. Thinking of mass as a measure of inertia surely brings more insights. Mass also tells the strength at which an object interacts with other bodies through the gravitational attraction. As you will see later, from the view of Albert Einstein, it is also possible to think of mass as a form of energy.}:

\begin{ilight}
	\centering \keypoint{Mass} is an intrinsic property of a body to resist any change in its state of motion.
\end{ilight}

\example{A box of 6.0 kg is being pushed along a horizontal surface with a force of 30 N. The resistive force acting is 21 N. What is the acceleration of the box?}

\begin{soln}
\begin{equation*}
	\fnet = F - f = ma \RA a = \frac{F-f}{m} = \frac{30-21}{6.0} = 1.5 \mpss 
\end{equation*}
\end{soln}

\example{A car of mass 800 kg is travelling at a speed of $20 \mps$. The driver then operates the brake pedal so a braking force of 2000 N gradually brings the car to stop. (a) What is the deceleration for the car? (b) What is the stopping distance?}


\begin{soln} Using Newton's second law and noticing braking force acts opposite to direction of motion:

{
	\centering
	
	$\fnet = ma \RA -2000 = 800 \times a \RA a = -2.5 \mpss$
	
	$2as = v^2 - u^2 \RA s = \frac{v^2-u^2}{2a} = \frac{0^2 - 20^2}{2\times(-2.5)} = 80\text{ m}$
	
}
\end{soln}


\begin{marginfigure}
	\vspace{-18pt}
	\begin{center}
		\begin{tikzpicture}[scale=1]
		\draw[thick] (-1,6) -- (1,6);
		\draw[thick] (0,6) -- (0,0);
		\draw[thick, ->, blue] (.5,2) -- (.5,0.5) node[midway, right]{\textcolor{black}{pull}};
		\draw[fill=gray!40] (0,3) circle[radius=.8];
		\end{tikzpicture}
	\end{center}
	\vspace{-20pt}
\end{marginfigure}
	
\example{A massive ball is suspended on a string. A second string is attached to the bottom of the ball. If one pulls the bottom string with a gradually increasing force, does the top string or the bottom string break first? What if the bottom string is jerked, which string breaks?}	
	
\begin{soln}
When tension gradually increases, system is always in equilibrium, the tensions in strings must have $T_\text{top} = T_\text{bottom} + mg$. The top string suffers a greater force, so it breaks first, however, when bottom string is jerked, the ball tends to remain at rest due to its large mass, preventing sudden change to the tension in top string so in this case bottom string is more likely to snap.
\end{soln}

\example{Leviation of a slinky - When the top end of the slinky is dropped, the information of the tension change must propagate to the bottom end before both sides begin to fall; the top of an extended slinky will drop while the bottom initially remains in its original position, compressing the spring.}

\subsection{Newton's third law}

Every force is part of a pair of interactions between one body and another, when one body exerts a force on another, the second body also exerts a reaction on the first.

\begin{ilight}
	\keypoint{Newton's third law}, also called the \keypoint{action-reaction principle}, states that action and reaction are always equal in magnitude, opposite in directions and of the same type.
\end{ilight}

\example{Suggest the action and reaction force in the following cases: (a) A person stands on a bathroom scale. (b) A helicopter hovers in air. (c) The earth orbits around the sun.}

\begin{soln}(a) The person exerts downward force on scale, scale exerts an upward reaction on the person.

(b) The rotors of helicopter push air downwards, air exerts an upward force on helicopter.

(c) The Sun pulls the earth through gravitational attraction, the Earth also attracts the sun in return. \end{soln}


\subsection{Force analysis \& free-body diagrams}

When doing mechanics problems, it is necessary to find all forces applied upon an object, to visualise all these forces, it is helpful to draw a \keypoint{free-body diagram} (FBD)\index{free-body diagram}. A FBD shows a simplified version of the body with arrows indicating forces applied, it is recommended to follow the routine stated below when solving a mechanics problem:

\begin{compactitem}
	\item[(1)] draw a FBD for the object in the problem
	
	\item[(2)] resolve and find the resultant force with aid of the FBD
	
	\item[(3)] apply Newton's laws to write down the equation of motion for the object
	
	\item[(4)] solve the equation(s) to find acceleration
	
	\item[(5)] use kinematic relations to deduce information about motion of the object
\end{compactitem}


% % % % % % % % % % % % % % % % % % % %

\section{Types of forces}

\subsection{Weight}\label{ch_weight}

All objects exert attractive forces of gravity upon each other.\footnote{You will learn more about gravitational forces in Year 13.} The \keypoint{weight} of a body is due to the gravitational pull from our planet -- the Earth\index{weight}.\\
The weight $W$ of any object is proportional to its mass $m$: \begin{empheq}[box=\tcbhighmath]{equation*}{W=mg}\end{empheq}

Where $g$ is strength of the gravitational field, or the gravitational acceleration constant, The value of $g \approx 9.81 \text{ N kg}^{-1}$ but this value for $g$ does not hold in a satellite orbit, on Mars, near a black hole, etc.


The concept of weight is different from mass in many aspects:
\begin{compactitem}
	\item[--] weight is a force, so it is a vector (always acting downwards still makes a direction)
	
	mass is a scalar, it has magnitude only
	
	\item[--] weight is measured in newtons, mass is measured in kilograms
	
	\item[--] weight of object depends on its mass but also strength of gravitational field
	
	mass is an intrinsic property of object, so does not depend on force fields
	
	same object can have different weights on different planets, but its mass will be the same\footnote{Here we do not take into account the effects of \emph{relativity}. A clever student who has learned Einstein's theories might suggest the mass of the same object increases with its velocity.}
\end{compactitem}

\example{An astronaut finds that he weighs 300 N on the surface of Mars, where the gravitational field strength is known to be $3.7 \text{ N kg}^{-1}$. Find his mass and hence estimate his weight if he returns to his home on the Earth.}

\begin{soln} Mass of astronaut: $m = \frac{W_M}{g_M} = \frac{300}{3.7} \approx 81.1 \text{ kg}$

Weight on earth: $W_E = mg_E = 81.1 \times 9.81 \approx 795 \text{ N}$ \end{soln}

\subsection*{Free fall}

All things on the earth fall because of the force of gravity, if we ignore the restraints such as air resistance and upthrust force on a falling object, say the object is under the influence of gravity only, then the object is in a state called \keypoint{free fall}. Assuming the object is subject to gravity only, the resultant force is simply its weight, applying the Newton's second law, we have $\fnet = W \RA ma = mg$ so acceleration of the freely-falling object is\footnote[][-5cm]{In the derivation, the mass terms cancel out. Rigorously speaking, these are two different masses. One is the measure of inertia, and the other is a measure of gravitational force. It is experimentally found that the inertia mass and the gravitational mass are equal. The fact that the two masses are equal has profound reasons. We have shown here acceleration of free fall equals gravitational field strength, but Albert Einstein’s suggests that it is actually impossible to distinguish between a uniform acceleration and a uniform gravitational field. This idea lies at the heart of his \emph{general theory of relativity}. Those who are interested in this topic are recommended to start from here and do some online researches.}: \begin{empheq}[box=\tcbhighmath]{equation*}{a=g}\end{empheq}

The above shows us that \keypoint{acceleration due to free fall} is simply equal to field strength $g$, so any object, regardless of its mass, has same acceleration due to free fall.\footnote{In \S\ref{ch_freefall} and \S\ref{ch:projectile}, the statement that acceleration of free fall is constant in absence of air resistance was asserted without further explanation. Now you know why.}


\subsection{Drag}

When a body moves through air, water or any fluid, it experiences resistance called drag force\index{drag}, factors that determine the value of fluid drag include:

\titem relative speed of the object to the fluid ($v\up \ra f\up$)

\titem cross section of the object ($A\up \ra f\up$)

\titem shape of the object (streamlined shape has smaller drag)

\titem density of the fluid ($\rho\up \ra f\up$)

but what determines the drag force is a complicated issue\footnote[][-4cm]{There are a few empirical formula for drag force, each of which is accurate under certain conditions.
	
For an object moving through a fluid at low speeds (\emph{laminar flow}, no turbulence occurs), the resistance it experiences is proportional to its speed: $f=bv$, where $b$ is some constant which depends on fluid viscosity and the effective cross-sectional area of the object.
	
If objects are moving at relative high speeds through the fluid such that \emph{turbulence} is produced behind the object, drag force is proportional to the speed squared: $ f = \frac{1}{2} \rho C_D A v^2$, where $\rho$ is the fluid's density, $A$ is the cross-sectional area, $C_D$ is a dimensionless quantity called the drag coefficient.}
in short, the drag force always acts in a direction to oppose relative motion of object through fluid.

\subsection*{Free fall through air}

\begin{marginfigure}
	\vspace*{-16pt}
	\centering
	\begin{tikzpicture}
	\draw[thick,blue,->] (0,0) -- (0,-2.5) node[right,black]{$mg$};
	\draw[thick,blue,->] (0,0) -- (0,1.5) node[right,black]{$f$};
	\draw[fill=gray!50] (0,0) circle (0.3);
	\draw[thick,purple,->] (1,0.5) -- (1,-0.5) node[right]{$v$};
	\end{tikzpicture}
	\vspace*{-16pt}
\end{marginfigure}

Let's consider an object falling through air from a very high tower, the forces acting are weight and air resistance (shown in the free-body diagram). The equation of motion for this falling object is:

{
	\centering
	
	$ \fnet = mg - f = ma $
	
}

As $v$ increases, air resistance $f$ increases, so net force $\fnet$ decreases, this means acceleration $a$ would decrease as object falls - i.e., speed will increase at a decreasing rate during the fall
\footnote[][7cm]{The velocity-time relation can be obtained for some simple models. Suppose air resistance is proportional to speed of the falling body, i.e., $f=bv$, then the equation of motion reads: $\fnet = m \frac{\dd v}{\dd t} = mg - bv$, where acceleration is written explicitly as the rate of change in velocity. With the initial conditions $v=0$ at $t=0$, we can solve this differential equation to obtain the speed of this falling object at any given time $t$:
\begin{equation*}
	\dd t = \frac{\dd v}{g - \frac{b}{m}v} \RA \int_0^t \dd t = \int_0^v \frac{\dd v}{g - \frac{b}{m}v} \end{equation*}
 \begin{equation*}\RA t = -\frac{m}{b} \ln \left(g - \frac{b}{m}v \right) \Bigg|_0^v = -\frac{m}{b} \ln \left( 1 - \frac{bv}{mg}\right)
\end{equation*}
	
Rearrange the terms, we find: 
\begin{equation*}v(t) = \frac{mg}{b} \left(1 - \ee^{-\frac{bt}{m}} \right)\end{equation*}}

%$v$-$t$ graph for this process is shown on the next page


\begin{figure*}[ht]
	\centering
	\begin{tikzpicture}[xscale = 1.6,yscale=1.65]
	\draw [thick, <->] (7.5,0) node[below]{$t$} -- (0,0) --(0,3.6) node[left]{$v$};
	\draw [thick,blue,domain=0:7,samples=20,smooth,variable=\x] plot (\x,{2.5*(1-exp(-0.8*\x))});
	\draw [thick,orange] (0,0) -- (1.8,3.6);
	\draw [thick,dashed] (0,2.5) node[left]{$v_t$} -- (7,2.5);
	\draw [->] (3,0.6) node[note]{when air drag is small at low speed,\\ acceleration $a$ is slightly less than $g$} -- (0.45,0.6);
	\draw [->] (5,1.6) node[note]{gravity is offset by the increasing air drag, \\ the velocity tends to a terminal value $v_t$} -- (5,2.4);
	\draw [->] (4.2,3.1) node[note]{in absence of air, an object falls freely \\ with constant acceleration $g\approx9.81\mpss$} -- (1.6,3.1);
	\end{tikzpicture}
	\caption{variation of velocity for a falling object through air}
\end{figure*}
at low speeds, air resistance is negligible, so $\fnet =ma \approx mg$,  acceleration of object at start of the fall is similar to $g$ but as $v$ increases, acceleration decreases so $a$ becomes less than $g$. After a sufficiently long time, acceleration gradually decreases to zero, velocity gradually increases and tends to a maximum value. At this stage, equilibrium is restored: $f=mg$, object no longer accelerates, this constant final velocity is known as the \keypoint{terminal Velocity}\index{terminal velocity}


\example{An object of 5.0 kg falls through the atmosphere from a very high altitude. After some time, it falls at a constant speed of $70 \mps$. Assume there is no significant change in gravitational field during the fall and the air resistance is proportional to speed: $f=bv$. (a) Find the value of the coefficient $k$. (b) Find the acceleration of the object when it is falling at $30\mps$.}

\begin{soln}equilibrium between weight and air drag when falling at terminal speed, so
\begin{equation*}
	mg = bv_t \RA b = \frac{mg}{v_t} = \frac{5.0\times9.81}{70} \approx 0.70 \text{ kg s}^{-1}
\end{equation*}

at any instant, equation of motion is: $\fnet = ma = mg - bv$

at $30\mps$, acceleration is: $a = \frac{mg-bv}{m} = \frac{5.0\times9.81-0.70\times30}{5.0} \approx 5.6 \mpss$ \end{soln}

\subsection*{Bubble rising in a liquid}

\begin{marginfigure}
	\vspace*{-16pt}
	\centering
	\begin{tikzpicture}
	\draw[thick,blue,->] (0.1,0) -- (0.1,-0.8) node[right,black]{$mg$};
	\draw[thick,blue,->] (-0.1,0) -- (-0.1,-1.6) node[left,black]{$f$};
	\draw[thick,blue,->] (0,0) -- (0,2.5) node[right,black]{$U$};
	\draw[fill=gray!50] (0,0) circle (0.3);
	\draw[thick,purple,->] (1,-0.5) -- (1,0.5) node[right]{$v$};
	\end{tikzpicture}
	\vspace*{-16pt}
\end{marginfigure}

let's now consider bubbles formed at the bottom of a fizzy drink. The forces acting on bubble are weight, water resistance and upthrust. The equation of motion for the rising bubble is:

{
	\centering
	
	$ \fnet = U - mg - f = ma $
	
}

As bubble moves faster, $f$ increases, then $\fnet$ decreases, so acceleration $a$ would gradually decrease to zero as bubble rises. The speed of the bubble increases and reaches a maximum value at terminal speed, $a\to0$, therefore one has: $$U=f+mg$$




\subsection{Normal contact}

When two objects are in contact, the interaction between them is called the \emph{contact force}. The 
\keypoint{normal contact force} is the component of contact force perpendicular to the contacting surface.\sidenote[][]{By definition, "normal" is always at right angles to the surface in question. Cf. refraction and reflection.}

Of course, we know that there are actually only 4 forces in the universe so the origin of normal contact is the \emph{electrostatic interaction} between atoms. It's convenient shorthand however to be able to talk about contact and non-contact forces and it's \textit{close enough.}\footnote{Pun intended.} When two objects are pressed against each other, surface atoms get close enough that the strong electrostatic repulsion between electron clouds of the atoms prevent them from penetrating through one another.

\newpage

\begin{marginfigure}
	\vspace*{-10pt}
	\centering
	\begin{tikzpicture}[scale=0.9]
	\draw[thick] (-2,0) -- (2,0);
	\draw[thick] (-1,0) rectangle (1,1.5);
	\draw[thick,blue,->] (0.04,0.75) --++ (0,-2.1) node[right,black]{$mg$};
	\draw[thick,blue,->] (-0.04,0) --++ (0,2.1) node[left,black]{$R$};
	\end{tikzpicture}
	\vspace*{-5pt}
\end{marginfigure}


\example{A box of mass $m=4.0 \text{ kg}$ is resting on a horizontal ground. What is the normal contact force acting?}

\begin{soln} Equilibrium between weight and normal contact, so
\begin{equation*}
	R-W=0 \RA R=mg = 4.0 \times9.81 \approx 39.2 \text{ N} 
\end{equation*} \end{soln}

\begin{marginfigure}
	\vspace*{-10pt}
	\centering
	\begin{tikzpicture}[scale=0.9]
	\draw[thick] (-1,0) rectangle (1,3);
	% \draw[fill, gray] (0,1.5) circle(0.25);
	% \draw[ultra thick, gray] (0,1.5) -- (0,0.8) -- (-0.2,0) (0,0.8) -- (0.2,0);
	% \draw[ultra thick, gray] (0,1.25) -- (-0.3,0.6) (0,1.25) -- (0.3,0.6);
\def\H{2.0}     % human height
  \draw[thick,gray,fill=gray] (0,\H+0.03) circle (0.15*\H) coordinate (H);
 \draw[thick,gray] (H)++(-90:0.15*\H) coordinate (N) to[out=-85,in=85]++ (0,-0.40*\H) coordinate (P);
  \draw[thick, gray,line cap=round] (N)++(-85:0.03) to[out=-110,in=90]++ (-0.08*\H,-0.4*\H);
  \draw[thick, gray, line cap=round] (N)++(-85:0.03) to[out=-60,in=-110]++ (0.15*\H,-0.12*\H)  to[out=70,in=-70]++ (-0.01*\H,0.24*\H);
  \draw[thick, gray] (P) to[out=-110,in=85] ($(H)+(-0.08*\H,-\H)$);
  \draw[thick, gray] (P) to[out=-80,in=108] ($(H)+(0.06*\H,-\H)$);

  
	\draw[blue,fill] (0,.9) circle(0.05);
	\draw[thick,blue,->] (0,.9) --++ (0,-1.5) node[right,black]{$mg$};
	\draw[thick,blue,->] (0,.9) --++ (0,1.5) node[left,black]{$R$};
	\draw[ultra thick] (0,3) --++ (0,0.8);
	\draw[purple,thick,->] (1.5,0.8) --++ (0,1.4) node[right]{$v$};
	\end{tikzpicture}
	\vspace*{-5pt}
\end{marginfigure}

\example{A man of $80$ kg stands in a lift. Find his apparent weight, i.e., the contact force, when the lift is (a) moving upwards at steady speed of $2.0 \mps$, (b) accelerating upwards at $2.0 \mpss$, (c) moving upwards but slowing down at a deceleration of $1.5 \mpss$.}

\begin{soln} The forces acting on man are weight and normal contact, for either case, equation of motion for the man reads:

{

\centering

$\fnet = ma = R - mg$

}

so normal contact force: $R = mg + ma$

when rising at steady speed, man is in equilibrium ($a=0$), so: $R=mg = 80\times9.81 \approx 785 \text{ N}$

when accelerating upwards ($a = +2.0 \mpss$): $R = 80 \times 9.81 + 80 \times 2.0 \approx 945 \text{ N}$

when decelerating upwards ($a = -1.5 \mpss$): $R = 80 \times 9.81 + 80 \times (-1.5) \approx 665 \text{ N}$ \end{soln}


\example{A sleigh of mass 15 kg lies at rest on an icy ground. The surface is frictionless. A force $P$ of 75 N is applied to the sleigh. Find the normal contact force and the acceleration of the sleigh if $P$ is acting (a) horizontally, (b) at an angle $\alpha$ to the horizontal where $\tan\alpha=\frac{3}{4}$.}

\begin{marginfigure}
	\vspace*{-20pt}
	\centering
	\begin{minipage}[b]{0.2\textwidth}
		\centering
		\begin{tikzpicture}[scale=0.84]
		\draw[thick,fill] (0,0) circle(0.05);
		\draw[blue,thick,->] (0,0) -- (0,-3) node[right]{$mg$};
		\draw[blue,thick,->] (0,0) -- (0,3) node[right]{$R$};
		\draw[blue,thick,->] (0,0) -- (1.5,0) node[right]{$P$};
		\node at (0,-4) {(a)};
		\end{tikzpicture}
	\end{minipage}\hfil
	\begin{minipage}[b]{0.2\textwidth}
		\centering
		\begin{tikzpicture}[scale=0.84]
		\draw[thick,fill] (0,0) circle(0.05);
		\draw[blue,thick,->] (0,0) -- (0,-3) node[right]{$mg$};
		\draw[blue,thick,->] (0,0) -- (0,2.1) node[right]{$R$};
		\draw[blue,thick,->] (0,0) -- (1.2,0.9) node[right]{$P$};
		\draw[dashed] (0,0) -- (1,0);
		\draw(0.5,0) arc(0:36.89:0.5);
		\node at (19:0.8) {$\alpha$};
		\node at (0,-4) {(b)};
		\end{tikzpicture}
	\end{minipage}
\end{marginfigure}

\begin{soln} 
Free-body diagrams for both cases are shown for both cases, no net force acts in vertical direction. The net force in horizontal direction provides acceleration

when $P$ acts horizontally:

	
	$ R = mg = 15\times 9.81 \approx 147 \text{ N}$
	
	$ P = ma \RA a = \frac{P}{m} = \frac{75}{15} = 5.0 \mpss $
	
when $P$ acts at angle $\alpha$:
	$R + P\sin\alpha = mg \RA R = 15\times9.81-75\times\frac{4}{5} \approx 87 \text{ N}$
	
	\eqyskip $ P\cos\alpha = ma \RA a = \frac{P\cos\alpha}{m} = \frac{75\times\frac{3}{5}}{15} = 3.0 \mpss $
\end{soln}


\subsection{Friction}

Friction is the component of contact force that is parallel to contact surfaces, when there is potential or actual sliding between surfaces, frictional force come into action\index{friction}

\titem for surfaces that \emph{tend} to move relative to each other, \keypoint{static friction} acts to oppose this tendency.

\titem if surfaces are already sliding over one another, then \keypoint{kinetic, or dynamic friction} opposes this motion.\footnote{\piste We dont need to know these in these terms - we often lump dynamic friction together with any other frictive forces, like air-resistance, and just call th resultant "Drag". Remember, a lot of Physics is knowing when to simplify.}

You could say that static friction $f_S$ is self-adjusting, in that it arises as an $N3^{rd}$ force and will always be equal to the applied force - up until the moment the capacity for static friction is exceeded\footnote{You'll learn about the coefficient of friction $\mu$ in maths}.
In other words, an object placed on a rough surface can stay at rest when acted by a small external force $F$ it can do so because $f_S$ equalises external force to maintain equilibrium, if no external force acts, then $f_S=0$

Friction, on microscopic level, is an \emph{electromagnetic force} in nature when two surfaces are in contact, irregularities on the surface touch each other surface atoms come very close and bonds are formed through electrostatic force in some sense, surface atoms get \emph{cold welded} to each other when surfaces try to move relative to each other, this electrostatic weld is origin of friction.

\subsection{Inclined slopes}\index{inclined slopes}

The inclined slope is probably the entry ticket into the business end of mechanics problems. There are \emph{hundreds} of examples of this notorious problem to be found in any physics textbook and any exam paper on mechanics.

\begin{marginfigure}
%	\vspace*{-12pt}
	\centering
 \begin{tikzpicture}[scale=1]
	\draw[thick] (0,0) -- (30:5) -- ++(0,-2.5) -- cycle;
	\draw[thick] (0.8,0) arc(0:30:0.8);
	\node at (15:1) {$\theta$};
	\draw[thick] (30:2.5) -- ++(120:1) -- ++(30:1) -- ++(-60:1);
	\draw[thick] (-1,0) -- (5,0);
	\draw[thick,blue,->] (2.348,1.933) --++ (0,-1.6) node[right,black]{$W$};
	\draw[thick,blue,->] (30:3) ++ (120:0.5) --++ (120:1.4) node[left]{$R$};
    \draw[thick,dashed,red,->] (30:3) ++ (120:0.5) --++ (120:-1.4) node[left]{$W_{\perp}$};
	\draw[thick,blue,->] (30:3) ++ (120:0.5) --++ (30:1.8) node[above]{$F$};
	\draw[thick,dashed,red,->] (30:3) ++ (120:0.5) --++ (210:0.8) node[left]{$W_{\parallel}$};
	\draw[purple,thick,->] (2.7,3) --++ (30:0.8) node[midway,above]{$v$};
	%	\draw[thick,blue,->] (2.348,1.933) --++ (210:0.8) node[left]{$W_\parallel$};
	%	\draw[thick,blue,->] (2.348,1.933) --++ (-60:1.38564) node[right]{$W_\perp$};;
	%	\draw[dashed] (2.348,1.933) ++ (210:0.8) --++ (-60:1.38564) --++ (30:0.8) ;
	\end{tikzpicture}
	\vspace*{-16pt}
\end{marginfigure}

The problem is about a mass $m$ placed on a plane inclined at angle $\theta$ to the horizontal, the mass could sit at rest on, slide down, or get pulled/pushed up the plane. The motion of the mass could be affected by weight, friction, normal contact, or other forces.\\
They almost all have the same trick - resolve the forces in directions parallel and perpendicular to the slope because this allows you to use the equations of motion to work out motion up or down the slope.
It is almost inevitable that you'll need to break weight into two components\footnote{We have already done that in Example \ref{ex:comp-of-W}.}

\titem component of weight parallel down the slope is: $W_\parallelslant = mg\sin\theta$

\titem component of weight perpendicular to slope is: $W_\perp = mg\cos\theta$

\begin{marginfigure}
	\centering
	\begin{tikzpicture}[scale=1.35]
	\draw[thick] (0,0) -- (20:3.5) -- ++(0,-1.1971) -- cycle;
	\draw[thick] (1,0) arc(0:20:1);
	\node at (10:1.2) {$\theta$};
	\draw[thick] (20:1.7) -- ++(110:0.6) -- ++(20:0.6) -- ++(-70:0.6);
	\draw[thick,blue,->] (20:2) ++ (110:0.3) --++ (0,-1.5) node[right,black]{$W$};
	\draw[thick,blue,->] (20:2) ++ (110:0.3) --++ (110:1.4) node[left]{$R$};
	\draw[thick,blue,->] (20:2) ++ (110:0.3) --++ (20:0.6) node[above]{$f$};
	\end{tikzpicture}
	\vspace*{-16pt}
\end{marginfigure}
\newpage

\example{A block of mass $m$ stays at rest on an inclined plane. The plane makes an angle $\theta$ with the horizontal. Find the normal contact force $R$ and the frictional force $f$ acting on the block.}

\begin{soln} The block is in equilibrium, so $\fnet=0$ in any direction.

Parallel to slope: $f = W_\parallelslant \RA f=mg\sin\theta$

Normal to slope: $R = W_\perp \RA R =mg\cos\theta$ \end{soln}


\begin{marginfigure}
	\centering
	\vspace*{-16pt}
	\begin{tikzpicture}[scale=1.1]
	\draw[thick] (0,0) -- (20:5) -- ++(0,-1.71) -- cycle;
	\draw[thick] (1,0) arc(0:20:1);
	\node at (10:1.2) {$\theta$};
	\draw[thick] (20:3) -- ++(110:1) -- ++(20:1) -- ++(-70:1);
	\draw[thick,blue,->] (20:3.5) ++ (110:0.5) --++ (0,-1.5) node[right,black]{$W$};
	\draw[thick,blue,->] (20:3.5) ++ (110:0.5) --++ (110:1.4) node[left]{$R$};
	\draw[purple,thick,->] (2,2.4) --++ (200:0.8) node[midway,above]{$v$};
	\end{tikzpicture}
	\vspace*{-16pt}
\end{marginfigure}

\example{A block of mass $m$ slides down a \emph{smooth} slope. The angle of the slope to the horizontal is $\theta$. Find the acceleration of the block.}

\begin{soln} only force acting along the slope is component of weight down the slope, so:

{
	\centering
	
	$ \fnet = ma = mg\sin\theta \RA a = g\sin\theta $
	
}

as $\theta\to0$, $a\to0$, this shows if plane becomes horizontal, the block simply stays put

as $\theta\to90^\circ$, slope becomes vertical, block would undergo free fall, so naturally $a\to g$ \end{soln}


\begin{marginfigure}
	\vspace*{25pt}
	\centering
	\begin{tikzpicture}[scale=1.1]
	\draw[thick] (0,0) -- (20:5) -- ++(0,-1.71) -- cycle;
	\draw[thick] (1,0) arc(0:20:1);
	\node at (10:1.2) {$\theta$};
	\draw[thick] (20:3) -- ++(110:1) -- ++(20:1) -- ++(-70:1);
	\draw[thick,blue,->] (20:3.5) ++ (110:0.5) --++ (0,-1.5) node[right,black]{$W$};
	\draw[thick,blue,->] (20:3.5) ++ (110:0.5) --++ (110:1.4) node[left]{$R$};
%	\draw[thick,blue,->] (20:3.5) ++ (110:0.5) --++ (20:1.8) node[above]{$F$};
	\draw[thick,blue,->] (20:3.5) ++ (110:0.5) --++ (20:0.375) node[below]{$f$};
	\draw[purple,thick,->] (2,2.4) --++ (200:0.8) node[midway,above]{$v$};
	%	\draw[thick,blue,->] (20:3.5) ++ (110:0.5) --++ (200:0.8) node[left]{$W_\parallelslant$};
	%	\draw[thick,blue,->] (20:3.5) ++ (110:0.5) --++ (-70:1.38564) node[right]{$W_\perp$};;
	%	\draw[dashed] (2.348,1.933) ++ (210:0.8) --++ (-60:1.38564) --++ (30:0.8) ;
	\end{tikzpicture}
%	\vspace*{-16pt}
\end{marginfigure}

\example{A block of mass 2.0 kg slides down a rough slope from rest. The slope is inclined at angle $\theta=20^\circ$ to the horizontal, and the block experiences a constant friction of 5.0 N. (a) What is the block's acceleration? (b) What is the distance travelled in 2.5 seconds?}

\begin{soln} Resolving along slope:

{
	\centering
	
	$\fnet = mg\sin\theta - f = ma$
	
$a = \frac{2.0 \times 9.81\times\sin20^\circ - 5.0}{2.0} \approx 0.855 \mpss $\\
	
}

 Distance travelled: $s = \cancelto{0}{ut} + \frac{1}{2}at^2 = \frac{1}{2}\times0.855\times2.5^2 \approx 2.67 \text{ m}$ \end{soln}


\begin{marginfigure}
%	\vspace*{16pt}
	\centering
	\begin{tikzpicture}[scale=1.1]
	\draw[thick] (0,0) -- (18:5) -- ++(0,-1.545) -- cycle;
	\draw[thick] (1,0) arc(0:20:1);
	\node at (10:1.2) {$\theta$};
	\draw[thick] (18:1.5) -- ++(108:1) -- ++(18:1) -- ++(-72:1);
	\draw[thick,blue,->] (18:2) ++ (108:0.5) --++ (0,-1.4) node[right,black]{$W$};
	\draw[thick,blue,->] (18:2) ++ (108:0.5) --++ (108:1.3) node[left]{$R$};
	%	\draw[thick,blue,->] (20:3.5) ++ (110:0.5) --++ (20:1.8) node[above]{$F$};
	\draw[thick,blue,->] (18:2) ++ (108:0.5) --++ (198:0.4) node[above]{$f$};
	\draw[purple,thick,<-] (2.8,2.4) --++ (198:0.8) node[midway,above]{$v$};
	%	\draw[thick,blue,->] (20:3.5) ++ (110:0.5) --++ (200:0.8) node[left]{$W_\parallelslant$};
	%	\draw[thick,blue,->] (20:3.5) ++ (110:0.5) --++ (-70:1.38564) node[right]{$W_\perp$};;
	%	\draw[dashed] (2.348,1.933) ++ (210:0.8) --++ (-60:1.38564) --++ (30:0.8) ;
	\end{tikzpicture}
	%	\vspace*{-16pt}
\end{marginfigure}


\example{A block of mass 3.0 kg is travelling up an inclined slope at an initial speed of $2.8 \mps$. The slope makes an angle of $18^\circ$ with the horizontal. A constant friction of 7.5 N acts on the block.\\ (a) What is the block's deceleration?\\ (b) How far does the block travel along the slope before its speed decreases to zero? \\(c) Suggest whether the block could stay on the slope.\\}

\begin{soln} Resolving along slope (take direction of initial velocity as positive direction): 

{
	\centering
	
	$\fnet = -mg\sin\theta - f = ma \RA a = \frac{-mg\sin\theta-f}{m} = \frac{-3.0\times9.81\times\sin18^\circ-7.5}{3.0} \approx - 5.53 \mpss$
	
	\eqyskip $v^2 - u^2 = 2as \RA s = \frac{v^2 - u^2}{2a} = \frac{0^2 - 2.8^2}{2\times(-5.5.3)} \approx 0.709 \text{ m}$ 
	
}

Note that component of weight down the slope is: $W_\parallelslant = mg\sin\theta = 3.0\times9.81\times18^\circ \approx 9.1 \text{ N}$

$W_\parallelslant > f$, so friction is not enough to prevent block from sliding back down the slope \end{soln}






\subsection{Many-body problems}

The problems we have been dealing with so far only involve one body, however a lot of the "resolving vectors" type problems involve several objects that mutually interact. We can often exploit the fact that if you take any individual part of a system and look into the \textit{internal} forces between the objects of interest you'll find that any force acting between objects \emph{within} system, has an equal but opposite reaction force.\\
The system can also be treated as a whole and we can analyse the \textit{net external force} acting on entire system and work out any resultant or combined force/acceleration/etc.


\example{Two boxes $A$ and $B$ are placed on a smooth surface. They are accelerated together by a horizontal force $F$ as shown. Find the acceleration and the contact force between them.}

\begin{figure}[ht]
	\centering

% HORIZONTAL ground - lift
\begin{tikzpicture}[scale=1.5]
  \def\W{3.6}  % ground width
  \def\D{0.2}  % ground depth
  \def\H{1.9}  % human height
  \def\h{0.8}  % mass height
  \def\w{1.0}  % mass width
  \def\mx{0.10*\W} % mass x coordinate
  
  % SETUP
  \draw[ground] (-0.5*\W,0) rectangle++ (\W+1,-\D);
  
  \draw (-0.5*\W,0) --++ (\W+1,0);
  \draw[mass] (\mx,0) rectangle++ (\w,\h) node[midway] {$A$};
  \draw[mass] (\mx+\w,0) rectangle++ (\w,\h) node[midway] {$B$};
  
  % PERSON
   \coordinate (H) at (0,0.75*\H);
  \draw[thick,line cap=round]
    (H)++(-165:0.3) to[out=-140,in=60]++ (-130:0.3)
    to[out=65,in=-90,looseness=1.0]++ (80:0.45) to[out=90,in=120,looseness=1.4]++ (20:0.2); % pony tail
  \draw[thick,fill=white] (H) circle (0.3);
  \draw[thick,line cap=round] (H)++(-140:0.3) to[out=80,in=-120,looseness=1.8]++ (40:0.6); % hair
  \draw[thick] (H)++(-120:0.3) coordinate (N) to[out=-115,in=40]++ (-130:0.36*\H) coordinate (P);
  \draw[thick,line cap=round] (N)++(-115:0.03) to[out=-60,in=170] (1.04*\mx,\h);
  \draw[thick,line cap=round] (N)++(-115:0.03) to[out=-70,in=170] (\mx,0.92*\h);
  \draw[thick,line cap=round] (P) to[out=-120,in=30,looseness=1.3] (-0.35*\W,0);
  \draw[thick,line cap=round] (P) to[out=-50,in=30,looseness=1.8] (-0.22*\W,0);
  
  % FORCES
 
    % \draw[<->] (\mx+1.4*\w,2.1*\h) node[below=2,left=0] {$y$}
    %   |-++ (0.7*\h,-0.7*\h) node[right=0] {$x$};
    % \draw[force] (\mx+0.3*\w,0.05*\h) --++ (0, 1.3*\h) node[above] {$\vbF_\mathrm{N}$};
    % \draw[force] (\mx+0.8*\w,0.50*\h) --++ (0,-1.3*\h) node[right=5,below=-3] {$-mg\vu{y}$};
     \draw[line width=1mm,blue,->] (\mx,0.950*\h) --++ ( 2*\h,0) node[above] {F};
    % \draw[force] (\mx+0.1*\w,0.10*\h) --++ (-0.9*\h,0) node[right=6,above=1] {$\vbF_\mathrm{f}$};
 
  
\end{tikzpicture}




 
	% \begin{tikzpicture}[scale=0.8]
	% \draw[thick, blue, ->] (-2,1) -- (0,1) node[midway, above]{$F$};
	% \draw (0,0) rectangle (2,2);
	% \draw (2,0) rectangle (4,2);
	% \node at (1,1) {$A$};
	% \node at (3,1) {$B$};
	% \draw (-1,0) -- (5,0);
	% \end{tikzpicture}	
\end{figure}


\begin{soln}
free-body diagrams for $A$, $B$, and entire system are given below
\end{soln}

\begin{figure*}[ht]
	\centering
	\begin{minipage}{0.3\textwidth}
		\centering
		\begin{tikzpicture}
		\draw[thick] (-0.5,-0.5) rectangle (0.5,0.5);
		\draw[blue,fill] (0,0) circle(0.06);
		\draw[blue,thick,->] (0,0) -- (1.5,0) node[above]{$F$};
		\draw[blue,thick,->] (0,0) -- (-1,0) node[above]{$R$};
		\node at (0,-1) {$A$};
		\end{tikzpicture}
	\end{minipage}\hfil
	\begin{minipage}{0.3\textwidth}
		\centering
		\begin{tikzpicture}
		\draw[thick] (-0.5,-0.5) rectangle (0.5,0.5);
		\draw[blue,fill] (0,0) circle(0.06);
		\draw[blue,thick,->] (0,0) -- (1,0) node[above]{$R$};
		\node at (0,-1) {$B$};
		\end{tikzpicture}
	\end{minipage}\hfil
	\begin{minipage}{0.3\textwidth}
		\centering
		\begin{tikzpicture}
		\draw[thick] (-0.5,-0.5) rectangle (0.5,0.5);
		\draw[blue,fill] (0,0) circle(0.06);
		\draw[blue,thick,->] (0,0) -- (1.5,0) node[above]{$F$};
		\node at (0,-1) {$A$ and $B$ combined};
		\end{tikzpicture}
	\end{minipage}	
\end{figure*}
\sidenote{In fact, only two of the three equations are independent. You can easily check that adding the equation for $A$ to that for $B$ would produce the equation for the system. To solve the two unknowns for this problem, any two of the three equations shall do the job.}
\begin{soln}equations of motion can be written down for each free-body diagram and solved
\begin{equation*}
\text{for $A$:}  F-R = M_A a 
\end{equation*}
\begin{equation*}
\text{for $B$:}  R = M_B a 
\end{equation*}
\begin{equation*}
\text{for system:} F=(M_A+M_B)a  
\end{equation*}
\begin{equation*}
a=\frac{F}{M_A+M_B}
\end{equation*}
\begin{equation*}
R= \frac{M_B}{M_A+M_B}F 
\end{equation*}
\end{soln}

\example{A vehicle of mass 1500 kg is towing a trailer of mass 500 kg by a light inextensible tow-bar. The engine of the vehicle exerts a driving force of 9600 N, and the tractor and the trailer experience resistances of 3600 N and 1800 N respectively. Find the acceleration of the vehicle and the tension in the tow-bar.}


    
free-body diagrams for trailer, vehicle and entire system are given below

\begin{figure*}[ht]
	\centering
	\begin{minipage}{0.15\textwidth}
		\centering
		\begin{tikzpicture}[scale=1.1]
		\draw[thick] (-0.5,-0.5) rectangle (0.5,0.5);
		\draw[blue,fill] (0,0) circle(0.08);
		\draw[blue,thick,->] (0,0) -- (0.855,0) node[above]{$T$};
		\draw[blue,thick,->] (0,0) --++ (-0.54,0) node[above right]{$f_t$};
		\node at (0,-1) {trailer};
		\end{tikzpicture}
	\end{minipage}\hfil
	\begin{minipage}{0.38\textwidth}
		\centering
		\begin{tikzpicture}[scale=1.1]
		\draw[thick] (-0.5,-0.5) rectangle (0.5,0.5);
		\draw[blue,fill] (0,0) circle(0.08);
		\draw[blue,thick,->] (0,0) -- (2.88,0) node[above]{$F_D$};
		\draw[blue,thick,->] (0,-0.05) --++ (-1.08,0) node[below]{$f_v$};
		\draw[blue,thick,->] (0,0.05) --++ (-0.855,0) node[above]{$T$};
		\node at (0,-1) {vehicle};
		\end{tikzpicture}
	\end{minipage}\hfil
	\begin{minipage}{0.38\textwidth}
		\centering
		\begin{tikzpicture}[scale=1.1]
		\draw[thick] (-0.5,-0.5) rectangle (0.5,0.5);
		\draw[blue,fill] (0,0) circle(0.08);
		\draw[blue,thick,->] (0,0) -- (2.88,0) node[above]{$F_D$};
		\draw[blue,thick,->] (0,-0.05) --++ (-1.08,0) node[below]{$f_v$};
		\draw[blue,thick,->] (0,0.05) --++ (-0.54,0) node[above right]{$f_t$};
		\node at (0,-1) {system};
		\end{tikzpicture}
	\end{minipage}	
\end{figure*}
\begin{soln}
equations of motion can be written down for each free-body diagram:
   
\begin{align*}
\text{for trailer:} & T - f_t = M_t a \RA T - 1800 = 500a\\
\text{for tractor:} & F_D - f_v - T = M_v a \RA 9600 - 3600 - T = 1500a \\
\text{for system:} & F_D - f_v - f_t = (M_v + M_t)a \RA  = 2000a \\
\end{align*}

solving simultaneous equations, we find
\begin{equation*}
a= 2.1 \mpss, \quad \text{and} \quad T=2850 \text{ N} 
\end{equation*}
\end{soln}

\subsection{Tension in ropes}
Tension, like the normal reaction force, is a useful shorthand that tells us a lot more about the situation than you might, as a new physicist, realise. The tension is a force along the length of a flexible medium, such as a rope or cable, it always acts both ways. To make that clearer, consider a person pulling a rope attached to a wall:
\begin{figure}
    \centering
    
\begin{tikzpicture}[scale=1.5]
  \def\WH{2.3}  % wall height
  \def\WT{0.3}  % wall thickness
  \def\GW{2.1}  % ground width
  \def\GD{0.25} % ground depth
  \def\h{0.6}   % mass height
  \def\w{0.7}   % mass width
  \def\H{2.0}   % human height
  \def\mx{0.5*\GW} % mass x coordinate
  
  % SETUP
  \rope{(0,0.6*\H) --++ (0.3*\GW,0) coordinate (RH)}
  \draw[ground] (0,0) -- (0,\WH) --++ (-\WT,0) --++ (0,-\WH-\GD) --++
                (\WT+\GW,0) -- (\GW,0) -- cycle;
  \draw (0,\WH) -- (0,0) -- (\GW,0);
  
  % PERSON
  \draw[thick] (0.7*\GW,\H) circle (0.3) coordinate (H);
  \draw[thick] (H)++(-98:0.3) coordinate (N) to[out=-98,in=75]++ (-0.07*\GW,-0.40*\H) coordinate (P);
  \draw[thick,line cap=round] (N)++(-98:0.03) to[out=-115,in=10] ([yshift=0.5]RH);
  \draw[thick,line cap=round] (N)++(-98:0.03) to[out=-100,in=0] ([yshift=-0.5]RH);
  \draw[thick] (P) to[out=-110,in=78] (0.46*\GW,0);
  \draw[thick] (P) to[out=-105,in=88] (0.52*\GW,0);
  
\end{tikzpicture}
 \caption{Pulling a rope from a wall}
    \label{fig:pullarope}
\end{figure}
The force on the rope from the person must be equal to the force on the person by the rope - from Newton's 3rd law. By the same logic, the force on the wall must be the same as the force exerted on the rope by the person.

\subsection*{Pulleys}

% \begin{figure}
% \begin{tikzpicture}
%   \def\W{2.8}  % ground width
%   \def\H{2.3}  % ground height
%   \def\D{0.3}  % ground depth
%   \def\h{0.50} % mass height
%   \def\w{0.75} % mass width
%   \def\L{0.7}  % rope length
%   \def\t{0.1}  % peg thickness
%   \def\R{0.30} % pulley radius
%   \def\F{0.8}  % force magnitude
%   \def\v{0.40} % velocity magnitude
%   \def\ang{-100}      % pull angle
%   \def\px{0.58*\W}    % pulley x position
%   \def\py{0.80*\H}    % pulley y position
%   \def\my{0.30*\H}    % mass y position
%   \def\rx{\px-1.0*\R} % rope-peg x position
%   \def\Hh{2.4}        % human height
%   \def\Hw{0.8}        % human width
%   \def\Hy{-0.4*\H}    % human y position
%   \def\Hx{0.21*\W}    % human x position
%   \coordinate (P)  at (\px,\py);    % pulley
%   \coordinate (M)  at (\px+\R,\my); % mass
%   \coordinate (T)  at ({\px+\R*sin(\ang)},{\py-\R*cos(\ang)}); % mass
%   \coordinate (RH) at ($(T)+(\ang:1.3)$); % right hand
  
%   % PULLEY & MASS
%   \rope{ %\draw[rope,line cap=round]
%     (RH) -- (T) arc(\ang+270:0:\R) --++ (0,\my-\py)}
%   \pic at (P) {pulley={\R}};
%   \pic at (P) {mount={90:\H-\py}};
%   \pic at (M) {weight={\w,0.5*\w,\h}};
%   \draw[mass] (M)++(-\w/2,0) rectangle++ (\w,-\h) node[midway] {$m$};
%   \draw[ground] (0,\H) rectangle++ (\W,\D);
%   \draw (0,\H) --++ (\W,0);
  
%   % PERSON
%   \draw[thick] (\Hx,\Hy+\Hh) circle (0.4) coordinate (H); % head
%   \draw[thick,line cap=round] % body
%     (H)++(-117:0.4) coordinate (N) to[out=-100,in=110]++ (0,-0.40*\Hh) coordinate (B);
%   \begin{pgfonlayer}{back} % draw on back
%     \draw[thick,line cap=round] % left arm
%       (N)++(-85:0.03) to[out=-70,in=230,looseness=1.2] ($(RH)!0.18!(T)+(0.03,0)$);
%   \end{pgfonlayer}
%   \draw[thick,line cap=round] % right arm
%     (N)++(-85:0.03) to[out=-105,in=220,looseness=1.25] ($(RH)!0.02!(T)+(0.03,0.02)$);
%   \draw[thick] (B) to[out=-65,in=72] (\Hx+0.04*\Hw,\Hy); % left leg
%   \draw[thick] (B) to[out=-55,in=95] (\Hx+0.39*\Hw,\Hy); % right leg
  
%   % FORCES
%   % \draw[vvec] (M)++(0.7*\w,-0.3*\h) --++ (0,\v) node[above=2,right=-2] {$v$};
%   % \ifthenelse{\boolean{showforces}}{
%   %   \draw[vvec] (T)++(\ang-16:0.4) --++ (\ang:\v) node[anchor=65,inner sep=1] {$v$};
%   %   \draw[force] (RH)++(\ang+165:0.6*\F) --++ (\ang:\F) node[below=-2] {$\vbT$}; % hand
%   %   \draw[force] (M)++(0.2*\w,0.3*\h) --++ (0,\F) node[right=-2] {$\vbT$};
%   %   \draw[force] (M)++(-0.3*\w,-0.8*\h) --++ (0,-\F) node[below=-3] {$-mg\vu{y}$};
%   %}{}
  
% \end{tikzpicture}
% \end{figure}


A \emph{pulley} is basically a wheel that carries a string/rope/cable, in this section, we only consider pulleys whose axis of rotation is fixed, such pulleys can be used to change direction of tension in a taut string. We also assume pulleys to be \emph{ideal}: they have no mass and no friction\footnote{I really need a spherical cow icon, this will have to do: 
    \includegraphics[scale=.1]{sphericalcow.png}}
for an ideal pulley, tensions on both sides are equal: $T_1 = T_2$.



\begin{marginfigure}
	\vspace{0pt}
	\centering
		\begin{tikzpicture}[scale=0.72]
		\draw[thick,gray] (0, 0.5) -- (0,4.5);
		\draw[thick,gray] (2, 1.5) -- (2,4.5);
		\draw[thick]  (-0.5,-0.5) rectangle (0.5,0.5);
		\draw[thick]  (1.5,0.5) rectangle (2.5,1.5);
		\node at (0,0) {$A$};
		\node at (2,1) {$B$};
		\draw[thick] (1,4.5) circle(1) circle(0.9);
		\draw[ultra thick] (0,6) -- (2,6) (1,6) -- (1,4.5);
		\draw [blue, thick, ->] (0, 0.5) -- (0,2.1) node[right]{$T$};
		\draw [blue, thick, ->] (2, 1.5) -- (2,3.1) node[right]{$T$};
		\draw [blue, thick, ->] (0, -0.5) -- (0,-2.5) node[right]{$m_Ag$};
		\draw [blue, thick, ->] (2, 0.5) -- (2,-0.83) node[right]{$m_Bg$};
		\draw [purple, thick, <-] (-1, -0.5) -- (-1,0) node[left]{$a$} -- (-1,0.5);
		\draw [purple, thick, <-] (3, 1.5) -- (3,1) node[right]{$a$} -- (3,0.5);
		\end{tikzpicture}
	\vspace{-10pt}
\end{marginfigure}

\example{Two blocks of mass $m_A$ and $m_B$ ($m_A>m_B$) are joined together by a light inextensible string. The string passes over a smooth pulley as shown. The two blocks are suddenly released from rest. Find the acceleration of each block and the tension in the string.}

\begin{soln}
    
Apply Newton's second law to each block:

{
	\centering
	
	$\left\{ \begin{array}{llll} \text{for $A$:} & m_Ag-T &=& m_Aa \\ \text{for $B$:} & T-m_Bg &=& m_Ba \end{array} \right.$
	
}

\vspace{0.4em} adding the two, one obtains equation of motion for whole system:

{
	\centering
	
	$ m_A g - m_B g = (m_A+m_B) a $
	
}

		
solving these equations, we find
\begin{equation*}
	a = \frac{m_A-m_B}{m_A+m_B}g \qquad T = \frac{2m_Am_Bg}{m_A+m_B} 
\end{equation*}
\end{soln} 


\example{A mass $M=4.0 \text{ kg}$ is attached to a block of mass $m=2.0 \text{ kg}$ through	a light string which passes over a frictionless pulley as shown. When both masses are released, find the acceleration and the tension in the string.}

\begin{marginfigure}
	\centering
	\vspace*{-10pt}
	\begin{tikzpicture}[scale=0.64]
	\draw[thick] (-6,0) rectangle (-4.5,1.6);
	\draw[very thick] (-7,0) -- (0,0) -- (0,-4);
	\draw[thick] (0.4,-1.5) rectangle (1.4,-2.5);
	\draw[ultra thick] (0,0) -- (0.6,0.5) ;
	\draw (0.6,0.5) circle(0.3) circle(0.25);
	\draw[thick,gray] (-4.5,0.8) -- (0.6,0.8) (0.9,0.5) -- (0.9,-1.5);
	\draw[thick,blue,fill] (-5.25,0.8) circle(0.05);
	\draw[thick,blue,->] (-5.25,0.8) --++ (1.4,0) node[above]{$T$};
	\draw[thick,blue,->] (-5.25,0.8) --++ (0,-2.1) node[right]{$Mg$};
	\draw[thick,blue,->] (-5.25,0.8) --++ (0,2.1) node[right]{$R$};
	\draw[thick,blue,fill] (0.9,-2) circle(0.05);
	\draw[thick,blue,->] (0.9,-2) --++ (0,1.4) node[right]{$T$};
	\draw[thick,blue,->] (0.9,-2) --++ (0,-2.1) node[right]{$mg$};
	\foreach \x in {-0.3,-0.8,...,-6.3} \draw (\x,0) --++ (-0.4,-0.4);
	\foreach \y in {-0.3,-0.8,...,-3.3} \draw (0,\y) --++ (-0.4,-0.4);
	\draw[thick,purple,->] (2,-1.2) -- (2,-2.8) node[midway,right]{$a$};
	\draw[thick,purple,->] (-4.2,2.1) --++ (1.6,0) node[midway,above]{$a$};
	\end{tikzpicture}
\end{marginfigure}

\begin{soln}
  Apply Newton's second law to each mass:

{
	\centering
	
	$\left\{ \begin{array}{lcll} \text{for $M$:} & T &=& Ma \\ \text{for $m$:} & mg-T &=& ma \end{array} \right.$
	
}

\vspace{0.4em} adding the two equations, we have:

{
	\centering
	
	$ mg = (M+m) a $
	
}

so acceleration is:

{
	\centering
	
	$a = \frac{mg}{M+m} = \frac{2.0\times9.81}{4.0+2.0} = 3.27 \mpss$
	
}

tension in string: $T = Ma = 4.0 \times 3.27 \approx 13.1 \text{ N}$ 
  
\end{soln}



	
	
	
\subsection{End-of-chapter questions}

\subsection*{Newton's first law}

\question{A little girl tries to lift a luggage bag of mass 25 kg. She pulls upwards with a force of 150 N. The bag does not move. What is the normal reaction from the floor?}

\question{To push a trolley around in a supermarket with constant velocity, you need to exert a steady force. How does this fact agree with Newton's first law, which suggests that motion with constant velocity requires no force?}

\question{A worker is pulling a wagon of mass of 40 kg across a lawn at a constant velocity. He applies a force of 200 N at an angle of $15^\circ$ above the horizontal. (a) Draw a free-body diagram for the wagon. (b) Find the frictional force. (c) Find the normal contact force.}

\subsection*{Newton's second law}

\question{(a) Forces of 3.0 N and 4.0 N act at right angles upon a mass of 160 g. What is the acceleration produced? (b) If the angle between the two forces are allowed to vary, what is the maximum possible acceleration they produce on the same mass? (c) What about the minimum possible acceleration?}
	
\question{Explain why it becomes increasingly easier for an rocket to accelerate as it travels through space. (Hint: consider the fuel carried by the rocket.)}

\question{Many cars are equipped with airbags which can inflate quickly in case of a collision event. Using Newton's second law, suggest why airbags could protect the driver and the passenger in the car during a car crash.}

\question{A rocket of mass $30,000$ kg is launched vertically upwards at uniform acceleration of $1.6 \mpss$. What is the minimum thrust force required?}

\question{A fire-fighter of mass 85 kg slides down a vertical pole. He descends through a distance of 6.0 m in 2.0 seconds. (a) Find the average acceleration. (b) Find the average frictional force acting on the fire-fighter.}

\question{A trolley has mass $m$. A person needs to push the trolley with force $F$ to produce an acceleration of $a$, and with force $2F$ to produce an acceleration of $3a$. Find, in terms of $m$ and $a$, the constant resistive force opposing the trolley's motion.}

\question{A girl stands onto a bathroom scale and finds the reading is 35.0 \text{ kg}. She then takes the scale into a lift, what mass reading would she observe if the lift (a) is going down at a constant speed, (b) is accelerating downwards at $2.1 \mpss$?}

\question{A pirate finds a box of gold coins at the bottom of a lake. The box and its contents have a total mass of 40 kg. The pirate pulls on the box by means of a cable, so that the box is made to rise vertically through the water. Meanwhile, the flow of water creates a constant horizontal force on the box, and the upthrust on the box is known to be 150 N. At one instant, the pirate applies a force of 380 N at an angle of $25^\circ$ to the upward vertical, and the acceleration of the box is found to be $0.80 \mpss$. Assume all the forces acting are coplanar. (a) Draw a free-body diagram for the box. (b) Find the horizontal force due to water flow. (c) Find the drag force on the box.}



\subsection*{Newton's third law}

\question{A book placed on your desk experiences two forces: its weight and the support force. Identify the associated reaction forces.}

\question{A student deduces that a rocket travelling in space can never accelerate because the propelling force acting on the rocket is cancelled by an equal and opposite force. Explain why this statement is incorrect.}

\question{A U-shaped magnet lies on a top-pan balance and a mass reading of 180 g is registered. A current-carrying wire is then placed above the magnet. The wire experiences an additional force of 0.30 N that acts upwards. What is the mass reading on the balance?}

\subsection*{Terminal velocity}

\question{A light ball and a heavy ball of the same size are released from a very high tower, state and explain whether they will reach the ground at the same time.}

\question{A stone is dropped from rest from a high tower. Air resistance is not negligible as the stone reaches terminal speed. Sketch two separate graphs to show the variation of its displacement and acceleration with time.}

\question{How does the terminal speed of a parachutist before opening the parachute compare to that after? Explain your reasons.}

\question{A ball is thrown horizontally from the top of a cliff. Effects of air resistance cannot be neglected. What happens to the horizontal and vertical components of the ball’s velocity?}

\question{A small sphere of mass 20.0 g is dropped from rest in a viscous liquid. When the sphere is moving at a speed of $v$, the viscous drag has a magnitude of $f = \alpha v^2$, where $\alpha=14 \text{ kg m}^{-1}$. (a) What is the sphere's acceleration at the instant when it is released? (b) What is the acceleration when it is moving at $5.0 \text{ cm s}^{-1}$? (c) What is the terminal velocity?}

\question{A stone is thrown with some initial velocity at an angle to the horizontal. Sketch on the same graph the path of the stone if (a) air resistance is negligible, (b) air resistance is significant.}

\subsection*{Inclined slopes}

\question{A 3.0 kg mass is placed on an inclined plane and it does not move. Given that the normal contact force acting on it is 28.0 N. (a) Find the angle of the plane to the horizontal. (b) Find the frictional force acting on the mass.}

\question{A small mass slides down a frictionless slope with an acceleration of $2.8 \mpss$. Determine the angle that the slope makes with the horizontal.}

\question{A car of mass 1400 kg is moving up a slope at a constant velocity of $13.5\mps$. The slope makes an angle of $6.0^\circ$ to the horizontal. Total resistive force of 650 N acts on the car. What is the driving force required to push the car up the slope?}

\question{A shopping trolley somehow loses control and runs down a straight slope from rest. The slope makes an angle of $3.0^\circ$ to the horizontal. The resistive force acting on the trolley is a constant 15 N. The trolley and its contents have a total mass of 40 kg. (a) Find the acceleration of the trolley. (b) Determine the time for the trolley to travel a distance of 4.0 m along the the slope. (c) Suggest why the slope in shopping malls are not made any steeper.}

\question{A heavy log of mass 240 kg is initially placed at a point $P$ at the bottom of a slope. A motor drags the log up the slope through a cable. The slope is inclined at an angle of $16^\circ$ to the horizontal. The motor provides a tension of 1200 N parallel to the slope. The friction that acts on the log is a constant 450 N. (a) Find the acceleration of the log. (b) Find the time taken to pull the log through a distance of 8.0 m to a point $Q$. (c) Find the velocity of the log at $Q$. (d) The cable breaks when the log reaches $Q$, find the distance moved beyond $Q$ until the log's speed becomes zero. (e) The log will then slide back down the slope. Find the time for the log to return to its starting position. (f) Sketch a $v$-$t$ graph for the log from the start at $P$ until it returns to $P$.}

\subsection*{Many-body problems}

\question{Block $A$ of mass 5.0 kg is connected by means of a light string to block $B$ of mass 3.0 kg. The two blocks are placed on a horizontal table. A force of 30 N is applied to pull on block $A$. Given that the friction on each block is 30\% of its own weight. (a) Find the acceleration of the blocks. (b) Find the tension in the string.}

\begin{marginfigure}
	\centering
	\vspace*{2pt}
	\begin{tikzpicture}[scale=0.64]
	\draw[thick] (-6,0) rectangle (-4.5,1.6);
	\node at (-5.25, 0.8) {$M$};
	\draw[very thick] (-7,0) -- (0,0) -- (0,-5) -- (2,-5);
	\draw[thick] (0.4,-1.5) rectangle (1.4,-2.5);
	\node at (0.9, -2) {$m$};
	\draw[ultra thick] (0,0) -- (0.6,0.5) ;
	\draw[thick] (0.6,0.5) circle(0.3) circle(0.25);
	\draw[thick,gray] (-4.5,0.8) -- (0.6,0.8) (0.9,0.5) -- (0.9,-1.5);
	\foreach \x in {-0.3,-0.8,...,-6.3} \draw (\x,0) --++ (-0.4,-0.4);
	\foreach \x in {0.3,0.8,...,1.8} \draw (\x,-5) --++ (-0.4,-0.4);
	\foreach \y in {-0.3,-0.8,...,-4.8} \draw (0,\y) --++ (-0.4,-0.4);
	\draw [<->] (0.9,0.05-5) -- (0.9,-0.05-2.5) node[midway,right]{$H$};
	\end{tikzpicture}
\end{marginfigure}

\question{A box of mass $M=3.6 \text{ kg}$ rests on a horizontal, rough surface. The box is connected to a block of mass $m=2.0 \text{ kg}$ through a light cord that passes over a frictionless pulley as shown. The box is released from rest. Given that the box experiences a frictional force of 12 N and the block is initially at a height of $H=0.80 m$ above the floor. (a) Find the acceleration of the block. (b) Determine the time taken for the block to hit the floor.}


\chapter{Mechanical Equilibrium}

In this chapter, we will study the mechanical equilibrium of objects. We often consider mechanical problems in the context of \emph{point objects}, in the case of a dimensionless point, there cannot be any turning moment and zero resultant force suffices for equilibrium. For \emph{rigid bodies}, any force many produce a turning effect called a \emph{moment}.

It follows then, that rigid bodies must satisfy another condition to stay in equilibrium...

\subsection{Moment of force}

\begin{ilight}
	The \keypoint{moment}\index{moment} of a force, is defined as the product of the force and the perpendicular distance from the pivot to the line of action. The word \keypoint{torque}\index{torque} is often used instead of moment, especially when we're \cancel{torquing} talking about the dynamic rotation of a motor and a change in angular momentum, rather than a statics problem. It's convenient to use $\tau$ as shorthand for moment: \begin{empheq}[box=\tcbhighmath]{equation*}{\tau = Fd_\perp} \end{empheq}
\end{ilight}

\begin{marginfigure}
	\centering
	\begin{tikzpicture}[scale=0.78,rotate=15]
	\draw[fill=gray!20,very thick] (-3.2,0.4) to [out=25,in=170] (-1,0.7) to [out=-10,in=150] (2.7,0.9) to [out=-30,in=35] (2.5,-1.3) to [out=215,in=30] (-0.5,-1.3) to [out=210,in=205] (-3.2,0.4);
	\draw[thick,blue,->] (2,-.2) --++ (-1.5,3) node[above]{$F$};
	\draw[fill] (-2.5,-.2) circle (0.05) node[left]{$O$};
	\draw (-2.6,-.3) --++ (-.6,-1) node[below]{pivot};
	\draw[thick,orange,<->] (-2.5,-.2) --++ (26.565:4.025) node[pos=0.56,above left,black]{$d_\perp$};
	\draw[thick] (-2.5,-.2) -- (2,-.2) node[midway,above]{$d$};
	\draw[thick] (1.5,-.2) arc(180:116.565:.5);
	\node at (1.4,0.2) {$\theta$};
	\draw[very thick,purple,->] (24:3) arc(10:42:4.5);
	\end{tikzpicture}
 \caption{Moment of a shape}
 \label{fig:moment}
\end{marginfigure}

The unit of torque\sidenote[][-10cm]{You'll not learn about angular momentum unless you take the sports science option} or moment is the Newton-meter : $[\tau] = \text{N m}$\\

The critical part of working out the moment of a force is that the displacement from the axis of rotation and the force are perpendicular. This means that if you have an object like \ref{fig:moment} then you will need to resolve \emph{either} the force \emph{or} the displacement. Which one you choose resolve is immaterial, so you might as well use whichever is most convenient!\\
\example{In the diagram, if we resolve the displacement; $$d_\perp = d\sin\theta$$ the moment of this force is: $$\tau = Fd\sin\theta$$\\
If we resolve the force: $$F_\perp = F\sin\theta$$ so moment of this force is: $$\tau = F\sin\theta d$$}

The moment of a force is treated\footnote{\piste Rigorously speaking, moment is a \emph{pseudovector}, which means that it does not transform quite like a normal vector although it does have a direction. In particular, if an object acted by a force is reflected across a plane, the moment of this force would not be reflected. Instead, it would be reflected and \emph{reversed} as a vector quantity}
\footnote[][]{\piste Using vector notation, moment of a force can be defined as a \emph{cross product}: $\overrightarrow{\tau} = \overrightarrow{r}\times\overrightarrow{F}$, where $\overrightarrow{r}$ is the position vector from the pivot to the point at which the force is applied.}
it can act in \emph{clockwise} or \emph{anti-clockwise} direction. 

The moment of a force produces turning effects,\footnote{Moment is like the rotational counterpart of a force: force changes the state of translational motion, moment changes the state of rotation.} if there exists a non-zero moment, the object will have clockwise or anti-clockwise \emph{angular acceleration}.\\
\titem Note that moment of a force depends on choice of pivot location. The moment of the same force with respect to different points can be very different. We often use this trick in questions to reduce the difficulty of the problem or the number of calculations required.

\subsection{Torque of couple}\label{ch:torque-of-couple}

Let's take a pair of equal but opposite forces acting at different positions on the same object, as in the figure \ref{fig:torquecouple}.

\begin{marginfigure}
	\centering
	\vspace*{-8pt}
	\begin{tikzpicture}
	\draw[thick,gray, rotate=30, rounded corners] (-2.4,-.3) rectangle (2.4,.3);
	\draw[thick,blue,->] (-1.73,-1) -- ++(0,2) node[left]{$F$};
	\draw[thick,blue,->] (1.73,1) -- ++(0,-2) node[right]{$F$};
	\draw[dashed,<->,purple] (-1.73,0) -- (0,0) node[midway,above]{$d_{\perp,1}$};
	\draw[dashed,<->,purple] (1.73,0) -- (0,0) node[midway,above]{$d_{\perp,2}$};
	\draw[fill] (0,0) circle(0.05);
	\draw[thick,<->] (-1.73,-1.8) -- (1.73,-1.8) node[midway,above]{$L_\perp$};
	\draw[thick] (-1.73,-1.9) -- (-1.73,-1.7) (1.73,-1.9) -- (1.73,-1.7);
	\end{tikzpicture}
	\vspace*{-16pt}
 \label{fig:torquecouple}
\end{marginfigure}

The two forces produce moments in the same direction, but because they have equal magnitudes and opposite directions, they produce no translational acceleration. The combined effect is called the torque of a couple. The resultant torque due to the couple is:

	$$ \tau = Fd_{\perp,1} + Fd_{\perp,2} = 2F(d_{\perp,1} + d_{\perp,2}) $$

$d_{\perp,1} + d_{\perp,2}$ is perpendicular distance $L_\perp$ between the pair, so:

$$ \tcbhighmath{\tau = F L_\perp} $$

\begin{ilight}
	\keypoint{Torque of a couple} can be therefore defined as the product of one force in the couple and the perpendicular distance between the pair
\end{ilight}
Note that: the torque of couple does not depend on choice of pivot. For same force pair, resultant moment is constant about any point.



\subsection{Moment of weight}

If you recall that weight is a force of gravity which is actually experienced by \emph{all} parts of the object, then you will realise  we need to sum up torques on each part of this object. This this brings a problem since the displacement of every part of the object will be different, fortunately, this calculation can be simplified using the idea of centre of gravity.

\begin{ilight}
	\keypoint{centre of gravity} is a point at which the entire weight of an object is considered to act\index{centre of gravity}
\end{ilight}

There is a similar concept called the centre of mass;

\keypoint{centre of mass} is the average position of all the mass that makes up the object.

Near the surface of the earth, mass and weight are directly proportional to each other so centre of mass is interchangeable with centre of gravity if we stay on earth\footnote{\piste You might like to explore this idea a bit - when and where would they be in a different part of a body?}.

For a regularly-shaped uniform object, the centre of gravity/mass is its geometrical centre, If an object is hung freely, centre of gravity/mass is vertically below the point of suspension, otherwise weight would produce a non-zero torque about the point of suspension, causing the object to rotate until torque becomes zero:

\begin{figure*}[ht]
	\centering
	\begin{minipage}[b]{0.45\textwidth}
		\centering
		\begin{tikzpicture}[scale=0.9]
		% triangle
		\draw[thick] (-1.4,0) -- (0,2.4) -- (1.4,0) --cycle;
		\draw[dashed] (0,0) -- (0,2.4) (0.7,1.2) -- (-1.4,0) (-0.7,1.2) -- (1.4,0);
		\draw[purple,fill] (0,0.8) circle (0.06);
		% rectangle
		\draw[thick] (-5,2) rectangle (-1.2,4);
		\draw[dashed] (-5,2) -- (-1.2,4) (-5,4) -- (-1.2,2);
		\draw[purple,fill] (-3.1,3) circle (0.06);	
		% circle
		\draw[thick] (-3.2,0.2) circle(1.2);
		\draw[purple,fill] (-3.2,0.2) circle (0.06);		
		\end{tikzpicture}
	centre of mass of uniform lamina
		
		is at the geometrical centre
	\end{minipage}\hfil
	\begin{minipage}[b]{0.45\textwidth}
		\centering
		\begin{tikzpicture}[scale=0.9]
		\draw[thick,rotate=30] (-1.5,-1) rectangle (1.5,1);
		\draw[fill] (0,1) circle (0.05);
		\draw[thick] (0,1) -- (0,3);
		\draw[very thick] (-0.8,3) -- (0.8,3);
		\draw[dashed] (0,1) -- (0,-1.8);
		\draw[purple, fill] (0,0) circle (0.05);
		\end{tikzpicture}
		
		centre of mass is vertically below
		
		the point of suspension
	\end{minipage}
\end{figure*}

A simple, practical way to find the centre of gravity/mass of an object is to suspended it from several positions. Each time we draw a vertical \emph{plumb-line} through the point of suspension. The centre of mass/gravity lies where the lines intersect.
\section{Stability}

By using the idea of torque, we can easily understand how, or why some items are stable, and others less-so. When an object is sat on a level surface, undisturbed, the resultnt forces must pass through the centre of mass, resulting in no torque:
% BLOCK - metastable
\def\H{3.1} % block height
\def\W{1.4} % block width
\def\R{sqrt(\W^2+\H^2)/2} % block half-diagonal
\def\D{0.25} % ground depth
\def\ang{atan(\H/\W)} % angle diagonal ccw
\def\angcw{atan2(\H,-\W)} % angle diagonal cw

\begin{figure}[ht]
\centering
\begin{tikzpicture}[scale=1.5]
  \coordinate (O) at (0,0);
  \coordinate (CM) at (0,\H/2);
  \coordinate (BL) at (-\W/2,0); % bottom left corner
  \coordinate (BR) at (\W/2,0); % bottom right corner
  \coordinate (TR) at (\W/2,0); % top right corner
  \draw[ground] (-0.8*\W,0) rectangle (1.4*\W,-\D);
  \draw[thick] (-0.8*\W,0) -- (1.4*\W,0);
  \draw[mass] (-\W/2,0) rectangle++ (\W,\H);
  \node[left] at (-\W/2,\H/2) {$W$};
  \node[above] at (0,\H) {$H$};
  \draw[CM] (CM) circle(0.06*\W);
  \draw[dashed] (BR) -- (CM) --++ ({\angcw}:0.2*\W);
  \draw[dashed] (BR)++(130:{\R}) arc(130:35:{\R});
  \draw pic[myarr,"$\theta$"scale=0.9,draw,angle radius=18,angle eccentricity=1.35] {angle=TR--BR--CM};
  \draw[force] (CM) --++ (0,-0.24*\H) node[midway,below=2,left=-.1] {$m\vb{g}$};
  \draw[force] (O) --++ (0,0.24*\H) node[midway,left=-.1] {$\vb{F}_\mathrm{N}$};
  \node[mypurple,above right=.1] at (BR) {$\vb*\tau=0$};
  %\draw[->] (F)++(42:1.9*\L) arc(70:0:0.5*\L) node[midway,above right=-2] {$\alpha$};
  \draw[rvec] (BR) -- (CM) node[below,right] {$\vb{r}$};
  \draw[rvec] (BR) -- (O) node[midway,above] {$\vb{r}_\mathrm{t}$};
  %\draw[rvec] (BR) -- (O) node[midway,right=2,below=-1.5] {\contour{white}{$\vb{r}_\mathrm{t}$}};
  \fill[blue!50!black] (BR) circle(0.03*\W);
\end{tikzpicture}
\end{figure}
As soon as the block is disturbed from this position, there will be a resultant torque. If the torque produced causes the block to return to it's origional position we would say the object is \emph{stable}. 
\begin{marginfigure}
\centering
% BLOCK - torque
\begin{tikzpicture}[scale=1.5]
  \def\ang{11}
  \coordinate (CM) at ({\angcw-\ang}:{\R});
  \coordinate (BR) at (0,0); % bottom right corner
  \coordinate (BL) at ({180-\ang}:\W); % bottom left corner
  \coordinate (TR) at ({\ang}:\W); % bottom left corner
  \coordinate (T) at ({\angcw-\ang-8}:{1.15*\R}); % torque 
  \draw[ground] (-1.3*\W,0) coordinate (L) rectangle (0.9*\W,-\D);
  \draw[thick] (-1.3*\W,0) -- (0.9*\W,0);
  \draw[mass,rotate around=({-\ang}:(BR))]
    (-\W,0) rectangle++ (\W,\H);
  \draw[CM] (CM) circle(0.06*\W);
  \draw[dashed] (BR) --++ (0,{1.6*\R});
  \draw[dashed] (BR)++(120:{\R}) arc(120:40:{\R});
  \draw[force] (CM) --++ (0,-0.24*\H) node[midway,below=2,left=-1] {$m\vb{g}$};
  \draw[force] (BR) --++ (0,0.24*\H) node[below right=0] {$\vb{F}_\mathrm{N}$};
  \pic[scale=1] at (T) {Tin};
  \node[mypurple,above=1] at (T) {$\vb*\tau$};
  \draw pic[myarr,draw,angle radius=35,angle eccentricity=1.1] {angle=BL--BR--L};
  \draw[rvec] (BR) -- (CM) node[midway,below left=-2] {$\vb{r}$};
  \draw[->] (110:1.1*\H) arc(110:148:0.4*\H) node[above left=-1] {$\alpha$};
  \fill[blue!50!black] (BR) circle(0.03*\W);
\end{tikzpicture}
\end{marginfigure}
As you can see - there can be more than one stable position for most shapes;\footnote{except the Gombok, which is only stable in one orientation}:


\begin{marginfigure}
% BLOCK - plot
\centering
\begin{tikzpicture}
  \def\xmax{3.4}
  \def\ymax{2.8}
  \def\angms{25}            % unstable angle
  \def\angus{atan(tan(\angms))} % unstable angle
  \def\angne{60}            % non-equilibrium
  \def\A{2.0}               % amplitude / yscale
  \def\yW{\A*sin(\angms)}   % y position phi =  0
  \def\yH{\A*cos(\angms)}   % y position phi = 90
  \def\xH{0.7*\xmax}        % x position phi = 90
  \def\xi{(90-\angms)/\om}  % x position unstable
  \def\om{(90/(\xH))}
  \def\h{0.12*\xmax} % height mini block
  \def\w{\h*tan(\angus)} % width mini block
  \coordinate (O) at (0,0);
  \coordinate (P-2) at ({(\angus-90)/\om},0.014);
  \coordinate (P-1) at ({(\angne-90)/\om},0.013);
  \coordinate (P+1) at ({(90-\angne)/\om},0.018);
  \coordinate (P+2) at ({(90-\angus)/\om},0.018);
  
  % MASS BLOCKS
  \begin{scope}[opacity=0.5]
    \draw[mass,rounded corners=0.9]
      ({-\xH-\w/2},0.018) rectangle++ ({\w},\h) % phi = -90
      (-\h/2,0.018) rectangle++ (\h,{\w})       % phi = 0
      ({\xH-\w/2},0.018) rectangle++ ({\w},\h); % phi = +90
    \draw[mass,rotate around={{\angus}:(P-2)}]
      (P-2) rectangle++ ({\w},\h);
    \draw[mass,rotate around={{\angne}:(P-1)}]
      (P-1) rectangle++ ({\w},\h);
    \draw[mass,rotate around={{-\angne}:(P+1)}]
      (P+1) rectangle++ (-{\w},\h);
    \draw[mass,rotate around={{-\angus}:(P+2)}]
      (P+2) rectangle++ (-{\w},\h);
    \draw[very thin]
      ({\xH+\w/2+0.3*\h},0.018) arc(0:90:0.3*\h) node[pos=0.5,above right,scale=0.5] {$\phi$}
      (P+1)++(-5:0.50*\h) arc(-5:{90-\angne}:0.50*\h) node[pos=0.7,right,scale=0.5] {$\phi$}
      (P+2)++(-5:0.30*\h) arc(-5:{90-\angus}:0.30*\h) node[pos=0.5,above right,scale=0.5] {$\phi$};
    \draw[very thin,mydashed]
      (P-2) --++ (0,1.5*\h)
      (P-1) --++ (0,0.8*\h)
      (P+1) --++ (0,1.2*\h)
      (P+2) --++ (0,1.5*\h);
  \end{scope}
  \fill[blue!50!black!50]
    ({-\xH-\w/2},\h+0.015) circle (0.03)
    (-\h/2,0.015) circle (0.03)
    (P-2)++({90+\angus}:\h) circle (0.03)
    (P-1)++({90+\angne}:\h) circle (0.03)
    (P+1) circle (0.03)
    (P+2) circle (0.03)
    ({\xH+\w/2},0.015) circle (0.03);
  
  % PLOT
  \draw[->,thick] (0,-0.1*\ymax) -- (0,1.05*\ymax) node[below=3,left] {$U=mgh$};
  \draw[->,thick] (-1.05*\xmax,0) -- (1.05*\xmax,0) node[below left] {$\phi$}; %$[\si{\degree}]$ %\theta
  \tick{-\xH,0}{90} node[below,scale=0.9] {$-90^{\text{o}}$};
  \tick{\xH,0}{90} node[below,scale=0.9] {$90^{\text{o}}$};
  \node[below,below left,scale=0.9] at (O) {$0^{\text{o}}$};
  \draw[dashed] (-\xH,0) --++ (0,1.1*\A);
  \draw[dashed] (\xH,0) --++ (0,1.1*\A);
  \tick{0,{\yH}}{0} node[right,above left,scale=0.9] {$mgH/2$};
  \tick{0,{\yW}}{0} node[right,below left,scale=0.9] {$mgW/2$};
  \draw[dashed] (-1.2*\xH,{\yH}) -- (1.24*\xH,{\yH});
  \draw[dashed] (-1.2*\xH,{\yW}) -- (1.24*\xH,{\yW});
  \draw[very thick,orange!80!black,samples=100,smooth,variable=\t]
    plot[domain=-\xmax:-\xH](\t,{\A*sin(\om*abs(\t)-\angms)}) --
    plot[domain=-\xH:0](\t,{\A*sin(\om*abs(\t)+\angms)}) --
    plot[domain=0:\xH](\t,{\A*sin(\om*\t+\angms)}) --
    plot[domain=\xH:\xmax](\t,{\A*sin(\om*\t-\angms)}); %node[right] {$x$}
    
  % CENTER OF MASS
  \draw[CM] (-\xH,{\yH+0.1}) circle(0.08) node[below=1,scale=0.85,fill=white,inner sep=1] {metastable};
  \draw[CM] (\xH,{\yH+0.1}) circle(0.08) node[below=1,scale=0.85,fill=white,inner sep=1] {metastable};
  \draw[CM] ({(90-\angms)/\om},{\A+0.1}) circle(0.08) node[above=.5,scale=0.85] {unstable};
  \draw[CM] ({(\angms-90)/\om},{\A+0.1}) circle(0.08) node[above=.5,scale=0.85] {unstable};
  \draw[CM] (0,{\yW+0.1}) circle(0.08) node[below=.5,scale=0.85] {stable};
  
\end{tikzpicture}
\end{marginfigure}


If we apply these ideas to the real world we can use them to understand that one of the more impressive looking tricks of the circus can be brought into the hands of everyone. 
Consider a tightrope walker:
% TIGHT ROPE arms
\def\L{3.2}   % length
\def\T{0.08}  % plank thickness
\def\H{2.2}   % human height
\def\CM{0.06} % CM radius
\begin{marginfigure}
\centering
   \begin{tikzpicture}
  \coordinate (O) at (0,0);
  \draw[rope] (O) circle (0.08);
  \draw[thick] (0,\H) circle(0.3) coordinate (H);
  \draw[thick] (H)++(-90:0.3) coordinate (N) to[out=-92,in=92]++ (0,-0.40*\H) coordinate (P);
  \draw[limb] (N)++(-120:0.03) to[out=-120,in=30]++ (-0.27*\H,-0.3*\H); % right arm
  \draw[limb] (N)++(-60:0.03) to[out=-60,in=140]++ (0.27*\H,-0.3*\H);   % left arm
  \draw[limb] (P) to[out=-95,in=95] (96:0.1); % right leg
  \draw[limb] (P) to[out=-85,in=88] (84:0.1); % left leg
  \draw[CM] (0,0.54*\H) circle(\CM) node[below right=0,scale=0.8] {CM};
\end{tikzpicture}
    \caption{Walking on a tight-rope unaided.}
    \label{fig:tightrope}
\end{marginfigure}

We can see that the Centre of Mass is substantially higher than the contact point and the situation is unstable. The artist has to be quite active in maintaining their CoM above the contact point - it's hard work!

One of the tricks that you can use however is a long pole to assist balance. Not only does this give you a mass that's easy to move around but it lowers the center of mass:
\begin{marginfigure}
% TIGHT ROPE - stick
\centering
\begin{tikzpicture}
  \coordinate (O) at (0,0);
  \draw[rope] (O) circle (0.08);
  \draw[thick] (0,\H) circle(0.3) coordinate (H);
  \draw[thick] (H)++(-90:0.3) coordinate (N) to[out=-91,in=92]++ (0,-0.40*\H) coordinate (P);
  \draw[limb] (P) to[out=-95,in=95] (96:0.1); % left leg
  \draw[limb] (P) to[out=-85,in=88] (84:0.1); % right leg
  \draw[myred,line width=1.5,line cap=round]
    (-\H,0.1*\H) to[out=40,in=180] (0,0.47*\H) to[out=0,in=140] (\H,0.1*\H);
  \draw[limb] (N)++(-95:0.03) to[out=-120,in=80]++ (-0.19*\H,-0.4*\H); % right arm
  \draw[limb] (N)++(-85:0.03) to[out=-60,in=100]++ (0.19*\H,-0.4*\H);   % left arm
  \draw[CM] (0,0.3*\H) circle(\CM) node[below right=0,scale=0.8] {CM};
\end{tikzpicture}
\end{marginfigure}
This might be clearer if we picture the kind of balance support given to beginners: 
\begin{marginfigure}
% TIGHT ROPE - weights
\begin{tikzpicture}
  \def\L{1.2}  % handle length
  \def\w{0.3}  % mass width
  \def\h{0.45} % mass height
  \def\r{0.02} % rope radius
  \coordinate (O) at (0,0);
  \coordinate (SL) at (-\L/2,0.42*\H); % stick left
  \coordinate (SR) at ( \L/2,0.42*\H); % stick right
  \coordinate (ML) at (-\L/2,-0.2*\H); % mass left
  \coordinate (MR) at ( \L/2,-0.2*\H); % mass right
  \draw[rope] (O) circle (0.08);
  \draw[thick] (0,\H) circle(0.3) coordinate (H);
  \draw[thick] (H)++(-90:0.3) coordinate (N) to[out=-91,in=92]++ (0,-0.40*\H) coordinate (P);
  \draw[limb] (P) to[out=-95,in=95] (96:0.1); % left leg
  \draw[limb] (P) to[out=-85,in=88] (84:0.1); % right leg
  \draw[myred,line width=1.5,line cap=round] (SL)++(-0.05,0) -- (SR) --++ (0.05,0);
  \draw[limb] (N)++(-95:0.03) to[out=-120,in=85]++ (-0.14*\H,-0.44*\H); % right arm
  \draw[limb] (N)++(-85:0.03) to[out=-60,in=95]++ (0.14*\H,-0.44*\H);   % left arm
  \draw[CM] (0,-0.15*\H) circle(\CM) node[below=.5,scale=0.8] {CM};
  \draw[rope] (SL)++(-\r,0.02) arc(180:0:\r) -- ($(ML)+(\r,0)$) --++ (-2*\r,0) -- cycle;
  \draw[rope] (SR)++(-\r,0.02) arc(180:0:\r) -- ($(MR)+(\r,0)$) --++ (-2*\r,0) -- cycle;
  \draw[mass] (ML)++(-\w/2,-\h) rectangle++ (\w,\h);
  \draw[mass] (MR)++(-\w/2,-\h) rectangle++ (\w,\h);
\end{tikzpicture}
\end{marginfigure}


\begin{marginfigure}
	\vspace*{-8pt}
	\centering
	\begin{tikzpicture}[scale=0.6]
	\draw[thick] (-4,-.3) rectangle (4,.3);
	\draw[thick] (-1,-.3) --++ (.4,-.6) --++ (-.8,0) node[midway,below]{$P$} -- cycle;
	\draw[thick,blue,->] (0,0) -- (0,-2) node[right]{$W$};
	\draw[thick,blue,->] (4,0) --++ (0,1.5) node[left]{$F$};
	\draw[thick,blue,->] (-4,0) --++ (0,-1.5) node[right]{$F$};
	\draw[<->] (-4,2.4) -- (-1,2.4)node[midway,above]{30 cm};
	\draw[<->] (-1,2.4) -- (4,2.4)node[midway,above]{50 cm};
	\foreach \x in {-4,-1,4} \draw (\x,2.2) --++ (0,.4);
	\end{tikzpicture}
	\vspace*{-16pt}
\end{marginfigure}


\example{The diagram shows a uniform beam of weight $W=20 \text{ N}$ and length 80 cm pivoted at point $P$. $P$ is 30 cm from one end. Two equal but opposite forces of magnitude $F = 12 \text{ N}$ are acting at the two ends of the beam as shown. What is the resultant moment about point $P$?}

\begin{soln} Moment of weight: $\tau_w = W d_w = 20 \times \left(0.50-\frac{1}{2}\times0.80\right) = 2.0 \text{ N m} \quad$ (clockwise)

Torque of couple: $\tau_c = F L = 12 \times 0.80 = 9.6 \text{ N m} \quad$ (anti-clockwise)

Resultant moment: $\tau_\text{net} = \tau_c - \tau_W = 9.6 - 2.0 = 7.6 \text{ N m} \quad$ (anti-clockwise) \end{soln}

\subsection{Principle of moments}

If there is no turning effect for an object, the total moment of all forces must vanish:

\begin{ilight}
	For a rigid body in equilibrium, sum of all clockwise moments must be equal to the sum of anti-clockwise moments \emph{about any point}, this is called the \keypoint{principle of moments} \index{principle of moments}
\end{ilight}

IE.an object in equilibrium has no turning effect about any point so resultant moment is zero about any point.\\
\footnote{As long as there is no resultant force, then zero resultant moment about any particular point would imply zero resultant moment about any point.
	
Mathematically, let's take a collection of forces $\overrightarrow{F_1}, \overrightarrow{F_2}, \cdots, \overrightarrow{F_n}$ acting at positions $\overrightarrow{r_1}, \overrightarrow{r_2}, \cdots, \overrightarrow{r_n}$ on an object with respect to some fixed point $O$. Suppose their resultant moment vanishes, i.e., $\sum \overrightarrow{\tau_i} \equiv \sum\overrightarrow{r_i} \times \overrightarrow{F_i} =0$, and also their resultant force vanishes, i.e., $\sum \overrightarrow{F_i} = 0$. If we focus on a different point $O'$ with a relative displacement $\overrightarrow{R}$ to point $O'$, then taking moments about $O'$, we will have:
\begin{equation*}
	\sum \overrightarrow{\tau_i'} = \sum \overrightarrow{r_i'}\times\overrightarrow{F_i} = \sum (\overrightarrow{r_i}+\overrightarrow{R})\times\overrightarrow{F_i} = \sum \overrightarrow{r_i}\times\overrightarrow{F_i} + \overrightarrow{R} \times \sum\overrightarrow{F_i} = 0
\end{equation*}
which shows zero resultant moment about one point together with zero resultant force guarantee resultant moment must be zero about any point in space.}

\begin{marginfigure}
	\vspace*{-8pt}
	\centering
	\begin{tikzpicture}[scale=0.6]
	\draw[thick] (-4.5,-.3) rectangle (4.5,.3);
	\draw[thick] (-1.5,-.3) --++ (.4,-.6) --++ (-.8,0) node[midway,below]{$P$} -- cycle;
	\draw[thick,blue,->] (0,0) -- (0,-2.5) node[right]{$mg$};
	\draw[<->] (-4.5,1.2) -- (-1.5,1.2)node[midway,above]{30 cm};
	\draw[<->] (-1.5,1.2) -- (4.5,1.2)node[midway,above]{60 cm};
	\foreach \x in {-4.5,-1.5,4.5} \draw (\x,1) --++ (0,.4);
	\draw[thick] (-4.5,0) --++ (0,-1) node[below, rectangle,draw]{3.0 kg};
	\end{tikzpicture}
	\vspace*{0pt}
\end{marginfigure}

\example{A uniform rod of length 90 cm is pivoted 30 cm from one end. It is balanced with a 3.0 kg load. Find the mass of the rod.}

\begin{soln}
    
 Take moments about $P$:
	
	$ 3.0 \times 9.81 \times 0.30 = m \times 9.81 \times \left(0.60-\frac{1}{2}\times0.90\right)$

so we find mass of rod: $m = 6.0 \text{ kg}$ 
\end{soln}


\example{A student balances a metre rule of mass 120 g supported on a fulcrum at the 40 cm mark. She then places a 20 g mass on the 70 cm mark and a 50 g mass on the 25 cm mark as shown. To balance the rule, what mass should she place on the 15 cm mark?}

\begin{figure}[ht]
	\centering
	\begin{tikzpicture}[scale=1]
	\draw[thick] (-5,-.3) rectangle (5,.3);
	\foreach \x in {10,20,...,90} \draw[thick] (\x/10-5,.1) node[below]{{\scriptsize \x}} --++ (0,.2) ; 
	\foreach \x in {5,15,...,95} \draw[thick] (\x/10-5,.2) --++ (0,.1) ; 
	\draw[thick] (-1,-.3) --++ (.3,-.5) --++ (-.6,0) -- cycle;
	\draw[fill] (2,.45) circle (0.15);
	\node[above] at (2,.5) {{\scriptsize 20 g}};
	\draw[fill] (-2.5,.45) circle (0.15);
	\node[above] at (-2.5,.5) {{\scriptsize 50 g}};
	\draw[gray,fill] (-3.5,.45) circle (0.15);
	\node[above] at (-3.5,.56) {{\scriptsize $m=?$}};
	\end{tikzpicture}
\end{figure}

\begin{soln}
    
 Take moments about the support:
\begin{equation*}
	mg \times (40-15) + 0.050g\times(40-25) = \end{equation*}
 \begin{equation*}
     0.12g\times(50-40) + 0.020g\times(70-40) \RA m=42 \text{ g} 
\end{equation*}
\end{soln}

\example{A cylinder of weight 100 N and diameter 50 cm rests against point $P$ of a curb of height 10 cm. What is the minimum force required to cause the cylinder to roll to the left?}

\begin{marginfigure}
	\vspace*{-15pt}
	\centering
	\begin{tikzpicture}[scale=0.4]
	\draw[thick] (0,0) circle (5) node[right]{$O$};
	\draw[thick] (-8,-3) -- (-4,-3) -- (-4,-5) --++ (9,0);
	\draw[thick,blue,->] (4,3) node[above right, black]{$A$} --++ (-3,4) node[above]{$F$};
	\draw[dashed] (4,3) -- (-4,-3) node[midway, above,rotate=36.87] {50 cm};
	\draw[dashed] (0,-3) node[right] {$B$}-- (-4,-3) node[below left] {$P$};
	\draw[dashed] (0,0) -- (0,-5) node[below] {$C$} -- (-8,-5);
	\draw[thick,blue,->] (0,0) -- (0,-4) node[left]{$W$};
	\draw[<->] (-8.5,-3) -- (-8.5,-5) node[midway, left] {10 cm};
	\end{tikzpicture}
	\vspace*{-0pt}
\end{marginfigure}


\begin{soln} To roll the cylinder, force applied must produce a torque no less than that of weight. The force is minimum if perpendicular displacement is greatest. 
Take moments about $P$ (see diagram):
	$ F_\tmin \times |PA| = W \times |PB| $


Note that $|PB| = \sqrt{|OP|^2 - |OB|^2}$, so
$|PB| = \sqrt{0.25^2 - (0.25-0.10)^2} = 0.20 \text{ m} $

Plug back into the equation above:
\begin{equation*}
	F_\tmin \times 0.50 = 100 \times 0.20 \RA F_\tmin = 40 \text{ N} 
\end{equation*}
\end{soln}

\subsection{Mechanical equilibrium}

Combining Newton's first law and principle of moments, we have the following statement:

\begin{ilight}
	For any mechanical system in equilibrium, two conditions must be satisfied:
	
	\begin{compactenum}
		\item[--] resultant force is zero in any direction: $\sum F = 0$
		
		\item[--] resultant moment is zero about any point: $\sum \tau = 0$
	\end{compactenum}
\end{ilight}

These two conditions allow for many problems to be solved.

\begin{marginfigure}
	\vspace*{-8pt}
	\centering
	\begin{tikzpicture}[scale=0.72]
	\draw[thick] (-4,-.2) rectangle (4,.2);
	\draw[thick] (-3.9,-.5) circle(.3);
	\draw[thick] (3.9,-.5) circle(.3);
	\draw[thick,blue,->] (0,0) -- (0,-1) node[right]{$W_p$};
	\draw[thick,blue,->] (-2,0.8) --++ (0,-3.5) node[right]{$W_m$};
	\draw[thick,blue,->] (-3.9,-.2) --++ (0,3) node[right]{$F_A$};
	\draw[thick,blue,->] (3.9,-.2) --++ (0,1.5) node[right]{$F_B$};
	\draw[<->] (-4,3.5) -- (-2,3.5)node[midway,above]{$x$};
	\draw[<->] (-2,3.5) -- (4,3.5)node[midway,above]{$L-x$};
	\foreach \x in {-4,-2,4} \draw (\x,3.3) --++ (0,.4);
	\node[below] at (-4,-.8) {$A$};
	\node[below] at (4,-.8) {$B$};
	\draw[fill] (-2,1.6) circle(0.2);
	\draw[ultra thick] (-2,1.5) -- (-2,0.9);
	\draw[ultra thick] (-2,0.9) -- (-2.25,0.2) (-2,0.9) -- (-1.75,0.2);
	\draw[ultra thick] (-2.3,0.8) -- (-2,1.4) -- (-1.7,0.8);
	\end{tikzpicture}
	\vspace*{0pt}
\end{marginfigure}


\example{A uniform plank of weight 100 N and length $L=6.0 \text{ m}$ rests horizontally on two supports $A$ and $B$. A man of weight 800 N stands a distance of $x=1.5 \text{ m}$ from end $A$. Determine the forces acting at the two supports.}

\begin{soln} Take moments about $A$: $ W_m x + W_p \cdot\frac{L}{2} = F_B L$

{
	\centering

	
	$ 800 \times 1.5 + 100 \times 3.0 = F_B \times 6.0 \ra F_B = 250 \text{ N} $
	
}

Take moments about $B$: $ W_m (L-x) + W_p \frac{L}{2} = F_A L$

{
	\centering

	$ 800 \times 4.5 + 100 \times 3.0 = F_A \times 6.0 \ra F_B = 650 \text{ N} $
	
}

One can check that: $F_A + F_B = W_m + W_p$, there must be no resultant force in vertical direction \end{soln}


\begin{marginfigure}
	\vspace*{0pt}
	\centering
	\begin{tikzpicture}[scale=1]
	\draw[thick] (0,5) -- (0,0) node[below left]{$O$} -- (4,0);
	\draw[ultra thick, gray] (0,4) node[left,black]{$A$} -- (3,0) node[below,black]{$B$};
	\draw[<->] (-1,0) -- (-1,4) node[above,midway,rotate=90]{$y=80\text{ cm}$};
	\draw[<->] (0,-1) -- (3,-1) node[below,midway]{$x=60\text{ cm}$};
	\draw[thick,blue,->] (0,4) --++ (1.2,0) node[above]{$R_A$};
	\draw[thick,blue,->] (3,0) --++ (0,1.5) node[above]{$R_B$};
	\draw[thick,blue,->] (3,0.05) --++ (-1.2,0) node[above]{$f_B$};
	\draw[thick,blue,->] (1.5,2) --++ (0,-1.5) node[left]{$W$};
	\end{tikzpicture}
	\vspace*{-16pt}
\end{marginfigure}


\example{A uniform ladder of weight 120 N rests on a rough ground against a smooth wall as shown. The dimensions are labelled on the diagram. (a) What is the contact force acting at $B$? (b) What is the contact force acting at $A$? (c) What is the frictional force at $B$?}


\begin{soln} Free-body diagram is drawn as shown

resolve vertically: $R_B = W \RA R_B = 120 \text{ N}$

take moments about $B$: $R_A y = W \frac{x}{2}$

{
	\centering
	
	$ R_A = \frac{120 \times 0.30}{0.80} = 45 \text{ N}$
	
}

resolve horizontally: $f_B = R_A \RA f_B = 45 \text{ N}$ 
\end{soln}
\begin{marginfigure}
	\vspace*{-12pt}
	\centering
	\begin{tikzpicture}[scale=0.8]
	\draw (0,-1.5) -- (0,4.5);
	\draw[ultra thick] (0,0) node[left]{$A$} -- (6,0) node[right]{$B$};
	\draw (6,0) -- (0,3.5) node[left]{$C$};
	\draw (5.4,0) arc(180:150:0.6);
	\node[left] at (5.3,0.3) {$30^\circ$};
	\draw (6,0) -- (6,-2);
	\draw (5.6,-2.7) rectangle (6.4,-2);
	\draw[thick,blue,->] (6,0) --++ (0,-0.8)node[right]{$W_b$};
	\draw[thick,blue,->] (3,0) --++ (0,-1.2)node[right]{$W_r$};
	\draw[thick,blue,->] (6,0) --++ (150:2.8)node[above]{$T$};
	\draw[thick,blue,->] (0,0) --++ (12.1*.2,3*.2)node[above]{$F_A$};
	\end{tikzpicture}
	\vspace*{-20pt}
\end{marginfigure}


\example{The diagram shows a uniform rod $AB$ of weight 60 N being held horizontally to a vertical wall by means of a light string. The string is attached to the rod at $B$, where a basket of weight 40 N is suspended. The other end of the string is fixed on the wall at $C$. The angle between the string and the rod is $30^\circ$. (a) Find the tension in the string. (b) Find the force acting on the rod at point $A$.}

\begin{soln} Take moments about $A$: $TL\sin\theta = W_b L + W_r  \frac{1}{2}L$

{
	\centering
	
	$ T \sin30^\circ = 40 + 60\times\frac{1}{2} \RA T=140 \text{ N} $
	
}

resolve horizontally: $F_{A,x} = T\cos\theta \RA F_{A_x}=140\cos30^\circ \approx 121 \text{ N}$

resolve vertically: $F_{A,y}+T\sin\theta = W_r + W_b \RA F_{A,y} = 60+40 - 140\sin30^\circ = 30 \text{ N} $

force at $A$: $F_A = \sqrt{F_{A,x}^2 + F_{A,y}^2} \RA F_A = \sqrt{121^2 + 30^2} \approx 125 \text{ N}$ \end{soln}


\subsection{Two forces in equilibrium}

The problem of two balanced forces is trivial

\begin{marginfigure}
	\vspace{-15pt}
	\begin{center}
		\begin{tikzpicture}[scale=0.6]
		\draw[fill] (0,0) circle(0.1);
		\draw[thick,blue,->] (0,0) -- (2,0) node[above]{$F_1$};
		\draw[thick,blue,->] (0,0) -- (-2,0) node[above]{$F_2$};
		\end{tikzpicture}
	\end{center}
	\vspace{-10pt}
\end{marginfigure}

Suppose two forces $F_1$ and $F_2$ are acting on an object in equilibrium

\begin{compactenum}
	\item[--] to have zero resultant force, $F_1$ and $F_2$ must be equal but opposite
	
	\item[--] to have zero resultant moment, $F_1$ and $F_2$ must act along same line
	
	otherwise they produce torque of couple that causes turning effects
\end{compactenum}


\subsection{Three forces in equilibrium}

\subsection{Force triangle}

When there are more than two forces, the situation becomes more complicated and we use \emph{vector diagrams} to solve the problems.

Suppose a set of forces $\overrightarrow{F_1}$, $\overrightarrow{F_2}$, $\cdots$, $\overrightarrow{F_n}$ are in equilibrium. If that's the case then no resultant force requires $\overrightarrow{F_1} + \overrightarrow{F_2} + \cdots + \overrightarrow{F_n} = 0$

Recall that resultant force is vector sum of all forces acting, and now this sum has to vanish, so if the force vectors are connected head to tail, they should form a closed $n$-polygon:

\begin{figure}[ht]
	\centering
	\begin{tikzpicture}[scale=0.7]
	\draw[fill] (0,0) circle(0.1);
	\draw[thick,blue,->] (0,0) -- (3,0) node[above]{$\overrightarrow{F_1}$};
	\draw[thick,blue,->] (0,0) -- (1,2) node[above]{$\overrightarrow{F_2}$};
	\draw[thick,blue,->] (0,0) -- (-2,2) node[above]{$\overrightarrow{F_3}$};
	\draw[thick,blue,->] (0,0) -- (-3,-1) node[above]{$\overrightarrow{F_4}$};
	\draw[thick,blue,->] (0,0) -- (1,-3) node[right]{$\overrightarrow{F_5}$};
	
	\node at (6,0) {\Huge{$\RA$}};
	
	\draw[thick,blue,->] (10,-2) -- ++(3,0) node[midway,below]{$\overrightarrow{F_1}$};
	\draw[thick,blue,->] (13,-2)-- ++(1,2) node[midway,right]{$\overrightarrow{F_2}$};
	\draw[thick,blue,->] (14,0)-- ++(-2,2) node[midway,above right]{$\overrightarrow{F_3}$};
	\draw[thick,blue,->] (12,2)-- ++(-3,-1) node[midway,above]{$\overrightarrow{F_4}$};
	\draw[thick,blue,->] (9,1) -- ++(1,-3) node[midway,left]{$\overrightarrow{F_5}$};
	\end{tikzpicture}
	\caption{an $n$-polygon formed by a set of $n$ balanced forces}\label{fig:force-polygon}
\end{figure}

In the case of three balanced forces, net force is zero means they should form a \keypoint{force triangle}. Any unknown forces can then be solved by cracking a geometric problem.

\begin{marginfigure}
	\vspace*{-16pt}
	\centering
	\begin{tikzpicture}[scale=0.7]
	\draw (-1.5,0) rectangle (1.5,-2);
	\draw[->,thick,blue] (0,-1) -- (0,-3) node[right]{$W$};
	\draw[->,thick,blue] (-0.5,0) -- ++(150:3) node[above]{$T_1$};
	\draw (-1,0) arc(180:150:0.5);
	\node[left] at (-1.1,0.25) {\scriptsize$\theta=30^\circ$};
	\draw (1,0) arc(0:30:0.5);
	\node[right] at (1.1,0.25) {\scriptsize$\theta=30^\circ$};
	\draw[->,thick,blue] (0.5,0) -- ++(30:3) node[above]{$T_2$};
	\node at (0,-4) {\huge{$\Downarrow$}};
	\draw[->,,thick,blue] (-1.0825,-5) -- ++(0,-2.5) node[midway,left]{$W$};
	\draw[->,,thick,blue] (-1.0825,-7.5) -- ++(30:2.5) node[midway,below]{$T_2$};
	\draw[<-,,thick,blue] (-1.0825,-5) -- ++(-30:2.5) node[midway,above]{$T_1$};
	\end{tikzpicture}
	\vspace*{-10pt}
\end{marginfigure}

\example{A painting of weight $W=20\text{ N}$ is supported by two strings as shown. Both strings form an angle $\theta=30^\circ$ to the horizontal. Find the tension in the strings.}



\begin{soln} By resolving forces, we have: 

{
	\centering
	
	$\Biggl\{\begin{array}{l}
	T_1 \cos\theta = T_2 \cos\theta\\
	T_1\sin\theta + T_2\sin\theta = W
	\end{array} $ 
	
}

we solve the equations to obtain:

{
	\centering
	
	$ T_1 = T_2  = \frac{W}{2\sin\theta} = \frac{20}{2\sin30^\circ}=20\text{ N} $
	
}


alternatively, we can construct the force triangle as shown

$T_1$, $T_2$ and $W$ form an equilateral triangle, so
\begin{equation*}
	T_1=T_2=W=20\text{ N} 
\end{equation*}
\end{soln}
\begin{marginfigure}
	\vspace*{-12pt}
	\centering
	\begin{tikzpicture}[scale=0.7]
	\draw (-1.5,0) rectangle (1.5,-2);
	\draw[->,thick,blue] (0,-1) -- (0,-3) node[right]{$W$};
	\draw[->,thick,blue] (-0.5,0) -- ++(150:3) node[above]{$T_1$};
	\draw (-1,0) arc(180:150:0.5);
	\node[left] at (-1.1,0.25) {\scriptsize$\theta_1=30^\circ$};
	\draw (1,0) arc(0:45:0.5);
	\node[right] at (0.9,0.32) {\scriptsize$\theta_2=45^\circ$};
	\draw[->,thick,blue] (0.5,0) -- ++(45:3) node[above]{$T_2$};
	\node at (0,-4) {\huge{$\Downarrow$}};
	\draw[->,,thick,blue] (-.95,-5) -- ++(0,-3) node[midway,left]{$W$};
	\draw[->,,thick,blue] (-0.95,-8) -- ++(45:2.69) node[midway,below]{$T_2$};
	\draw[<-,,thick,blue] (-.95,-5) -- ++(-30:2.20) node[midway,above]{$T_1$};
	\draw (-0.95,-5.5) arc(-90:-30:0.5);
	\draw (-0.95,-7.5) arc(90:45:0.5);
	\node at (-0.55,-5.65) {{\scriptsize $60^\circ$}};
	\node at (-0.6,-7.3) {{\scriptsize $45^\circ$}};
	\end{tikzpicture}
	\vspace*{-18pt}
\end{marginfigure}

\example{The same painting of weight $W=20\text{ N}$ is supported by two strings at different angles $\theta_1=30^\circ$ and $\theta_2=45^\circ$ as shown. Find the forces in the two strings.}


\begin{soln} By resolving forces, we have: 

{
	\centering
	
	$ \left\{\begin{array}{l}
		T_1 \cos\theta_1 = T_2 \cos\theta_2\\
		T_1\sin\theta_1 + T_2\sin\theta_2 = W
	\end{array} \right.
	\RA
	\left\{\begin{array}{l}
	\frac{\sqrt{3}}{2}T_1 = \frac{\sqrt{2}}{2}T_2\\
	\frac{1}{2}T_1 + \frac{\sqrt{2}}{2}T_2 = 20
	\end{array} \right.
	$
	
}


solving this, we find: $T_1 \approx 14.6\text{ N}\, , \, T_2 \approx 17.9\text{ N} $

alternatively, we construct the force triangle as shown

$T_1$ and $T_2$  are related to $W$ by the the law of sine:

{
	\centering
	
	$\frac{W}{\sin75^\circ} = \frac{T_1}{\sin45^\circ} = \frac{T_2}{\sin60^\circ} $
	
}

from this we get the same result: $T_1 \approx 14.6\text{ N}\, , \, T_2 \approx 17.9\text{ N} $ \end{soln}



\subsection{Concurrent forces}

Tor three forces in equilibrium, they must produce zero resultant moment about any point.

\begin{marginfigure}
	\vspace*{-15pt}
	\centering
	\begin{tikzpicture}[xscale=0.4, yscale=0.6]
	\draw[fill=gray!20] (-3.2,0.4) to [out=25,in=170] (-1,0.7) to [out=-10,in=150] (2.7,0.9) to [out=-30,in=35] (2.5,-1.3) to [out=215,in=30] (-0.5,-1.3) to [out=210,in=205] (-3.2,0.4);
	\draw[thick, blue, ->] (1,0) --++ (4,0) node[above]{$F_1$};
	\draw[thick, blue, ->] (240:1) --++ (240:2) node[left]{$F_3$};
	\draw[thick, blue, ->] (150:0.6) --++ (150:3.46) node[left]{$F_2$};
	\draw[dashed] (0,0) -- (1,0);
	\draw[dashed] (0,0) -- (240:1);
	\draw[dashed] (0,0) -- (150:0.6);
	\draw[fill] (0,0) circle(0.09 and 0.06);
	\node at (-60:0.5) {$P$};
	\end{tikzpicture}
	\vspace*{-15pt}
\end{marginfigure}

Suppose lines of action of $F_1$ and $F_2$ meet at point $P$, the moment of $F_1$ and moment of $F_2$ about $P$ are both zero. To produce zero resultant moment about $P$, then moment of $F_3$ about $P$ must vanish which suggests that the line of action of $F_3$ must pass through point $P$.
The lines of action of $F_1$, $F_2$ and $F_3$ therefore, must pass through the same point\footnote{In the case of three parallel forces in equilibrium, we can introduce the notion of an ideal point at infinity, so that parallel lines could meet at that point(!).}

Such three forces are said to be \emph{concurrent}

\subsection*{Summary for three forces in equilibrium}

For three forces in equilibrium, we can now conclude:

\begin{compactenum}
	\item[--] the three force vectors must be able to form a force triangle
	
	this is a consequence of zero resultant force
	
	\item[--] the lines of action for the three forces must pass through same point
	
	this is a consequence of zero resultant moment
\end{compactenum}






\subsection{End-of-chapter questions}


\subsection*{Mechanical equilibrium}

\question{Is it possible for an object to be in equilibrium if only one force is acting on it?}

\question{If three forces are in equilibrium, suggest and explain whether the lines of action must lie in the same plane.}

\include{momentum}
\include{work-energy}
\chapter{Materials}

In this chapter, we study how materials behave under a range of conditions. An important notion is the elasticity of materials, for \keypoint{elastic} materials, when external force is removed, the material returns to its original shape. If the material cannot restore to original shape, it is said to be \keypoint{inelastic}, or \keypoint{plastic}

\section{Springs}

\subsection{Hooke's law}

When a force $F$ is applied to a spring, it is stretched from original length $L_0$ to some length $L$. This extension of the spring, $x = L - L_0$, is proportional to the force applied.

\begin{ilight}
	Extension of an ideal spring is directly proportional to the load applied (within a certain range), this is called \keypoint{Hooke's law}\index{Hooke's law}.
\end{ilight}

\begin{figure}[ht]
	\centering
	\begin{tikzpicture}[xscale=1.5]
	\draw[thick,decorate, decoration={coil,amplitude=4pt, segment length=5pt}] (0,0) -- (3,0);
	\draw[fill] (-0.1,-0.5) rectangle (0,0.5);
	\draw[dashed] (4.2,-2.5) -- (4.2,-3);
	\draw[thick,decorate, decoration={coil,amplitude=4pt, segment length=6pt}] (0,-1.5) -- (3.6,-1.5);
	\draw[fill] (-0.1,-2) rectangle (0,-1);
	\draw[very thick,blue,->] (3.6,-1.5) -- (4.4,-1.5) node[right]{$F$};
	\draw[thick,decorate, decoration={coil,amplitude=4pt, segment length=7pt}] (0,-3) -- (4.2,-3);
	\draw[fill] (-0.1,-3.5) rectangle (0,-2.5);
	\draw[very thick,blue,->] (4.2,-3) -- (5.8,-3) node[right]{$2F$};
	\draw[dashed] (0,0) -- (0,-3);
	\draw[dashed] (3,0.5) -- (3,-3.5);
	\draw[dashed] (3.6,-1.5) -- (3.6,-.5);
	\draw[<->] (0,-.8) -- (3,-.8) node[midway,above]{$L_0$};
	\draw[<->,red] (3,-.8) -- (3.6,-.8) node[midway,above]{$x$};
	\draw[dashed] (4.2,-2) -- (4.2,-3);
	\draw[<->] (0,-2.3) -- (3,-2.3) node[midway,above]{$L_0$};
	\draw[<->,red] (3,-2.3) -- (4.2,-2.3) node[midway,above]{$2x$};
	\end{tikzpicture}
\end{figure}

Hooke's law can be summarised by the equation: \begin{empheq}[box=\tcbhighmath]{equation*}{F=kx}\end{empheq}

The proportionality constant $k$ is called the \keypoint{spring constant}, a larger $k$ means a greater force is required to extend the spring by same amount. A spring with a large $k$ is said to be \emph{stiff}. 
Note that we usually fall into the trap of thinking about materials under tension, however, Hooke's law also holds if spring is being compressed - i.e., if a spring is pushed to have compression of $x$, we still have force applied $F=kx$.

The linear relationship between $F$ and $x$ is only true up to the limit of the material capability. The point beyond which Hooke's law no longer holds is called the \keypoint{limit of proportionality}. Beyond this point it will still stretch, (and return to the origin if unloaded) but not in a linear fashion.
If the force on the spring is greater still then, spring might be overstretched and no longer exhibits elastic behaviour - it it may become permanently deformed. The point beyond which spring cannot return to original length is called the \keypoint{elastic limit} and the deformation called \keypoint{permanent set.}
\footnote{The elastic limit of a spring and the limit of proportionality are two different but always confused concepts. These two limits are usually very close to one another, but they are conceptually different.}

\begin{figure*}[ht]
	\centering
	\begin{tikzpicture}[yscale=0.96]
	\draw[<->] (9.6,0) node[below]{$x$} -- (0,0) -- (0,6.4) node[left]{$F$};
	\draw[thick] (0,0) -- (3,4) [out=53.13,in=185] to (9,6); \fill (3,4) circle(0.08); \fill (4.2,5.1) circle(0.08);
	\draw[<-] (1.8,2) -- ++(2.5,-0.8) node[note]{linear region for which \\ Hooke's law $F=kx$ holds};
	\draw[<-] (3.2,3.8) -- ++(1,-1) node[note]{limit of \\porportionality};
	\draw[<-] (3.9,5.2) -- ++(-1.5,.4) node[note]{elastic limit};
	\draw[<-] (6.4,5.5) -- ++(.8,-1.2) node[note]{non-linear region\\ where spring\\behaves plastically};
	\end{tikzpicture}
	\caption{force-extension graph for a typical spring under load}
\end{figure*}

\example{A spring has a natural length of 20.0 cm. When a mass of 250 g is suspended from the spring, the new length of the spring is 26.0 cm. Find the spring constant. }

\begin{soln}\begin{equation*}
	k = \frac{F}{x} = \frac{mg}{L-L_0} = \frac{0.250\times9.81}{(26.0-20.0)\times10^{-2}} \RA k \approx 40.9 \text{ N m}^{-1} 
\end{equation*}
\end{soln}

\example{A spring has a spring constant of $270 \text{ N m}^{-1}$. A mass of 1.2 kg is hung from the spring. When the mass is released from a position where the spring has an extension of 5.0 cm, what is the acceleration of the mass?}

\begin{soln}\begin{equation*}
	\fnet = F - mg = kx - mg = ma \RA 
\end{equation*}
\begin{equation*}
 a = \frac{kx-mg}{m} = \frac{270\times5.0\times10^{-2} - 1.2 \times 9.81}{1.2} \approx 1.44 \mpss 
\end{equation*}
\end{soln}

\subsection{Elastic potential energy in a spring}\label{ch-spring-energy}

We stretch a spring by exherting a force over a distance. It follows then that to stretch or compress a spring, work must be done. The work done stored is called \emph{elastic potential energy}. \index{energy!elastic potential}


\begin{marginfigure}
	\vspace{10pt}
	\centering
		\begin{tikzpicture}[scale=1]
		\draw[gray!30,fill] (0,0) -- (3.2,0) -- (3.2,2.4) -- cycle;
		\draw[<->] (4,0) node[below]{$x$} -- (0,0) -- (0,3) node[left]{$F$};
		\draw (0,0) -- (3.6,2.7) node[above]{$F=kx$};
		\draw[dashed] (0,2.4) -- (3.2,2.4) -- (3.2,0);
		\node at (2.25,0.8) {$W=\frac{1}{2}Fx$};
		\end{tikzpicture}
\end{marginfigure}

Note that force in spring varies as spring is stretched so to find work in stretching a spring by $x$, we cannot just multiply it by $F$ as one would do with a constant force doing work. Instead we compute the area under $F$-$x$ graph\footnote{Mathematically, we can also integrate over the total extension to find this work done $W=\int_0^x F \dd x = \int_0^x k x \dd x = \frac{1}{2}kx^2$, which of course gives the same result.}. The graph is a right-angled triangle so:
$$ W = \frac{1}{2}Fx = \frac{1}{2}kx^2 $$

The formula for the elastic potential energy stored in a spring is:

{
	\begin{empheq}[box=\tcbhighmath]{equation*}{E_p=\frac{1}{2}kx^2} \end{empheq}
	
}

\example{A steel spring has a spring constant of $20 \text{ N cm}^{-1}$. How much work is needed to stretch it from an extension of 3.0 cm to an extension of 5.0 cm?}

\begin{soln} Work done needed equals the increase in elastic potential energy:
\begin{equation*}
	W = \Delta E_p = \frac{1}{2}kx_2^2 - \frac{1}{2}kx_1^2 = \frac{1}{2} \times 2000 \times (0.050^2 - 0.030^2) = 1.6 \text{ J} 
\end{equation*}
\end{soln}

\example{A trolley of 400 g can travel freely along a horizontal surface. It is pushed against a spring buffer. Suppose the spring is initially compressed by 5.0 cm under a 20 N force. When the trolley is released, it accelerates until it becomes detached. What is the trolley's final speed?}

\begin{soln} Elastic potential energy in spring $\RA$ kinetic energy of trolley
\begin{equation*}
\frac{1}{2}Fx = \frac{1}{2}mv^2 \RA \frac{1}{2}\times20\times0.050 = \frac{1}{2}\times0.40\times v^2 \RA v\approx 1.58\mps 
\end{equation*}
\end{soln}

\example{The same spring is now set between two trolleys $A$ and $B$ of mass 400 g and 600 g. Initially the spring is again compressed by 5.0 cm under a force of 20 N. After both trolleys are released, what are their final speeds?}

\begin{figure}[ht]
	\centering
	\begin{tikzpicture}[scale=1.35]
	\draw[thick,decorate, decoration={coil,amplitude=4pt, segment length=5pt}] (-.8,.4) -- (.8,.4);
	\draw[thick] (-2.4,0) rectangle (-.8,.8);
	\draw[thick] (.8,0) rectangle (2.4,.8);
	\foreach \c in {-2,-1.2,1.2,2}
	{  \draw[thick] (\c,-0.2) circle(0.2);
		\draw[thick] (\c,-0.2) circle(0.15);  }
	\draw[thick] (-4,-.4) -- (4,-.4);
	\node at (1.6,0.4) {$B$}; \node at (-1.6,0.4) {$A$};
	\end{tikzpicture}
\end{figure}

\begin{soln} Elastic potential energy in spring transforms into kinetic energy of the two trolleys:
\begin{equation*}
\frac{1}{2}Fx=\frac{1}{2}m_Av_A^2 + \frac{1}{2}m_Bv_B^2 \RA 20\times0.050 = 0.40v_A^2 + 0.60v_B^2
\end{equation*}

Total momentum for trolley $A$ and $B$ as a whole is conserved:
\begin{equation*}
m_Bv_B - m_Av_A = 0 \RA 0.40v_A = 0.60v_B
\end{equation*}

solving the simultaneous equations, we find: $v_A \approx 1.22\mps, v_B \approx 0.82 \mps$ \end{soln}

\subsection{Spring combinations}

So far we have discussed the properties of a single spring, next we investigate how a set of springs respond to a given load.

\subsection*{Parallel springs}

\begin{marginfigure}
	\vspace{-15pt}
	\centering
	\begin{tikzpicture}[scale=1]
		\draw[thick,decorate, decoration={coil,amplitude=4pt, segment length=5pt}] (0,-0.5) -- (3,-0.5) node[midway,yshift=10pt]{$k_2$};
		\draw[thick,decorate, decoration={coil,amplitude=4pt, segment length=5pt}] (0,0.5) -- (3,0.5) node[midway,yshift=10pt]{$k_1$};
		\draw[fill] (-0.2,-0.8) rectangle (0,0.8);
		\draw[fill=black!50] (3,-0.8) rectangle (3.05,0.8);
		
		\draw[thick,decorate, decoration={coil,amplitude=4pt, segment length=6pt}] (0,-3) -- (3.6,-3) node[midway,yshift=10pt]{$k_2$};
		\draw[thick,decorate, decoration={coil,amplitude=4pt, segment length=6pt}] (0,-2) -- (3.6,-2) node[midway,yshift=10pt]{$k_1$};
		\draw[fill] (-0.2,-3.3) rectangle (0,-1.7);
		\draw[fill=black!50] (3.6,-3.3) rectangle (3.65,-1.7);
		\draw[very thick,blue,->] (3.6,-2.5) -- (5,-2.5) node[right]{$F$};
		\draw[dashed] (3.025,-0.8) -- (3.025,-3.5);
		\draw[dashed] (3.625,-1.7) -- (3.625,-1);
		\draw[<->] (3.025,-1.4) -- (3.625,-1.4) node[midway,above]{$x$};
	\end{tikzpicture}
%	\caption{two springs in parallel}
	\vspace{-10pt}
\end{marginfigure}

Let's take two springs connected in parallel, when the combination is stretched under a load of $F$, extension in each spring should be the same:

{
	\centering
	
	$x_1=x_2=x$
	
}

the forces in each spring will in general be different, but sum of these must be equal to load: 

{
	\centering
	
	$ F=F_1+F_2 $
	
}


Divide both sides by $x$, we have: $\frac{F}{x} = \frac{F_1}{x_1} + \frac{F_2}{x_2}$.

Recall the Hooke's law, this becomes: $k=k_1+k_2$

Generalize for $n$ springs in parallel connection, the combined spring constant is
\begin{empheq}[box=\tcbhighmath]{equation*}{k=k_1+k_2+\cdots+k_n}
\end{empheq}

\subsection*{Series springs}

Let's now take two springs in series

\begin{figure}[ht]
	\centering
	\begin{tikzpicture}[scale=1]
	\draw[thick,decorate, decoration={coil,amplitude=4pt, segment length=5pt}] (0,0) -- (2,0) node[midway,yshift=10pt]{$k_1$};
	\draw[thick,decorate, decoration={coil,amplitude=4pt, segment length=5pt}] (2.5,0) -- (4.5,0) node[midway,yshift=10pt]{$k_2$};
	\draw (2,0) -- (2.5,0);
	\draw[fill] (-0.2,-0.5) rectangle (0,0.5);
	\draw[fill=black!50] (4.5,-0.5) rectangle (4.55,0.5);
	
	\draw[thick,decorate, decoration={coil,amplitude=4pt, segment length=5pt}] (0,-1.5) -- (2.5,-1.5) node[midway,yshift=10pt]{$k_1$};
	\draw[thick,decorate, decoration={coil,amplitude=4pt, segment length=5pt}] (3,-1.5) -- (5.5,-1.5) node[midway,yshift=10pt]{$k_2$};
	\draw (2.5,-1.5) -- (3,-1.5);
	\draw[fill] (-0.2,-2) rectangle (0,-1);
	\draw[fill=black!50] (5.5,-2) rectangle (5.55,-1);
	\draw[very thick,blue,->] (5.5,-1.5) -- (7,-1.5) node[right]{$F$};
	\draw[dashed] (4.525,-0.5) -- (4.525,-1);
	\draw[dashed] (5.525,-0.5) -- (5.525,-1);
	\draw[<->] (4.525,-0.8) -- (5.525,-0.8) node[midway,above]{$x$};
	\end{tikzpicture}
%	\caption{two springs in series}
\end{figure}


The force in each spring is the same: $F_1=F_2=F$

But total extension is the sum of individual extensions: $x=x_1+x_2$.

Divide by the same $F$, we have: $\frac{x}{F} = \frac{x_1}{F_1} + \frac{x_2}{F_2}$.

For combined spring constant, we find: $\frac{1}{k} = \frac{1}{k_1} + \frac{1}{k_2}$

For $n$ springs connected in series, the combined spring constant is therefore given by
\begin{empheq}[box=\tcbhighmath]{equation*}
{\frac{1}{k} = \frac{1}{k_1} + \frac{1}{k_2} + \cdots \frac{1}{k_n}} 
\end{empheq}

If we set up $n$ springs in parallel, we actually make it thicker, it then requires a stronger force to stretch it by the same extension so we indeed see the combined spring constant is greater than that of any individuals. If we set up $n$ springs in series, we make it longer, instead of pulling one spring at a time, the same force now stretches $n$ springs simultaneously, this gives rise to a greater total extension so the combined spring constant must be less than that of any individual.

\example{(a) A spring with $k_1 = 20 \text{ N cm}^{-1}$ is connected in series with a second spring with $k_2 = 30 \text{ N cm}^{-1}$. When a force of 60 N is applied, what is the total extension of the combination? (b) The same two springs are now connected in parallel. When a force of 50 N is applied on the combination, what is the extension?}

\begin{soln} For series connection: $k = \left(\frac{1}{k_1} + \frac{1}{k_1}\right)^{-1} = \left( \frac{1}{20} + \frac{1}{30} \right)^{-1} = 12 \text{ N cm}^{-1} \RA x = \frac{F}{k} = \frac{60}{12} = 5.0 \text{ cm}$

\eqskip or sum up extension of each spring: $x = x_1 + x_2 = \frac{F}{k_1} + \frac{F}{k_2} = \frac{60}{20} + \frac{60}{30} = 5.0 \text{ cm}$

for parallel connection: $k = k_1 + k_2 = 20 + 30 = 50 \text{ N cm}^{-1} \RA x = \frac{F}{k} = \frac{50}{50} = 1.0 \text{ cm} $ \eoe
\end{soln}
\begin{marginfigure}
	\vspace*{-5pt}
	\centering
	\begin{tikzpicture}
	\draw[thick] (-1.5,0) -- (1.5,0);
	\foreach \x in {-1,0,1}
	\draw[thick,decorate, decoration={coil,amplitude=4pt, segment length=5pt}] (\x,0) --++ (0,-1.5);
	\draw[thick] (-1.5,-1.5) -- (1.5,-1.5);
	\foreach \x in {-0.5,0.5}
	\draw[thick,decorate, decoration={coil,amplitude=4pt, segment length=5pt}] (\x,-1.5) --++ (0,-1.5);
	\draw[thick] (-1,-3) -- (1,-3);
	\draw[thick,decorate, decoration={coil,amplitude=4pt, segment length=5pt}] (0,-3) --++ (0,-1.5);
	\draw[thick] (-0.5,-4.5) -- (0.5,-4.5);
	\draw[thick] (0,-4.5) --++ (0,-0.5);
	\draw[thick] (-0.5,-6) rectangle (0.5,-5);
	\node at (0,-5.5) {$6.0 \text{ N}$};
	\end{tikzpicture}
	\vspace*{-16pt}
\end{marginfigure}


\example{A set of identical springs are set up as shown. Each individual spring extends by 1.0 cm under a load of 1.0 N. Assume the limit of proportionality is not exceeded, what is the total extension for this combination when a load of 6.0 N is applied?}

\begin{soln} For a single spring: $k = 1.0 \text{ N cm}^{-1}$, the combined spring constant: $k_\text{total} = \left( \frac{1}{3k} + \frac{1}{2k} + \frac{1}{k} \right)^{-1}$

{
	\centering
	
	\eqyskip $ k_\text{total} = \left(\frac{1}{3.0} + \frac{1}{2.0} + \frac{1}{1.0} \right)^{-1} = \frac{6}{11} \approx 0.545 \text{ N cm}^{-1}$
	
}

Total extension: $x = \frac{F}{k_\text{total}} = \frac{6.0}{0.545} = 11 \text{ cm}$

Alternatively, we can find and add extensions of each layer. In particular, the top layer withstands a force of 6.0 N shared by three springs, so each spring has a force of 2.0 N, extension is 2.0 cm. Similarly, the other two layers extend by 3.0 cm and 6.0 cm respectively, hence, total extension: $x = 2.0 + 3.0 + 6.0 = 11 \text{ cm}$ \end{soln}

\subsection{Stress, strain \& Young modulus}

From daily experience, it is easier to stretch a longer wire than a shorter one, the same tensile force also produce greater effects on a thinner material than on a thicker one. To accurately compare materials, we need better quantities to describe what's going on. To study a material's response to a tensile force, several new quantities are to be introduced:

\subsection{Stress \& Strain}

\begin{ilight}
	\centering \keypoint{Tensile strain}\index{strain} is defined as the ratio of the extension of a wire to its natural length: \begin{empheq}[box=\tcbhighmath]{equation*}{\epsilon = \frac{x}{L}}	\end{empheq}
\end{ilight}

\begin{ilight}
	\centering \keypoint{Tensile stress}\index{stress} is defined as the force applied per unit cross-sectional area: 
	    \begin{empheq}[box=\tcbhighmath]{equation*}{\sigma = \frac{F}{A}}\end{empheq}
\end{ilight}

Strain is the ratio of two lengths, so it is unit free. It gives us a measure comparable to extension that allows us to compare materials of different length.

Stress is comprable to pressure - by considering the force per unit area we have a measure that allows us to compare materials of different (uniform) cross-section. The units of stress, like pressure: $[\sigma] = \text{N m}^{-2} = \text{Pa}$ (pascal)

The elastic behaviour of any material, though the physical mechanism might be different, can be accurately modelled as a spring. Hooke's law says $F \propto x$, it then follows that $\frac{F}{A} \propto \frac{x}{L}$, so $\sigma \propto \epsilon$

i.e., stress and strain should also be proportional to each other within certain limit

\example{A copper wire has a cross-sectional area of $1.5 \times 10^{-6} \text{ m}^2$. The breaking stress of the wire is $2.0 \times 10^8 \text{ Pa}$. Find the breaking force.}

\begin{soln}
    
\begin{equation*}
\sigma = \frac{F}{A} \RA F = \sigma A = 2.0 \times 10^8 \times 1.5 \times 10^{-6} = 3.0 \times 10^2 \text{ N} 
\end{equation*}
\end{soln}

\subsection{Young modulus}

\begin{ilight}
	\centering The ratio of stress to strain of a material is called the \keypoint{Young modulus}\index{Young modulus}:\begin{equation*} 
 \tcbhighmath{E = \frac{\sigma}{\epsilon}} \end{equation*}
\end{ilight}


Because one of the parts of this is dimensionless, and the other pascals, the unit of measurementis: $[E] = \text{Pa}$ (pascal)

Young modulus is a property of the material, for the same material, Young modulus is a constant, no matter in what shape it takes. i.e., it does not depend on the length or the cross section of the object.\\
Young modulus values can look ``wrong'' in answers to problems. This is because typical value of Young's modulus for metals: $E_\text{metal} \sim 10^{11} \text{ Pa}$

Young modulus is a measure of the stiffness of a material - to produce same strain, greater stress is required for a material with greater Young modulus. You can get back to the spring constant from $E = \frac{\sigma}{\epsilon} = \frac{\sfrac{F}{A}}{\sfrac{x}{L}} = \frac{FL}{xA}$, we rearrange to get: $F = \frac{EA}{L}x$\\

Compare this with Hooke's law, we can identify the force constant to be given by: $k = \frac{EA}{L} $

\titem $E \up \ra k \up$, stiffer material makes stiffer springs

\titem $A \up \ra k \up$, more difficult to stretch a thick spring (think about parallel springs)

\titem $L \up \ra k \down$, easier to stretch a long spring (think about series springs)
\newpage

\begin{figure*}[ht]
    \begin{center}
        \begin{tikzpicture}[scale=1]
            \draw[->] (-0.5,0) -- (10,0) node[anchor=north] {$\epsilon$};
            \draw[->] (0,-0.5) -- (0,6) node[anchor=east] {$\sigma$};
            \draw (0,0) -- (1.5,3) .. controls (2,3.5) .. (2.5,3) .. controls (3,2.8) .. (4,4) .. controls (5.5,5.5) and (8,5.5) .. (10,4.5);
            \fill (1.5,3) circle (2pt) node[anchor=south east] {A};
            \fill (2,3.4) circle (2pt) node[anchor=south] {B};
            \fill (7,5.2) circle (2pt) node[anchor=south] {C};
            \fill (10,4.5) circle (2pt) node[anchor=south] {D};
            \node[draw, text width = 4cm] at (13,3) {
            \begin{description}
                \item[A] The limit of proportionality
                \item[B] The yield point
                \item[C] The ultimate stress (maximium stress)
                \item[D] The breaking point
            \end{description}
            };
        \end{tikzpicture}
    \end{center}
    \caption{Stress-strain curve for a ductile material}
    \label{stress-strain}
\end{figure*}

Figure \ref{stress-strain} shows an example stress-strain curve. Note that the limit of proportionality is often a good approximation of the elastic limit of a metal. The ``breaking stress'' usually refers to the ultimate stress, i.e. the maximum stress the material can withstand, rather than the stress at the breaking point.

\example{A 200 N tensile force is applied on a steel wire of 1.5 m, the wire extends by 5.0 mm. The diameter of the cross section is 0.60 mm. What is the Young modulus of the steel wire?}

\begin{soln}\begin{equation*}
	E = \frac{\sfrac{F}{A}}{\sfrac{x}{L}} = \frac{FL}{Ax} = \frac{200 \times 1.5}{\pi \times (0.30\times10^{-3})^2 \times 5.0 \times 10^{-3}} \RA E \approx 2.1 \times 10^{11} \text{ Pa} 
\end{equation*}
\end{soln}
\example{A copper wire of length 2.0 m is under a stress of $7.8 \times 10^7 \text{ Pa}$. Given that the Young modulus of copper is $1.2 \times 10^{11} \text{ Pa}$, what is (a) the strain of the wire, (b) the extension of the wire?}

\begin{soln} strain: $\epsilon = \frac{\sigma}{E} = \frac{7.8 \times 10^7}{1.2 \times 10^{11}} \RA \epsilon = 6.5 \times 10^{-4} = 0.065 \%$

extension: $x = \epsilon L = 6.5 \times 10^{-4} \times 2.0 \RA x = 1.3 \times 10^{-3} \text{ m}$ \end{soln}

\example{Several blocks of steel are used to support a bridge. Each block has a height of 30 cm and a cross section of 15 cm $\times$ 15 cm. The steel block is designed to compress 2.0 mm when the maximum load is applied. Given that the Young modulus of steel is $2.1 \times 10^{11} \text{ Pa}$, what is the maximum load that can be supported by one block?}

\begin{soln} \begin{equation*}
	E = \frac{FL}{Ax} \RA F = \frac{EAx}{L} = \frac{2.1\times10^{11} \times (0.15\times0.15) \times2.0\times10^{-3}}{0.30} \end{equation*}
 
 \begin{equation*}
 \RA F \approx 3.15 \times 10^7 \text{ N} 
\end{equation*}
\end{soln}

\example{Two metal wires $A$ and $B$ are of the same length and they extend by the same amount under the same load. Given that Young modulus of wire $A$ is twice of $B$, what is the ratio of their diameters?}

\begin{soln}\begin{equation*}
	E = \frac{FL}{Ax} \RA A = \frac{1}{4}\pi d^2 = \frac{FL}{Ex} \RA 
 \end{equation*}
 \begin{equation*}
 d^2 \propto \frac{1}{E} \RA \frac{d_A}{d_B} = \sqrt{\frac{E_B}{{E_A}}} = \frac{1}{\sqrt{2}} 
\end{equation*}
\end{soln}

\example{A full-size crane is ten times greater than a model crane in all linear dimensions. If they are made of the same materials, what is the ratio of the cable's extension?}

\begin{soln} same material means same density $\rho$ and same Young's modulus $E$, so:
\begin{equation*}
E = \frac{FL}{Ax} \RA x = \frac{FL}{AE} = \frac{mgL}{\pi r^2 E} = \frac{\rho V g L}{\pi r^2 E} \RA
\end{equation*}
\begin{equation*}
x \propto \frac{VL}{r^2} \propto \frac{l^3 l}{l^2} \propto l^2 \RA \frac{x_\text{full}}{x_\text{model}} = 10^2 = 100 
\end{equation*}
\end{soln}
\newpage
\subsection{Measurement of Young modulus}\label{ch-measure-E}

The measurement of the Young's modulus of a metal wire is a required practical with a lot of detail. The experimental setup can be laid out as shown:


\begin{figure*}[ht]
	\centering
	\begin{tikzpicture}
	\draw[thick] (0,-0.1) -- (0,1.5);
	\foreach \y in {0.2,0.4,...,1.4} \draw[thick] (0,\y) --++ (-0.2,-0.2);
	\draw[thick] (-0.3,-0.1) -- (11.1,-0.1) -- (11.1,-3.5);
	\draw[thick,fill=white] (11.2,0) circle (0.78);
	\draw[very thick,gray] (0,0.8) -- (11.2,0.8) arc (90:0:0.8) -- (12,-1.5);
	\draw[thick,fill] (11.25,0) --++ (0.05,-0.05) --++ (-.9,-.9) --++ (-0.05,0.05) -- cycle; 
	\draw (5,0.9) --++ (1,1) node[above]{metal wire};
	\draw[Green,fill=blue] (7.9,0.9) -- (8.1,0.9) -- (8,0.6) -- cycle;
	\draw (8,1) --++ (1,1) node[above]{marker};
	\draw (11.4,0.6) --++ (1,1) node[above]{pulley};
	\draw[thick,fill=white] (7.5,0) rectangle (9,-0.2);
	\draw[thick,fill=white] (8,0.05) circle(0.16);
	\draw (8,-0.3) --++ (-1,-1) node[below,twoline]{travelling\\microscope};
	\foreach \y in {-2.9,-2.25,-1.6} {
		\draw[thick] (11.84,\y-0.04) arc (180:0:0.16) [out=-90, in=90] to (12, \y-0.3);
		\draw[thick,fill=white] (11.6,\y-0.6) rectangle (12.4,\y-0.3);
	}
	\draw (12.5,-2) --++ (.8,.8) node[above]{masses};
	\end{tikzpicture}
\end{figure*}

The original length $L$ (up to the marker) of the wire is measured with \emph{metre rule} $\pm 1$mm. The diameter $d$ of the wire is measured with \emph{micrometer} $\pm$0.01mm, then cross-sectional area is: $A = \frac{1}{4}\pi d^2$.\\

For a recorded mass $m$ attached to the wire, then force applied is $F=mg$\\

The extension $x$ of the wire is taken to be distance moved out by the marker. In the lab we typically use a meter rule but in some textbooks you'll see this can be measured with a \emph{travelling microscope}. A travelling microscope is basically a microscope that can move back and forth along a rail. The position of the microscope can be varied by turning a screw. This position can be read off a vernier scale. So in short, a travelling microscope can be used to measure the change in length with a very high resolution (typically to a precision of 0.01mm or 0.02 mm).\\
Once the data is collected, the stress can be calculated by $\sigma = \frac{F}{A}$, and strain can be calculated by $\epsilon = \frac{x}{L}$.\\

A graph of stress against strain can be plotted, a best fit line can be drawn and by taking the gradient of the straight-line section, the Young modulus can be obtained.

\newpage
\subsection{Stress-strain curves}

Stress-strain curve for a material can be obtained using the methods in \S\ref{ch-measure-E} 

\begin{figure}[ht]
	\centering
	\begin{tikzpicture}[scale=0.85]
	\draw[<->] (9.6,0) node[below]{strain} -- (0,0) -- (0,6.4) node[left]{stress};
	\draw[thick] (0,0) -- (3,4) [out=53.13,in=185] to (9,6); \fill (3,4) circle(0.08);
	\fill (4.2,5.1) circle(0.08);
	\fill (9,6) circle(0.08);
	\draw[<-] (1.8,2) -- ++(2.3,-0.9) node[note]{linear\\region};
	\draw[<-] (3.2,3.8) -- ++(1,-1) node[note]{$L_P$: limit of \\porportionality};
	\draw[<-] (3.9,5.2) -- ++(-1.5,.4) node[note]{$L_E$: limit\\ of elasticity};
	\draw[<-] (6.5,5.6) -- ++(.2,-1.5) node[note]{plastic flow};
	\draw[<-] (9,5.8) -- ++(.1,-1.2) node[note]{breaking\\point};
	\end{tikzpicture}
	\caption{stress-strain graph for a typical metal}
\end{figure}

Ductile and plastic materials show some plastic deformation before breaking. Strictly speaking, ductility is continued plastic deformation at lower forces.
Brittle materials on the other hand, break before any plastic deformation.


\begin{figure}[ht]
	\centering
	\begin{tikzpicture}[scale=0.85]
	\draw[<->] (6,0) node[below]{strain} -- (0,0) -- (0,6.4) node[left]{stress};
	\draw[thick] (0,0) -- (3,4);
    \draw[thick] (0,0) -- (2,4);
    \fill (3,4) circle(0.08);
    \fill (2,4) circle(0.08);
    \draw[<-] (2,4) -- ++(1,1) node[note]{stiffer\\material};
	\draw[<-] (1.8,2) -- ++(2.3,-0.9) node[note]{linear\\region};
	\draw[<-] (3.2,3.8) -- ++(1,-1) node[note]{breaking\\point};
		\end{tikzpicture}
	\caption{stress-strain graph for two brittle materials}
\end{figure}

For all materials, up to limit of proportionality $L_E$, stress is proportional to strain and the Young modulus can be calculated by the gradient of the line before $L_P$. In addition, consider the product of stress and strain: $\sigma \cdot \epsilon \sim \frac{F}{A} \frac{x}{L} \sim \frac{Fx}{AL} \sim \frac{W}{V}$. The area under stress-strain curve therefore, gives the work done per unit volume to stretch the wire. Up to limit of elasticity $L_E$, metal wire can return to original length when it is unloaded.

\begin{marginfigure}
	\vspace*{-15pt}
	\centering
	\begin{tikzpicture}[scale=0.62]
	\draw[fill=yellow!30] (0,0) -- (3,4) [out=53.13,in=190] to (7,5.7)  -- (2.725,0) -- cycle;
	\draw[fill=green!30] (7,5.7)  -- (2.725,0) -- (7,0) -- cycle;
	\draw[<->] (8,0) node[below]{strain} -- (0,0) -- (0,6.4) node[above]{stress};
	\draw[thick,blue,postaction={decorate},decoration={markings,mark=at position 0.45 with {\arrow{>}}}] (0,0) -- (3,4) [out=53.13,in=190] to (7,5.7);
	\draw[thick,blue,postaction={decorate},decoration={markings,mark=at position 0.6 with {\arrow{>}}}] (7,5.7) -- (2.725,0);
	\draw[thick] (7,5.7) [out=10,in=185] to (8,5.85);
	\fill (4.0,4.92) circle(0.08) node[above left] {$L_E$};
	\node[blue,rotate=53.13, above] at (2.4,3.2) {loading};
	\node[blue,rotate=53.13, below] at (4.6,2.5) {unloading};
	\node at (3.5,2.7) {{\Large $X$}};
	\node at (6,2) {{\Large $Y$}};
	\end{tikzpicture}
	\vspace*{-16pt}
\end{marginfigure}

If a wire is stretched beyond the elastic limit, it follows a different path when force is removed and displays permanent set:

\begin{figure}[ht]
	\centering
	\begin{tikzpicture}[scale=0.85]
	\draw[<->] (7,0) node[below]{strain} -- (0,0) -- (0,6.4) node[left]{stress};
	\draw[thick] (0,0) -- (3,4) [out=53.13,in=185] to (6.5,5.8);
    \draw[thick] (2,0) -- (6.5,5.8);
    \fill (6.5,5.8) circle(0.08);
 \fill (3,4) circle(0.08);
	\fill (4.2,5.1) circle(0.08);
	\draw[<-] (1.5,0) -- ++(3,0.9) node[note]{permanent\\set};
	\draw[<-] (3.2,3.8) -- ++(1,-1) node[note]{$L_P$: limit of \\porportionality};
	\draw[<-] (3.9,5.2) -- ++(-1.5,.4) node[note]{$L_E$: limit\\ of elasticity};
 \draw[<-] (5,5.6) -- ++(1.5,1) node[note]{undergoing plastic\\deformation};
 \draw[<-] (2,3) -- ++(-1.5,.4) node[note]{Loading};
	\draw[<-] (6.5,5.6) -- ++(.2,-1.5) node[note]{unloaded\\ here};
	\end{tikzpicture}
	\caption{stress-strain graph for a material displaying permanent set}
\end{figure}

As discussed above, the area under stress-strain graph relates to energy, so the work done in stretching wire is given by $X+Y$. The energy returned when wire contracts is $Y$. The difference $X$ gives energy loss as heat in the wire.


\subsection{Metals}
Metals consist of positive ions in a sea of delocalised electrons. During the elastic phase of deformation the spaces between the ions get larger and smaller. The metallic bonds resist this change from their equilibrium length and act like small springs acting to return the spacing to its original length.

An initial expectation of the plastic phase of deformation in a metallic lattice may be that the planes of ions slip past one another; however an analysis of the forces required for such movement gives an answer hundreds of times higher than the measured yield stress. Instead, the plastic deformation of metals must be explained in terms of \emph{dislocations}. A dislocation occurs when there is a gap in the metallic lattice. Dislocations occur naturally in materials and enable plastic deformation to occur through the breaking of individual bonds in succession, rather than all at once.

\begin{figure*}[h]\begin{center}
    \begin{tikzpicture}
        % Left hand
        \foreach \y in {5,6,7} {
            \foreach \x in {0,1,2,3,4} {
                \draw (\x,\y) circle (4mm);
            }
        }
        \foreach \x in {0.1,1.15,2.8,3.8} {
                \draw (\x,4) circle (4mm);
            }
        \foreach \x in {0.3,1.3,2.6,3.6} {
            \draw (\x,3) circle (4mm);
        }
        \foreach \x in {0.5,1.4,2.4,3.35} {
            \draw (\x,2) circle (4mm);
        }
        \draw[red, very thick] (0,7) -- (0.1,4) -- (0.5,2);
        \draw[red, very thick] (1,7) -- (1,5) -- (1.4,2);
        \draw[red, very thick] (2,7) -- (2,5);
        \draw[red, very thick] (3,7) -- (3,5) -- (2.4,2);
        \draw[red, very thick] (4,7) -- (4,5) -- (3.4,2);
         % right hand
        \foreach \y in {5,6,7} {
            \foreach \x in {7,8,9,10,11} {
                \draw (\x,\y) circle (4mm);
            }
        }
        \foreach \x in {7.1,8.15,9.2,10.8} {
                \draw (\x,4) circle (4mm);
            }
        \foreach \x in {7.3,8.3,9.3,10.6} {
            \draw (\x,3) circle (4mm);
        }
        \foreach \x in {7.5,8.4,9.4,10.35} {
            \draw (\x,2) circle (4mm);
        }
        \draw[red, very thick] (7,7) -- (7.1,4) -- (7.5,2);
        \draw[red, very thick] (8,7) -- (8,5) -- (8.4,2);
        \draw[red, very thick] (9,7) -- (9,5) -- (9.4,2);
        \draw[red, very thick] (10,7) -- (10,5);
        \draw[red, very thick] (11,7) -- (11,5) -- (10.4,2);
    \end{tikzpicture}
    \caption{The movement of a dislocation}\label{disloc}
\end{center}\end{figure*}

Figure \ref{disloc} shows a dislocation moving within a metal which would allow the metal to deform by moving one atom at a time. As the movement of dislocations is the dominant mode of plastic deformation, changes to the ability of dislocations to move through the metal have significant effects on it properties. For example:
\begin{description}
    \item[Work Hardening] As a metal is deformed, the dislocations move through the structure. Slowly the dislocations reach grain-boundaries or other dislocations and are no longer able to move. The metal therefore becomes less ductile and more brittle. This may be a desired property in order to harden a metal, or the additional brittleness may be undesirable.
    \item[Alloying] The addition of alloying atoms to the lattice can `pin' a dislocation in place (as shown in figure \ref{alloying}. The metal is therefore no longer able to deform by the movement of dislocations so the metal has a greater yield stress and is less ductile. Examples include adding carbon to iron to produce steel or adding zinc to copper to produce brass.
\end{description}
\begin{figure}[h]\begin{center}
    \begin{tikzpicture}
        % Left hand
        \foreach \y in {5,6,7} {
            \foreach \x in {0,1,2,3,4} {
                \draw (\x,\y) circle (4mm);
            }
        }
        \foreach \x in {0.1,1.15,2.8,3.8} {
                \draw (\x,4) circle (4mm);
            }
        \foreach \x in {0.3,1.3,2.6,3.6} {
            \draw (\x,3) circle (4mm);
        }
        \foreach \x in {0.5,1.4,2.4,3.35} {
            \draw (\x,2) circle (4mm);
        }
        \fill [red] (2,4) circle (2mm);
    \end{tikzpicture}
    \caption{Alloying atom pinning a dislocation}\label{alloying}
\end{center}\end{figure}

The effect on a stress-strain graph can be seen in figure \ref{stress-strain-alloy} below.

\begin{figure}[ht]
    \begin{center}
        \begin{tikzpicture}[scale=0.5]
            \draw[->] (-0.5,0) -- (10,0) node[anchor=north] {$\epsilon$};
            \draw[->] (0,-0.5) -- (0,8) node[anchor=east] {$\sigma$};
            \draw (0,0) -- (1.5,3) .. controls (2,3.5) .. (2.5,3) .. controls (3,2.8) .. (4,4) .. controls (5.5,5.5) and (8,5.5) .. (10,4.5) node[anchor=north] {Pure metal};
            \fill (10,4.5) circle (1mm);
            \draw (0,0) -- (1.3,5.5) .. controls (1.5,5.8) .. (2,6) node[anchor=south] {Alloy};
            \fill (2,6) circle (1mm);
        \end{tikzpicture}
    \end{center}
    \caption{Stress-strain curve for a ductile material}
    \label{stress-strain-alloy}
\end{figure}

\subsection{Polymers}

Polymers consist of long chain molecules\begin{marginfigure}
    \centering
    \includegraphics[]{polyethelene.png}
    \caption{Long chain molecule of polyethylene, a common material.}
\end{marginfigure} weakly held together by inter-molecular forces. Initially the molecules are likely to be tangled-up together. As force is applied it is initially difficult to move the polymer chains from this state (\textbf{A}). As the chains begin to unravel they straighten out by bond rotation, requiring relatively little force for a large increase in strain (\textbf{B}). As the polymer chains become straight it becomes much more difficult to extend the material any further without damaging the material (\textbf{C}).

\begin{figure}[ht]
    \begin{center}
        \begin{tikzpicture}[scale=.75]
            \draw[thick, ->] (-0.5,0) -- (11,0) node[anchor=north] {$\epsilon$};
            \draw[thick, ->] (0,-0.5) -- (0,9) node[anchor=east] {$\sigma$};
            \draw (0,0) .. controls (3,7) and (6,1) .. (10,8);
            \draw (1.5,1) node{\textbf{A}};
            \draw (5,1) node{\textbf{B}};
            \draw (8.5,1) node{\textbf{C}};
            \draw[dashed] (3,0) -- (3,3.6);
            \draw[dashed] (7.2,0) -- (7.2,4.7);
        \end{tikzpicture}
    \end{center}
    \caption{Stress-strain curve for a polymer}
    \label{stress-strain-polymer}
\end{figure}

\section{Rubber}
When a rubber is exposed to stress or strain energy, internal rearrangements such as rotation and extension of the polymer chains occur. 
\begin{marginfigure}
    \includegraphics[]{figures/naturalrubber.png}
    \caption{The twisted structure of natural rubber}
\end{marginfigure} These changes occur as a function of the energy applied and the duration and rate of application, as well as the temperature at which the energy is applied.
The difference between the two curves is the hysteresis losses.
\begin{figure}
	\vspace*{-10pt}
	\centering
	\begin{tikzpicture}[scale=0.54]
	\draw[<->] (6.4,0) node[below]{strain} -- (0,0) -- (0,6.4) node[left]{stress};
	\draw[thick,blue,postaction={decorate},decoration={markings,mark=at position 0.5 with {\arrow{>}}}] (0,0) [out=60,in=210] to (3,4) [out=30,in=225] to (6,6);
	\draw[thick,blue,postaction={decorate},decoration={markings,mark=at position 0.5 with {\arrow{>}}}] (6,6) [out=240,in=30] to (3,3) [out=210,in=50] to (0,0);
	\end{tikzpicture}
	\vspace*{-16pt}
\end{figure}
	
\section{End-of-chapter questions}


\subsection*{Springs \& Hooke's law}

\question{A spring has a force constant of $500 \text{ N m}^{-1}$. The spring has a stretch length of 30 cm when a 40 N weight is hung from it. Find the natural length of the spring.}

\question{Two springs, one with spring constant $k_1 = 8.0 \text{ N cm}^{-1}$ and the other with $k_2 = 12 \text{ N cm}^{-1}$. (a) When they are connected in parallel, what is the total extension when the combination supports a load of 60 N? (b) If they are connected in series, what is the total extension when they support the same load?}

\question{A weight of 100 N is placed on top of a spring. The spring is compressed by 2.0 cm. Assume the spring obeys Hooke's law, how much strain energy is stored in the spring?}

\question{A spring with a force constant $800 \text{ N m}^{-1}$ is supported and stands vertically. A ball of mass 60 g falls vertically onto it. The ball has a speed of $3.0\mps$ as it makes contact with the spring. Assume no energy loss, what is the maximum compression of the spring?}

\question{A trolley of mass 160 g is placed on a frictionless track. The trolly pushed against a spring buffer of a force constant $90 \text{ N m}^{-1}$. The spring buffer is compressed by 7.0 cm. The trolley is then released from rest. (a) What is the initial acceleration of the trolley? (b) Assume no loss of energy, what is the final speed of the trolley along the track?}

\question{
	Spring $A$ is stiffer than spring $B$, i.e., $k_A>k_B$. On which spring is more elastic potential energy stored if they are stretched by (a) the same extension, (b) the same force?
}

\begin{marginfigure}
	\vspace*{-5pt}
	\centering
	\begin{tikzpicture}
	\draw[help lines, gray!50, very thin, step=0.1] (0,0) grid (5,4);
	\draw[help lines, step=1] (0,0) grid (5,4);
	\foreach \y in {0,1,2,3, 4} \node[left] at (0,\y) {$\y$};
	\foreach \x in {0,1,2,3,4,5} \node[below] at (\x,0) {$\x$};
	\node at (-1,3.5) {load/kg};
	\node at (4,-0.8) {extension/cm};
	\draw[thick,blue] (0,0) -- (5,3.2);
	\end{tikzpicture}
	\vspace*{-25pt}
\end{marginfigure}

\question{A load is attached to the lower end of a spring. The extension of the spring is measured when the load increases. The variation with extension of the load is shown. (a) Suggest whether the spring obeys Hooke's law. (b) Find the spring constant. (c) How much elastic potential energy is stored in the spring when the extension is 5.0 cm? (d) How much work is done to increase the extension from 2.0 cm to 4.0 cm?}


\question{
	Suppose you have a spring, a ruler, a mass hanger and a set of masses. Suggest how the apparatus may be used to determine the load on the spring at (a) the limit of elasticity, (b) the limit of proportionality?
}


\subsection*{Stress, Strain, \& Young modulus}

\question{A cable of diameter 1.5 mm is under a tension of 200 N. Find the stress in this cable.}

\question{Estimate the stress in your neck when it supports your head in a vertical position.}

\question{A metal wire of natural length 1.8 m and diameter 0.70 mm is fixed to the ceiling at one end. When a mass of 6.5 kg is hung from the lower end, the wire extends by 2.7 mm. (a) Find the strain of the wire. (b) Find the Young modulus of the metal.}

\question{A wire has a diameter of 0.50 mm with a Young modulus of 0.18 TPa. The length of the wire is increased by 0.20\% by a force $F$. (a) Find the stress in the wire. (b) Find the force $F$.}


\question{Two metal wires $A$ and $B$ are made of different materials. The diameter of wire $A$ is twice that of wire $B$, and the Young modulus of wire $A$ is three times that of wire $B$. If the wires are extended by the same strain, what is the ratio of the tension in wire $A$ to tension in wire $B$?}



\question{Two rods $A$ and $B$ of the same diameter are joined end to end and hug vertically. Rod $B$ is twice as long as rod $A$ but has half the Young modulus. When a mass is hung from the combination, what is the ratio of the extension of rod $A$ to the extension of rod $B$?}

\question{Bones of different animals have very similar compositions. Suggest why heavier animals appear to have thicker bones?}


\begin{marginfigure}
	\vspace*{5pt}
	\centering
	\begin{tikzpicture}
	\draw[help lines, gray!50, very thin, step=0.1] (0,0) grid (5,5);
	\draw[help lines, step=1] (0,0) grid (5,5);
	\foreach \y in {0,1,2,3,4,5} \node[left] at (0,\y) {$\y$};
	\foreach \x in {0,0.1,0.2,0.3,0.4,0.5} \node[below] at (10*\x,0) {$\x$};
	\node at (-1,4.5) {$F$/N};
	\node at (4,-0.8) {$x$/mm};
	\draw[thick,blue] (0,0) -- (5,2.4);
	\end{tikzpicture}
	\vspace*{-25pt}
\end{marginfigure}

\question{A load $F$ is suspended from a copper wire. The graph shows the load-extension relation. (a) State two quantities other than $F$ and $x$ that are required to determine the Young modulus of copper. (b) Suggest how the two quantities may be measured. (c) Use the graph to find the energy stored in the wire when a load of 2.0 N is applied. (d) Given that steel has about twice the Young modulus as copper, sketch the variation with $x$ of $F$ for a steel wire that has the same dimensions as the copper wire.}

\begin{marginfigure}
	\vspace*{-10pt}
	\centering
	\begin{tikzpicture}[scale=0.54]
	\draw[<->] (6.4,0) node[below]{strain} -- (0,0) -- (0,6.4) node[left]{stress};
	\draw[thick,blue,postaction={decorate},decoration={markings,mark=at position 0.5 with {\arrow{>}}}] (0,0) [out=60,in=210] to (3,4) [out=30,in=225] to (6,6);
	\draw[thick,blue,postaction={decorate},decoration={markings,mark=at position 0.5 with {\arrow{>}}}] (6,6) [out=240,in=30] to (3,3) [out=210,in=50] to (0,0);
	\end{tikzpicture}
	\vspace*{-16pt}
\end{marginfigure}

\question{
	In an experiment, a speciman of a rubber compound is being stretched and relaxed. The stress-strain curve is plotted. (a) State and explain whether this rubber compound behaves elastically. (b) The tyres on a vehicle are made of this rubber compound. Explain why the tyres become warm as the car travels on a road.
}

%\include{fluids}
\chapter{Waves}

Wave is a way of transferring and storing \emph{energy} without the transport of matter. Most familiar examples are surface waves on water, sound waves, light waves. In the next two chapters, we are going to look at the basics of wave motion and some of the most important wave phenomena.

%Wave motion is the propagation of disturbance -- deviations from a state of equilibrium -- from one place to another. 
%
%Waves consist of continuous oscillations around some fixed locations.
%
%There are mainly two types of waves, \keypoint{mechanical waves} and \keypoint{electromagnetic waves}. Mechanical waves propagate as oscillations of matter, thus it must travel through a \emph{medium}. Electromagnetic waves can travel in free space.
%
%A wave can either be \keypoint{transverse} or \keypoint{longitudinal} depending on the direction of its oscillation. The direction of propagation of a transverse wave is perpendicular to the direction of its oscillation, and the wave motion of a longitudinal wave is parallel to its oscillation.

\subsection{wave terminologies}

wave motion is the propagation of disturbance from one place to another

\begin{figure*}[ht]
	\centering
	\begin{tikzpicture}[xscale=1.2, yscale=0.95]
	\draw [->] (0,-2) -- (0,2.4) node[midway,left]{$O$} node[above]{displacement};
	\draw [->] (0,0) -- (3.1*pi,0) node[below]{position};
	\draw[very thick,->] (1.8*pi,2) -- (2.6*pi,2);
	\node[above,twoline] at (2.2*pi, 2.1) {direction of\\wave motion};
	\foreach \t/\labelcolor in {1/{blue}, 2/{Green}, 3/{purple}, 4/{gray}} {
		\draw[\labelcolor] ({0.4*pi+(\t-1)*0.6)}, 1.7) --++ (0.2,0.5) node[above]{$t_\t$};
		\draw [thick, \labelcolor, domain=0:3*pi, samples=30, smooth] plot (\x,{1.6*sin((\x+0.6-0.6*\t)*1.25 r)});
	}
	\end{tikzpicture}
	
	\caption{wave pattern at different times ($t_1<t_2<t_3<t_4$) as wave travels in space}
\end{figure*}


to describe the wave motion and particle vibrations, we can define the following quantities:


\cmt as a wave moves out, each point oscillates back and forth about their rest positions

distance from a particle's equilibrium position is called \keypoint{displacement} of the particle \index{displacement}

\cmt greatest displacement for a particle is called the \keypoint{amplitude} ($A$) \index{amplitude}

\cmt wave pattern repeats itself over a certain distance

distance between two adjacent points undergoing exactly same motion is the \keypoint{wavelength} ($\lambda$) \index{wavelength}

one can think of wavelength as crest-to-crest distance, trough-to-trough distance, etc.\footnote{For now, we take for granted that a wave is transverse. There are also longitudinal waves for which terms like crest and trough do not apply. We will get into that in \S\ref{ch-Lwaves}.}

\cmt each point also repeats its vibrational motion over a certain time interval

time for a particle to complete a full oscillation cycle is the  \keypoint{period} ($T$) \index{period}

\cmt number of oscillations for a particle per unit time is the \keypoint{frequency} ($f$) \index{frequency}

frequency can also be defined as the number of crests passing a given point per unit time

frequency of a wave is related to its period by \begin{empheq}[box=\tcbhighmath]{equation*}{f=\frac{1}{T}}\end{empheq}

unit of frequency: $[f] = \text{Hz}$ (hertz), where $1 \text{ Hz} = 1 \text{ s}^{-1}$

\begin{figure*}[ht]
	\centering
	\begin{minipage}{0.48\textwidth}
		\centering
		\begin{tikzpicture}[scale=0.55]
		\draw [->] (0,-2) -- (0,2) node[midway,left]{$O$} node[above]{displacement};
		\draw [->] (0,0) -- (3*pi,0) node[below]{position};
		\draw [very thick,blue,domain=0:3*pi,samples=30,smooth,variable=\x] plot (\x,{1.6*sin(\x*1.25 r)});
		\draw[<->, thick, red] (.4*pi,1.6) -- (.4*pi,0) node[midway,right]{$A$};
		\draw[<->, thick, red] (.4*pi,1.8) -- (2*pi,1.8) node[midway,above]{$\lambda$};
		\end{tikzpicture}
		
		wave pattern of all particles at one instant
	\end{minipage}\hfil
	\begin{minipage}{0.48\textwidth}
		\centering
		\begin{tikzpicture}[scale=0.55]
		\draw [->] (0,-2) -- (0,2) node[midway,left]{$O$} node[above]{displacement};
		\draw [->] (0,0) -- (3*pi,0) node[below]{time};
		\draw [very thick,blue,domain=0:3*pi,samples=30,smooth,variable=\x] plot (\x,{1.6*sin(\x*1.25 r)});
		\draw[<->, thick, red] (.4*pi,1.6) -- (.4*pi,0) node[midway,right]{$A$};
		\draw[<->, thick, red] (.4*pi,1.8) -- (2*pi,1.8) node[midway,above]{$T$};
		\end{tikzpicture}
		
		vibration of one specific particle at all times
	\end{minipage}
\end{figure*}

\cmt wave energy is transferred along the direction of wave motion at a certain speed $v$

in one period, wave moves forward by a distance of one wavelength

so \keypoint{wave speed} \begin{empheq}[box=\tcbhighmath]{equation*}{v=\frac{\lambda}{T}}\end{empheq}, or in terms of frequency, \begin{empheq}[box=\tcbhighmath]{equation*}{v=\lambda f}\end{empheq}

\example{When a wave travels on a water surface, the maximum depth of water is 21 cm and the minimum depth is 18 cm. What is the amplitude of the wave?}

\begin{soln} amplitude is half the end-to-end distance: $ A = \frac{1}{2}(21-18) = 1.5 \text{ cm} $ \end{soln}


\example{A wave travelling at $4.0 \mps$ has a wavelength of 50 cm, what is its period?}

\begin{soln} \begin{equation*}
	v = \frac{\lambda}{T} \RA T = \frac{\lambda}{v} = \frac{0.50}{4.0} = 0.125 \text{ s} 
\end{equation*}
\end{soln}

\subsection{transverse \& longitudinal waves}

a wave can either be \emph{transverse} or \emph{longitudinal}, depending on the direction of its oscillation

\subsection{transverse waves}

\begin{ilight}
	\centering a \keypoint{transverse} wave has vibrations at right angle to its direction of energy transfer \index{transverse wave}
\end{ilight}

\cmt examples of transverse waves: wave on a string, surface wave on water, light wave, etc.

\cmt for a transverse wave, greatest displacement in positive direction is called a \emph{crest}, or a \emph{peak}

greatest displacement in negative direction is called a \emph{trough}


\subsection{longitudinal waves} \label{ch-Lwaves}

\begin{ilight}
	\centering a \keypoint{longitudinal} wave has vibrations in parallel direction to energy transfer \index{longitudinal wave}
\end{ilight}

\cmt examples of longitudinal waves: sound waves, wave along a stretched slinky, etc.

\cmt for a longitudinal wave, if medium gets squeezed, we say this region is a \emph{compression}

if a medium expands, we say this region is a \emph{rarefaction}

\cmt wavelength of a longitudinal wave can be defined as compression-to-compression distance

\begin{figure*}[ht]
	\centering
	\begin{tikzpicture}[xscale=1, yscale=1.7]
	\foreach \x in {-0.75,-0.5,...,10.75}
	\draw[thick] ({\x} , -1) --++ (0,1);
	\end{tikzpicture}
	
	all particles at their equilibrium position before a pressure wave is set up
\end{figure*}

\begin{figure*}[ht]
	\centering
	\begin{tikzpicture}[xscale=1, yscale=1.7]
	\foreach \x in {-0.75,-0.5,...,11} 
		\draw[thick] ({\x+0.5*cos(\x*0.5*pi r)} , -1) --++ (0,1);
	\foreach \x in {-0.5,-0.25,...,11.25} 
		\draw[thick] ({\x+0.5*cos((\x-1)*0.5*pi r)} , -2.5) --++ (0,1);
	\foreach \x in {-0.5,-0.25,...,11} 
		\draw[thick] ({\x+0.5*cos((\x-2)*0.5*pi r)} , -4) --++ (0,1);
		\foreach \x in {3,7} \node[Green] at (\x,-1.2) {\phantom{p}rarefaction\phantom{p}};
		\foreach \x in {1,5,9} \node[blue] at (\x,-1.2) {compression};
		\foreach \x in {0,4,8} \node[Green] at (\x,-2.7) {\phantom{p}rarefaction\phantom{p}};
		\foreach \x in {2,6,10} \node[blue] at (\x,-2.7) {compression};
		\foreach \x in {1,5,9} \node[Green] at (\x,-4.2) {\phantom{p}rarefaction\phantom{p}};
		\foreach \x in {3,7} \node[blue] at (\x,-4.2) {compression};
		\node at (-1.5,-0.5) {{\large $t_1$}};
		\node at (-1.5,-2) {{\large $t_2$}};
		\node at (-1.5,-3.5) {{\large $t_3$}};
		\draw[very thick,->] (2.8,0.3) -- (7.2,0.3);
		\node[above,twoline] at (5, 0.4) {direction of wave motion};
		\draw[thick,red,<->] (1,-0.5) -- (5,-0.5) node[above,pos=0.55]{$\lambda$};
	\end{tikzpicture}
	
	\caption{compression and rarefaction regions at different times ($t_1<t_2<t_3$)	as a longitudinal pressure wave travels in space}
\end{figure*}



\example{What is the distance between a compression and a rarefaction for a sound wave of frequency 550 Hz? (Speed of sound in air is about $330 \mps$.)}

\begin{soln}\begin{equation*}
	d = \frac{1}{2}\lambda = \frac{v}{2f} = \frac{330}{2\times 550} \RA d = 0.30 \text{ m} 
\end{equation*}
\end{soln}

\subsection{sound waves}

sound waves propagate via the compression and rarefaction of air (or other medium) \index{sound}

molecules near vibrating source is pushed away from rest positions and into their neighbours

neighbouring molecules then in turn push into their neighbours, and so on

the disturbance is transferred through the medium, forming a sound wave

\cmt sound waves are longitudinal
	
\begin{figure*}[ht]

\begin{tikzpicture}[scale=1]
\pgfmathsetmacro{\noe}{6}

\foreach \t in {0,0.1,...,14}{
	\pgfmathsetmacro{\np}{pow({1+cos((\t-2)*pi/2.5 r)}, 4)*\noe/4}
	\foreach \idx in {0,1,...,\np} {
		\draw[fill] (\t+rnd*0.1-0.05,rnd*1.8) circle (0.025);	% wave 1
	}
}
\draw[<->,thick] (7,2.4) -- (12,2.4) node[midway,above]{{\footnotesize one wavelength}};
\draw[->,thick] (1,2.4) -- (4,2.4) node[midway,above]{{\footnotesize wave motion}};
\node[below] at (2,-.2) {{\footnotesize compression}};
\node[below] at (7,-.2) {{\footnotesize compression}};
\node[below] at (12,-.2) {{\footnotesize compression}};
\node[below] at (4.5,-.2) {{\footnotesize rarefaction}};
\node[below] at (9.5,-.2) {{\footnotesize rarefaction}};
\end{tikzpicture}
\end{figure*}

\cmt propagation of sound waves require medium (air, water, steel, etc.)
	
sound cannot travel in vacuum
	
\cmt speed of sound is material-dependent, but not frequency-dependent
	
	sound in general travels faster in denser medium
	
	\titem $v_\text{air} \approx 340 \mps$ (under standard atmospheric pressure and room temperature)
	
	\titem $v_\text{water} \approx 1500 \mps$
	
	\titem $v_\text{steel} \approx 5000 \mps$
	
\cmt \emph{pitch} of a note is related to frequency of sound wave

rapid vibrations of sound source at high frequencies produce a high pitch

\cmt \emph{loudness} of sound mainly depends on amplitude of vibration

a greater amplitude means the wave is more energetic so it sounds louder
	

\subsection*{measurement of sound waves}

two key apparatuses for sound measurement are the \keypoint{microphone} and the \keypoint{oscilloscope} \index{oscilloscope}

sound waves can be captured by a \emph{microphone}, which converts sound into electrical signals

electrical signals can then be sent into an \emph{oscilloscope} (see figure\footnote{The beautiful figure of the oscilloscope was created by \emph{Hugues Vermeiren}, who generously shared the source codes on TeXample: \url{http://www.texample.net/tikz/examples/textronics-oscilloscope/}}) for measurement

oscilloscope can be thought as an upgraded voltmeter showing how voltage varies with time

\begin{figure}[ht]
	\centering
	\includegraphics{oscilloscope.pdf}
	\caption{the display and controls of a typical cathode-ray oscilloscope}
\end{figure}

\cmt oscilloscope displays the variation of voltage ($y$-axis) with time ($x$-axis)

\cmt horizontal scale (time axis) is given by \emph{time-base} settings

vertical scale (voltage axis) is given by \emph{voltage gain}, or \emph{Y-sensitivity} settings

\cmt period $T$ of the sound wave can be found using time-base settings

then frequency of sound wave is calculated: $f=\frac{1}{T}$

\cmt voltage amplitude can be found using the voltage gain


\begin{marginfigure}
	\vspace*{-12pt}
	\centering
	\begin{tikzpicture}[scale=0.8]
	\draw[fill=gray!80, rounded corners] (-0.2,-3.2) rectangle (7.2,3.2);
	\draw[Green!50, fill] (0,-3) rectangle (7,3);
	\draw[Green!20, ultra thick, domain=0:7, samples=100] plot (\x, {2*sin((\x-0.6)*120)});
	\draw[black!80, step=1] (0,-3) grid (7,3);
	\end{tikzpicture}
	\vspace*{-16pt}
\end{marginfigure}

\example{A sound wave is detected by a microphone and the trace is displayed on an oscilloscope as shown. If the time-base is set at 0.5 ms div$^{-1}$ and the voltage gain is set at 2 V div$^{-1}$. What is the frequency and the amplitude of the signal?}

\begin{soln} period: $T = 3 \times 0.5 \text{ ms} = 1.5 \times10^{-3} \text{ s}$

\eqyskip frequency: $f = \frac{1}{T} = \frac{1}{1.5 \times10^{-3} \text{ s}} \approx 667 \text{ Hz}$

\eqyskip amplitude: $A = 2 \times 2 \text{ V} = 4.0 \text{ V}$ \end{soln}




\subsection{electromagnetic waves}

a wave can be either \emph{mechanical} or \emph{electromagnetic}, depending on whether it requires a medium\index{electromagnetic wave}

\cmt waves that require a \emph{medium} to travel are called \keypoint{mechanical waves}

examples are sound waves, water waves, wave on a string, etc.

mechanical waves involve vibration of matter particles

\cmt no medium is needed for propagation of \keypoint{electromagnetic waves} i.e., they can travel in vacuum



\subsection{properties of electromagnetic waves}

\cmt electromagnetic waves can travel in free space

\cmt electromagnetic waves involve vibrations of electric and magnetic fields
	
an altering electric field can generate an altering magnetic field, which then further produces a new altering electric field, so on and so forth

electric and magnetic fields then permeate through space, transferring energy and information
	
\cmt electromagnetic waves are all transverse
	
vibration of electric fields and magnetic fields are both perpendicular to wave motion

\begin{figure*}[ht]
	\centering
	\begin{tikzpicture}[xscale=0.72]
	\draw[Green, thick, ->] (0,-2.1) -- (0,2.1) node[left, twoline] {electric\\field};
	\draw[blue, thick, dotted] (0 ,0) -- (0.7, 1.4);
	\draw[blue, thick, ->] (0, 0) -- (-0.7,-1.4) node[left, twolinecap] {magnetic\\field};
	\draw[thick, ->] (-0.3,0) -- (14,0) node[right, twoline]{direction\\ of energy\\transfer};
	\draw[Green, thick, domain=0:13.5, samples=150] plot (\x, {1.6*sin(2*\x r)});
	\draw[blue, thick] plot [smooth] coordinates {
		(0, 0)
		(0.0007, -0.1987)
		(0.0053, -0.3894)
		(0.0177, -0.5646)
		(0.0413, -0.7174)
		(0.0793, -0.8415)
		(0.134, -0.932)
		(0.2073, -0.9854)
		(0.3002, -0.9996)
		(0.4131, -0.9738)
		(0.5454, -0.9093)
		(0.6958, -0.8085)
		(0.8623, -0.6755)
		(1.0422, -0.5155)
		(1.2325, -0.335)
		(1.4294, -0.1411)
		(1.6292, 0.0584)
		(1.8278, 0.2555)
		(2.0213, 0.4425)
		(2.2059, 0.6119)
		(2.3784, 0.7568)
		(2.5358, 0.8716)
		(2.6758, 0.9516)
		(2.7968, 0.9937)
		(2.8981, 0.9962)
		(2.9795, 0.9589)
		(3.0417, 0.8835)
		(3.0864, 0.7728)
		(3.1156, 0.6313)
		(3.1323, 0.4646)
		(3.1397, 0.2794)
		(3.1415, 0.0831)
		(3.1417, -0.1165)
		(3.1442, -0.3115)
		(3.1529, -0.4941)
		(3.1715, -0.657)
		(3.2032, -0.7937)
		(3.2506, -0.8987)
		(3.316, -0.9679)
		(3.4007, -0.9985)
		(3.5053, -0.9894)
		(3.6296, -0.9407)
		(3.7727, -0.8546)
		(3.9328, -0.7344)
		(4.1075, -0.5849)
		(4.2939, -0.4121)
		(4.4886, -0.2229)
		(4.6876, -0.0248)
		(4.8872, 0.1743)
		(5.0832, 0.3665)
		(5.272, 0.544)
		(5.4499, 0.6999)
		(5.6139, 0.8278)
		(5.7614, 0.9228)
		(5.8905, 0.9809)
		(6, 1)
		(6.0896, 0.9792)
		(6.1597, 0.9193)
		(6.2114, 0.8228)
		(6.2468, 0.6935)
		(6.2683, 0.5366)
		(6.2791, 0.3582)
		(6.2828, 0.1656)
		(6.2832, -0.0336)
		(6.2842, -0.2315)
		(6.2899, -0.4202)
		(6.304, -0.5921)
		(6.3298, -0.7404)
		(6.3704, -0.8592)
		(6.4282, -0.9437)
		(6.5047, -0.9906)
		(6.601, -0.998)
		(6.7172, -0.9657)
		(6.8526, -0.8948)
		(7.0059, -0.7883)
		(7.1749, -0.6503)
		(7.3568, -0.4864)
		(7.5484, -0.3031)
		(7.7461, -0.1078)
		(7.946, 0.0919)
		(8.144, 0.2879)
		(8.3362, 0.4724)
		(8.5191, 0.6381)
		(8.6892, 0.7784)
		(8.8438, 0.8876)
		(8.9807, 0.9614)
		(9.0985, 0.9969)
		(9.1963, 0.9927)
		(9.2744, 0.9488)
		(9.3336, 0.8672)
		(9.3755, 0.751)
		(9.4024, 0.6048)
		(9.4173, 0.4346)
		(9.4235, 0.247)
		(9.4248, 0.0495)
		(9.4251, -0.1499)
		(9.4283, -0.3433)
		(9.4385, -0.5231)
		(9.459, -0.682)
		(9.4932, -0.8137)
		(9.5435, -0.9129)
		(9.6121, -0.9758)
		(9.7001, -0.9998)
		(9.808, -0.9839)
		(9.9356, -0.9288)
		(10.0817, -0.8367)
		(10.2444, -0.7112)
		(10.4213, -0.5573)
		(10.6094, -0.3813)
		(10.805, -0.19)
		(11.0044, 0.0089)
		(11.2037, 0.2073)
		(11.3988, 0.3976)
		(11.586, 0.5719)
		(11.7617, 0.7235)
		(11.9231, 0.8462)
		(12.0676, 0.9352)
		(12.1935, 0.9869)
		(12.2996, 0.9993)
		(12.3859, 0.9718)
		(12.4528, 0.9056)
		(12.5016, 0.8033)
		(12.5345, 0.6689)
		(12.5539, 0.5079)
		(12.5633, 0.3266)
		(12.5662, 0.1324)
		(12.5664, -0.0672)
		(12.568, -0.2641)
		(12.5748, -0.4504)
		(12.5906, -0.6188)
		(12.6187, -0.7626)
		(12.6621, -0.8759)
		(12.7229, -0.9543)
		(12.8027, -0.9946)
		(12.9023, -0.9954)
		(13.0218, -0.9564)
		};
	\end{tikzpicture}
	\caption{variation in electric and magnetic fields for an electromagnetic wave}
\end{figure*}

	
\cmt all electromagnetic waves travel at a constant speed $c=3.0\times10^8\mps$ in free space

i.e., speed of light in vacuum is constant\footnote{The speed of light in vacuum is actually a \emph{universal} physical constant. According to Einstein's special relativity, this is the upper limit for the speed at which matter and information can travel. This speed is also independent of the inertial reference frame one chooses, i.e., the speed of light in vacuum is the same for all observers, regardless of the motion of the source or the observer.}



\subsection{electromagnetic spectrum}

electromagnetic waves come in a wide range of wavelengths and frequencies 

distribution of electromagnetic radiation according to wavelength or frequency is the \keypoint{electromagnetic spectrum}\index{electromagnetic spectrum} (see diagram)



\begin{figure*}[!ht]
	\centering
	\begin{tikzpicture}[yscale=1.04]
	% define visible spectrum shading
	\pgfdeclareverticalshading{visible}{100bp}
	{color(0bp)=(violet!70); color(25bp)=(violet!70); color(35bp)=(blue!70);
		color(45bp)=(green!70); color(55bp)=(yellow!70); color(65bp)=(orange);
		color(75bp)=(red); color(100bp)=(red)}
	% grayish background for spectrum labels
	\draw[gray!15,fill] (-1.5,-8) rectangle (1.5,7.5);
	% axes
	\draw (-1.5,-8) -- (-1.5,7.5) (1.5,-8) -- (1.5,7.5);
	\node at (-3,7.7) {wavelength};
	\node at (3,7.7) {frequency};
	% spectrum labels
	\shade[shading=visible] (-1.5,-1.15) rectangle (1.5,-0.75);
	\foreach \x/\xlabel in {1.8/{radio waves}, -2/{microwaves}, -4.6/{infrared}, -6.35/{visible light}, -7.8/{ultraviolet}, -11/{X-rays}, -14.8/{$\gamma$-rays} } \node at (0,{\x*0.7+3.5}) {\xlabel};
	% values of wavelengths
	\foreach \x in {5,3,1,...,-15} {
		\draw (-1.5,{\x*0.7+3.5}) --++ (-0.5,0);
		\node at (-3,{\x*0.7+3.5}) {$10^{\x}$ m};
	}
	% values of frequencies
	\foreach \x in {3,5,...,23} {
		\draw (1.5,{-\x*0.7+8.8}) --++ (0.5,0);
		\node at (3,{-\x*0.7+8.8}) {$10^{\x}$ Hz};
	}
	% dashed lines as separators
	\foreach \x in {-1,-3,-9,-13} \draw[dashed,thick] (-1.5,{\x*0.7+3.5}) --++ (3,0);
	% extension for visible spectrum
	\draw[thick, gray!75] (1.5,-1.15) --++ (3,0) -- (5,-3);
	\draw[thick, gray!75] (1.5,-0.75) --++ (3,0) -- (5,1);
	\shade[shading=visible] (5,-3) rectangle (7,1);
	\node at (6, 1.3) {visible light};
	\node[right] at (7.2, 0.8) {700 nm};
	\node[right] at (7.2, -2.8) {400 nm};
	\end{tikzpicture}
	\caption{the electromagnetic spectrum}
\end{figure*}

\newpage

\cmt the table below shows electromagnetic spectrum in somewhat more precise details \footnote[][-2cm]{There are no precise accepted boundaries between different ranges in the electromagnetic spectrum. The boundaries are actually somewhat ambiguous. The ranges of different portions tend to overlap. For example, a large portion of the ranges of X-rays overlap with that of $\gamma$-rays, and microwaves are considered by many people as a subdivision of radio waves. Therefore, the values given here are merely supposed to give you some rough idea about the order of magnitudes for electromagnetic wavelengths and frequencies. So the point I want to make here is: do not take the borderlines too seriously.}
\vspace{3cm}
\begin{figure*}
\begin{center}
	\begin{tabular}{|C{4cm}|C{4cm}|C{4cm}|}
		\hline type of radiation & range of wavelength & range of frequency \\ 
		\hline radio waves  & 10$^{-1}$ $\sim$ 10$^{6}$ m & 10$^{2}$ $\sim$ 10$^{9}$ Hz \\ 
		\hline microwaves & 10$^{-3}$ $\sim$ 10$^{-1}$ m & 10$^{9}$ $\sim$ 10$^{11}$ Hz \\ 
		\hline infra-red & 7$\times$10$^{-7}$ $\sim$ 10$^{-3}$ m & 10$^{11}$ $\sim$ 10$^{14}$ Hz \\ 
		\hline visible light  & 4$\times$10$^{-7}$ $\sim$ 7$\times$10$^{-7}$ m & 10$^{14}$ $\sim$ 10$^{15}$ Hz \\ 
		\hline ultraviolet & 10$^{-9}$ $\sim$ 4$\times$10$^{-7}$ m & 10$^{15}$ $\sim$ 10$^{17}$ Hz \\ 
		\hline X-rays & 10$^{-13}$ $\sim$ 10$^{-9}$ m & 10$^{17}$ $\sim$ 10$^{19}$ Hz \\ 
		\hline $\gamma$-rays & 10$^{-16}$ $\sim$ 10$^{-11}$ m & 10$^{19}$ $\sim$ 10$^{24}$ Hz\\ 
		\hline
	\end{tabular} 
\end{center}
\end{figure*}

\cmt each type of electromagnetic waves has important applications in some area
\footnote{The entries listed here only include a teeny-weeny part of the uses of electromagnetic radiation, somewhat based on my personal taste. I also included a handful of explanations for the examples that I found interesting (otherwise I would not choose them), as you will see a huge load of footnotes in the next few pages. You are encouraged to do some researches as well, I can guarantee that you will not be disappointed.}

\begin{compactitem}
	\item[--] radio waves
	
	\xskip telecommunication (TV/radio broadcast, satellite communication) \footnote{Having the longest wavelengths of all radiation, radio waves have the best ability to diffract around obstacles in city buildings and mountains, therefore a large area can be covered.}

	\item[--] microwaves
	
	\xskip telecommunication (mobile phones, WiFi, Bluetooth, satellite communication)
	\footnote{Microwaves do not diffract sufficiently as radio waves, but they can transmit more information per unit time because they have higher frequencies. Microwaves are used in short-range telecommunication, including mobile phones, wireless networks, and bluetooth connections.}
		
		
	\xskip heating food (microwave ovens)
	\footnote{Frequency of microwaves are close to the natural frequencies of the rotational motion of water molecules. When food is exposed to microwaves, the water molecules in the food resonate and vibrate more violently, causing a rise in the food's temperature.}
	
	\item[--] infrared (IR)
	
	\xskip IR thermography (temperature monitors, thermographic cameras)
	\footnote{All objects emit electromagnetic radiation based on their temperatures. According to the law of black-body radiation, objects near room temperature emit thermal energy as infrared radiation, so variations in the temperature can be detected.}
	
	\xskip night-vision devices
	\footnote{Night-vision devices convert photons (just think of them as particles of light for now) into electrons, which are amplified by a chemical and electrical process and then converted back into visible light. Infrared sources can be used to augment the available light, increasing the visibility in the dark.}
	
	\xskip IR heating (IR heat lamps, IR saunas)
	
	
	\xskip IR data transmission (remote controls, optical-fibre communication)
	
	\item[--] visible light
	
	\xskip human vision
	\footnote{Human eyes are only sensitive to a small fraction of the electromagnetic spectrum. The wide variety of colours that we see is actually built up from the relative intensities of red, green and blue light collected by the three colour detectors in our eyes.}
	
	\item[--] ultraviolet (UV)
	
	\xskip UV sterilising (drinking water treatment, disinfection of medical facilities, etc.)
	\footnote{Short-wavelength UV light can damage the DNA's in living organisms. A microorganism exposed to germicidal UV light might not be able to reproduce, and becomes harmless. For the same reason, overexposure to UV radiation present in the sunlight can cause sunburn, or even skin cancer.}
		
	\xskip fluorescent dyes (black light fluorescent paint, UV watermarks)
	\footnote{UV radiation can cause many substances to glow through chemical reactions. UV watermarks that are visible under UV light are used to prevent counterfeiting of currency, or forgery of important documents such as passports and ID cards.}
	
	\item[--] X-rays
	
	\xskip medical imaging (X-ray imaging, CT scans)
	\footnote{X-rays are very energetic and thus very penetrating. They can pass through human body easily to form an image giving information about the structures of tissues and bones.}
	
	\xskip security check (luggage scanners)
	
	\xskip X-ray crystallography
	\footnote{X-rays can be diffracted by the lattice of atoms in a crystal. The diffraction pattern gives information about the structure of the lattice. X-ray crystallography is a very important experimental technique to study the microscopic structures of new materials.}
	
	\item[--] $\gamma$-rays
	
	\xskip radiation therapies (cancer treatment)
	\footnote{$\gamma$-rays have extremely high frequencies. They are even more energetic and penetrating. They can be used to damage the DNA of cancerous tissue, and hence kill the cancerous cells as a treatment.}
	
	\xskip medical imaging (PET scans)
	\footnote{Positron emission tomography (PET) uses a radioactive tracer to produce $\gamma$-rays within the tissues of interest. The energy and location of these $\gamma$-rays can be detected and sent to a computer to build up a 3D image of the body part.}
\end{compactitem}

\example{A beam of electromagnetic radiation is known to have a frequency of 25 THz in vacuum. What is its wavelength?}

\begin{soln}
    
\begin{equation*}
	\lambda = \frac{c}{f} = \frac{3.00\times10^8}{25\times10^12} \RA \lambda = 1.2 \times 10^{-5} \text{ m} \quad \text{(infra-red)} 
\end{equation*}
\end{soln}

	

\subsection{wave intensity}

energy can be transmitted along a wave

the degree to which energy is concentrated is called the intensity of the wave

here concentration has two meanings -- concentration in time and concentration in area $S$\footnote{To avoid confusion, I reserved letter `$A$' for wave amplitude and chose `$S$' to represent an area.}

\begin{ilight}
	\centering \keypoint{intensity} of a wave is defined as the power $P$ per unit area on a cross section $S$: \begin{empheq}[box=\tcbhighmath]{equation*}{I=\frac{P}{S}}\end{empheq}\index{wave intensity}
\end{ilight}

\cmt unit of wave intensity: $[I] = \text{ W m}^{-2}$

\cmt intensity is proportional to square of its amplitude: $\boxed{I \propto A^2}$.

\cmt intensity of a wave decreases as it spreads out in space

intensity at distance of $r$ from a point source is inversely proportional to $r^2$: $\boxed{I \propto \frac{1}{r^2}}$
\footnote{As a wave produced from a point source travels out by a distance $r$ away from the source, the energy it carries is spread uniformly over the surface area of the sphere of radius $r$, that is: $S=4\pi r^2$. So the intensity of a wave obeys an inverse square law: $I \propto \frac{1}{r^2}$.}

\example{If the amplitude of an incoming wave is increased by 50\%, what is the increase in the wave intensity?}

\begin{soln} $I \propto A^2 \RA \frac{I'}{I} = \left( \frac{A'}{A} \right)^2 = \left( 1 + 50\% \right)^2 =2.25 \RA $ so intensity is increased by 125\% \end{soln}

\example{Two observers $A$ and $B$ are at a distance $r_A$ and $r_B$ from a point source where $r_A=2r_B$. (a) Find the ratio of their intensities $\frac{I_A}{I_B}$. (b) Find the ratio of their amplitudes $\frac{A_A}{B_B}$.}

\begin{soln}

{
	\centering
	
	$ I \propto \frac{1}{r^2} \RA \frac{I_A}{I_B} = \left(\frac{r_B}{r_A}\right)^2 = \left(\frac{1}{2}\right)^2 \RA \frac{I_A}{I_B} = \frac{1}{4} $
	
	\vspace*{0.4em} $ I \propto A^2 \RA \frac{A_A}{A_B} = \sqrt{\frac{I_A}{I_B}} = \sqrt{\frac{1}{4}} \RA \frac{A_A}{A_B} = \frac{1}{2} $
	
}

\end{soln}




\subsection{polarisation}

\subsection{plane polarisation}

\begin{ilight}
	a wave is \keypoint{plane-polarised}, or simply \keypoint{polarised}, if the vibration is in one single direction at right angle to the direction of propagation of energy\index{polarisation}
\end{ilight}

\begin{figure}[ht]
	\centering
	\begin{minipage}{0.45\textwidth}
		\centering
		\begin{tikzpicture}[scale=0.64]
		\draw[very thick, blue, ->] (210:2) -- (30:8) node[twoline, above, black]{direction of\\energy transfer};
		\foreach \t in {0,45,90,135} {
			\draw[red, ultra thick, <->] (\t:1.5) --++ ({\t+180}:3);
			\draw[red, ultra thick, <->] (30:3.5) ++ (\t:1.5) --++ ({\t+180}:3);
		}
		\end{tikzpicture}
		
		an unpolarised wave
		
		(wave vibrations in multiple directions)
	\end{minipage}\hfil
	\begin{minipage}{0.45\textwidth}
		\centering
		\begin{tikzpicture}[scale=0.64]
		\draw[very thick, blue, ->] (210:2) -- (30:8) node[twoline, above, black]{direction of\\energy transfer};
		\foreach \t in {90} {
			\draw[red, ultra thick, <->] (\t:1.5) --++ ({\t+180}:3);
			\draw[red, ultra thick, <->] (30:3.5) ++ (\t:1.5) --++ ({\t+180}:3);
		}
		\end{tikzpicture}
		
		a plane-polarised wave
		
		(wave vibrations in one direction only)
	\end{minipage}
\end{figure}

\cmt polarisation\footnote[][-10cm]{Apart from plane polarisation where the vibrations are fixed in a single direction, there are also \emph{circular} or \emph{elliptical} polarisation, for which the direction of vibrations \emph{rotate} in a plane as the wave travels.} is a phenomenon associated with \emph{transverse} waves only

longitudinal waves (such as sound waves) cannot be polarised

this is because longitudinal vibrations are always parallel to direction of energy transfer

\cmt most important examples are polarisation of \emph{electromagnetic waves}

regarding vibration of electromagnetic waves, we usually focus on variations of electric field

\cmt sunlight, light emitted from incandescent bulbs are unpolarised

they are emitted randomly from their sources, so they contain all planes of polarisation


\cmt a few of the great many applications of polarisation are:

\begin{compactitem}
	\item[--] polarising sunglasses \& photography \sidenote[][-10cm]{When unpolarised light is reflected at the boundary of two media, the reflected beam would become partially polarised. Polarising sunglasses use this effect to reduce the sunlight reflected from shiny surfaces, such as a lake (for sightseeing tourists) or road surfaces (for drivers). For the same reason, photographers often place a polarising filter in front of the camera lens to darken skies and increase the contrast.}
	
	\item[--] polarised 3D films\footnote[][-6.5cm]{The eyeglasses worn by viewers contain a pair of polarising filters with mutually perpendicular axes. When two images are projected onto the same screen, each filter restricts the light reaching each eye by passing only one of the images that is polarised in the same direction. This creates an 3D illusion.}
	
	\item[--] liquid-crystal display (LCD) technology\footnote[][-4cm]{Liquid crystal arrays can be realigned by applying an electric field. This causes the axis of polarization to rotate, hence backlight may be allowed to pass or be blocked, forming dark patterns of the display.}
	
	\item[--] radio transmission and reception\footnote[][-2cm]{Since electric currents flow in certain directions only, so the radio waves and microwaves produced from aerials are intrinsically polarised. This means the reception of signals would be sensitive to a particular direction, but totally insensitive to the normal direction. This explains why altering the orientation of antennas could greatly enhance the quality of reception.}
\end{compactitem}

\subsection{polarisers \& Malus's law}

a plane-polarised light can be produced from unpolarised light using a \keypoint{polariser}

a polariser blocks vibrations in all planes except the plane of polarisation\footnote[][-2cm]{What we are discussing here is just one type of polariser, often called a \emph{linear absorptive polariser}, which is basically a synthetic transparent plastic sheet made of certain crystals. The thin sheet contains long-chain organic molecules aligned parallel to each other. When an unpolarised light passes through the polariser, there is a strong absorption of the electric field parallel to the alignment of molecules, so vibration of the electric field is allowed to pass in one direction only.}

direction along which vibration is allowed to pass is called the axis of the polariser

\begin{figure*}[ht]
	\centering
	\begin{tikzpicture}[scale=0.65]
	\draw[very thick, blue, ->] (0,0) -- (10,0);	
	\draw[fill=cyan!30] (0,0) ++ (0,2) ++ (160:2) --++ (-20:4) --++ (0,-4) --++ (160:4) --++ (0,4);
	\foreach \t in {-1.5,-1,...,1.5} \draw[very thin, gray!80, decorate, decoration={snake, amplitude=0.6}] (0,0) ++ (0,\t) ++ (160:1.6) --++ (-20:3.2);
	\draw[thick, Green, dashed] (0,1.8) -- (0,-1.8);
	\draw[Green] (0.1,1.2) -- (1.6,3) node[black, right, twoline]{axis of\\transmission};
	\draw[gray!80] (-1.4,1.6) -- (-3,3) node[black, left, twoline]{direction of alignment\\of molecules};
	\draw[very thick, blue] (-10,0) -- (0,0);
	\foreach \t in {-20,40,90,125} {
		\draw[red, ultra thick, <->] (-4,0) ++ (\t:1.5) --++ ({\t+180}:3);
		\draw[red, ultra thick, <->] (-8,0) ++ (\t:1.5) --++ ({\t+180}:3);
	}
	\foreach \t in {4,8} \draw[red, ultra thick, <->] (\t,-1.5) --++ (0,3);
	\node at (-6,-3.2) {unpolarised light};
	\node at (6,-3.2) {plane-polarised light};
	\node at (0,-3.2) {polariser};
	\end{tikzpicture}
	
	\caption{unpolarised light goes through a polariser and becomes plane-polarised}
\end{figure*}

\cmt polarised light can be sent through a second polarising filter (sometimes called an \emph{analyser})


\begin{figure*}[ht]
	\centering
	\begin{minipage}{0.45\textwidth}
		\centering
		\begin{tikzpicture}[scale=0.4]
		\draw[very thick, blue, ->] (3,0) -- (8.4,0);	
		\draw[fill=orange!30] (3,0) ++ (0,2) ++ (160:2) --++ (-20:4) --++ (0,-4) --++ (160:4) --++ (0,4);
		\foreach \t in {-1.5,-1,...,1.5} \draw[very thin, gray!80, decorate, decoration={snake, amplitude=0.6}] (3,0) ++ (0,1.5) ++ (160:\t) --++ (0,-3);	
		\draw[very thick, blue] (-3,0) -- (3,0);	
		\draw[fill=cyan!30] (-3,0) ++ (0,2) ++ (160:2) --++ (-20:4) --++ (0,-4) --++ (160:4) --++ (0,4);
		\foreach \t in {-1.5,-1,...,1.5} \draw[very thin, gray!80, decorate, decoration={snake, amplitude=0.6}] (-3,0) ++ (0,\t) ++ (160:1.6) --++ (-20:3.2);	
		\draw[very thick, blue] (-8.4,0) -- (-3,0);
		\foreach \t in {-20,40,90,125} {
			\draw[red, very thick, <->] (-6.5,0) ++ (\t:1.5) --++ ({\t+180}:3);
		}
		\foreach \t in {0} \draw[red, very thick, <->] (\t,-1.5) --++ (0,3);
		\node at (-3,-3.5) {polariser};
		\node at (3,-3.5) {analyser};
		\end{tikzpicture}
		
		no transmission of light if polariser and 
		
		analyser are aligned at right angles
	\end{minipage}\hfil
	\begin{minipage}{0.45\textwidth}
		\centering
		\begin{tikzpicture}[scale=0.4]
		\draw[very thick, blue, ->] (3,0) -- (8.4,0);	
		\draw[fill=orange!30] (3,0) ++ (0,2) ++ (160:2) --++ (-20:4) --++ (0,-4) --++ (160:4) --++ (0,4);
		\foreach \t in {-1.5,-1,...,1.5} \draw[very thin, gray!80, decorate, decoration={snake, amplitude=0.6}] (3,0) ++ (0,\t) ++ (160:1.6) --++ (-20:3.2);	
		\draw[very thick, blue] (-3,0) -- (3,0);	
		\draw[fill=cyan!30] (-3,0) ++ (0,2) ++ (160:2) --++ (-20:4) --++ (0,-4) --++ (160:4) --++ (0,4);
		\foreach \t in {-1.5,-1,...,1.5} \draw[very thin, gray!80, decorate, decoration={snake, amplitude=0.6}] (-3,0) ++ (0,\t) ++ (160:1.6) --++ (-20:3.2);	
		\draw[very thick, blue] (-8.4,0) -- (-3,0);
		\foreach \t in {-20,40,90,125} {
			\draw[red, very thick, <->] (-6.5,0) ++ (\t:1.5) --++ ({\t+180}:3);
		}
		\foreach \t in {0,6.5} \draw[red, very thick, <->] (\t,-1.5) --++ (0,3);
		\node at (-3,-3.5) {polariser};
		\node at (3,-3.5) {analyser};
		\end{tikzpicture}
		
		polarised light is unaffected if polariser
		
		and analyser are aligned in parallel
	\end{minipage}
\end{figure*}

\cmt if polarised light of initial intensity $I_0$ has an angle $\theta$ to the axis of analyser

\keypoint{Malus's law}\index{Malus's law} states that transmitted intensity is given by: ${I = I_0 \cos^2 \theta}$

this is because only electric field parallel to axis of analyser is transmitted

so transmitted amplitude\footnote{This is actually the amplitude of the electric field strength.} satisfies: $A = A_0 \cos\theta$, where $A_0$ is the initial amplitude

recall that intensity is proportional to square of amplitude, so $I = I_0 \cos^2 \theta $

\example{A beam of light polarised in the vertical direction has an amplitude $A$ and intensity $I$. It passes through a polarising filter whose axis of polarisation is at $45^\circ$ to the vertical. (a) What is the amplitude and the direction of polarisation of the emerging beam? (b) If the emerging beam then enters another filter whose axis of polarisation is at $75^\circ$ to the vertical, what is the amplitude and intensity of the emerging beam?}

\begin{soln} through first filter: 
$A_1 = A \cos\theta_1 = A \cos 45^\circ \RA A_1 = \frac{1}{\sqrt{2}} A $
light is polarised in a direction at $60^\circ$ to the vertical

through second filter:
$A_2 = A_1 \cos\theta_2 = \frac{1}{\sqrt{2}} A \times \cos (75^\circ - 45^\circ) \RA A_2 = \frac{\sqrt{6}}{4} A $

$I_2 = \left(\frac{\sqrt{6}}{4} \right)^2 I \RA I_2 = \frac{3}{8} I$

 or, $I_2 = I \cos^2\theta_1 \cos^2 \theta_2^2 = I \times \cos^2 45^\circ \times \cos^2 (75^\circ-45^\circ) \RA I_2 = \frac{3}{8} I$ 

\end{soln}





\subsection{diffraction}

a wave has the ability to bend around obstacles or pass through narrow gaps

\begin{ilight}
	when a wave passes through a aperture/slit/gap/hole or encounters an obstacle, it spreads out/around the corners, this is known as \keypoint{diffraction}\index{diffraction} of waves
\end{ilight}

\cmt diffraction is a general property of all waves

some examples of diffraction are

\begin{compactitem}
	\item[--] sound waves can diffract through an open door
	
	so you can hear people in the next room talking, even though you cannot see them
	
	\item[--] light ray can bend when it goes through/around a small hole/small particles
	
	sky appears red at sunset as red light is diffracted most by dust particles in atmosphere
\end{compactitem}

\cmt a wave is diffracted most when its wavelength is close to size of aperture/obstacle

a wave with longer wavelength usually diffracts more than a wave of shorter wavelength

\begin{figure*}[htp]
	\centering
	\begin{tikzpicture}[scale=1.25]
	\draw[ultra thick] (0,2.4) -- (0,0.25) (0,-0.25) -- (0,-2.4);
	\foreach \x in {0.25,0.75,1.25,1.75,2.25}{
		\draw[blue,thick] (-\x,-2.4) -- (-\x,2.4);
		\draw[blue,thick] (-65:\x) arc(-65:65:\x);
	}
	\draw[ultra thick] (6,2.4) -- (6,0.8) (6,-0.8) -- (6,-2.4);
	\foreach \x in {0.25,0.75,1.25,1.75,2.25}{
		\draw[blue,thick] (-\x+6,-2.4) -- (-\x+6,2.4);
		\draw[blue,thick] (6,-0.8) ++ (-25:\x) to[out=65,in=270] (6+\x,-0.8-\x/6) -- (6+\x,0.8+\x/6) to[out=90,in=-65] (6+0.9063*\x,0.8+0.4226*\x);
	}
	\end{tikzpicture}
	
	\caption{diffraction of a wave as it passes through an aperture of width that is (a) close to the wavelength, (b) greater than the wavelength}
\end{figure*}

\newpage


\subsection{Doppler effect}

\begin{ilight}
	relative motion between wave source and the observer causes a change in observed frequency, this is known as the \keypoint{Doppler effect}\index{Doppler effect}
\end{ilight}

\cmt any type of wave can exhibit Doppler effect

\cmt when wave source moves towards observer, a higher frequency is observed

when wave source moves away from observer, a lower frequency is observed

\cmt Doppler effect also occurs if source is at rest but observer is in motion

as long as there is \emph{radial} motion between source and observer, there is shift in frequency

\cmt Doppler effect finds its use in many areas, some examples are

\begin{compactitem}
	\item[--] \emph{Doppler radars}: used to measure velocity of moving target (speeding cars, tennis balls, etc.)
	
	\item[--] \emph{Doppler ultrasonography}: used to image blood flow in human bodies
	
	\item[--] \emph{astronomy}: used to study motion of stars and galaxies
\end{compactitem}

\newpage

\subsection*{explanation for Doppler effect}

suppose wave source is moving towards the observer at speed $v_s$

at time $t=0$, the source emits a wavefront which travels at speed $v$

\begin{figure*}[ht]
	\centering
	\begin{tikzpicture}[scale=0.95]
	\draw[thick] (0,0) circle (0.2);
	\draw[fill] (12,0.4) circle (0.2);
	\draw[ultra thick] (12,0.2) -- (12,-0.1) -- (11.8,-0.6) (12,-0.1) -- (12.2,-0.6) (11.75,-0.1) -- (12,0.15) -- (12.25,-0.1);
	\draw[blue,thick] (0.4, 0.62513) arc (10:-10:3.6);
	\node at (0,-1) {source};
	\node at (12,-1) {observer};
	\end{tikzpicture}
\end{figure*}

after one period, wavefront travels forward by a distance of $vT$

at same time, source moves forward by a distance of $v_s T$ and emits a new wavefront

\begin{figure*}[ht]
	\centering
	\begin{tikzpicture}[scale=0.95]
	\draw[thick,gray!80] (0,0) circle (0.2);
	\draw[thick] (2,0) circle (0.2);
	\draw[gray, ->] (0.3,0) -- (1.7,0);
	\draw[fill] (12,0.4) circle (0.2);
	\draw[ultra thick] (12,0.2) -- (12,-0.1) -- (11.8,-0.6) (12,-0.1) -- (12.2,-0.6) (11.75,-0.1) -- (12,0.15) -- (12.25,-0.1);
	\draw[Green,thick] (2.4, 0.62513) arc (10:-10:3.6);
	\draw[blue,thick] (6.4, 0.62513) arc (10:-10:3.6);
	\draw[gray!80,dashed,thick] (0.4, 0.62513) arc (10:-10:3.6);
	\node at (0,-1) {source};
	\node at (12,-1) {observer};
	\draw[<->] (2.55,-1.5) -- (6.5,-1.5) node[above, midway]{$\lambda'$};
	\draw[<->] (0.5,-1.5) -- (2.45,-1.5) node[above, midway]{$v_s T$};
	\draw[<->] (0.5,-2.4) -- (6.5,-2.4) node[above, midway]{$\lambda = vT$};
	\end{tikzpicture}
\end{figure*}

if source is at rest, observer simply perceives a wavelength $\lambda = vT$

now source is moving closer, a shorter wavelength $\lambda'$ is perceived

since frequency is inversely proportional to wavelength, so higher frequency $f'$ is observed

similar discussion follow for the case where source moves away from observer

\cmt change in observed frequency is due to change in wavelength caused by relative motion

if source moves towards/away from observer, apparent wavelength becomes shorter/longer

\subsection*{Doppler effect equation}

we are now ready to derive an equation for the shift of observed frequency

quantitatively, we can write: $\lambda' = \lambda \mp v_s T \quad$ ("$-$"/"$+$" if source is moving closer/away)

using $v = \lambda f$ and $f=\frac{1}{T}$, this becomes: $ \frac{v}{f'} = \frac{v}{f} \mp \frac{v_s}{f} $

rearranging, one finds observed frequency is given by: $\boxed{f' = f\frac{v}{v \mp v_s}}$

\example{A police car moving towards you at $16 \mps$ sirens at 500 Hz. Given that the speed of sound in air is $340\mps$, at what frequency do you hear the siren?}

\solc\begin{equation*}
	f' = f\frac{v}{v - v_s} = 500 \times \frac{340}{340-16} \RA f' \approx 525 \text{ Hz} 
\end{equation*}

\example{A star emits an $H_\alpha$ line of wavelength 656 nm. (a) What is the frequency of this $H_\alpha$ line? (b) An observer on earth detects a wavelength of 680 nm, what can we say about motion of the star? (c) Find the relative speed of the star with respect to the earth.}

\sol original frequency of $H_\alpha$: $f = \frac{c}{\lambda} = \frac{3.00\times10^8}{656 \times 10^{-9}} \approx 4.57 \times 10^{14} \text{ Hz}$

observed wavelength is longer (redshift), this means the star is moving away, or receding

to find speed of star, we can consider observed frequency: $f' = \frac{c}{\lambda'} = f\frac{c}{c + v_s}$
\begin{equation*}
	\frac{3.00\times10^8}{680\times10^{-9}} = 4.57 \times 10^{14} \times \frac{3.00\times10^8}{3.00\times10^8 + v_s} \RA v_s \approx 1.10 \times 10^7 \mps
\end{equation*}
	
alternatively, we can consider observed wavelength: $\lambda' = \lambda + v_s T$, or: $\Delta \lambda = \lambda' - \lambda = \frac{v_s}{f}$
\begin{equation*}
	(680-656)\times10^{-9} = \frac{v_s}{4.57 \times 10^{14}} \RA  v_s \approx 1.10 \times 10^7 \mps 
\end{equation*}



% \chapter{Waves} \label{chapter:Waves}


% \section{Superposition}\syl{state, explain and use the principle of superposition}
% When two or more waves arrive at the same place at the same time they combine with each other according to the principle of superposition.

% % \begin{marginfigure}\superpositionpic{60}{0}{1.5}{1.5}\caption{Constructive interference}\end{marginfigure}

% This states that the resultant displacement at any point is equal to the sum of the displacements of the individual waves.  Note that a displacement may be negative so addition of two displacements may result in a resultant, which is smaller than either of the original displacements.//
% When two waves of equal amplitude and wavelength meet in phase they interfere constructively to produce a new wave of twice the amplitude but the same wavelength\syl{state the conditions necessary for two-source interference to be
% observed, i.e. constant phase difference, vibrations in the same line}
% .

% When two waves of equal amplitude and wavelength meet in antiphase they interfere destructively to cancel each other out.
% % \begin{marginfigure}\superpositionpictwo{60}{180}{1.5}{-1.5}\caption{destructive interference}\end{marginfigure}

% When two waves of equal amplitude and wavelength meet with any other phase difference the resultant will have an intermediate amplitude.
% % \begin{marginfigure}{\superpositionpic {60}{60}{1.5}{1.5}}\caption{Two waves, out of phase}\end{marginfigure}

% \section{Coherence}

% Coherence\syl{give examples of coherent and incoherent sources,} is an essential condition for the interference of waves.  Two sources of waves are described as coherent if they emit waves with a constant phase difference (note that the waves do not necessarily have to be in phase).  Two waves arriving at a point are said to be coherent if there is a constant phase difference between them as they pass that point.

% \section{Interference} Interference cannot be observed with light - or any form of electromagnetic radiation - from two separate light sources.  This is because the phase difference between light waves from the two sources changes randomly so the points of cancellation and reinforcement move about at random.  The two sets of waves need to be produced from a single source, either:
% \begin{itemize} \item by dividing the wavefront and arranging for the divided wavefronts to overlap, as in the double slit experiment (see below) or \item dividing the amplitude of the waves using a partial reflector and arranging for the separated waves to overlap, as in thin film interference.
% \end{itemize}
% \subsection{Path difference and Interference} 
% \newif\ifstartcompletesineup
% \newif\ifendcompletesineup
% \pgfkeys{
%     /pgf/decoration/.cd,
%     start up/.is if=startcompletesineup,
%     start up=true,
%     start up/.default=true,
%     start down/.style={/pgf/decoration/start up=false},
%     end up/.is if=endcompletesineup,
%     end up=true,
%     end up/.default=true,
%     end down/.style={/pgf/decoration/end up=false}
% }
% \pgfdeclaredecoration{complete sines}{initial}
% {
%     \state{initial}[
%         width=+0pt,
%         next state=upsine,
%         persistent precomputation={
%             \ifstartcompletesineup
%                 \pgfkeys{/pgf/decoration automaton/next state=upsine}
%                 \ifendcompletesineup
%                     \pgfmathsetmacro\matchinglength{
%                         0.5*\pgfdecoratedinputsegmentlength / (ceil(0.5* \pgfdecoratedinputsegmentlength / \pgfdecorationsegmentlength) )
%                     }
%                 \else
%                     \pgfmathsetmacro\matchinglength{
%                         0.5 * \pgfdecoratedinputsegmentlength / (ceil(0.5 * \pgfdecoratedinputsegmentlength / \pgfdecorationsegmentlength ) - 0.499)
%                     }
%                 \fi
%             \else
%                 \pgfkeys{/pgf/decoration automaton/next state=downsine}
%                 \ifendcompletesineup
%                     \pgfmathsetmacro\matchinglength{
%                         0.5* \pgfdecoratedinputsegmentlength / (ceil(0.5 * \pgfdecoratedinputsegmentlength / \pgfdecorationsegmentlength ) - 0.4999)
%                     }
%                 \else
%                     \pgfmathsetmacro\matchinglength{
%                         0.5 * \pgfdecoratedinputsegmentlength / (ceil(0.5 * \pgfdecoratedinputsegmentlength / \pgfdecorationsegmentlength ) )
%                     }
%                 \fi
%             \fi
%             \setlength{\pgfdecorationsegmentlength}{\matchinglength pt}
%         }] {}
%     \state{downsine}[width=\pgfdecorationsegmentlength,next state=upsine]{
%         \pgfpathsine{\pgfpoint{0.5\pgfdecorationsegmentlength}{0.5\pgfdecorationsegmentamplitude}}
%         \pgfpathcosine{\pgfpoint{0.5\pgfdecorationsegmentlength}{-0.5\pgfdecorationsegmentamplitude}}
%     }
%     \state{upsine}[width=\pgfdecorationsegmentlength,next state=downsine]{
%         \pgfpathsine{\pgfpoint{0.5\pgfdecorationsegmentlength}{-0.5\pgfdecorationsegmentamplitude}}
%         \pgfpathcosine{\pgfpoint{0.5\pgfdecorationsegmentlength}{0.5\pgfdecorationsegmentamplitude}}
% }
%     \state{final}{}
% }



% One way of demonstrating interference effects is to take a single source of waves (s) and allow the waves to travel by two different paths to a detector (D).

% \begin{figure*}
% \begin{center}
% \begin{tikzpicture}

% \tikzset{interface/.style={
%         % The border decoration is a path replacing decorator. 
%         % For the interface style we want to draw the original path.
%         % The postaction option is therefore used to ensure that the
%         % border decoration is drawn *after* the original path.
%         postaction={draw,decorate,decoration={border,angle=-45,
%                     amplitude=0.3cm,segment length=2mm}}}}
                    
% \tikzset{waveo/.style={decorate, draw=red,
%         decoration={complete sines,amplitude=20pt, segment length=10pt, start up, end down}}}
% \tikzset{wavet/.style={decorate, draw=blue,
%         decoration={complete sines,amplitude=20pt,start up, end up, segment length=10pt}}}
%         \tikzset{wavetr/.style={decorate, draw=blue,
%         decoration={complete sines,amplitude=20pt,start up, end down, segment length=10pt}}}
% \draw[dashed] (0,0)-- (5,-4)-- (10,0)node[anchor=west,text width=4cm]{Detector, D};
% \draw[waveo] node [anchor=east]{Source, S}(0,0) -- node[midway, fill=white]{Path $\vec{SR}$}(4.7,-4) (5.3,-4)  node[anchor=south, fill=white]{Mirror, R}-- (10,0)node[midway, sloped, fill=white]{Path $\vec{RD}$};
% \draw[wavet] (0,0)-- (10,0);
% \draw[thick] (10,2)--(10,-3);
% \draw[dashed] (0,0)-- (10,0) node[midway,sloped, fill=white]{Path $\vec{SD}$};
% \draw[thick, interface](4,-4)--(6,-4);
% \end{tikzpicture}
% \caption{Illustrating path difference.}
% \end{center}

% \end{figure*}

% This can be done as shown below by having one path directly from S to D and another path via a reflector (R).  The path difference\syl{show an understanding of path difference, phase difference, and
% coherence} is the difference between the distances travelled, SD and SRD.

% As R is moved perpendicular to the direction $\vec{SD}$ a point will be reached where the direct wave and the reflected wave arrive at D in phase.  At this point constructive interference occurs at D and a maximum amplitude is detected.  (Note that this occurs when the path difference has a value like $\frac{1}{2}$ or $\frac{3}{2}$ wavelengths as a phase change of $\frac{1}{2}$ cycle occurs on reflection.)

% If R is moved away from the line $\vec{SD}$ the path difference increases and a point is reached where the direct and reflected waves arrive in antiphase.   A minimum amplitude is now detected.  The amplitude will not in general be zero as the amplitude of the reflected wave will be less than that for the direct wave as it has travelled further so complete cancellation will not occur.

% Further movement of R away from $\vec{SD}$ will restore the maximum amplitude as the two waves will again be in phase.  The increase in the length $\vec{SRD}$ that occurs in moving from one maximum to the next is one wavelength and this provides a method for measuring wavelength.

% \section{Single slit diffraction}
% In general it is considered that light only travels in straight lines - evidence for example is the production of shadows. When water waves pass through a gap or past an object they spread into the "shadow" region as shown below (A). This is known as "diffraction". The width of the opening controls the degree of diffraction. A narrow opening causes more diffraction (B).

% \subsection{Diffraction through a single narrow slit:}

% If  the slit is more than a few millimetres wide, a sharp rectangular shadow is formed (C). 
% If the slit is narrowed (D) the pattern begins to change - a wide central band is observed with a series of bright fringes either side. The central band is twice the width of the outer fringes. As the slit approaches its narrowest, the pattern broadens and becomes less bright (due to cutting down of light able to pass through slit).

% If white light is used the central band is white, however, the fringes are coloured.
% If monochromatic (single wavelength/colour) light is used the central band and fringes are the same colour. Red light has a more spread out pattern than blue - it also has a larger wavelength than blue - diffraction increases with wavelength.

% Intensity distribution of pattern\syl{recall the shape of the intensity pattern from a single slit and its effect
% on double-slit and diffraction grating patterns}:
% The diagram below (fig: \ref{single}) shows a typical diffraction pattern obtained when monochromatic light passes through a single narrow slit. Also plotted is the intensity of the pattern.

% \begin{figure}
% \label{single}
% \includegraphics[]{Single.png}
% \caption{Single slit diffraction pattern and geometry. }
% \end{figure}

% The central peak has the greatest intensity, with decreasing intensity with each fringe either side. Each band/fringe has its greatest intensity in its centre. This is the point where light is most likely to fall in its range.
% The fringes are caused by the edges of the slit acting like two wave sources, the waves from the two wave "sources" interfere superimposing an interference pattern over the diffraction pattern.	 

% \section{Youngs experiment} 
% The modern double-slit experiment is a demonstration that light and matter can display characteristics of both classically defined waves and particles; moreover, it displays the fundamentally probabilistic nature of quantum mechanical phenomena. A simpler form of the double-slit experiment was performed originally by Thomas Young in 1801. He believed it demonstrated that the wave theory of light was correct and his experiment is sometimes referred to as Young's experiment or Young's slits. The experiment belongs to a general class of "double path" experiments, in which a wave is split into two separate waves that later combine into a single wave. Changes in the path lengths of both waves result in a phase shift, creating an interference pattern.
% \begin{figure*}
% \begin{tikzpicture}
% \tikzset{waveo/.style={decorate, draw=red,
%         decoration={complete sines,amplitude=20pt, segment length=10pt, start up, end up}}}
% \tikzset{wavet/.style={decorate, draw=blue,
%         decoration={complete sines,amplitude=20pt,start up, end up, segment length=10pt}}}
%         \tikzset{wavetr/.style={decorate, draw=blue,
%         decoration={complete sines,amplitude=20pt,start up, end down, segment length=10pt}}}
% \draw (0,0) -- (0,-.5);
% \draw (0,0) -- (0,.5);
% \draw (0,.8)-- (0,1.5);
% \draw (0,-.8)-- (0,-1.5);
% \draw[dashed] node [anchor=north east]{A} (0,.65)-- (10,0)node[anchor=west,text width=4cm]{Constructive interference. Path lengths $=n\lambda$};
% \draw[waveo] (0,.65)-- (10,0);
% \draw[waveo] (0,-.65)-- (10,0);
% \draw[thick] (10,2)--(10,-3);
% \draw[dashed] (0,.65)-- (10,-2);
% \draw[dashed] node [anchor=south east]{B}(0,-.65)-- (10,-2);
% \draw[wavetr] (0,.65)-- (10,-2) node[anchor=west, text width=4cm]{Destructive interference. Path lengths differ by $\frac{n}{2}\lambda$};
% \draw[wavet] (0,-.65)-- (10,-2);
% \draw[dashed] (0,-.65)-- (10,0) node[midway, fill=white]{Path};

% \end{tikzpicture}
% \caption{Constructive and destructive interference as a result of path differences.}
% \end{figure*}

% The position of the bright and dark fringes is determined by the geometry of the experimental set-up.The distance from center to bright spot n is:
%  \[y = \frac{n \lambda D}{d}\]
% We know that the maxima will be at angles such that the path difference will be whole number of wavelengths. The geometric derivation is shown below.\syl{describe an experimental demonstration of two-source interference for light, appreciating the historical importance of Young's experiment, and be familiar with experiments which demonstrate two source interference for water waves, sound waves and microwaves;} 

% \begin{figure}[!ht]
% \includegraphics[]{twin.PNG}
% \caption{Beam pattern and geometric workings for Young's slits}
% \end{figure}

% Note that this makes use of the small angle approximation, a useful simplification of the basic trigonometric functions which is approximately true in the limit where the angle approaches zero.\begin{marginfigure}
% \includegraphics[]{small.png}
% \caption{Illustrating the small angle approximation}
% \end{marginfigure} They are truncations of the Taylor series for the basic trigonometric functions to a second-order approximation.

% \section{Diffraction grating}

% A diffraction grating is usually a piece of glass or plastic with closely-spaced lines on it. A transmission grating has thin clear spaces/slits where light can pass through, while a reflection gratings have shiny surfaces between the lines that reflect light.
% The tracks on a compact disc are very close together - the shiny gaps between the lines act as a reflection grating, leading to a colourful light spectrum being formed when light reflects off the surface.
% Transmission gratings are used to measure wavelengths of light to a very high degree of accuracy (spectrometer) - all substances have characteristic absorption/emission spectra.

% \subsection{Explaining diffraction gratings:}

% The individual lines act like wave sources or thin slits - they produce diffracting waves which will then interfere. Destructive interference leads to cancellation in nearly all directions - there are a few very clearly defined directions where constructive interference occurs, leading to bright fringes. The direction of these fringes depends on the wavelength of light - the wavelength can thus be determined!

% \paragraph{Diffraction grating formula} 

% Parallel rays of monochromatic light of wavelength ``$\lambda$'' fall on a diffraction grating with slit separation/grating spacing ``d''.

% \begin{marginfigure}
% \includegraphics[scale=.5]{grating.jpg}
% \end{marginfigure}

% If the grating has $N$ lines per metre, the grating spacing is given by; $$d=\frac{1}{N}$$\syl{use the equation $d sin\theta = n\lambda$ for a diffraction grating}

% For constructive interference light from adjacent slits must be in phase (with light from all slits/lines).

% $\Delta l= d sin \theta$, where $\theta$ is the angle of diffraction

% Hence $$d sin\theta = n\lambda$$


% \section{Standing Waves}
% Standing or stationary waves\syl{understand that a stationary wave can be regarded as a superposition of two progressive waves of equal amplitude and frequency,travelling in opposite directions and that the internodal distance is $\frac{\lambda}{2}$}
%  are a superposition effect, which occurs when two identical waves (implying waves with a similar amplitude and frequency) travelling in opposite directions interfere.  The most common example of this is where incident and reflected waves interfere, for example on a spring or in a pipe in a musical instrument.

% With a string, spring or pipe of fixed length there will be certain frequencies where a standing wave pattern is produced as shown below.  These frequencies will be those which give an integral number of half wavelengths within the fixed length.  The standing wave is produced by interference of the waves travelling in both directions caused by reflection at the ends. See the diagram on the right.  

% It is characterised by nodes, which are points of zero displacement, and antinodes, which are points of maximum displacement.  The first few possible patterns for transverse waves on a string or spring are shown below.
% \begin{figure*}
% \begin{tikzpicture}[x={(1cm,0.5cm)},z={(0cm,1cm)},y={(1cm,-0.2cm)}]

%         %repere
%         \draw[->] (0,-pi,0) --++ (6,0,0) node[above right] {Frequency};
%         \draw[->] (0,-pi,0) --++ (0,6.5,0) node[right] {Time};
%         \draw[->] (0,-pi,0) --++ (0,0,1.5);        
%             %sinusoides
%             \draw[blue] plot[domain = -pi:+pi, samples = 300] 
%             (1,\x,{0.3*cos(10*.5/10*(\x) r)});
%             \draw[blue] (1,-pi-0.15,0) node [left]{$f_{0}$};
           
%             \draw[blue,dashed] plot[domain = -pi:+pi, samples = 300] 
%             (1,\x,{-0.3*cos(10*.5/10*(\x) r)});
%             \draw[blue] (0,-pi-0.15,0);
            
%             \draw[blue] plot[domain = -pi:+pi, samples = 300] 
%             (3,\x,{0.3*sin(10*1/10*(\x) r)});
%             \draw[blue] (3,-pi-0.15,0) node [left]{$f_{1}$};
%              \draw[blue,dashed] plot[domain = -pi:+pi, samples = 300] 
%             (3,\x,{-0.3*sin(10*1/10*(\x) r)});
%             \draw[blue] (3,-pi-0.15,0);
            
%             \draw[blue] plot[domain = -pi:+pi, samples = 300] 
%             (5,\x,{0.3*cos(10*1.5/10*(\x) r)});
%             \draw[blue] (5,-pi-0.15,0) node [left]{$f_{2}$};
%             \draw[blue,dashed] plot[domain = -pi:+pi, samples = 300] 
%             (5,\x,{-0.3*cos(10*1.5/10*(\x) r)});
%             \draw[blue] (5,-pi-0.15,0);
             
       
%     \end{tikzpicture}
% \end{figure*}  
% The separation of adjacent nodes or antinodes is half a wavelength.  The different numbers of half wavelengths on the spring or string correspond to different allowed frequencies.  The lowest of these (on the left) is called the fundamental.  The others are whole number multiples of the fundamental frequency and are called harmonics.  Each subsequent frequency is known as an overtone. Because of the phase change on reflection at the fixed ends the ends are both nodes.

% The wavelength can be found by measuring the node separation and doubling it.  This gives a method for finding the speed of sound or electromagnetic waves if the frequency is known, using v = f.  For example, if a standing wave is set up using 1 GHz radio waves and the node separation is found to be 0.15 m the wavelength will be 0.30 m and the wave speed is given by:

% $$c = f \times \lambda= (1 \times 10^{9} Hz) \times  0.30 m = 3 \times 10^{8} ms^{-1}$$

% \section{Standing Waves in Musical Instruments}

% Stringed instruments such as violins and pianos produce standing waves as described above.  Hitting, plucking or bowing the string produces waves, which travel off in both directions.  The pitch of the fundamental and therefore of the harmonics may be increased by reducing the length of the string, by increasing the tension in it or by using a lighter string.

% Wind instruments work by setting up longitudinal standing waves in tubes.  Examples are organ pipes, trumpets and clarinets.

 
% If the tube is open at one end and closed at the other there will be a node at the closed end and an antinode at the open end.  At the fundamental frequency the length of the tube is therefore 1/4 of a wavelength.  Harmonics occur at frequencies where the length of the tube is 3/4, 5/4, 7/4 etc wavelengths.closed
% \begin{marginfigure}
% \includegraphics[]{closedpipe.png}
% \end{marginfigure}
% When any musical instrument is played the note you hear is at the fundamental frequency but a number of harmonics are also present.  The relative intensity of the different harmonics determines the quality of the note you hear.  It is this, which enables us to distinguish the same note played on different instruments.

% If the tube is open at both ends both ends have an antinode and the fundamental occurs when the tube length is half a wavelength.  Harmonics are at whole number multiples of the fundamental frequency. \begin{marginfigure}
% \includegraphics[]{openpipe.png}
% \end{marginfigure}

% \begin{figure}
% \includegraphics[]{pipes.jpg}
% \end{figure}
 
% \section{Refraction of Light}

% The Snell's Law equation is written in such a way as to emphasise symmetry. It is used to predict the change in direction that results from a change in velocity.\\
% \begin{figure*}
% \centering
% \begin{tikzpicture}[
%     media/.style={font={\footnotesize\sffamily}},
%     wave/.style={
%         decorate,decoration={snake,post length=1.4mm,amplitude=2mm,
%         segment length=2mm},thick},
%     interface/.style={
%         % The border decoration is a path replacing decorator. 
%         % For the interface style we want to draw the original path.
%         % The postaction option is therefore used to ensure that the
%         % border decoration is drawn *after* the original path.
%         postaction={draw,decorate,decoration={border,angle=-45,
%                     amplitude=0.3cm,segment length=2mm}}},
%    scale=1.2]
%     % Round rectangle
%     \fill[gray!10,rounded corners] (-4,-3) rectangle (4,0);
%     % Interface
%     \draw[blue,line width=.5pt,interface](-4,0)--(4,0);
%     % Vertical dashed line
%     \draw[dashed,gray](0,-3)--(0,3);
%     % Coordinates system
%     \draw(0,0.15)node[above]{$x$};
%     \draw[<->,line width=1pt] (1,0) node[above]{$y$}-|(0,-1) node[left]{$z$};
%     % Incidence
%     \draw[->,wave]
%          (150:4.5cm)--(150:3.5cm)node[right]{Incident ray};
%     \draw[gray](0:0cm)--(155:2cm);
%     \path (0,0)++(113:1cm)node{$\phi$};
%     \draw[->](0,0.75)arc(90:155:.75cm);
%     % Transmission
%     \draw[->,wave]
%          (-60:2.5cm)--(-60:3.5cm)node[right]{Refracted ray};
%     \draw[gray](0:0cm)--(-60:2cm);
%     \path (0,0)++(-70:1cm)node{$\theta$};
%     \draw[->] (0,-0.75) arc (-90:-60:.75cm);
  
%     % Media names
%     \path[media] (-3,.6)  node {media 1}
%                  (-3,-.6) node {media 2};

%     % $x$ axis
%     \filldraw[fill=white,line width=1pt](0,0)circle(.12cm);
%     \draw[line width=.6pt] (0,0)
%                           +(-135:.12cm) -- +(45:.12cm)
%                           +(-45:.12cm) -- +(135:.12cm);
%     % Interface pointer
%     \draw[-latex,thick](3.2,0.5)node[right]{$\mathsf{S_{1,2}}$}
%          to[out=180,in=90] (3,0);
         
%        \node at (0,-4){\begin{large}$\text{Snell's Law}=\frac{\sin\phi}{\sin\theta} = \frac{v_1}{v_2} = \frac{n_2}{n_1}$\end{large}};
%     % To-paths are really useful for drawing curved lines. The above
%     % to path is equal to:
%     %
%     % \draw[-latex,thick](3.2,0.5)node[right]{$\mathsf{S_{1,2}}$}
%     %      ..controls +(180:.2cm) and +(up:0.25cm) .. (3,0);
%     % Internally the to path is translated to a similar bezier curve,
%     % but the to path syntax hides the complexity from the user. 
% \end{tikzpicture}
% \caption{Refraction from low refractive index to higher refractive index (eg. Air to glass)}
% \end{figure*}
% Looking at the equation for snell's law we see that there must be a point at which the incident ray causes a refraction at, or above 90 degrees. This is called the \textbf{\textit{critical angle}}. At shallower angles of incidence, total internal reflection occurs.
% The only application of total internal reflection required to be learnt is the step-index multimode (`thick core') optical fibre. Step-index means simply that the core is glass of one refractive index and the cladding is glass of a lower index (students need to know why it has to be lower), with an abrupt change in index at the interface.
% \begin{figure*}
% \centering
% \begin{tikzpicture}[scale=1]
% \shade[left color=black!40!white,right color=black!10!white](0,-2.5)--(10,-2.5)--(10,2.5)--(0,2.5)--cycle;
% \draw[dashed, thin, red](-1,0)--(11,0);
% \draw[thin](0,-2.50)--(0,2.5);
% \draw(0,2)--(10,2);
% \draw(0,2.5)--(10,2.5);
% \draw(0,-2)--(10,-2);
% \draw(0,-2.5)--(10,-2.5);
% \draw[thick, blue,->](-1.8,-1)--(0,0)--(2,2)node(A){}--(6,-2)--(9.5,2);
% \draw[dashed, thin](2,1)--(2,3)node[]{};
% \node at (1,2.25){$n_2$};
% \node at (1,-1.75){$n_1$};
% \draw[->](2,1.4)arc(-90:-125:.75cm)node[anchor=north,midway]{$\theta_i$};
% \draw[->](2,2.75)arc(90:-40:.75cm)node[anchor=south,midway]{$\theta_r$};
% \node at (7,3.25){$n_1 \sin{\theta_i}=n_2 \sin{\theta_r}$};
% \node at (9,-1.75){\tiny Core};
% \node at (9,-2.35){\tiny Cladding};
% \begin{scope}
% \path[clip,decoration={random steps, segment length=6pt, amplitude=3pt},decorate] (9.9,-2.6)--(9.9,2.6)--(11,2.6)--(11,-2.6)-- cycle; 
% \shade[left color=white,right color=white](9.9,-3)--(9.9,3)--(11,3)--(11,-3)-- cycle;
% \end{scope}
% \end{tikzpicture}
% \caption{Total internal reflection in a single fibre}

% \end{figure*}
% While such fibres are fine for conveying light for illumination, students need to know that they can't be used for transmitting a rapid sequence of data over a long distance. Multimode dispersion is the problem.
% \begin{figure*}
% \centering
% \pgfmathdeclarefunction{gauss}{2}{%
%   \pgfmathparse{1/(#2*sqrt(2*pi))*exp(-((x-#1)^2)/(2*#2^2))}%
% }
% \begin{tikzpicture}[scale=1]

% \shade[left color=black!40!white,right color=black!10!white](0,-2.5)--(10,-2.5)--(10,2.5)--(0,2.5)--cycle;
% \draw[dashed, thin, red](-1,0)--(10,0);
% \draw[thin](0,-2.50)--(0,2.5);
% \draw(0,2)--(10,2);
% \draw(0,2.5)--(10,2.5);
% \draw(0,-2)--(10,-2);
% \draw(0,-2.5)--(10,-2.5);
% \draw[thick, blue,->](-1.8,-1)--(0,0)--(2,2)--(6,-2)--(9.5,2);
% \draw[thick, cyan,->](-2.2,-1)--(0,0)--(2.3,2)--(6.6,-2)--(10.9,2);
% \draw[dashed, thin](2,1)--(2,3)node[]{};
% \draw[dashed, thin](2.3,1)--(2.3,3)node[]{};
% \node at (1,2.25){$n_2$};
% \node at (1,-1.75){$n_1$};
% \node at (9,-1.75){\tiny Core};
% \node at (9,-2.35){\tiny Cladding};
% \node at (12,-3.25){\tiny{Broad pulse}};
% \node at (-3,-3.25){\tiny{Defined pulse}};
% \begin{scope}
% \path[clip,decoration={random steps, segment length=6pt, amplitude=3pt},decorate] (9.9,-2.6)--(9.9,2.6)--(11,2.6)--(11,-2.6)-- cycle; 
% \shade[left color=white,right color=white](9.9,-3)--(9.9,3)--(11,3)--(11,-3)-- cycle;
% \end{scope}

% \begin{axis}[scale=.5,
%     anchor=origin,
%     at={(-3cm,-3cm)},
%     style={samples=51,smooth},
%     hide axis
% ]
% \addplot[mark=none,thick] {gauss(0,.4)};
% \end{axis}
% \begin{axis}[scale=.5,
%     anchor=origin,
%     at={(12cm,-3cm)},
%     style={samples=51,smooth},
%     hide axis
% ]
% \addplot[mark=none,thick] {gauss(0,2)};
% \end{axis}

% \draw[thick, blue,->](-1.8,-1)--(0,0)--(2,2)--(6,-2)--(9.5,2);
% \draw[thick, cyan,->](-2.2,-1)--(0,0)--(2.3,2)--(6.6,-2)--(10.9,2);
% \end{tikzpicture}
% \caption{Multi-mode dispersion in a fibre resulting from the two different length paths}
% \end{figure*}

% Light travelling at an angle to the axis of the fibre will travel further for a given axial length of fibre, than light travelling parallel to, or at a smaller angle to, the axis, and so will arrive later. Thus the arrival time of an element of data encoded in the light is smeared out. The element could start to arrive (by the shortest route) earlier than the previous element has finsished arriving by its longest route. Even worse confusion can occur. 
% There are two ways round the problem of multimode dispersion\ldots
% \begin{itemize}
% \item Make fibres with graded index cores. This means cores which have a progressively lower index as we go out from the axis towards the interface with the cladding. The lower the index the faster the light travels so, if the grading is correctly calculated, the longer, more zigzaggy, paths cash in on the `faster medium' and take no longer than the short, axial, route. Clever stuff, but note that graded index fibres are not in the WJEC specification. 
% \item Make the core very thin. Its diameter must  be no more than a few wavelengths of the light (or infrared) being carried. Such fibres are monomode. Light travels parallel to the axis. There are no zigzag modes. Students are required to know this. You are not required to know why very thin fibres are monomode. This is just as well, because it cannot be shown by ray optics, nor even by simple application of Huygens Principle. Electromagnetic wave theory is needed.
% \end{itemize}
% The website\\ \url{http://www.techoptics.com/pages/Fiber\%20Optics\%20-\%20Optical\%20Fiber.html}   \\ 
% gives an excellent summary of fibres for data transmission, with some facts and figures.

% \section{Photons}

% Everything in nature seems to come in lumps or quanta (singular: quantum). For example, ordinary matter is made of atoms, and electric charge comes in units of e. This lumpiness was only becoming fully accepted a hundred years ago. But in 1905 Einstein made the bold suggestion that light, too, was `lumpy'. Light quanta are now called photons.
% A photon is a discrete packet of electromagnetic radiation energy. The energy of a photon is given by
% $$ E=hf $$
% In which $f$ is the frequency of the light and $h$ is a constant called Planck's constant. $h = 6.6 \times 10^{-34} \text{Js}$ \sidenote{For interest only \ldots The constant h had first arisen in the earlier (1900) work of Max Planck on the radiation inside a cavity with hot walls. Planck had shown that the energies of oscillating particles in the wall seemed to be quantised.}

% Einstein suggested some experiments in which the quantisation of light should reveal itself. The simplest to understand involved the photoelectric effect, a phenomenon known about since the late 1880s.

% \section{Laser Physics}
% Lasers, by now, are in nearly every household as they were once in every self-respecting science fiction film. Although, the death ray mystique will appeal to most A-level students there is real value in studying the physics which forms the foundation of laser construction. Unlike other subjects such as relativity or nuclear physics which can only be touched upon within an A-level course, the essential physics underpinning lasers can be taught reasonably well in a few lessons. Also, these fundamentals of lasers can be taught at the right level while not oversimplifying the subject.

% LASER is an acronym and stands for Light Amplification by Stimulated Emission of Radiation. Which leads us nicely onto what is stimulated emission?

% \subsection{The Three Important Atomic Processes}
% These three processes are:
% \begin{enumerate}
% \item Absorption of light
% \item Spontaneous emission of light
% \item Stimulated emission of light
% \end{enumerate}
% Absorption of light by an atom is shown in the diagram below - a photon of the correct energy is absorbed by the atom and an electron gains enough energy to move from the ground state to the excited state (Note: for the moment we are only considering the ground state and the first excited state only).

% \begin{figure*}
% \usetikzlibrary{shadings}
% \centering
% \begin{tikzpicture} [scale=.7, wave/.style={
%         decorate,decoration={snake,post length=1.4mm,amplitude=2mm,
%         segment length=2mm},thick},
%     interface/.style={
%         % The border decoration is a path replacing decorator. 
%         % For the interface style we want to draw the original path.
%         % The postaction option is therefore used to ensure that the
%         % border decoration is drawn *after* the original path.
%         postaction={draw,decorate,decoration={border,angle=-45,
%                     amplitude=0.3cm,segment length=5mm}}} ]
   
% \pgfdeclareverticalshading{rainbow}{100bp} 
%  {color(0bp)=(red); color(25bp)=(red); color(35bp)=(yellow); 
%   color(45bp)=(green); color(55bp)=(cyan); color(65bp)=(blue); 
%   color(75bp)=(violet); color(100bp)=(violet)}
%     \begin{scope}[]
%         \shade[shading=rainbow]
%             (0,0) rectangle (3,8); 
%     \draw[very thick,black] (0,3) -- (3,3);
%     \draw[very thick,black] (0,6.3) -- (3,6.3);
%      \end{scope}
% \draw[thick, black] (3.5,-1)--(10,-1);
% \draw[thick, black] (3.5,3)--(10,3);
% \draw[thick, black] (3.5,6.3)--(10,6.3);
% \draw[dashed,<-]  (5,6.3)--(5,-1);
% \draw[dashed,<-]  (7,3)--(7,-1);
% \draw[wave,blue,->](-2,8.5)--(5,3.5);
% \draw[wave,green,->](-1,6)--(7,1);
% \draw [fill=black] (5,-1) circle (2mm);
% \draw [fill=black] (7,-1) circle (2mm);
%     \end{tikzpicture}
%     \caption{Absorption of light. The photon energy has to correspond to a possible energy level transition for the energy to be absorbed by an electron.}
% \end{figure*}
 
% Spontaneous emission is the reverse process - an electron drops spontaneously (and randomly) from the excited state to the ground state and emits a photon of the same energy. These photons have random phase and random direction.

% \begin{figure*}
% \usetikzlibrary{shadings}
% \centering
% \begin{tikzpicture} [scale=.7, wave/.style={
%         decorate,decoration={snake,post length=1.4mm,amplitude=2mm,
%         segment length=2mm},thick},
%     interface/.style={
%         % The border decoration is a path replacing decorator. 
%         % For the interface style we want to draw the original path.
%         % The postaction option is therefore used to ensure that the
%         % border decoration is drawn *after* the original path.
%         postaction={draw,decorate,decoration={border,angle=-45,
%                     amplitude=0.3cm,segment length=5mm}}} ]
   
% \pgfdeclareverticalshading{rainbow}{100bp} 
%  {color(0bp)=(red); color(25bp)=(red); color(35bp)=(yellow); 
%   color(45bp)=(green); color(55bp)=(cyan); color(65bp)=(blue); 
%   color(75bp)=(violet); color(100bp)=(violet)}
%     \begin{scope}[]
%         \shade[shading=rainbow]
%             (0,0) rectangle (3,8); 
%     \draw[fill=black] (0,0) rectangle (3,3);
%     \draw[fill=black] (0,3.1)rectangle (3,5); 
% \draw[fill=black] (0,5.1)rectangle (3,8); \end{scope}
% \draw[thick, black] (3.5,0)--(10,0)node[anchor=west]{Ground State};
% \draw[thick, black] (3.5,3)--(10,3)node [anchor=west]{$2.32eV$};
% \draw[thick, black] (3.5,5)--(10,5) node [anchor=west] {$2.75eV$};
% \draw [fill=black] (5,5) circle (2mm);
% \draw [fill=black] (7,3) circle (2mm);
% \draw[dashed,->]  (5,5)--(5,0);
% \draw[dashed,->]  (7,3)--(7,0);
% \draw[wave,blue,->](5,2)--(8,2)node[anchor=west]{$\lambda=?$};
% \draw[wave,green,->](7,1)--(10,1)node[anchor=west]{$\lambda=?$};
 
%     \end{tikzpicture}
%     \caption{Emission spectra and corresponding energy levels. Note: the frequency and wavelength can be calculated from the change in energy of the electron using $E=hf$}
% \end{figure*}
% However, there is also a third process [which was originally proposed by Einstein in 1917]. This process is known as stimulated emission - an electron is `stimulated' to drop from its excited state by an incoming photon. 


% \begin{figure}
% \includegraphics[scale=.5]{stimem.png}
% \caption{Stimulated emission - an incoming photon of an energy matching the transition can cause an electron to drop to a lower energy level. In order to do this it must release a photon identical to the first.}
% \end{figure}


% The reason that the electron is stimulated to drop is that the incoming photon is an electromagnetic wave and its e-m field will exert an oscillating force on the excited electron. If the incoming photon is of the correct frequency, this oscillating force will cause the excited electron to drop and both photons will exit with the same frequency, phase and direction. Note: again, the incoming photon needs to be of the correct energy.

% \subsection{Inverting The Population}
% In order to get as much light out of a system as is possible we need to get as many atoms excited as is possible.\begin{marginfigure}
% \includegraphics[]{popin.jpg}
% \caption{Population inversion - when more electrons are in an excited state than in the ground state. This is the opposite of a `normal' situation}
% \end{marginfigure} Obviously, the more electrons we have in an excited state the more will drop and emit photons (either spontaneously or through stimulation). However, there is one serious problem that arises when we produce a lot of light - the very photons that we produce are the actual photons that can be absorbed (they have the correct energy to produce both effects). If we have photons being absorbed all the time then our laser beam isn't getting any stronger.

% Forget, for the moment about spontaneous emission (we are allowed to but we'll explain why later). When a photon arrives at an atom one of three things can happen:
% \begin{enumerate}
% \item It can pass by and do nothing.
% \item It can be absorbed (if the atom is in the ground state).
% \item It can cause stimulated emission (if the atom is in the excited state).
%  \end{enumerate}
% When it comes to producing a laser beam with a high intensity the three options above will have the following effect on the beam.
% \begin{enumerate}
% \item No change in the beam.
% \item Net loss of one photon from the beam.
% \item Net gain of one photon in the beam.
% \end{enumerate}
% We need to arrive at a situation where stimulated emission is more likely than absorption so that the laser beam increases in intensity. Since stimulated emission occurs if the electrons are in the upper level and absorption when electrons are in the lower level we need to get more electrons into the upper, excited level. This is called population inversion (or $N_{2} > N_{1}$ as stated in the syllabus, where $N_2$ and $N_1$ are the number of electrons in the excited state and the ground state respectively).

% Unfortunately, this goes against what happens in nature - lower energy levels are always more heavily populated than higher energy levels when we have thermal equilibrium. There's only one thing for it - get rid of this thermal equilibrium. How do we do this? We continue to pump energy into exciting electrons to higher energy levels to maintain a population inversion and to break the conditions of thermal equilibrium.

% Population inversion is not usually possible if we only have two energy levels (if pumping is carried out by light).
%  \newpage
% \subsection{The 3 Energy Level Laser System}

% \begin{figure}
% \includegraphics[]{lasers.jpg}
% \end{figure}
% \sidenote{Note:\begin{itemize} \item E3 (to E2) has to have a short lifetime because E3 cannot start to fill up - pumping won't then be possible. Also, we don't want the electrons to stay in E3 and have them stimulated to drop back to E1 by the pumping light - that's back to the 2-level system again which wasn't quite good enough.

% \item More than half the electrons from E1 must be pumped to E2 (via E3) in order to obtain a population inversion - that's a lot of electrons!\end{itemize}}
% \begin{enumerate}
% \item Pumping. Electrons are promoted from the ground state (E1) to E3 usually by using an external light source or by electron collisions.
% \item Electrons drop quickly (because E3 is chosen to have a short lifetime of the order of nanoseconds) to the metastable (E2). Calling E2 metastable means that it has a long lifetime and electrons stay there for a long time (not that long really around a millisecond but that's a very long time for an electron).
% \item This is the transition that produces the laser photons so we must have N2 > N1. Note that, although stimulated emission still reduces our population inversion, the pumping is at a different wavelength. We have to make sure that the pumping [1] exceeds the stimulated emission [3] to maintain a population inversion.
% \end{enumerate}

% \subsection{The 4 Energy Level Laser System}

% \begin{figure*}
% \includegraphics[]{fourlevel.jpg}
% \end{figure*}
% \begin{enumerate}
% \item Pumping again.
% \item Quick drop to the metastable state E2.
% \item This is the laser light producing transition so this time N2 > N1. However, because E0 is the ground state, E1 is practically empty initially so obtaining population inversion is far, far easier (definitely no need to pump half the electrons!).
% \item Another quick transition so E1 has a short lifetime. This is because we want E1 to be empty so that we have a population inversion (if N1 is small it's easier for N2 to be larger than N1).							
%  \end{enumerate}
% \subsection{Laser Construction}

% \begin{figure*}
% \includegraphics[scale=.5]{lasercav.PNG}
% \end{figure*}

% In order to ensure that the laser produces light of a high enough intensity, the above set up is used. The amplifying medium is the region where the population inversion exists. This means that the conditions are right in the amplifying medium for stimulated emission. Under these conditions one photon has the potential to produce two photons and these can produce 4 photons, then 8 photons etc. Like a chain reaction, this process will lead to an exponential increase in output energy. Laser physicists aren't happy with this, they go even further - they use mirrors to ensure that this exponential increase happens many times. Because only 1\% of the light exits each time it reflects back and forth between the mirrors, on average, the beam will pass through the amplifying medium a hundred times before it exits. Now, considering that each time the beam passes through the amplifying medium it is increasing exponentially, this factor of 100 makes an enormous difference. 

% This all leads to very high light intensities inside the amplifying medium and this is why (as was said earlier) we can forget about spontaneous emission. Imagine that you're an excited electron sitting happily in your higher energy level. Normally, you'll just drop down spontaneously when your time is up. But, inside a laser, there's so much light that you never drop spontaneously because before your time's up you've been disturbed by another photon, stimulated to join in with all the other light and join in coherently as well!

% \subsection{Efficiency}
% Usually, lasers are very inefficient beasts. Because of the large energies required to maintain a population inversion, their efficiencies are generally far below 1\%. Some reasons for this:
% \begin{itemize}
% \item The pumping energy is considerably larger than the output photon energy.
% \item High intensity pumping combined with the high intensity of the laser beam means that the amplifying medium will get very hot. So, there will be large heat losses. To make this matter worse, we need to cool the amplifying medium usually so that it, or its container, doesn't melt. By cooling the system we just transfer more heat and increase our losses but better this than destroy a \textsterling 50 000 laser!
% \end{itemize}
 
% \subsection{Semiconductor Lasers}
% The basic structure of a standard `edge emitting' semiconductor laser is shown below. The whole block shown below is a semiconductor chip with dimensions approximately 0.5mm $\times$ 0.5 mm $\times$ 1 mm.
% \begin{figure*}
% \includegraphics[]{diode.png}
% \end{figure*}

 
% The above laser fits the basic shape of a normal laser, the mirrors, however, are far from the 100\% and 99\% reflecting ideals discussed earlier. The mirrors are simply due to the semiconductor-air boundary at the edges of the chip. [This in fact gives 40\% reflection only (at both sides).] This would be disastrous for highly inefficient gas lasers but not for our semiconductor laser. The reason why:
% \begin{itemize}
% \item The population inversion inside the semiconductor sandwich area is millions of times higher than in gas lasers [$10^{25}$ electrons/$m^{3}$].
% \item The exponential increase in light intensity (i.e. 1 photon becoming two, becoming four etc.) occurs far more quickly because of the higher population inversion.
% \item So the fact that we lose 60\% of the light at each reflection is compensated for by having huge gains between the mirrors.
%  \end{itemize}

% \subsection{Advantages and Uses of Laser Diodes}
% These are straightforward and can be summarised as follows:

% Advantages:
% \begin{itemize}
% \item Cheaper
% \item Smaller
% \item More efficient
% \item Easy to mass produce	
% \end{itemize}
% Some uses:
% \begin{itemize}
% \item Inside DVD and CD players
% \item Bar-code readers
% \item Telecommunications (via optical fibres)
% \item Image scanning
% \item Laser surgery
% \end{itemize}
% The usefulness of laser diodes is `reflected' in the number of them produced annually - around 1 billion (109) laser diodes are produced worldwide per year!
 
% For more than enough further reading see:
% \begin{itemize}
% \item \url{http://members.aol.com/WSRNet/tut/ut1.htm}
% \item Wikipedia \url{http://en.wikipedia.org}  then type `laser' or `semiconductor laser'
% \item Google search `laser theory'
% \end{itemize}


	
\subsection{end-of-chapter questions}



\subsection*{wave terminologies}

\question{
	A bottle floating on a water surface is seen to bob up and down over a distance of 20 cm twice each second. What is its amplitude and its frequency?
}

\question{A horn produces a note of frequency of 512 Hz. Given that sound travels at $330 \mps$ in air, find the wavelength of this sound wave.}

\begin{marginfigure}
	\vspace*{-12pt}
	\centering
	\begin{tikzpicture}[scale=0.78]
	\draw[help lines, gray!50, very thin, step=0.2] (0,-2) grid (6,2);
	\draw[help lines, step=1] (0,-2) grid (6,2);
	\foreach \y in {-2,-1,0,1,2} \node[left] at (0,\y) {$\y$};
	\foreach \x in {0,0.1,0.2,0.3,0.4,0.5,0.6} \node[below] at (10*\x,-2) {$\x$};
	\node at (-1.2,1.5) {$y$/cm};
	\node at (5,-2.8) {$x$/m};
	\draw [thick, blue, domain=0:6, samples=30, smooth] plot (\x,{1.6*sin(75*\x)});
	\draw [thick, Green, domain=0:6, samples=30, smooth] plot (\x,{-0.8*cos(120*\x)});
	\node[blue] at (1.8, 1.6) {$P$};
	\node[Green] at (4.5, 1.1) {$Q$};
	\end{tikzpicture}
	\vspace*{-18pt}
\end{marginfigure}

\question{The graph shows the variation of displacement $y$ with the distance $x$ of wave $P$ and wave $Q$ at some instant. (a) For each of the two waves, state the amplitude and the wavelength. (b) Given that both waves have a time period of 0.30 s, find the frequency of the waves, and hence find the speed for each wave.}


\subsection*{transverse \& longitudinal waves}

\question{
	An ultrasound waves of frequency 2.0 MHz travels with a speed of $1600 \mps$ through water. What is the shortest distance from a point of maximum pressure to a point of	minimum pressure?
}

\begin{marginfigure}
	\vspace*{-12pt}
	\centering
	\begin{tikzpicture}[xscale=0.6,yscale=0.75]
	\draw [->] (0,-2) -- (0,2.4) node[midway,left]{$O$} node[left]{$y$};
	\draw [->] (0,0) -- (2.5*pi,0) node[below]{$x$};
	\draw[very thick,->] (1.6*pi,2) -- (2.4*pi,2);
	\node[above,twoline] at (2*pi, 2.2) {direction of\\wave motion};
	\draw [thick, blue, domain=0:2.4*pi, samples=30, smooth] plot (\x,{1.6*sin(\x*1.2 r)});
	\draw[fill, red] (5*pi/12,1.6) circle (0.08) node[above, black]{$P$};
	\draw[fill, red] (10*pi/12,0) circle (0.08) node[above right, black]{$Q$};
	\draw[fill, red] (15*pi/12,-1.6) circle (0.08) node[below, black]{$R$};
	\draw[fill, red] (20*pi/12,0) circle (0.08) node[above left, black]{$S$};
	\end{tikzpicture}
	\vspace*{-16pt}
\end{marginfigure}

\question{
	The graph shows the variation with distance $x$ of the displacement $y$ of a wave on a stretched rope at
	time $t = 0$. The wave has a frequency of 1.0 Hz. (a) State and explain whether this wave is transverse or longitudinal. (b) For the points $P$, $Q$, $R$, and $S$ labelled on the graph, which one of them is moving upwards with greatest speed at $t=0$? (c) Sketch the wave pattern at $t = 0.50 \text{ s}$. (c) For the points $P$, $Q$, $R$, and $S$, what are their positions at $t = 0.50 \text{ s}$?
}


\question{
	Given a slinky toy, suggest how you can demonstrate (a) a transverse wave, (b) a longitudinal wave to your classmates.
}

\subsection*{sound waves}

\question{
	If some alien civilisation blows up the moon one day, can we hear it?
}

\question{An oscilloscope is used to measure a sound wave. When the time base is set at 5.0 ms div$^{-1}$, 3 complete waveforms are observed over 5 divisions. Calculate, (a) the period of the wave, (b) the frequency of the wave.}

\question{Describe the changes to the wave pattern displayed on an oscilloscope if the sound being measured has (a) a higher pitch, (b) a greater volume.}


\subsection*{electromagnetic waves}

\question{
	Green light has a wavelength of 500 nm in vacuum. What is its frequency?
}

\question{
	An electromagnetic wave has a period of 1.0 ps. (a) What is the wavelength of this wave? (b) What is the number of wavelengths in a distance of 1.0 m?
}

\question{
	(a) What is the frequency of the longest-wavelength ultraviolet wave? (b) What is the frequency of the shortest-wavelength infra-red radiation?
}

\question{
	State a typical value for the wavelength for the following radiation in vacuum: (a) infra-red, (b) X-ray, (c) microwave, (d) ultraviolet.
}

\question{
	A student argues that the speed of an electromagnetic wave is proportional to its frequency because we have $c=\lambda f$. State and explain whether the student's argument is correct.
}

\question{
	(a) Suggest as many properties of electromagnetic waves as you can. (b) Among these properties you have listed, which property of electromagnetic wave is distinct from other transverse waves?
}

%\question{
%	Suppose a sound wave and an electromagnetic wave have the same wavelength, compare their frequencies.
%}

\subsection*{intensity of waves}

\question{
	A sound wave in air has an amplitude $A_0$ and intensity $I_0$. If the amplitude increases to $3A_0$, what is the new intensity?
}

\question{
	The intensity of radiation arriving normally on a solar panel was $500 \text{ W m}^{-2}$. (a) If the panel has an effective area of $4.0 \text{ m}^2$, how much energy would arrive in one hour? (b) If the radiation has double the amplitude, how much energy would arrive in one hour?
}

\question{
	The intensity $I$ of a sound wave can be given by the formula: $I = k v \rho f^2 A^2$, where $v$ is the speed of sound, $\rho$ is the density of the medium, $f$ is the frequency and $A$ is the amplitude. Show that $k$ is a unit-free constant.
}



\question{
	A point source gives out spherical waves. How does the wave amplitude $A$ vary with the distance $r$ from the source?
}

\subsection*{polarisation}

\question{
	A horizontally polarised beam of light of intensity $I_0$ passes through a polarising filter whose plane of polarisation is at $30^\circ$ to the horizontal. What is the transmitted intensity?
}

\question{
	When polarized light is sent through some chemical solution, the plane of polarization is rotated. If the intensity is $I$ without the solution being placed in the beam but $0.80I$ after the sugar is placed in the beam, what is the angle of rotation?
}

\question{
	A beam of unpolarised light is sent through two polarising filters $A$ and $B$. Filter $A$ has its axis of
	polarisation in the vertical direction, and filter $B$ has its axis of polarisation in the horizontal direction. (a) What is the emergent intensity? (b) A third polarising filter $C$ is inserted between the filters $A$ and $B$, such that filter $C$ has its plane of polarisation at an angle of $45^\circ$ to the vertical. What is the new emergent intensity?
}

\question{
	Suggest how you may use a stretched elastic rope to demonstrate the phenomenon of polarisation.
}

\subsection*{diffraction}

\question{
	A hill stands between a house and a radio transmitter, but radio signals sent from the transmitter are still received at the house. How can this happen?
}

\question{
	Diffraction can be demonstrated by observing how water waves in a ripple tank go through a gap. Suggest whether the following action would lead to greater or less obvious diffraction: (a) increasing the frequency of the water waves, (b) increasing the amplitude of the waves, (c) increasing the width of the gap, (d) increasing the speed of the waves by increasing the depth of water.
}

\question{
	Microwave ovens use microwaves of frequency of around 2.4 GHz to heat food. The front door of many microwave ovens is made of glass with a metal grid, where gap in the grid are a few mm across. By reference to the wavelength of the microwave, explain how the front door keeps microwaves in but could let light out.
}

% \subsection*{Doppler effect}

% For the questions below, take the speed of sound to be $340 \mps$. 

% \question{A car travels with a constant speed along a straight road. The car's horn is known to sound at a frequency of 500 Hz, but an observer standing on the roadside hears a frequency of 450 Hz. What is the magnitude and the direction of the car's velocity relative to the observer?}


% \question{
% 	An ambulance travelling at a steady speed of $25 \mps$ passes close to a stationary observer. The warning signal on the ambulance has a frequency of 1500 Hz. What is the overall change in the frequency heard by the observer as the ambulance goes by?
% }

% \question{
% 	If the motion of the source is at right angles to an observer, state and explain whether there will be a Doppler effect?
% }


% \question{
% 	A girl sits on a horizontal platform that is rotating about a vertical axis. The girl moves in a circular path with a constant speed of $8.0 \mps$. The girls starts blowing a whistle which emits a sound of frequency 1200 Hz. To an observer standing at a distance, (a) find the maximum and the minimum frequency heard. (b) Describe the variation in the frequency of the sound heard by the observer.
% }

% \question{
% 	A spectral line of a particular wavelength is emitted from a distant star. The wavelength observed is found to vary periodically from a maximum value to a minimum value and back to the maximum value. What can you say about the motion of the star?
% }



% \question{
% 	[This question is beyond CAIE syllabus] (a) A source emits sound wave of frequency $f$. If the observer, rather than the source, is in motion, suggest whether there is any change in the apparent frequency $f^\prime$ heard by the observer. (b) If the observer is moving towards the source at a velocity of $v_o$, show that the observed frequency $f'$ is: $f' = f \frac{v+v_o}{v}$. (c) If the observer is moving away from the source, derive a similar expression for the observed frequency. 
% }

\section{Superposition of Waves}

When two of more waves meet together, the resultant motion is a combination of the individuals. They can form a \emph{resultant wave} when they overlap, after which they cross one another and neither is affected. In this chapter, we will study the principle of superposition, and look at two important consequences: the phenomenon of \emph{interference} and \emph{stationary waves}.

\subsection{Superposition of waves}



\subsection{Principle of superposition}

\begin{ilight}
	 When two (or more) waves meet to form a resultant wave, the \keypoint{principle of superposition} states that the resultant displacement is the vector sum of each individual displacement\index{superposition of waves!principle of superposition}
\end{ilight}

\example{Displacement–time graphs for two waves $A$ and $B$ meeting at some point, together with the resultant wave formed at that point.}

\begin{figure*}[ht]
	\centering
	\begin{minipage}{0.45\textwidth}
		\centering
		\begin{tikzpicture}[scale=0.8]
		\draw [->] (0,-2) -- (0,2) node[above]{displacement};
		\draw [->] (0,0) -- (6.4,0) node[below]{$t$};
		\draw [Green, dashed] (0,0) --++ (0.8,1) --++ (1.6, -2) --++ (1.6,2) --++ (1.6, -2) --++ (0.4,0.5) node[below right]{$A$};
		\draw [purple, dashed] (0,-0.8) -- (1.6,-0.8) -- (2.4, 0.4) -- (4 ,0.4) -- (4.8, -0.5) -- (5.6, -0.5) -- (6, 0.4) node[right]{$B$}; 
		\draw [thick, blue] (0,-0.8) -- (0.8, 0.3) -- (1.6, -0.8) -- (2.4, -0.6) -- (4, 1.4) -- (4.8, -0.5) -- (5.6, -1.5) -- (6,0.1);
		\end{tikzpicture}
	\end{minipage}\hfil
	\begin{minipage}{0.45\textwidth}
		\centering
		\begin{tikzpicture}[scale=0.8]
		\draw [->] (0,-2) -- (0,2) node[above]{displacement};
		\draw [->] (0,0) -- (6.4,0) node[below]{$t$};
		\draw [Green,dashed,domain=0:6,samples=30,smooth] plot (\x,{0.6*sin(\x*3 r)}) node[below right]{$A$};
		\draw [purple,dashed,domain=0:6,samples=30,smooth] plot (\x,{1.2*sin(\x*1.5 r)}) node[above right]{$B$};
		\draw [thick,blue,domain=0:6,samples=30,smooth] plot (\x,{1.2*sin(\x*1.5 r) + 0.6*sin(\x*3 r)});
		\end{tikzpicture}
	\end{minipage}
\end{figure*}

Two important situations are constructive superposition and destructive superposition\index{superposition of waves!constructive superposition}\index{superposition of waves!destructive superposition}

\keypoint{Constructive superposition} occurs when resultant wave has greatest possible amplitude. This happens when peaks of the two waves meet together (also trough meets trough)\footnote{This statement is implicitly for transverse waves. For longitudinal waves to superpose constructively, the regions of compression (or rarefaction) must overlap.} 	peaks of both waves must arrive with a time difference $\Delta t = 0, \, T, \, 2T, \cdots$
	
\keypoint{Destructive superposition} occurs when amplitudes of each wave cancel out one another this happens when peak of one wave overlaps with trough of the other wave peaks of both waves must arrive with a time difference $\Delta t = \frac{1}{2}T, \, \frac{3}{2}T, \, \frac{5}{2}T, \cdots$

\begin{figure*}[ht]
	\centering
	\noindent\begin{minipage}{0.45\linewidth}
		\centering
		\begin{tikzpicture}[xscale=0.8, yscale=0.75]
		\draw [->] (0,-2) -- (0,2) node[above]{displacement};
		\draw [->] (0,0) -- (6.4,0) node[below]{$t$};
		\draw [Green, dashed,domain=0:6,samples=30,smooth] plot (\x,{0.6*sin(\x*1.2 r)});
		\draw [purple, dashed,domain=0:6,samples=30,smooth] plot (\x,{1.2*sin(\x*1.2 r)});
		\draw [thick,blue,domain=0:6,samples=30,smooth,variable=\x] plot (\x,{1.8*sin(\x*1.2 r)});
		\end{tikzpicture}
	\end{minipage}\hfil
	\begin{minipage}{0.45\linewidth}
		\centering
		\begin{tikzpicture}[xscale=0.8, yscale=0.75]
		\draw [->] (0,-2) -- (0,2) node[above]{displacement};
		\draw [->] (0,0) -- (6.4,0) node[below]{$t$};
		\draw [Green,dashed,domain=0:6,samples=30,smooth] plot (\x,{-1.2*sin(\x*1.2 r)});
		\draw [purple ,dashed,domain=0:6,samples=30,smooth] plot (\x,{1.6*sin(\x*1.2 r)});
		\draw [very thick,blue,domain=0:6,samples=30,smooth] plot (\x,{0.4*sin(\x*1.2 r)});
		\end{tikzpicture}
	\end{minipage}\hfil
	\caption{constructive and destructive superposition of two transvers waves}
\end{figure*}
\vspace{2cm}
Note: longitudinal waves are often represented by transverse diagrams. Transverse are much clearer and easier to read:
\vspace{1cm}
 

\begin{figure*}[htp]
	\centering
	\noindent\begin{minipage}{0.5\linewidth}
		\centering
		\includegraphics[width=.95\textwidth]{longitudinal+I.pdf}
	\end{minipage}\hfil
	\begin{minipage}{0.5\linewidth}
		\centering
		\includegraphics[width=.95\textwidth]{longitudinal-I.pdf}
	\end{minipage}
	\caption{constructive and destructive superposition of two sound waves} 
\end{figure*}


\subsection{Path difference}

The \keypoint{path difference}\index{path difference} is the difference in distance that two waves must travel from their sources to a given point. It plays a crucial role in wave interference, where waves can either constructively or destructively superpose depending on their path difference. Path difference is generally expressed in multiples of wavelength. To tell whether two waves superpose constructively or destructively:

{

\centering
\vspace*{6pt}
\begin{tcolorbox}[colframe=red, colback=yellow!30, width=0.65\textwidth]
	\setlength{\baselineskip}{\baselineskip}%
	
	\centering
	
	if ${\Delta L = 0, \, \lambda, \, 2\lambda, \cdots}$, then superposition is constructive
	
	if ${\Delta L = \frac{\lambda}{2}, \, \frac{3\lambda}{2}, \, \frac{5\lambda}{2}, \cdots}$, then superposition is destructive
	
\end{tcolorbox}\vspace*{-4pt}

}

\newpage

\example{Two loudspeakers are wired to produce identical sound signals in unison. The sound wave produced has a wavelength of 80 cm. Describe the volume you hear when you are at a distance of (a) 10 m from both speakers, (b) 10 m from one speaker and 12 m from the other?}

\begin{soln} path difference for case (a): $\Delta L = 10 - 10  = 0$

so constructive superposition, resultant amplitude is large, a loud sound is heard

path difference for case (b): $\Delta L = 12 - 10 = 2.0 \text{ m} \RA \Delta L = \frac{5}{2}\lambda$

so destructive superposition, resultant amplitude is small, sound is quiet \end{soln}

\subsection{Phase difference}

When two waves or two vibrating particles are compared, it is also useful to describe how much one is out of step with the other in terms of their \keypoint{phase difference}\index{phase difference} ($\Delta \phi$)
\footnote{More logically, we should first define the notion of \emph{phase angle} before talking about the \emph{phase difference}, which is simply the difference between the phase angles of two waves.
	
	Wave motion can generally be described by sine (or cosine) functions. Since the initial displacement of a wave is not necessarily zero, one can introduce a constant term that appears in the argument of the sine function. This term shifts the origin of the sine, and we call that the \emph{phase angle} $\phi$ of the wave, measured in radians (or degrees). One can think of $\phi$ as a number giving the fraction in a complete oscillation.
	
	In particular, the displacement $y$ of a transverse wave along the direction $x$ of energy transfer can be given by: $y = A \sin\left( \frac{2\pi x}{\lambda} - \frac{2\pi t}{T} + \phi \right)$. You can verify that this indeed describes a sinusoidal wave of wavelength $\lambda$ and period $T$ travelling at speed $v = \frac{\lambda}{T}$, while $\phi$ determines when and where the peaks show up.}

Phase difference is measured in radians (rad) or degrees ($^\circ$),  if two waves have a phase difference $\Delta \phi = 0,2\pi,4\pi,\cdots$, the two waves are said to be \keypoint{in phase}\footnote{If $\Delta \phi$ is given in degrees, then two waves are in phase if $\Delta \phi = 0, \, 360^\circ, \, 720^\circ, \cdots$}. In this case, peak of one wave overlaps with peak of the other. If two waves are not in phase, then they are said to be \keypoint{out of phase}
If $\Delta \phi = \pi, 3\pi, 5\pi, \cdots$, then two waves are completely out of phase, or \keypoint{anti-phase}\footnote{If $\Delta \phi$ is given in degrees, then two waves are anti-phase if $\Delta \phi = 180^\circ, \, 540^\circ, \, 900^\circ, \cdots$} in this case, peak of one wave meets the trough of the other. 

We can tell whether two waves superpose constructively or destructively from their phase difference as follows:

{
	
	\centering
	\vspace*{6pt}
	\begin{tcolorbox}[colframe=red, colback=yellow!30, width=0.65\textwidth]
		\setlength{\baselineskip}{\baselineskip}%
		
		\centering
		
		if ${\Delta \phi = 0, \, 2\pi, \,4\pi, \cdots}$, then superposition is constructive
		
		if ${\Delta \phi = \pi, \, 3\pi, \, 5\pi, \cdots}$, then superposition is destructive
		
	\end{tcolorbox}\vspace*{-4pt}
	
}






\newpage

Phase difference can also be related to path difference by: \begin{empheq}[box=\tcbhighmath]{equation*}{ \frac{\Delta \phi}{2\pi} = \frac{\Delta L}{\lambda} = \frac{\Delta t}{T} } \end{empheq}

\example{The diagrams below each shows the displacement of a particular point in two transverse waves. If both waves are travelling in the same direction with a wavelength of 30 cm, what is the phase difference and path difference between each wave?}

\begin{figure*}[htp]
	\centering
	\begin{minipage}{0.32\textwidth}
		\centering
		\begin{tikzpicture}[xscale=0.56,yscale=0.6]
		\draw [thick, ->] (0,0) --(7.5,0) node[below]{$t$};
		\draw [thick, ->] (0,-3) --(0,2.5) node[left]{$x$};
		\draw [thick,blue,domain=0:7,smooth] plot (\x,{2*sin(\x r)}) node[right]{$B$};
		\draw [thick,Green,domain=0:7,smooth] plot (\x,{2*sin(((\x - pi/2) r)}) node[right]{$A$};
		\end{tikzpicture}
		
		(a)
	\end{minipage}
	\begin{minipage}{0.32\textwidth}
		\centering
		\begin{tikzpicture}[xscale=0.56,yscale=0.6]
		\draw [thick, ->] (0,0) --(7.5,0) node[below]{$t$};
		\draw [thick, ->] (0,-3) --(0,2.5) node[left]{$x$};
		\draw [thick,blue,domain=0:7,smooth] plot (\x,{2*sin(\x r)}) node[right]{$B$};
		\draw [thick,Green,domain=0:7,smooth] plot (\x,{2*sin(((\x - pi) r)}) node[right]{$A$};
		\end{tikzpicture}
		
		(b)
	\end{minipage}
	\begin{minipage}{0.32\textwidth}
		\centering
		\begin{tikzpicture}[xscale=0.56,yscale=0.6]
		\draw [thick, ->] (0,0) --(7.5,0) node[below]{$t$};
		\draw [thick, ->] (0,-3) --(0,2.5) node[left]{$x$};
		\draw [thick,blue,domain=0:7,smooth] plot (\x,{2*sin(\x r)}) node[right]{$B$};
		\draw [thick,Green,domain=0:7,smooth] plot (\x,{2*sin(((\x - pi/3) r)}) node[below right]{$A$};
		\end{tikzpicture}
		
		(c)
	\end{minipage}
	\caption{examples of phases differences between two waves}
\end{figure*}

\begin{soln}
(a) $\Delta t = \frac{1}{4}T \RA \Delta \phi = \frac{1}{4}\times 2\pi = \frac{1}{2}\pi, \quad \Delta L = \frac{1}{4} \lambda = \frac{1}{4}\times30 = 7.5 \text{ cm}$

\eqyskip (b) $\Delta t = \frac{1}{2}T \RA \Delta \phi = \frac{1}{2}\times 2\pi = \pi, \quad \Delta L = \frac{1}{2} \lambda = \frac{1}{2}\times30 = 15 \text{ cm}$

\eqyskip (c) $\Delta t = \frac{1}{6}T \RA \Delta \phi = \frac{1}{6}\times 2\pi = \frac{1}{3}\pi, \quad \Delta L = \frac{1}{6} \lambda = \frac{1}{6}\times30 = 5.0 \text{ cm}$ \end{soln}


\subsection*{Brief summary}

The conditions for constructive and destructive superposition can be summarised as follows:

\begin{center}
{\renewcommand{\arraystretch}{1.25}
\begin{tabular}{|c|c|c|}
	\hline
	& constructive superposition & destructive superposition \\
	\hline
	simple criteria & peak meets peak & peak meets trough \\
	\hline
	time difference & $\Delta t = n \cdot T$ & $\Delta t = \left(n+\frac{1}{2}\right) \cdot T$ \\
	\hline
	path difference & $\Delta L = n \cdot \lambda$ & $\Delta L = \left(n+\frac{1}{2}\right) \cdot \lambda$ \\
	\hline
	phase difference & $\Delta \phi = n\cdot 2\pi\,$ or $\,n\cdot 360^\circ$ & $\Delta \phi = \left(n+\frac{1}{2}\right)\cdot 2\pi \,$ or $\,\left(n+\frac{1}{2}\right)\cdot 360^\circ$ \\
	\hline
\end{tabular}
}
\end{center}

where $n$ is a whole number: $n=0, 1, 2, 3 \cdots$

Remember that $\Delta t$, $\Delta L$ and $\Delta \phi$ are just different ways to describe how much one wave is ahead of/lag behind another wave, so they must be closely interrelated to each other.




\newpage

\subsection{Interference}

When two (or more) waves of constant phase difference meet, stable regions of constructive and destructive superposition are produced alternately, this phenomenon is called \keypoint{interference}\index{interference}



\begin{figure*}[htp]
	\centering
	\includegraphics[width=1\textwidth]{interference.pdf}
\caption{constructive (solid lines) and destructive (dashed lines) interference of waves produced from two coherent sources $A$ and $B$}
\end{figure*}
\vspace{1cm}
Regions of constructive interference are often called \emph{maxima}

\begin{compactitem}
	\item[--] points on line C$_0$ are of equal distance to wave sources, i.e., path difference $\Delta L = 0$
	
	so maxima are observed along C$_0$
	
	\item[--] points on C$_1$, C$_2$ and C$_3$ have path difference $\Delta L = \lambda, \, 2\lambda, \, 3\lambda$
	
	so maxima are also formed along C$_1$, C$_2$ and C$_3$
	
	\item[--] equivalently, points on C$_n$ have phase difference $\Delta \phi = n\cdot 2\pi$, where $n=0,1,2,3,\cdots$
\end{compactitem} 

Regions of destructive interference are often called \emph{minima}

\begin{compactitem}
	\item[--] points on D$_1$, D$_2$ and D$_3$ have path difference $\Delta L = \frac{1}{2}\lambda, \, \frac{3}{2}\lambda, \, \frac{5}{2}\lambda$
	
	so maxima are also formed along D$_1$, D$_2$ and D$_3$
	
	\item[--] similarly, points on D$_n$ have phase difference $\Delta \phi = \left( n+\frac{1}{2} \right) \cdot 2\pi$, where $n=1,2,3,\cdots$
\end{compactitem} 






\newpage


Stable interference means regions of maxima/minima remain maxima/minima can be easily seen. This requires \keypoint{coherent} wave sources, which means constant phase difference over time\index{coherence}.

The easiest way to produce coherent waves is to use the the same source, such as

\begin{compactitem}
	\item[--] dippers driven by a common vibrating beam in a ripple tank
	
	\item[--] loudspeakers driven by the same signal generator
 \item[--] a laser beam shone on two slits.
\end{compactitem}

\begin{figure}[ht]
	\centering
	\begin{tikzpicture}[scale=1.2]
	% water in ripple tank
	\draw[cyan!30, fill] (2,-0.3) -- (2,-1.5) -- (-1,-3.5) -- (-5,-3.5) -- (-5,-2.3) -- (-2,-0.3) -- cycle;
	\draw[white, fill, opacity=0.3] (2,0) -- (2,-0.3) -- (-1,-2.3) -- (-5,-2.3) -- (-5,-2) -- (-1, -2) -- cycle;
	% outline for ripple tank
	\draw[gray!40] (-2,-0.3) -- (-2,-1.5) -- (2,-1.5) (-2,-1.5) -- (-5,-3.5);
	\draw[thick] (2,0) -- (2,-1.5) -- (-1,-3.5) -- (-5,-3.5) -- (-5,-2) -- (-2,0) -- cycle;
	\draw[thin, gray] (2,-0.3) -- (-1, -2.3) -- (-5, -2.3) -- (-2,-0.3) -- cycle;
	\draw[thick] (-5,-2) -- (-1, -2) -- (-1, -3.5) (-1, -2) -- (2,0);
	\draw[thick] (-2,0) -- (-2,-0.3);
	% dippers
	\foreach \x in {-0.8,0.9} {
		\draw[thick, fill=orange!40] ({-0.6+\x}, 0.2)  circle (0.4);
		\draw[ultra thick] ({-0.6+\x}, 0.6) -- ({-0.6+\x},1.0);
		\draw[<->] ({-0.6+\x}, -0.25) --++ (0,-0.35);
	}
	\draw[ultra thick] (-1.4, 0.98) -- (0.3, 0.98);
	\draw[ultra thick] (-0.55,1) --++ (0,0.5) --++ (1.2,0.5) node[below right]{to vibrator};
	\draw (0.5, 0.4) --++ (1.2,0.5) node[right]{dipper};
	% ripples
	\draw[gray] (-1.62,-0.4) [out=-75 , in=180] to (-1.45,-0.52) [out=0, in=210] to (-1.0,-0.4);
	\draw[gray] (-1.85,-0.4) [out=-75 , in=180] to (-1.55,-0.65) [out=0, in=210] to (-0.7,-0.4);
	\draw[gray] (-2.08,-0.4) [out=-75 , in=180] to (-1.6,-0.8) [out=0, in=210] to (-0.4,-0.4);
	\draw[gray] (0.08,-0.4) [out=-75 , in=180] to (0.25,-0.52) [out=0, in=210] to (0.7,-0.4);
	\draw[gray] (-0.15,-0.4) [out=-75 , in=180] to (0.15,-0.65) [out=0, in=210] to (1.0,-0.4);
	\draw[gray] (-0.38,-0.4) [out=-75 , in=180] to (0.1,-0.8) [out=0, in=210] to (1.3,-0.4);
	\end{tikzpicture}
	
	\caption{Demonstrating interferece of water waves in a ripple tank}
\end{figure}


Conversely, we do not see interference from beams of light from different lamps or laser sources because light from different sources are usually \emph{incoherent}\footnote{Emission of light is associated with the electronic transitions between energy levels in an atom, which is a completely random process. So in general separate light sources are not coherent.}

This can be overcome by dividing a single beam into several beams using a number of slits. common apparatuses of doing so are the \emph{double-slit} and \emph{diffraction gratings}

\example{Two coherent waves meet in space, one of which has an amplitude of 0.30 cm and the other of 0.20 cm. How does the intensity of maxima compare with the intensity of minima?}

\begin{soln} Amplitude of maxima: $A_\tmax = A_1 + A_2 = 0.30 + 0.20 = 0.50 \text{ cm}$

Amplitude of minima: $A_\tmin = A_1 - A_2 = 0.30 - 0.20 = 0.10 \text{ cm}$

Ratio of intensities: $\frac{I_\tmax}{I_\tmin} = \left( \frac{A_\tmax}{A_\tmin} \right)^2 = \left( \frac{0.50}{0.10} \right)^2 \RA \frac{I_\tmax}{I_\tmin} = 25$ \end{soln}



\newpage


\subsection{Double-slit interference}\index{double-slit experiement}

Let's take a beam of light being guided through two narrow slits, light waves diffracting through the slits act as coherent wave sources and when they meet on a screen, interference pattern can be seen. Since this experiment is carried out with light, alternating bright and dark fringes are formed. 

This is known as \emph{Thomas Young}'s \keypoint{double-slit experiment}
\footnote[][-4cm]{The nature of light has been argued since the history of human civilisation. It has been long debated whether light is a \emph{wave} or it is made of \emph{particles}. It was until the early 1800s when Thomas Young carried out the famous double-slit experiment that now bears his name, people became certain that light has wavelike properties. We know now of course that it has particle -like properties too. For a detailed review, check out the Wikipedia article: \url{https://en.wikipedia.org/wiki/Light}.}



\colorlet{wall}{blue!30!black}
\colorlet{myblue}{blue!70!black}
\colorlet{myred}{red!70!black}
\colorlet{mydarkred}{red!50!black}
\colorlet{mylightgreen}{green!60!black!70}
\colorlet{mygreen}{green!60!black}
\colorlet{myredgrey}{red!50!black!80}
\colorlet{myshadow}{blue!30!black!90}
\tikzstyle{wave}=[myblue,thick]
\tikzstyle{mydashed}=[black!70,dashed,thin]
\tikzstyle{mymeas}=[{Latex[length=3,width=2]}-{Latex[length=3,width=2]},thin]
\tikzstyle{mysmallarr}=[-{Latex[length=3,width=2]}]


\newcommand\rightAngle[4]{
  \pgfmathanglebetweenpoints{\pgfpointanchor{#2}{center}}{\pgfpointanchor{#3}{center}}
  \coordinate (tmpRA) at ($(#2)+(\pgfmathresult+45:#4)$);
  \draw[white,line width=0.6] ($(#2)!(tmpRA)!(#1)$) -- (tmpRA) -- ($(#2)!(tmpRA)!(#3)$);
  \draw[mydarkred] ($(#2)!(tmpRA)!(#1)$) -- (tmpRA) -- ($(#2)!(tmpRA)!(#3)$);
}
\newcommand\lineend[2]{
  \def\w{0.1} \def\c{30}
  \draw[mygreen] (#1)++(#2:\w) to[out=#2-180-\c,in=#2+\c] (#1)
                               to[out=#2+\c-180,in=#2-\c]++ (#2-180:\w);
}
\def\tick#1#2{\draw[thick] (#1) ++ (#2:0.1) --++ (#2-180:0.2)}

% INTERFERENCE FADING
\begin{tikzfadingfrompicture}[name=interference]
  \def\lambd{0.5} % wavelength
  \foreach \r in {1,...,15}
    \foreach \j in {1,...,25}
       \path [line width=\lambd*\j,draw=transparent!0,opacity=0.04]
          (0,0) circle (\lambd*\r); %(0:\r) arc (0:180:\r);
\end{tikzfadingfrompicture}


\begin{figure*}
	\centering
 \begin{tikzpicture}[
    nodal/.style={mylightgreen,dashed,very thin},
    declare function={
      %xnode(\n,\dn,\lam,\f) = sqrt( (\n^2+(\n+\dn)^2)*\lambd^2/2 - (\n^2-(\n+\dn)^2)^2*\lambd^4/(4*\a^2) - \a^2/4 );
      xnode(\n,\dn,\lam,\f) = \lam/\f*sqrt( \n^2*(\f^2-\dn^2)+\n*\dn*(\f^2-\dn^2)+\dn^2*\f^2/2-(\f^4+\dn^4)/4 );
      ynode(\n,\dn,\lam,\a) = (2*\n*\dn+\dn^2)*\lam/(2*\f);
      intensity(\y,\lam,\a,\L) = cos(180*\a*\y/(2*\lam*sqrt(\L*\L+\y*\y)))^2;
    }
  ]
  
  \def\L{3.8}       % distance between walls
  \def\H{5.4}       % total wall height
  \def\h{2.8}       % plane wave height
  \def\t{0.15}      % wall thickness
  \def\a{1.15}      % slit distance
  \def\d{0.20}      % slit size
  \def\N{21}        % number of waves
  \def\lambd{0.20}  % wavelength
  \def\R{\N*\lambd} % wave radius
  \def\Nlines{3}    % number of nodal lines
  \def\A{1.6}       % amplitude
  %\def\r{0.06}      % point source radius
  %\def\nmax{10}
  \def\nsamples{100}
  \def\ang{62}
  
  \begin{scope}
    \clip (-\t/2,-\H/2) rectangle (\L,\H/2);
    %\clip (-\t/2,0.7*\a) -- (0.6*\L,\H/2) -- (\L,\H/2) --
    %      (\L,-\H/2) -- (0.6*\L,-\H/2) -- (-\t/2,-0.7*\a) -- cycle;
    
    % NODAL LINES
    \draw[nodal]
      (0.08*\N*\lambd,0) -- (1.06*\R,0);
    \coordinate (NP0) at (\L,0);  % to avoid "Dimension too large error"
    \foreach \dn [evaluate={
                   \f=\a/\lambd;
                   \nmin=2.5+0.2*\dn; %0.501*(-\dn+\f)
                   \nmax=10; %(NP0)
                   \c=int(\dn<\f);
                   \y=\L/sqrt((\a/(\lambd*\dn))^2-1);
                 }] in {1,...,\Nlines}{
      \coordinate (NP+\dn) at (\L,\y);  % to avoid "Dimension too large error"
      \coordinate (NP-\dn) at (\L,-\y); % to avoid "Dimension too large error"
      \ifnum\c=1
        \draw[nodal,variable=\n,samples=\nsamples,smooth]
          plot[domain=\nmin:\nmax] ({xnode(\n,\dn,\lambd,\f)},{ynode(\n,\dn,\lambd,\f)})
          -- (NP+\dn);
        \draw[nodal,variable=\n,samples=\nsamples,smooth]
          plot[domain=\nmin:\nmax] ({xnode(\n,\dn,\lambd,\f)},{-ynode(\n,\dn,\lambd,\f)})
          -- (NP-\dn);
      \fi
    }
    
    % WAVES
    \foreach \i [evaluate={\R=\i*\lambd;}] in {1,...,\N}{
      \ifodd\i
        \draw[myblue,line width=0.8] (0,\a/2)++(\ang:\R) arc (\ang:-\ang:\R);
        \draw[myred,line width=0.8] (0,-\a/2)++(\ang:\R) arc (\ang:-\ang:\R);
      \else
        \draw[myblue!80,line width=0.1] (0,\a/2)++(\ang:\R) arc (\ang:-\ang:\R);
        \draw[myred!80,line width=0.1] (0,-\a/2)++(\ang:\R) arc (\ang:-\ang:\R);
      \fi
    }
  \end{scope}
  
  % PLANE WAVES
  \foreach \i [evaluate={\x=-\i*\lambd;}] in {0,...,5}{
    \ifodd\i
      \draw[myblue,line width=0.8] (\x,-\h/2) -- (\x,\h/2);
    \else
      \draw[myblue,line width=0.1] (\x,-\h/2) -- (\x,\h/2);
    \fi
  }
  
  % WALL
  \fill[wall]
    (\t/2,\a/2-\d/2) rectangle (-\t/2,-\a/2+\d/2)
    (\t/2,\a/2+\d/2) rectangle (-\t/2,\H/2)
    (\t/2,-\a/2-\d/2) rectangle (-\t/2,-\H/2)
    (\L,-\H/2) rectangle (\L+\t,\H/2);
  
  % SHADES
  \begin{scope}[shift={(1.08*\L,0)}]
    \def\yz{\L/sqrt((\a/\lambd)^2-1)} % m = +- 1/2
    \def\yZ{\L/sqrt((\a/\lambd/2)^2-1)} % m = +- 1
    \clip (0,-\H/2) rectangle (1.1*\A,\H/2);
    \fill[white] (0,-\H/2) rectangle++ (\A,\H); % to fill seams
    \foreach \i [evaluate={\n=0.5*\i;\yn=\L/sqrt((\a/(2*\lambd*\n))^2-1);
                 }] in {1,...,\Nlines}{
      \ifodd\i % if even
        \fill[myshadow] (0,{-\yn-0.1}) rectangle++ (\A,0.2); % to fill seams
        \fill[myshadow] (0,{ \yn-0.1}) rectangle++ (\A,0.2); % to fill seams
      \fi
    }
    \path[left color=myshadow,right color=myshadow,middle color=white,shading angle={180}]
      (0,{-\yz}) rectangle (\A,{\yz});
    \foreach \i [evaluate={
                  \n=0.5*\i;
                  \m=0.5*(\i+1);
                  \yn=\L/sqrt((\a/(2*\lambd*\n))^2-1);
                  \ym=\L/sqrt((\a/(2*\lambd*\m))^2-1);
                  \dang=mod(\i,2)*180;
                 }] in {1,...,\Nlines}{
      \path[left color=myshadow,right color=white,shading angle={\dang}]
        (0,\yn) rectangle (\A,\ym);
      \path[left color=myshadow,right color=white,shading angle={180+\dang}]
        (0,-\yn) rectangle (\A,-\ym);
    }
  \end{scope}
  
  % INTENSITY
  \begin{scope}[shift={(1.1*\L+1.1*\A,0)}]
    \draw[->,thick] (-0.08*\A,0) -- (1.3*\A,0) node[right=-2] {$\expval{I}$}; % I axis
    \draw[->,thick] (0,-0.52*\H) -- (0,0.54*\H) node[right] {$y$}; % y axis
    \draw[nodal] (NP0) --++ (0.15*\L+2.1*\A,0); % green nodal lines
    \foreach \i [evaluate={\y=\L/sqrt((\a/(\lambd*\i))^2-1)}] in {1,...,\Nlines}{ % green nodal lines
      \draw[nodal] (NP+\i) --++ ({0.15*\L+1.1*\A+\A*intensity(\y,\lambd,\a,\L)},0);
      \draw[nodal] (NP-\i) --++ ({0.15*\L+1.1*\A+\A*intensity(\y,\lambd,\a,\L)},0);
    }
    \draw[myred,thick,variable=\y,samples=\nsamples,smooth,domain=-\H/2:\H/2]
      plot({\A*intensity(\y,\lambd,\a,\L)},\y);
    \foreach \i [evaluate={ % ticks
        \modd=\i; %int(\i);
        \meven=int(\i-1);
        \y=\L/sqrt((\a/(\lambd*\i))^2-1);
                }] in {1,...,\Nlines}{
      \ifodd\i
        \tick{0,-\y}{180} node[right=0,scale=0.85] {$m=-\frac{\modd}{2}$};
        \tick{0,\y}{180} node[right=0,scale=0.85] {$m=+\frac{\modd}{2}$};
      \else
        \tick{0,-\y}{180} node[right=0,scale=0.85] {$m=-\meven$};
        \tick{0,\y}{180} node[right=0,scale=0.85] {$m=+\meven$};
      \fi
    }
  \end{scope}
  
\end{tikzpicture}
	\caption{Young's double-slit interference experiment}
\end{figure*}


When waves from the two slits meet with path difference $\Delta L = 0, \, \lambda, \, 2\lambda, \cdots$, maxima is formed, ie bright fringes are observed.\\
When waves from the two slits meet with path difference $\Delta L = \frac{1}{2}\lambda, \, \frac{3}{2}\lambda, \, \frac{5}{2}\lambda, \cdots$, minima is formed, ie dark fringes are observed.

We can use the geometry of the slit separation, the distance to the screen, and the spacing of the fringes to calculate the wavelength of the light - at least we can if the slit-to-screen distance is much larger than the slit separation, i.e., $D \gg d$.\\

In this case, bright fringes are nearly equally spaced with a separation of: \begin{empheq}[box=\tcbhighmath]{equation*}{x = \frac{\lambda D}{d}}\end{empheq}

where $d$ is separation of the two slits, $D$ is slit-to-screen distance

\begin{figure*}
	\centering
	\begin{tikzpicture}[scale=0.55]
	\coordinate (A) at (0,1); \node[left] at (A) {$A$};
	\coordinate (B) at (0,-1); \node[left] at (B) {$B$};
	\coordinate (O) at (19.5,0);
	\coordinate (P) at (19.5,4);
	% double slit and screen
	\draw[thick] (0,4) -- (0,1.1) (0,0.9) -- (0,-0.9) (0,-1.1) -- (0,-4) node[below]{double slit};
	\draw (P) -- (0,0) node[left]{$M$} -- (O) node[midway,below]{$D$};
	\draw[thick] (19.5,5) -- (P) node[right] {$P$} -- (O) node[right] {$O$} node[midway,left]{$y$} -- (19.5,-5) node[below]{screen};
	% fringes
	\foreach \z in {-4, -2.667, -1.333, 0 , 1.333 ,2.667, 4}{
		\draw[blue,fill] (21.5,\z-0.2) rectangle (22.5,\z+0.2);
		\draw (22.8,\z) -- (23.2,\z);	
	}
	\foreach \z in {-4, -2.667, -1.333, 0 , 1.333 ,2.667}
	\draw[<->] (23,\z+0.03) -- (23,\z+1.273) node[midway,right]{$x$};
	\draw (22,-4.88) node[below] {fringes};
	% paths
	\draw[<->] (-1,1) -- (-1,-1) node[midway,left]{$d$};
	\draw[thick,blue] (A) -- (P) -- (B);
	\draw (0,1) -- (0.4706,-0.8824);
	\draw (0.4706,-0.8824) ++ (-0.05,0.2) -- ++(0.2,0.05) -- ++(0.05,-0.2) node[below]{$H$};
	\draw[ultra thick, red] (B) -- (0.4706,-0.8824);
	% zoom-in view
	\draw (3,0) arc(0:11.3:3); \node[right] at (3,0.32) {$\theta$};
	\draw[densely dotted] (0,0) circle(2);
	\draw[double] (-30:2) -- (10.5,-3);
	\coordinate (A') at (12,-2.5);
	\coordinate (B') at (12,-5.5);
	\draw (12,-1.5) -- (A') node[left] {$A$} -- (B') node[left]{$B$} node[midway,left]{$d$} -- (12,-6);
	\draw[blue] (A') -- ++(30:2.5);
	\draw[blue] (B') -- ++(30:3.6);
	\draw (A') -- ++(-60:2.598);
	\draw[dashed] (A') -- ++(2,0);
	\draw (A') ++ (.8,0) arc(0:30:.8);
	\draw (A') ++ (15:1.1) node{\small$\theta$};
	\draw[dashed] (B') -- ++(2,0);
	\draw (B') ++ (.8,0) arc(0:30:.8);
	\draw (B') ++ (15:1.1) node{\footnotesize$\theta$};
	\draw (A') ++ (0,-.8) arc(270:300:0.8);
	\draw (A') ++ (285:1.1) node{\footnotesize$\theta$};
	\draw[ultra thick,red] (B') -- ++(30:1.5);
	\draw (B') ++(30:1.3) -- ++(120:0.2) -- ++(30:0.2) node[below right]{$H$};
	\draw[red,->] (B') ++(30:.75) to[out=120,in=180] (13,-4.2) to[out=0,in=135] (14.4,-5);
	\draw[red] (14,-5.4) node[right] {\footnotesize{$\Delta L \approx d\sin\theta$}};
	\draw[densely dotted] (10.5,-6.4) rectangle (17.5,-1);
	\end{tikzpicture}
	\label{fringey}
 \caption{Deriving the Young slits formulae.}
\end{figure*}

The derivation of this formulae is one you should know and follows: 
Consider a point $P$ on the screen where a bright fringe is formed (see figure \ref{fringey})

The path difference between the slits must be a whole number of wavelength: $$ \Delta L = \big|PB-PA\big| = n\lambda $$

Because $D \gg d$, paths $PA$ and $PB$ can be considered approximately parallel, so path difference is length of the segment $HB$, then we have: $$\Delta L \approx d\sin\theta$$

In terms of $\theta$, bright fringes are found where: $$ d\sin\theta = n\lambda $$

To convert $\theta$ into the coordinate $y$, we have: $$\tan \theta = \frac{y}{D}$$

Because $\theta$ is very small, the small-angle approximation ($\sin \theta \approx \tan \theta$) can be applied, so positions where the bright fringes show up are: $$ \tcbhighmath{y_n =  n \times \frac{\lambda D}{d}} $$\begin{marginfigure}
\begin{tikzpicture}[scale=.7]
\begin{axis}[
  grid=both,
  xmin=-.1,
  xmax=.6,
  ymin=-.5,
  ymax=1,
  xlabel=$x$,
  ylabel=$y$,
  axis lines=center,
  >=stealth
]
  \addplot[
    domain=-.1:.6,
    blue,
    ultra thick,
    samples=100,
  ] plot[smooth] {sin(deg(x))};
  \addplot[
    domain=-.1:.6,
    red,
    ultra thick,
    samples=100,
  ] plot[smooth] {tan(deg(x))};
\end{axis}
\end{tikzpicture}

\caption{For small values of x, tan(x) [red] and sin(x) [blue] can be considered equal.} 
\end{marginfigure}


Here the coordinate subscript $n$ labels the order of the bright fringes. This shows that the $(n+1)$-th fringe is at a fixed distance from the $n$-th fringe, therefore the separation between neighbouring bright fringes is: $$ {x=\frac{\lambda D}{d}} $$

Aaltering width of slit could cause a change in \emph{brightness} of fringes.

\example{A teacher demonstrates the double-slit experiment with a beam of red light produced from a helium-neon laser. The beam of wavelength of 632 nm is sent through a double-slit separated by 0.30 mm and passed onto a wall at about 2.0 m from the slits. (a) What is the fringe separation? (b) If the experiment is carried out using a green laser, what changes do you expect?}

\begin{soln} Fringe separation: $ x = \frac{\lambda D}{d} = \frac{632 \times 10^{-9} \times 2.0 }{0.30 \times 10^{-3}} \RA x \approx 4.2 \times 10^{-3} \text{ m}$

if replaced by green light, wavelength becomes shorter, so smaller fringe separation \end{soln}


\example{Coherent light passes through a double slit. Initially, the light intensity from each slit is the same. State the change to the appearance of the fringes if (a) the separation between the slits is increased, (b) The width of one of the slits is reduced.}

\begin{soln} (a) From $ x = \frac{\lambda D}{d} $, shorter $\lambda$ results in smaller fringe separation

\hspace*{1.25em} but brightness of the fringes remains unchanged

(b) Reducing width of slit causes amplitude of light from that slit to decrease

\hspace*{1.25em} at maxima, resultant amplitude decreases, so bright fringes become less bright

\hspace*{1.25em} at minima, amplitudes do not cancel completely, so dark fringes become a bit brighter

\hspace*{1.25em} but no change to separation between fringes \end{soln}



\subsection{Multi-slit interference: diffraction gratings}

Passing light through even more slits also produces nice interference patterns, the \keypoint{diffraction grating} is such an apparatus, with hundreds or thousands of slits per mm\index{diffraction grating}

\begin{figure*}[ht]
	\centering
	\begin{tikzpicture}[xscale=1, yscale=0.84]
	% screen & fringes
	\draw[thick, fill=gray!20] (30:7.5) ++ (-3.2,1) --++ (6.4,0) --++ (0,-2) --++ (-6.4,0) --++ (0,2);
	\foreach \x/\order in {-2.5/2,-1/1,0/0,1/1,2.5/2} \draw[very thick, blue, fill] (30:7.5) ++ (\x, 0) circle (0.05) node[black, above] {{\footnotesize $n=\order$}}-- (30:2);
	% diffraction grating
	\draw[thick,fill=orange!35, opacity=0.75] (30:2) ++ (-0.6,0.6) --++ (1.2,0) --++ (0,-1.2) --++ (-1.2,0) --++ (0,1.2);
	\foreach \x in {0.2,0.16,...,-0.2}
	\draw[ultra thin] (30:2) ++ (-\x-0.01,0.4) --++ (0.02,0) --++ (0,-0.8) --++ (-0.02,0) --++ (0,0.8);
	\draw[very thick, blue] (30:2) -- (0,0);
	% laser source
	\draw[white, fill, opacity=0.75] (-0.6, 0.4) rectangle (0.6,-0.4);
	\draw[thick] (-0.6, 0.4) -- (0.6, 0.4) -- (0.6,-0.4);
	\draw[thick] (-0.6, 0.4) --++ (210:2);
	\draw[thick] (0.6,-0.4) --++ (210:2);
	\draw[thick] (0.6, 0.4) --++ (210:2.4);
	% nodes
	\draw (1.2,-1) node {laser source};
	\draw (30:7.5) ++ (0, 1.1) node[above] {screen};
	\draw (30:2) ++ (0, 0.7) node[above, twoline] {diffraction\\grating};
	\end{tikzpicture}
	\caption{multi-slit interference as light passes through a diffraction grating}
\end{figure*}

Compared to double-slit, maxima produced by diffraction grating are \emph{sharper} and \emph{brighter}, maxima are formed when waves from not two but many slits interfere constructively. The location of each maxima from a diffraction grating is given by: \begin{empheq}[box=\tcbhighmath]{equation*}{d \sin \theta = n\lambda}\end{empheq}

where $d$ is the separation between adjacent slits, and light is incident \emph{normally}.

\begin{figure*}[htp]
	\noindent
	\centering
	\begin{minipage}{0.64\textwidth}
		\centering
		\begin{tikzpicture}[scale=1.15]
		\draw[blue,thick] (-3,0) -- (3,0) node[right] {\textcolor{black}{$n=0$}};
		\foreach \x in {1,2,3}{
			\pgfmathsetmacro{\y}{asin(\x*0.3)}
			\draw[blue,thick] (-\y:3) node[right] {\textcolor{black}{$n=\x$}} -- (0,0) -- (\y:3) node[right] {\textcolor{black}{$n=\x$}};
			\draw[->] ({2-\x*0.5}, 0) arc (0:\y:{2-\x*0.5});
		}
		\draw[fill=gray!50, opacity=0.75] (-0.1,-0.4) rectangle (0.1,0.4);
		\foreach \y in {0.36,0.32,...,-0.36} \draw (-0.1,\y) -- (0.1,\y);
		\node at (8:1.75) {{\footnotesize $\theta_1$}};
		\node at (27:1.25) {{\footnotesize $\theta_2$}};
		\node at (51:0.8) {{\footnotesize $\theta_3$}};
		\end{tikzpicture}
  
  \end{minipage}\hfill
	\begin{minipage}{0.35\textwidth}
		\centering
		\begin{tikzpicture}[scale=1.3]
		\foreach \y in {0,1,2,3,4}
		\draw[ultra thick] (0,\y+0.1) -- (0,\y+0.9);
		\foreach \y in {1,2,3,4}{
			\draw[blue] (0,\y) -- ++(30:2);
			\draw (-0.3,\y) -- (-0.7,\y);
		}
		\foreach \y in {2,3,4}{
			\draw[<->] (-0.5,\y-0.97) -- (-0.5,\y-0.03) node[midway,left]{$d$};
			\draw (0,\y) -- ++(-60:.866);
			\draw[very thick, red] (0,\y-1) -- ++(30:.5);
			\draw[red, decoration={markings,mark=at position 1 with {\arrow[scale=2]{>}}},
			postaction={decorate},dashed] (0,\y-1) -- ++(30:.25) to[out=-45,in=135] (2,1);
			\draw (0,\y-.3) arc(270:300:.3);
			\draw (0,\y) ++ (285:.5) node {\footnotesize$\theta$};
		}
		\node[below] at (2,1) {$\Delta L = d\sin\theta$};
		\end{tikzpicture}
	\end{minipage}
 \caption{Diffraction gratings}
 \end{figure*}

You do not need to be able to reproduce this derivation, but it is offered here in contrast to the above. For constructive interference, path difference between adjacent slits must satisfy: $\Delta L = n\lambda $. Since slits are so close together, we consider the light rays passing through the slits are nearly parallel. Using the same small angle trick as before, we obtain: $d \sin \theta = n\lambda$ - this is known as the \emph{diffraction grating equation}.

Note that the greatest angle through which a beam of light can be diffracted is 90$^\circ$ which puts a constraint on the highest order possible:
\begin{equation*}
n_\tmax = \frac{d\sin\theta_\tmax}{\lambda} < \frac{d \sin 90^\circ}{\lambda} \RA \boxed{n_\tmax < \frac{d}{\lambda}}
\end{equation*}

The value of $\frac{d}{\lambda}$ should be rounded \emph{down} to the nearest whole number to give $n_\tmax$

If white light is passed through diffraction gratings, \emph{dispersion} is observed, this is the cause of the pretty colours seen in a dvd or cd:

\begin{figure*}[!ht]
	\centering
	\begin{tikzpicture}[scale=1.6]
	% background
	\draw[fill=gray!50] (-3.5,-2.7) rectangle (4.2,2.7);
	% zeroth order
	\draw[white,ultra thick] (-3,0) -- (3,0);
	% first and second order
	\foreach \WL/\yanse in {6.4/red, 6.0/orange, 5.6/yellow, 5.2/green, 4.8/blue, 4.4/violet}{
		\pgfmathsetmacro{\ystart}{asin(\WL/18)}
		\pgfmathsetmacro{\yend}{asin((\WL+0.4)/18)}
		\draw[\yanse!80, fill] (0,0) -- (\ystart:3) arc(\ystart:\yend:3) -- cycle;
		\draw[\yanse!80, fill] (0,0) -- (-\ystart:3) arc(-\ystart:-\yend:3) -- cycle;
		\pgfmathsetmacro{\zstart}{asin(\WL/9)}
		\pgfmathsetmacro{\zend}{asin((\WL+0.4)/9)}
		\draw[\yanse!50, fill] (0,0) -- (\zstart:3) arc(\zstart:\zend:3) -- cycle;
		\draw[\yanse!50, fill] (0,0) -- (-\zstart:3) arc(-\zstart:-\zend:3) -- cycle;
	}
	% order labels
	\foreach \t/\nlabel in {0/0, 5.4/1, 10.8/2, -5.4/1, -10.8/2}{
		\pgfmathsetmacro{\y}{asin(\t/18)}
		\node at (\y:3.5) {$n=\nlabel$};
	}
	% diffraction grating
	\foreach \y in {0.36,0.32,...,-0.36} \draw (-0.1,\y) -- (0.1,\y);
	\draw[fill=gray!20, opacity=0.7] (-0.1,-0.4) rectangle (0.1,0.4);
	\end{tikzpicture}
	\caption{dispersion of white light through a diffraction grating}
\end{figure*}



\newpage
Some notes on the peculiarities of this effect:
\begin{compactitem}
	\item[--] At zeroth order, diffraction grating equation $d\sin\theta_0 = 0\cdot\lambda$ is satisfied at $\theta_0 = 0$ for any $\lambda$, so we observe a collection of all colours, i.e., a white central maxima.
	
	\item[--] At first order, $d\sin\theta_1 = 1\cdot\lambda$, maxima for long wavelengths appear at greater angles, so the spectrum spreads into a rainbow band, with red light at outer end and violet closer to the centre.
	
	\item[--] At higher orders, the spectrum would be more spread out as the frequency dependence of diffracted angle becomes more pronounced. The spectra of different orders may even overlap to give complicated combinations.
\end{compactitem}

The diffraction grating is an import device in analysis of light - it is widely used in measurement of wavelength of light and in telecommunications.

\example{Light of 632 nm wavelength produces first-order maxima at 16$^\circ$ when passing through a diffraction grating at right angles. How many lines per millimetre are there in the diffraction grating?}

\begin{soln} slit separation: $d = \frac{n\lambda}{\sin\theta} = \frac{1\times632\times10^{-9}}{\sin 16^\circ} \RA d \approx 2.29 \times 10^{-6} \text{ m}$

number of lines in 1 mm: $N = \frac{1 \text{ mm}}{2.29 \times 10^{-6} \text{ m}} = \frac{1 \times 10^{-3}}{2.29 \times 10^{-6}} \RA N \approx 436$ \end{soln}

\example{Green light of 510 nm wavelength is incident normally on a diffraction grating with slit separation of 2.0 $\mu$m. (a) What is the highest order seen? (b) how many maxima are produced? (c) If blue light is used for the same diffraction grating experiment, suggest the effect on the diffraction pattern.}

\begin{soln} to find highest order, we have: $n_\tmax < \frac{d\sin 90^\circ}{\lambda} = \frac{2.0\times10^{-6}}{510\times10^{-9}} \approx 3.9 \RA n_\tmax = 3$

1st, 2nd, 3rd-order on either side and 0th-order at centre, so $ 2\times 3 + 1 = 7$ maxima are formed

blue light has shorter wavelength, so more maxima will be produced with blue light, and separation between each maxima would be closer \end{soln}

\example{Light of wavelength 630 nm produces a second-order maxima at angle of $60^\circ$ when it is directed at a diffraction grating. If visible light of another wavelength can also give a maximum at the same angle, what is this wavelength?}

\begin{soln} using diffraction grating equation: $d \sin\theta = n\lambda$, we find $n\lambda$ is constant for fixed $\theta$

from information given in the question, this constant is $ 2\times630 = 1260 \text{ nm}$

if $n=1$, $\lambda = 1260 \text{ nm}$, which would be infra-red

if $n=3$, $\lambda = 420 \text{ nm}$, which would be visible (violet)

if $n\geq 4$, $\lambda \leq 315 \text{ nm}$, which would be ultraviolet

so the desired wavelength is 420 nm \end{soln}


\subsection{Diffraction}

A wave has the ability to bend around obstacles or pass through narrow gaps

\begin{ilight}
	When a wave passes through a aperture/slit/gap/hole or encounters an obstacle, it spreads out/around the corners, this is known as \keypoint{diffraction}\index{diffraction} of waves.
\end{ilight}

Diffraction is a general property of all waves. Some examples of diffraction are:

\begin{compactitem}
	\item[--] sound waves can diffract through an open door
	
	so you can hear people in the next room talking, even though you cannot see them
	
	\item[--] light ray can bend when it goes through/around a small hole/small particles
	
	sky appears red at sunset as red light is diffracted most by dust particles in atmosphere
\end{compactitem}

Diffraction is a scale-dependent effect - a wave is diffracted most when its wavelength is close to size of aperture/obstacle. Because a single object is close to an ``infinite gap'', longer wavelengths usually diffracts more around an edge than a wave of shorter wavelength.

\begin{figure*}[htp]
	\centering
	\begin{tikzpicture}[scale=1.25]
	\draw[ultra thick] (0,2.4) -- (0,0.25) (0,-0.25) -- (0,-2.4);
	\foreach \x in {0.25,0.75,1.25,1.75,2.25}{
		\draw[blue,thick] (-\x,-2.4) -- (-\x,2.4);
		\draw[blue,thick] (-65:\x) arc(-65:65:\x);
	}
	\draw[ultra thick] (6,2.4) -- (6,0.8) (6,-0.8) -- (6,-2.4);
	\foreach \x in {0.25,0.75,1.25,1.75,2.25}{
		\draw[blue,thick] (-\x+6,-2.4) -- (-\x+6,2.4);
		\draw[blue,thick] (6,-0.8) ++ (-25:\x) to[out=65,in=270] (6+\x,-0.8-\x/6) -- (6+\x,0.8+\x/6) to[out=90,in=-65] (6+0.9063*\x,0.8+0.4226*\x);
	}
	\end{tikzpicture}
	
	\caption{diffraction of a wave as it passes through an aperture of width that is (a) close to the wavelength, (b) greater than the wavelength}
\end{figure*}

If we go beyond the GCSE understanding of diffraction as above we can we can see that the light passing through a single slit forms a diffraction pattern somewhat different from those formed by double slits or diffraction gratings, which we discuss earlier in the chapter on interference.
\begin{figure*}
    \centering
\begin{tikzpicture}[
    nodal/.style={mylightgreen,dashed,very thin},scale=1.5
  ]
  
  \def\L{3.8}       % distance between walls
  \def\H{5.4}       % total wall height
  \def\h{2.8}       % plane wave height
  \def\t{0.15}      % wall thickness
  \def\a{0.8}       % slit distance
  \def\d{0.20}      % slit size
  \def\N{21}        % number of waves
  \def\dN{1}        % shift radius
  \def\A{1.6}       % amplitude
  \def\lambd{0.20}  % wavelength
  \def\R{\N*\lambd} % wave radius
  \def\Nfringes{2}  % number of nodal lines
  \def\nshades{70}  % number of shades
  \def\nsamples{40}
  \def\ang{40}
  
  % PLANE WAVES
  \foreach \i [evaluate={\xp=-2*(\i-0.5)*\lambd; \xm=-2*(\i-1)*\lambd;}] in {1,...,3}{
    \draw[myblue,line width=0.8] (\xp,-\h/2) -- (\xp,\h/2);
    \draw[myblue,line width=0.1] (\xm,-\h/2) -- (\xm,\h/2);
  }
  
  % WAVES
  \begin{scope}
    \clip (-\t/2,-\H/2) rectangle (\L,\H/2);
    \foreach \i [evaluate={\Rp=2*(\i+\dN)*\lambd; \Rm=2*(\i+\dN+0.5)*\lambd;}] in {1,...,\N}{
      \draw[myblue,line width=0.8] (-2.2*\lambd,0)++(\ang:\Rp) arc (\ang:-\ang:\Rp);
      \draw[myblue!80,line width=0.1] (-2.2*\lambd,0)++(\ang:\Rm) arc (\ang:-\ang:\Rm);
    }
    \foreach \m [evaluate={\ddy=0.1*\m;}] in {1,2,2.7}{ % mask at destructive interference
      \foreach \i [evaluate={\ym=\L/sqrt((\a/(\lambd*\m))^2-1); \dy=0.55*\i/\nshades;}]
        in {1,...,\nshades}{
        \fill[white,opacity=0.045] (0.05,0.005+\ddy) -- (\L,{\dy+\ym}) --++ (0,-2*\dy) -- (0.05,-0.005+\ddy);
        \fill[white,opacity=0.045] (0.05,0.005-\ddy) -- (\L,{\dy-\ym}) --++ (0,-2*\dy) -- (0.05,-0.005-\ddy);
      }
    }
  \end{scope}
  
  % WALL
  \fill[wall]
    (\t/2,\a/2) rectangle (-\t/2,\H/2)
    (\t/2,-\a/2) rectangle (-\t/2,-\H/2)
    (\L,-\H/2) rectangle (\L+\t,\H/2);
  
  % SHADES
  \begin{scope}[shift={(1.08*\L,0)}]
    \clip (0,-\H/2) rectangle (1.1*\A,\H/2);
    \fill[myshadow] (0,-\H/2) rectangle (\A,\H/2); % to fill seams
    \def\yz{\L/sqrt((\a/\lambd)^2-1)} %\L/sqrt((\a/\lambd)^2-1)
    \path [left color=myshadow,right color=myshadow,middle color=white,shading angle={180}]
      (0,{-\yz}) rectangle (\A,{\yz});
    \foreach \i [evaluate={
                  \n=\i;
                  \m=\i+1;
                  \yn=\L/sqrt((\a/(\lambd*\n))^2-1); %\L*\n*\lambd/\a
                  \ym=\L/sqrt((\a/(\lambd*\m))^2-1);
                  \dang=mod(\i,2)*180;
                 }] in {1,...,\Nfringes}{
      \path [left color=myshadow,right color=myshadow,middle color=myshadow!50,shading angle={\dang}]
        (0,\yn) rectangle (\A,\ym);
      \path [left color=myshadow,right color=myshadow,middle color=myshadow!50,shading angle={180+\dang}]
        (0,-\yn) rectangle (\A,-\ym);
    }
  \end{scope}
  
  % INTENSITY
  \begin{scope}[shift={(1.2*\L+1.1*\A,0)}]
    \draw[->,thick] (-0.08*\A,0) -- (1.3*\A,0) node[right=-2] {$\expval{I}$};
    \draw[->,thick] (0,-0.52*\H) -- (0,0.54*\H) node[right] {$y$};
    \draw[myred,thick,variable=\y,samples=\nsamples,smooth,domain=-\H/2:\H/2]
      plot({\A*int_one(\y,\lambd,\a,\L)},\y);
    \foreach \i [evaluate={\y=\L/sqrt((\a/(\lambd*\i))^2-1);}] in {1,...,\Nfringes}{
      \tick{0,-\y}{180} node[right=-1,scale=0.9] {$m=-\i$};
      \tick{0,\y}{180} node[right=-1,scale=0.9] {$m=+\i$};
    }
  \end{scope}
  
\end{tikzpicture}
    
    \caption{Single slit diffraction}
    \label{single slit}
\end{figure*}

Figure \ref{single slit} Shows a single-slit diffraction pattern as we would obseve with eg. a laser. Note that the central maximum is larger than maxima on either side and that the intensity decreases rapidly on either side. In contrast, a diffraction grating produces evenly spaced lines that dim slowly on either side of the center.

\begin{figure*}
\centering
% DIFFRACTION - ONE SLIT - PATH DIFFERENCE close-up
\begin{tikzpicture}[scale=1.5]
  \message{Diffraction, path difference (close-up)^^J}
  
  \def\L{6.15}      % distance between walls
  \def\l{4.0}       % path length
  \def\H{4.3}       % total wall height
  \def\f{0.9}       % fractional height of projection point
  \def\t{0.15}      % wall thickness
  \def\a{2.5}       % slit distance
  \def\ang{26}      % angle
  \coordinate (T) at (0,\a/2);
  \coordinate (B) at (0,-\a/2);
  \coordinate (L) at (0,0);
  \coordinate (R) at (\L,0);
  \coordinate (I) at ({\a/2/tan(\ang)},0);
  
  % LINES
  \draw[mygreen,thick] (T) --++ (\ang:.8*\l) coordinate (PT); %node[midway,below=1,above left=-2] {$r_1$};
  \draw[mygreen,thick] (B) --++ (\ang:\l) coordinate (PB); %node[midway,left=2,below right=-1] {$r_2$};
  \draw[mygreen,thick] (L) --++ (\ang:.9*\l) coordinate (PR);
  \draw[dashed] (T) -- (B);
  \draw[mydarkred,dashed] (T) -- ($(L)!(T)!(PR)$) coordinate (M);
  %\draw[black!60,dashed] (T) --++ (.20*\L,0) coordinate (TR);
  \draw[black!60,dashed] (L) --++ (.6*\L,0) coordinate (LR);
  %\draw[black!60,dashed] (B) --++ (.20*\L,0) coordinate (BR);
  
  % LINE END
  \lineend{PT}{\ang+70}
  \lineend{PR}{\ang+70}
  \lineend{PB}{\ang+70}
  
  % ANGLES
  \draw pic[mysmallarr,"$\theta$",mydarkred,draw=mydarkred,angle radius=20,angle eccentricity=1.3]
    {angle = B--T--M}; %\contour{white}{}
  \draw pic[mysmallarr,"$\theta$",mydarkred,draw=mydarkred,angle radius=24,angle eccentricity=1.2]
    {angle = LR--L--PR};
  \draw pic[mysmallarr,"$\theta$",mydarkred,draw=mydarkred,angle radius=24,angle eccentricity=1.2]
    {angle = LR--I--PB};
  \rightAngle{T}{M}{PR}{0.3}
  
  % MEASURES
  \draw[<->,black] (-2.1*\t,0) --++ (0,\a/2) node[midway,fill=white,inner sep=0.1] {$a/2$};
  \draw[<->,black] (-2.1*\t,0) --++ (0,-\a/2) node[midway,fill=white,inner sep=0.1] {$a/2$};
  \draw[<->,black] ([shift={({\ang-90}:0.1)}]L) -- ([shift={({\ang-90}:0.1)}]M)
    node[midway,left=2,below right=0]{$\frac{a}{2}\!\sin\theta$};
  
  % WALL
  \fill[wall]
    (0,\a/2) rectangle (-\t,\H/2)
    (0,-\a/2) rectangle (-\t,-\H/2);
  
\end{tikzpicture}
\label{singleslitmaths}
\caption{Analysis of single slit diffraction}
 %\forceversofloat
\end{figure*}

The analysis of single-slit diffraction is illustrated in Figure \ref{singleslitmaths} in which the light arrives at the slit, illuminating it uniformly and is in phase across its width. We then consider light propagating onwards from different parts of the same slit. 

Assuming the screen is very far away compared with the size of the slit, rays heading toward a common destination are nearly parallel. When they travel straight ahead, they remain in phase, and we observe a central maximum. However, when rays travel at an angle $\theta$ relative to the original direction of the beam, each ray travels a different distance to a common location, and they can arrive in or out of phase. The ray from the bottom travels a distance of one wavelength farther than the ray from the top. Thus, a ray from the center travels a distance $\lambda/2$ less than the one at the bottom edge of the slit, arrives out of phase, and interferes destructively. A ray from slightly above the center and one from slightly above the bottom also cancel one another. In fact, each ray from the slit interferes destructively with another ray. In other words, a pair-wise cancellation of all rays results in a dark minimum in intensity at this angle. By symmetry, another minimum occurs at the same angle to the right of the incident direction (toward the bottom of the figure) of the light.

% \begin{tikzpicture}
%   \message{Diffraction, path difference^^J}
  
%   \def\L{5.9}       % distance between walls
%   \def\H{4.3}       % total wall height
%   \def\f{0.95}      % fractional height of projection point
%   \def\ang{atan((\f*\H+\a/2)/\L/2)} % theta
%   \def\t{0.15}      % wall thickness
%   \def\a{2.5}       % slit distance
%   \coordinate (T) at (0,\a/2);
%   \coordinate (B) at (0,-\a/2);
%   \coordinate (L) at (0,0);
%   \coordinate (R) at (\L,0);
%   \coordinate (P) at (\L,\f*\H/2);
%   \coordinate (M) at ($(L)!(T)!(P)$);
  
%   % LINES
%   \draw[mygreen,thick] (T) -- (P); %node[midway,above=-1] {$r_1$};
%   \draw[mygreen,thick] (B) -- (P); %node[midway,below=3,right=6] {$r_2$}; %right=6,below right=-4
%   \draw[mygreen,thick] (L) -- (P);
%   \draw[dashed] (T) -- (B);
%   \draw[dashed,black!60] (L) -- (R);
%   \draw[mydarkred,dashed] (M) -- (T);
  
%   % ANGLES
%   \draw pic[mysmallarr,"$\theta'$",mydarkred,draw=mydarkred,angle radius=25,angle eccentricity=1.2]
%     {angle = B--T--M};
%   \draw pic[mysmallarr,"$\theta$",mydarkred,draw=mydarkred,angle radius=31,angle eccentricity=1.14]
%     {angle = R--L--P};
%   \rightAngle{T}{M}{P}{0.3}
  
%   % MEASURES
%   \draw[<->,black] (0,-0.47*\H) --++ (\L,0) node[midway,fill=white,inner sep=1] {$L$};
%   \draw[<->,black] (-2.1*\t,0) --++ (0,\a/2) node[midway,fill=white,inner sep=0.1] {$a/2$};
%   \draw[<->,black] (-2.1*\t,0) --++ (0,-\a/2) node[midway,fill=white,inner sep=0.1] {$a/2$};
%   \draw[<->,black] ([shift={({\ang-100}:0.1)}]L) -- ([shift={({\ang-100}:0.1)}]M)
%     node[midway,left=3,below right=-2]{$\frac{a}{2}\sin\theta'$};
%   \draw[<->,black] ([shift={(2.1*\t,0)}]P) -- ([shift={(2.1*\t,0)}]R)
%     node[midway,fill=white,inner sep=1]{$y$};
  
%   % WALL
%   \fill[wall]
%     (0,\a/2) rectangle (-\t,\H/2)
%     (0,-\a/2) rectangle (-\t,-\H/2)
%     (\L,-\H/2) rectangle (\L+\t,\H/2);
%   \fill[mygreen!80!black] (P) circle (0.3*\t) node[right=1,above left=-1] {P};
  
% \end{tikzpicture}


As seen in the figure, the difference in path length for rays from either side of the slit is  $a \sin{\theta}$  and we see that a destructive minimum is obtained when this distance is an integral multiple of the wavelength.

Thus, to obtain destructive interference for a single slit,

$$ a \sin{\theta} = m \lambda$$ 
Where:
m is the order of the minimum.\\
a is the slit width,\\
$\lambda$  is the light’s wavelength,\\
$\theta$  is the angle relative to the original direction of the light\\

% \begin{figure}
%     \centering
%     % TWO SPLIT: interference x diffraction
% \begin{tikzpicture}
%  \def\A{3.5}
%  \def\d{3.4}     % slit distance
%  \def\a{0.9}     % slit size
%  \def\lambd{0.2} % wavelength
%  \def\k{15}      % x -> theta
%  \def\xmax{4.6}
%  \def\ymax{\A}
%  \def\N{4}
%  \def\D{2.3*\xmax}
%  \def\L{4.0}
%  \def\nsamples{80}
 
%  \draw[->,thick,black] (-1.05*\xmax,0) -- (1.1*\xmax,0) node[below] {$\theta$};
%  \draw[->,thick,black] (0,-0.1*\A) -- (0,1.12*\A);
%  \draw[wave,variable=\t,samples=\nsamples,smooth,domain=-\xmax:\xmax]
%     plot(\t,{\A*int_one_ang(\k*\t,\lambd,\a)});
 
%  \foreach \i [evaluate={\t=asin(\i*\lambd/\a)/\k;}] in {1,...,\N}{
%     \tick{-\t,0}{90} node[below=-1,scale=0.7] {$m=-\i$};
%     \tick{\t,0}{90} node[below=-1,scale=0.7] {$m=+\i$};
%  }
 
% \end{tikzpicture}
%     \caption{Caption}
%     \label{fig:enter-label}
% \end{figure}


\newpage




\subsection{Stationary waves}

\subsection{Formation of stationary waves}

Stationary waves are often formed when a travelling wave is \emph{reflected}. Forward-moving wave and backward-moving reflected wave combine to form stationary wave.

\begin{ilight}
	When two waves of same frequency and same speed travel in opposite directions and meet together, they superpose to form stable regions of constructive and destructive interference, this forms a \keypoint{stationary wave}, also known as a \keypoint{standing wave}\index{stationary wave}
\end{ilight}

\begin{figure*}[htp]
	\centering
	\noindent\begin{minipage}{0.25\linewidth}
		\centering
		\begin{tikzpicture}[scale=0.5]
		\foreach \y in {0, -3.5, -7.5}
		\draw[thin, dashed] (-pi, \y) -- (pi, \y);
		\draw [thick,blue,domain=-pi:pi,samples=24,smooth,variable=\x] plot (\x,{sin(\x r)});
		\draw[->,blue] (-0.5,1.5) -- (0.5,1.5);
		\draw [thick,Green,domain=-pi:pi,samples=24,smooth,variable=\x] plot (\x,{sin(-\x r)-3.5});
		\draw[<-,Green] (-0.5,1.5-3.5) -- (0.5,1.5-3.5);
		\draw [very thick,red,domain=-pi:pi,samples=5,smooth,variable=\x] plot (\x,{-7.5});
		\node at (0,-11) {$t=0$};
		\node[white] at (2,-11) {$\frac{1}{1}$};
		\end{tikzpicture}
	\end{minipage}\hfil
	\begin{minipage}{0.25\linewidth}
		\centering
		\begin{tikzpicture}[scale=0.5]
		\foreach \y in {0, -3.5, -7.5}
		\draw[thin, dashed] (-pi, \y) -- (pi, \y);
		\draw [thick,blue,domain=-pi:pi,samples=24,smooth,variable=\x] plot (\x,{-cos(\x r)});
		\draw[->,blue] (-0.5,1.5) -- (0.5,1.5);
		\draw [thick,Green,domain=-pi:pi,samples=24,smooth,variable=\x] plot (\x,{-cos(\x r)-3.5});
		\draw[<-,Green] (-0.5,1.5-3.5) -- (0.5,1.5-3.5);
		\draw [very thick,red,domain=-pi:pi,samples=24,smooth,variable=\x] plot (\x,{-2*cos(\x r)-7.5});
		\node at (0,-11) {$t=\frac{T}{4}$};
		\end{tikzpicture}
	\end{minipage}\hfil
	\begin{minipage}{0.25\linewidth}
		\centering
		\begin{tikzpicture}[scale=0.5]
		\foreach \y in {0, -3.5, -7.5}
		\draw[thin, dashed] (-pi, \y) -- (pi, \y);
		\draw [thick,blue,domain=-pi:pi,samples=24,smooth,variable=\x] plot (\x,{-sin(\x r)});
		\draw[->,blue] (-0.5,1.5) -- (0.5,1.5);
		\draw [thick,Green,domain=-pi:pi,samples=24,smooth,variable=\x] plot (\x,{-sin(-\x r)-3.5});
		\draw[<-,Green] (-0.5,1.5-3.5) -- (0.5,1.5-3.5);
		\draw [very thick,red,domain=-pi:pi,samples=5,smooth,variable=\x] plot (\x,{-7.5});
		\node at (0,-11) {$t=\frac{T}{2}$};
		\end{tikzpicture}
	\end{minipage}\hfil
	\begin{minipage}{0.25\linewidth}
		\centering
		\begin{tikzpicture}[scale=0.5]
		\foreach \y in {0, -3.5, -7.5}
		\draw[thin, dashed] (-pi, \y) -- (pi, \y);
		\draw [thick,blue,domain=-pi:pi,samples=24,smooth,variable=\x] plot (\x,{cos(\x r)});
		\draw[->,blue] (-0.5,1.5) -- (0.5,1.5);
		\draw [thick,Green,domain=-pi:pi,samples=24,smooth,variable=\x] plot (\x,{cos(\x r)-3.5});
		\draw[<-,Green] (-0.5,1.5-3.5) -- (0.5,1.5-3.5);
		\draw [very thick,red,domain=-pi:pi,samples=24,smooth,variable=\x] plot (\x,{2*cos(\x r)-7.5});
		\node at (0,-11) {$t=\frac{3T}{4}$};
		\end{tikzpicture}
	\end{minipage}
		\caption{The formation of a stationary wave (red) due to the superposition of one wave travelling to the right (blue) and another wave travelling to the left (green)}
\end{figure*}

from this one can visualise how the pattern of a stationary wave varies with time (see Figure \ref{fig:statwave}. There are points in the wave where destructive interference always occurs, these points are called \keypoint{nodes}, which have zero amplitudes. \index{node}

The points oscillating with the greatest amplitudes are called \keypoint{anti-nodes}. \index{anti-node} Points other than nodes and anti-nodes oscillate with different amplitudes. 

\begin{figure*}[htp]
	\centering
	\begin{tikzpicture}[scale=0.7]
	\foreach \amp in {-2,-1,1,2}
	\draw [thick,red,domain=-3*pi:3*pi,samples=60,smooth,variable=\x] plot (\x,{\amp*sin(\x r)});
	\draw [thick,blue] (-3*pi,0) -- (3*pi,0);
	\foreach \idx in {-1,0,1}{
		\draw[<->] (\idx*2*pi-pi+0.03,-2.5) -- ++(2*pi-0.06,0) node[midway,above]{$\lambda$};
		\draw (\idx*2*pi-pi,-2.8) -- ++(0,0.6) ++(2*pi,0) -- ++(0,-0.6);
	}
	\foreach \n in {-3,-2,...,3} \node at (\n*pi,0.7) {N};
	\foreach \a in {-2.5,-1.5,...,2.5} \node at (\a*pi,2.25) {A};
	\end{tikzpicture}
	\caption{variation of a stationary wave pattern with time}
 \label{fig:statwave}
\end{figure*}
The distance between two adjacent nodes is half of one wavelength. This allows us to say:: \begin{empheq}[box=\tcbhighmath]{equation*}{\text{the wavelength }\lambda = 2 \times \text{node-to-node distance}}\end{empheq}



\subsection{Stationary wave pattern \& boundary conditions}

Patterns of stationary waves depend heavily on the boundary conditions
The end can be \emph{fixed}, also called a \emph{closed} end, to form a \emph{node} - such as with a wave on a string.

The end can also be \emph{free} to oscillate, known as an \emph{open} end, to form \emph{an anti-node} - such as you might find in a flute.



\subsection*{Stationary wave between two fixed ends}

With both ends fixed, the two ends are both nodes, three patterns with the longest wavelengths and lowest frequencies are shown:

\begin{figure*}[htp]
	\centering
	\noindent\begin{minipage}{0.3\linewidth}
		\centering
		\begin{tikzpicture}[scale=0.7]
		\draw [thick,blue,domain=-pi:pi,samples=24,smooth,variable=\x] plot (\x,{cos(\x/2 r)});
		\draw [thick,dashed,blue,domain=-pi:pi,samples=24,smooth,variable=\x] plot (\x,{-cos(\x/2 r)});
		\node at (0,-2) {fundamental mode};
		\end{tikzpicture}
	\end{minipage}\hfil
	\begin{minipage}{0.3\linewidth}
		\centering
		\begin{tikzpicture}[scale=0.7]
		\draw [thick,blue,domain=-pi:pi,samples=24,smooth,variable=\x] plot (\x,{sin(\x r)});
		\draw [thick,dashed,blue,domain=-pi:pi,samples=24,smooth,variable=\x] plot (\x,{-sin(\x r)});
		\node at (0,-2) {2nd harmonic};
		\end{tikzpicture}
	\end{minipage}\hfil
	\begin{minipage}{0.3\linewidth}
		\centering
		\begin{tikzpicture}[scale=0.7]
		\draw [thick,blue,domain=-pi:pi,samples=24,smooth,variable=\x] plot (\x,{cos(3*\x/2 r)});
		\draw [thick,dashed,blue,domain=-pi:pi,samples=24,smooth,variable=\x] plot (\x,{-cos(3*\x/2 r)});
		\node at (0,-2) {3rd harmonic};
		\end{tikzpicture}
	\end{minipage}\hfil
\end{figure*}

The longest-wavelength mode is called the \emph{fundamental mode}, it also has lowest frequency. Modes with higher frequencies are usually referred to as \emph{excited modes}, or \emph{harmonics}\footnote{These modes are closely related to the notes played on musical instruments.}
in many texts different modes are labelled as 1st, 2nd, 3rd \emph{harmonics}, etc.

If two ends are separated by a distance of $L$, then wavelengths for different modes are:
\begin{equation*}
\lambda_1 = 2L \qquad \lambda_2 = L \qquad \lambda_3 = \frac{2L}{3} \qquad \cdots \qquad \lambda_n = \frac{2L}{n}
\end{equation*}

Making comparison with the fundamental wavelength $\lambda_1$ $\lambda_1$, we have
\begin{equation*}
\lambda_2 = \frac{\lambda_1}{2} \qquad \lambda_3 = \frac{\lambda_1}{3} \qquad \cdots \qquad \lambda_n = \frac{\lambda_1}{n}
\end{equation*}

Since wave speed remains constant as wave travels through the same medium, the frequency of each mode is inversely proportional to its wavelength, hence,  
\begin{equation*}
f_2 = 2f_1 \qquad f_3 = 3f_1 \qquad \cdots \qquad f_n=n f_1
\end{equation*}

\example{\emph{Melde's experiment} demonstrates a stationary wave formed on a vibrating string using an oscillator set at a particular frequency. Given that the length of the string is 80 cm. (a) What is the wavelength of this wave? (b) If one wants to construct a stationary of longer wavelength, suggest what changes can be made?}

\begin{figure*}[ht]
	\centering
	\begin{tikzpicture}[scale=1.1]
	\draw [thick,domain=0:10,samples=90,smooth] plot (\x,{.6*sin(\x*90)});
	\draw [thick,domain=0:10,samples=90,smooth] plot (\x,{-.6*sin(\x*90)});
	\draw[fill] (0,0) circle (0.1);
	\draw[fill] (10,0) circle (0.1);
	\draw[thick] (10,-1) -- (10,1);
	\foreach \y in {0.9,0.6,...,-0.7} \draw (10,\y) --++ (0.2,-0.2);
	\draw (-0.03, 0) rectangle (0.03, -1);
	\draw[thick] (-.8, -1) rectangle (.8,-2);
	\node at (0,-1.5) {oscillator};
	\node[twoline, right] at (10.2, 0) {fixed\\point};
	\end{tikzpicture}
	\caption{Melde's experiment: demonstration of stationary waves on a tense string}
\end{figure*}

\begin{soln} Five loops so five half-wavelengths $\RA L = \frac{5}{2}\lambda \RA \lambda = \frac{2}{5}L = \frac{2}{5} \times 80 \RA \lambda = 32\text{ cm}$

to have larger wavelength, one could do either of the following:
use a string with greater length (while keeping frequency of oscillator unchanged)
reduce the frequency of oscillation (constant wave speed, so $\lambda$ increases)

\titem increase tension in string (wave speed increases \end{soln}\footnote{The speed $v$ of a progressive wave travelling along a string of length $L$ and mass $m$ can be given by the equation: $v = \sqrt{\frac{TL}{m}}$, where $T$ is the tension in the string.} for same frequency, so $\lambda$ increases)




\subsection*{Stationary wave between one fixed end and one open end}

In this case, we have a node at one end and an anti-node at the other end, again we give the three patterns with the longest wavelengths and lowest frequencies.

\begin{figure*}[htp]
	\centering
	\noindent\begin{minipage}{0.3\linewidth}
		\centering
		\begin{tikzpicture}[scale=0.7]
		\draw [thick,blue,domain=0:2*pi,samples=24,smooth,variable=\x] plot (\x,{sin(\x/4 r)});
		\draw [thick,dashed,blue,domain=0:2*pi,samples=24,smooth,variable=\x] plot (\x,{-sin(\x/4 r)});
		\node at (pi,-2) {fundamental mode};
		\end{tikzpicture}
	\end{minipage}\hfil
	\begin{minipage}{0.3\linewidth}
		\centering
		\begin{tikzpicture}[scale=0.7]
		\draw [thick,blue,domain=0:2*pi,samples=24,smooth,variable=\x] plot (\x,{sin(3*\x/4 r)});
		\draw [thick,dashed,blue,domain=0:2*pi,samples=24,smooth,variable=\x] plot (\x,{-sin(3*\x/4 r)});
		\node at (pi,-2) {2nd harmonic};
		\end{tikzpicture}
	\end{minipage}\hfil
	\begin{minipage}{0.3\linewidth}
		\centering
		\begin{tikzpicture}[scale=0.7]
		\draw [thick,blue,domain=0:2*pi,samples=24,smooth,variable=\x] plot (\x,{sin(5*\x/4 r)});
		\draw [thick,dashed,blue,domain=0:2*pi,samples=24,smooth,variable=\x] plot (\x,{-sin(5*\x/4 r)});
		\node at (pi,-2) {3rd harmonic};
		\end{tikzpicture}
	\end{minipage}\hfil
%	\caption{The three lowest frequency modes of wave patterns formed between one fixed end and one open end}
\end{figure*}

\cmt longest-wavelength/lowest-frequency mode is also called the \emph{fundamental mode}

excited modes are called \emph{harmonics} as before

\cmt wavelengths for allowed patterns under the boundary conditions are
\begin{equation*}
\lambda_1 = 4L \qquad \lambda_2 = \frac{4L}{3} \qquad \lambda_3 = \frac{4L}{5} \qquad \cdots \qquad \lambda_n = \frac{4L}{2n-1}
\end{equation*}

comparing with the fundamental wavelength $\lambda_1$, we find
\begin{equation*}
\lambda_2 = \frac{\lambda_1}{3} \qquad \lambda_3 = \frac{\lambda_1}{5} \qquad \cdots \qquad \lambda_n = \frac{\lambda_1}{2n-1}
\end{equation*}

again the wave speed should be the same for all modes, so the allowed frequencies are
\begin{equation*}
f_2 = 3f_1 \qquad f_3 = 5f_1 \qquad \cdots \qquad f_n=(2n-1) f_1
\end{equation*}

\begin{marginfigure}
	\vspace*{-5pt}
	\centering
	\begin{tikzpicture}[xscale=0.64, yscale=0.85]
	\draw[<->] (9,1.8) -- (0,1.8) node[midway,above]{30 cm}; 
	\draw[very thick] (9,1) -- (0,1) -- (0,-1) -- (9,-1);
	\draw [thick,domain=0:9,samples=90,smooth] plot (\x,{0.95*sin(\x*30)});
	\draw [thick,domain=0:9,samples=90,smooth] plot (\x,{-0.95*sin(\x*30)});
	\end{tikzpicture}
	\vspace*{-16pt}
\end{marginfigure}


\example{The graph shows a standing sound wave produced in an air column in a closed pipe of length 30 cm. The frequency of this sound wave is 825 Hz. (a) What is the wavelength of the sound wave? (b) Calculate the speed of sound. (c) What is the frequency of the fundamental mode for sound wave in this air column? (d) Describe the motion of an air molecule near the open end.}
	
\begin{soln} (a) wavelength of this excited mode: $\lambda = \frac{4}{3}L = \frac{4}{3}\times 30 \times 10^{-2} \RA \lambda = 0.40 \text{ m}$

(b) wave speed: $v = \lambda f = 0.40 \times 825 \RA v = 330 \mps$

(c) wavelength of fundamental mode: $\lambda = 4L = 4\times 30\times10^{-2} \RA \lambda = 1.2 \text{ m}$

\hspace*{1.25em} fundamental frequency: $f = \frac{v}{\lambda} = \frac{330}{1.2} \RA f = 275 \text{ Hz}$

(d) note that an anti-node is formed near open end, also sound wave is longitudinal

\hspace*{1.25em} so air molecule near open end vibrate \emph{horizontally} with greatest amplitude \end{soln}



\begin{marginfigure}
	\vspace*{-8pt}
	\centering
	\begin{tikzpicture}[scale=1.5]
	\draw[cyan!40, fill] (-0.5,0) rectangle (0.5, 2.5);
	\draw[cyan!40, fill] (-0.8,0.3) rectangle (-0.5, 0.6);
	\draw[thick] (-0.8, 0.3) -- (-0.5, 0.3) -- (-0.5, 0) -- (0.5, 0) -- (0.5, 4); 
	\draw[thick] (-0.8, 0.6) -- (-0.5, 0.6) -- (-0.5, 4); 
	\draw[ultra thick] (-0.4, 4.5) -- (0.3, 4.5) arc(90:-90:0.1) -- (-0.4, 4.3);
	\draw[ultra thick] (0.4, 4.4) --++ (0.4,0);
	\draw (0.6, 4.5) --++ (0.4, 0.3) node[right, twoline]{tuning\\fork};
	\draw (0.3, 2) --++ (0.6, 0.4) node[right]{water};
	\draw[thick,->] (-0.9, 0.45) --++ (-0.4,0);
	\draw[very thick] (-1,0) -- (1,0);
	\end{tikzpicture}
	\vspace*{10pt}
\end{marginfigure}


\example{\label{recpracwave}A tuning fork is held above a tall cylinder which is initially filled with water. The water level, measured from the bottom of the cylinder, is lowered at a constant rate. A loud sound is first heard when the water level is at 60.0 cm, and the next loud sound is heard when the water level is at 26.8 cm. It is known that the speed of sound in air is $340 \mps$. What is the frequency of the sound wave produced by the tuning fork?}

\begin{soln} loud sound is heard when stationary wave is set up

node is formed at surface of water

node-to-node distance is $60.0 - 26.8 = 33.2 \text{ cm}$

wavelength of this wave: $\lambda = 2 \times 33.2 \text{ cm} = 0.664 \text{ m}$

frequency: $f = \frac{v}{\lambda} = \frac{340}{0.664} \RA f \approx 512 \text{ Hz}$ 
\end{soln}

\subsection*{Stationary wave between two open ends}

This is left as an exercise for the reader to prove:

You should verify that the allowed wavelengths are: $\lambda = 2L, \, L, \,\frac{2}{3}L, \cdots, \frac{2L}{n}, \cdots$

Frequencies of the fundamental and excited modes go as $f = f_1, \, 2f_1, \, 3f_1, \cdots, nf_1 \cdots$


\subsection{Measurement of speed of sound with stationary waves}
The required practical version of this is covered in Example \ref{recpracwave}, however there is a version that often appears in exam questions as follows;

\begin{figure}[ht]
	\centering
	\begin{circuitikz}
		% loudspeaker
		\draw (-0.6,0.55) to[loudspeaker, name=L] (-0.6,-0.55);
%		\node [waves, scale=0.7, right] at(L.north) {};
		\node[above] at (-0.6,0.65) {loudspeaker};
		% microphone
		\draw[thick] (5.22,0) rectangle (5.38,-0.8);
		\draw[thick, fill=white] (5.3,0) circle (0.2);
		\draw[thick] (5.3, -0.8) [out=-90, in=60] to (3.5,-2); 
		\node[above] at (5.3, 0.3) {microphone};
		\draw[thick, ->] (5.6,0) --++ (0.4,0);
		\draw[thick, ->] (5.0,0) --++ (-0.4,0);
		% oscilloscope
		\draw[thick, rounded corners] (1.9,-2) rectangle (5,-4);
		\draw[rounded corners] (2, -2.1) rectangle (4, -3.9);
		\foreach \y in {-2.4, -3, -3.6} {
			\draw[thick] (4.5, \y) circle (0.2);
		}
		\foreach \x in {2.2, 2.4, ..., 3.8} \draw[very thin, gray!40] (\x, -2.1) -- (\x, -3.9);
		\foreach \y in {-2.3, -2.5, ..., -3.7} \draw[very thin, gray!40] (2, \y) -- (4, \y);
		\draw [thick,Green, domain=2:4,samples=80,smooth] plot (\x,{0.65*sin((\x-2.7)*540) - 3});
		\node[below] at (3.5, -4.1) {oscilloscope};
		% metal plate
		\draw[thick, fill=gray!20] (9,0) ++ (-0.4,1) --++ (0.8,-0.5) --++ (0,-2) --++ (-0.8,0.5) --++ (0,2);
		\node[above] at (9, 1) {metal plate};
		% wave pattern
		\draw [gray, dashed, domain=0:9,samples=80,smooth] plot (\x,{0.3*sin(\x*80)});
		\draw [gray, dashed, domain=0:9,samples=80,smooth] plot (\x,{-0.3*sin(\x*80)});
	\end{circuitikz}
	\caption{arrangement of apparatuses for setting up a stationary sound wave}
\end{figure}

\cmt To set up a stationary sound wave, we need two travelling waves in opposite directions, this sound wave can be produced by a \emph{loudspeaker} (coupled to a \emph{signal generator}). The forward-travelling wave is reflected at a \emph{metal plate} and the reflected wave superposes with wave from loudspeaker to give a stationary wave.
The wave intensity at any point can be detected by a \emph{microphone} and observed on an \emph{oscilloscope}. To verify whether stationary wave is formed, we check whether nodes are formed - the positions of the metal plate and the microphone are adjusted carefully and if there exists a location where no signal is detected, this means microphone is at a node and a stationary wave has been formed between loudspeaker and reflector.
To then find wavelength of sound wave, we slowly move microphone to the next node, the wavelength of sound wave: $\lambda = 2\times \text{distance between neighbouring nodes}$

To find frequency of sound wave, we move microphone to anywhere that is not a node, from time base of the oscilloscope, we can work out period $T$ of the wave. Frequency is then given by the formula: $f=\frac{1}{T}$.

Finally,  speed of the sound wave\footnote{We should bear in mind that stationary waves do not propagate through space, so rigorously speaking, it is incorrect to say the speed of a stationary way. The wave speed calculated in this way refers to the speed of either of the two travelling waves that give rise to the stationary wave through superposition.} can be calculated: $v=\lambda f$




\subsection{Stationary waves \& progressive waves}

There are quite a few differences between a stationary wave and a progressive wave:

A progressive wave can transfer energy from one place to another but for stationary waves, vibrational energy is not transferred due to existence of nodes. 

In a stationary wave, all part have the same amplitude. Different points in general have different amplitudes over a stationary wave. Nodes have zero amplitude, anti-nodes have greatest amplitudes. Other points have various amplitudes depending on their positions.

In a progressive wave the displacements of each point at a particular moment can be all different, but the largest displacements they can reach are the same (if no loss of energy). This means that the phase difference between any two points in a progressive wave can take any value and different points reach their peaks at different times for a progressive wave so $\Delta \phi$ can be anything from 0 to $2\pi$ (or anything from 0 to $360^\circ$).

Phase difference between any two points in a stationary wave however, is either 0 or $\pi$ (or $180^\circ$). For points between adjacent nodes, they all reach their highest at same time despite differences in amplitudes, these points are all in phase, so $\Delta \phi=0$.
For points separated by one node, they are anti-phase, because when one is at its peak, the other is at its trough, so $\Delta \phi = \pi$ (or $180^\circ$)


	
\section{End-of-chapter questions}


\subsection*{Principle of superposition}

\question{
	Noise reduction headphones can cancel out external sound waves by producing their own sound waves. What is the phase difference between the external sound wave and the wave produced in the headphone?
}

\question{
	A wave of frequency of 10 Hz travels at a speed of $4.0 \mps$. What is the phase difference between two points 0.50 m apart?
}

\question{
	Two loudspeakers $A$ and $B$ are connected to the same signal generator. They give out sound signals of wavelength 0.50 m. A person stands in a position which is at a distance of 20.0 m from $A$ and 21.0 m from $B$. Initially, only loudspeaker $A$ is emitting sound. What happens to the intensity of the sound heard by the person when loudspeaker $B$ is also switched on?
}

\question{
	A students thinks that the superposition of waves is about the superposition of the wave amplitudes. Give a counter-example why he is wrong.
}

\question{
	Two separate loudspeakers generate sound waves of slightly different frequencies. (a) Explain why a stable interference is not observed. (b) If an observer stands at a fixed location, state and explain the variation of the intensity of the sound that he hears.
}

\question{
	Two waves $A$ and $B$ are of the same type. Given that they have a phase difference of $120^\circ$, and the intensity of wave $A$ is twice that of wave $B$. On the same graph, sketch the displacement against time to illustrate the two waves.
}


\question{
	Two waves $A$ and $B$ of the same amplitude have a phase difference of $90^\circ$. (a) On the same graph, sketch the displacement against time to illustrate the two waves. (b) Sketch the variation of the resultant displacement if $A$ and $B$ superpose together.
}



\subsection*{Double-slit experiment}

\question{
	A laser emits a yellow light of 580 nm. It is sent through a set of double slits with a separation of 0.40 mm. An interference pattern is seen on a screen at a distance of 2.5 m from the slits. What is the distance between a bright fringe and its closest dark fringe?
}


\question{
	A double-slit interference pattern is produced on a screen using a red laser of 630 nm wavelength. The screen is 4.0 m from the double slit, and the fringes are seen to have a spacing of 15 mm. What is the separation of the slits?
}

\newpage


\question{
	Light is incident normally on two narrow slits which are of 2.00 mm apart. A screen is at a distance of 2.5 m from the slits. The distance between the central maximum and the 4$^\text{th}$ bright spot is found to be 3.2 mm. What is the wavelength of the light?
}

\question{
	In a double-slit experiment, suggest by what factor does the separation between bright fringes change if (a) the distance between the slits and the screen is doubled, (b) the slit separation is doubled, (c) the frequency of incident light is doubled?
}

\question{
	A light of wavelength $\lambda$ passes through two narrow slits $S_1$ and $S_2$ and forms a series of bright and dark fringes on a screen. The $n^\text{th}$ dark fringe from the central bright fringe is observed at point $X$ on the screen. Find an expression, in terms of $\lambda$ and $n$, for $\big|S_2X - S_1 X\big|$.
}

\question{
	A beam of white light enters a double-slit. Describe and explain the appearance of (a) the central bright fringe, (b) the first bright fringe from the central fringe.
}

\question{
	A beam of light passes through a double slit such that the emergent intensity from each slit is initially the same. Describe the change to the brightness of the fringes if (a) the intensity of light through both slits are increased, (b) the intensity of light through one of the slit is increased.
}




\subsection*{Diffraction gratings}

\question{
	A beam of light is incident normally on a diffraction grating that has 1000 lines per mm. The angle between the first order maxima is $60^\circ$. Find the wavelength of the light.
}

\question{
	Light of wavelength 590 nm is directed through a diffraction grating. The first order maximum is formed at an angle of $20.0^\circ$. (a) What is the slit separation for the diffraction grating? (b) What is the angular separation between the first and second order maxima?
}

\question{
	A beam of light of 480 nm wavelength is incident normally on a diffraction grating which has a slit separation of 3.0 $\mu$m. How many intensity maxima can be observed?
}

\question{
	A diffraction pattern is observed when monochromatic light falls on a diffraction grating. If another diffraction grating with twice as many lines per millimetre is used, suggest the possible effects to (a) the total number of diffraction maxima, (b) the angle between the first and second order.
}

\question{
	A beam of light contains two wavelengths 420 nm and 630 nm. A diffraction grating of $6.0 \times 10^5$ lines per metre produces a pattern such that the second order of one of these wavelengths overlaps with the third order of the other wavelength. What is the angle $\theta$ at which the overlap occurs?
}

\question{
	A diffraction grating is used with several wavelengths of light. The angle $\theta$ of the second order maxima is measured for each wavelength $\lambda$. If a graph of $\sin\theta$ against $\lambda$ is plotted, what does the gradient of the graph represent?
}

\question{
	A beam of monochromatic light is incident normally on a diffraction grating. If the grating is rotated about the axis parallel to the incident beam, describe the changes to the positions of diffraction maxima.
}

\question{
	Light of unknown wavelength $\lambda$ is emitted from a distant star. Suggest the apparatuses you need and what measurements you need take in order to determine $\lambda$. 
}


\subsection*{Stationary waves}

\question{
	How many nodes, including the end points, are there in a standing wave that is three wavelengths long? What about anti-nodes?
}

\question{
	A stretched rope of length 150 cm is fixed at the two ends. At what wavelengths can stationary waves be set up on this rope?
}

\question{
	(a) A pipe of length $l$ is closed at one end but open at the other. Notes are produced by the pipe if stationary waves are set up. In terms of $l$ and the speed of sound $v$ in the air column, what is the frequency of the lowest note produced in the pipe? (b) If the pipe is open at both ends, what is the lowest frequency?
}

\question{
	Suggest and explain how you can lower the pitch of a note on a guitar by altering the length of the string.
}

\question{
	A glass cylinder stands upright and is initially full of water. A loudspeaker emits a sound wave of frequency 850 Hz into the cylinder from above. The speed of sound in air is measured to be $340 \mps$. A tap at the bottom of the cylinder is opened so that the water level slowly decreases. The first loud sound is heard when the water level is at a height of 70.0 cm. When the wave level drops to a height of $y$ cm, a second loud sound is heard. (a) What is the wavelength of the sound wave? (b) Explain why loud sounds are heard for these specific water levels. (c) Find the value of $y$.
}


\question{
	A glass tube, closed at one end, has a layer of fine dust sprinkled inside on its lower side. A loudspeaker placed near the open end emits sound waves at a particular frequency of 1700 Hz such that a stationary wave is produced inside the tube. The dust is seen to form small heaps of equal spacing as shown below. (a) Explain why heaps of dust are formed. (b) Label the nodes of the stationary wave on the graph. (c) Find the speed of sound in the pipe.
}

\begin{figure}[ht]
	\centering
	\begin{circuitikz}
		% loudspeaker
		\draw (-0.8,0.55) to[loudspeaker, name=L] (-0.8,-0.55);
		% pipe
		\draw[thick] (0,-0.75) --++ (10,0) --++ (0,1.5) --++ (-10,0);
		\draw [gray!50, fill, domain=0:10,samples=80,smooth] plot (\x,{-0.1*cos(\x*180)-0.6}) --++ (0,-0.05) --++ (-10,0) -- cycle;
		% labels
		\draw (8,0.8) --++ (0.6,0.5) node[above]{tube};
		\draw (-0.75,0.7) --++ (0.2,0.5) node[above]{loudspeaker};
		\draw (3,-0.6) --++ (0.4,1.9) node[above]{dust heap};
		\draw[dashed] (3,-0.8) --++ (0,-0.8);
		\draw[dashed] (9,-0.8) --++ (0,-0.8);
		\draw[<->] (3, -1.4) -- (9,-1.4) node[midway, above]{30 cm};
	\end{circuitikz}
\end{figure}
\newpage

\question{
	A stationary wave is set up between the two ends $X$ and $Y$ of an elastic cord. The graph shows the wave pattern at a time $t=t_0$. Given that the frequency of the vibration is 50 Hz. (a) Sketch the wave pattern between $X$ and $Y$ at time $t_1 = t_0 + 10 \text{ ms}$. (b) Sketch the wave pattern at time $t_2 = t_0 + 25 \text{ ms}$. (c) State the phase difference between the displacement of point $A$ and that of point $B$.
}

\begin{figure}[ht]
	\centering
	\begin{tikzpicture}[scale=1.25]
	\draw [thick, domain=0:8,samples=80,smooth] plot (\x,{0.8*sin(\x*90)});
	\draw[fill] (0,0) circle (0.05) node[left]{$X$};
	\draw[fill] (8,0) circle (0.05) node[right]{$Y$};
	\draw (5.5, 0.75*0.8) --++ (0.5,0.5) node[right]{elastic cord};
	\draw[fill, red] (0.5, {0.8*sin(45)}) circle(0.05) node[above left, black]{$A$};
	\draw[fill, red] (3, -0.8) circle(0.05) node[below, black]{$B$};
	\end{tikzpicture}
\end{figure}


\chapter{Cosmology}

\section{Stephan's law}
An objects luminosity depends on two factors:
Its surface temperature
Its surface area
The relationship between these is known as Stefan's Law or the Stefan-Boltzmann Law, which states:
The total energy emitted by a black body per unit area per second is proportional to the fourth power of the absolute temperature of the body

So Stefan's Law shows that the luminosity of a star is directly proportional:
\begin{itemize}
    \item To its radius $L \propto r^2$
    \item To its surface area $L \propto 4 \pi r^2$
    \item To its surface absolute temperature $L \propto T^4$
\end{itemize}


Stefan's Law equation is given by:
\begin{equation*}
    L = 4 \pi r2\sigma T^{4}
\end{equation*}

Where:
L = luminosity of the star (W)\\
r = radius of the star\\
$\sigma$ = the Stefan-Boltzmann constant\\
T = surface temperature of the star (K)\\
The surface area of a star (or other spherical object) can be calculated using: $A = 4\pi r^2$\\
Where r = radius of the star

\example{The surface temperature of Proxima Centuri, the nearest star to Earth, is $3000$K and its luminosity is $6.506 \times 10^{23} $W.

Calculate the radius of Proxima Centuri in kilometres and show your working clearly.}
\begin{soln}
List the known quantities: 

Surface temperature, $T = 3000 K$\\
Luminosity, $L = 6.506 \times 10^{23} $W\\
Stefan's constant, $\sigma = 5.67 \times 10^{-8} W m^{-2} K^{-4}$\\
Write down Stefan's Law \& sub in values\\
The radius of Proximal Centuri is 106 200 km (4 s.f.)
\end{soln}
\section{Wein's law}

\subsection{Black Body Radiator}
An ideal black body radiator is one that absorbs and emits all wavelengths.\\
A black body is a theoretical object, however, stars are the best approximation there is.\\
The radiation emitted from a black body has a characteristic spectrum that is determined by the temperature alone.
% CUSTOM COLORS
% See https://tikz.net/blackbody_color/
\definecolor{1000K}{rgb}{1,0.0337,0}
\definecolor{2000K}{rgb}{1,0.2647,0.0033}
\definecolor{3000K}{rgb}{1,0.4870,0.1411}
\definecolor{4000K}{rgb}{1,0.6636,0.3583}
\definecolor{5000K}{rgb}{1,0.7992,0.6045}
\definecolor{6000K}{rgb}{1,0.9019,0.8473}
\definecolor{8000K}{rgb}{0.7874,0.8187,1}
\definecolor{10000K}{rgb}{0.6268,0.7039,1}
\pgfdeclareverticalshading{rainbow}{100bp}{
  color(0bp)=(red); color(25bp)=(red); color(35bp)=(yellow);
  color(45bp)=(green); color(55bp)=(cyan); color(65bp)=(blue);
  color(75bp)=(violet); color(100bp)=(violet)
}
\colorlet{myred}{red!70!black}
\colorlet{mygreen}{green!70!black}
\colorlet{mydarkgreen}{green!55!black}
\pgfmathdeclarefunction{planck}{2}{%
  \pgfmathparse{1.191042972e26/(#1^5)/(exp(0.01439e9/(#1*#2))-1)}%
}
\pgfmathdeclarefunction{rayleighjeans}{2}{%
  \pgfmathparse{8.278160269e18*#2/(#1^4)}%
}
\pgfmathdeclarefunction{wien}{2}{%
  \pgfmathparse{1.191042972e26/(#1^5)*exp(-0.01439e9/(#1*#2))}%
}
\pgfmathdeclarefunction{lampeak}{1}{% % Wien's displacement law
  \pgfmathparse{2.898e6/#1}%
}
\begin{figure*}

\begin{tikzpicture}[scale=1.5]
  \def\N{60}
  \def\xmax{3100}
  \def\ymax{1.36e10}
  \def\tick#1#2{\draw[thick] (#1+.01*\ymax) -- (#1-.01*\ymax) node[below=-.5pt,scale=0.75] {#2};}
  \begin{axis}[
      every axis plot/.style={
        mark=none,samples=\N,domain=5:\xmax,smooth},
      xmin=(-.05*\xmax), xmax=(1.05*\xmax),
      ymin=(-.08*\ymax), ymax=(1.08*\ymax),
      restrict y to domain=0:\ymax,
      axis lines=middle,
      axis line style=thick,
      %enlargelimits=upper, % extend the axes a bit to the right and top
      tick style={black,thick},
      ticklabel style={scale=0.8},
      %xtick style={draw=none},xticklabels=none,
      max space between ticks=26,
      xlabel={Wavelength $\lambda$ [nm]},
      ylabel={Power $P$ [kW/sr\,m$^2$\,nm]},
      xlabel style={at={(rel axis cs:0.5,0)},below=-1pt,font=\small},
      ylabel style={at={(rel axis cs:-0.11,0.5)},rotate=90},
      width=9cm, height=7cm,
      %clip=false
      tick scale binop=\times,
      every y tick scale label/.style={at={(rel axis cs:0,1)},anchor=south}]
    ]
    
    % RAINBOW
    \shade[shading=rainbow,shading angle=90,opacity=0.5] (380,0) rectangle (740,\ymax);
    \node[above=-1pt,scale=0.8] at (200,\ymax) {\strut UV}; % 10 - 400 nm
    \node[above=-1pt,scale=0.8] at (570,\ymax) {\strut optical}; % 380 - 740 nm
    \node[above=-1pt,scale=0.8] at (920,\ymax) {\strut IR}; % 740 - 1050 nm
    
    % PLANCK
    \addplot[very thick,red]    {planck(x,3000)};
    \addplot[very thick,orange] {planck(x,4000)};
    \addplot[very thick,samples=3*\N,blue] {planck(x,5000)};
    %\addplot[dashed,thick,red,domain=1000:4000]    {rayleighjeans(x,3000)};
    %\addplot[dashed,thick,orange,domain=1000:4000] {rayleighjeans(x,4000)};
    %\addplot[dashed,thick,blue,domain=1000:4000]   {rayleighjeans(x,5000)};
    %\addplot[dashed,thick,red,domain=1000:4000]    {wien(x,3000)};
    %\addplot[dashed,thick,orange,domain=1000:4000] {wien(x,4000)};
    %\addplot[dashed,thick,blue,domain=1000:4000]   {wien(x,5000)};
    
    %% MAXIMUM (Wien's displacement law)
    \addplot[mygreen,thin,variable=T,domain=2500:6000]({lampeak(T)},{planck(lampeak(T),T)});
    
    % LABELS
    \node[above=0pt,scale=0.75,red] at (1150,{planck(1150,3000)}) {3000K};
    \node[above right=-1pt,scale=0.75,orange!80!black] at (740,{planck(740,4000)}){4000K};
    \node[above right=-1pt,scale=0.75,blue] at (800,{planck(800,5000)}) {5000K};
    
    %% TICKS
    %\tick{500,0}{500}
    %\tick{1000,0}{1000}
    %\tick{1500,0}{1500}
    %\tick{2000,0}{2000}
    %\tick{2500,0}{2500}
    %\tick{3000,0}{3000}
  \end{axis}
\end{tikzpicture}
\end{figure*}

The intensity-wavelength graph shows how thermodynamic temperature links to the peak wavelength for three different bodies.

\subsection {Wien's Displacement Law}
Wien’s displacement law relates the observed wavelength of light from an object to its surface temperature, it states:
The black body radiation curve for different temperatures peaks at a wavelength that is inversely proportional to the temperature

This relation can be written as:

\begin{equation*}
    \lambda_{\mathrm{max}}=\frac{W}{T}
\end{equation*}

Where:\\
$\lambda_{\mathrm{max}}$ = the maximum wavelength emitted by an object at the peak intensity (m) \\
T = the surface temperature of an object (K) \\
Wien's constant is $W=2.9\times 10^{-3} mK $

This equation shows that the higher the temperature of a body, the shorter the wavelength at the peak intensity. Hotter objects tend to be white or blue, and cooler objects tend to be red or yellow.//


\example{
The spectrum of the star Rigel in the constellation of Orion peaks at a wavelength of 263 nm, while the spectrum of the star Betelgeuse peaks at a wavelength of 828 nm.
Determine which of these two stars, Betelgeuse or Rigel, is cooler.}
\begin{soln}
List the known quantities
Maximum emission wavelength of Rigel = 263 nm = 263 × 10-9 m
Maximum emission wavelength of Betelgeuse = 828 × 10-9 m

$$T = \frac{2.9e-3}{\lambda_{max}}$$
Calculate the surface temperature of each star:
$$\text{Rigel: }T=\frac{2.9e{-3}}{263e{-9}}=11026 = 11 000K$$

$$\text{Betelgeuse: }T=\frac{2.9e{-3}}{828e{-9}}=3502 = 3500K$$

Step 4: Write a concluding sentence

Betelgeuse has a surface temperature of 3500 K, therefore, it is much cooler than Rigel
\end{soln}

\chapter{Electrical Quantities \& Components}

\section{Electrical quantities}

\subsection{Electric charges}\index{charge}

Electric charge is the property of matter that causes it to experience a force in an electric field.  Electric charges come in two types, \emph{positive} charges (+ve) and \emph{negative} charges (-ve) \footnote{Important to realise that the +ve and -ve signs are to distinguish each type. A -ve charge is not an deficit of charge in the way that a -ve bank account is a debt.}. You should be familiar with the fact that like charges repel each other, and opposite charges attract. The interaction between electric charges plays a central role in all electrostatic, chemical and electrical phenomena.

Electric charges are measured in coulombs: $[Q] = \text{C}$\\

All matter is made of atoms, which consist of a nucleus (+ve) and electrons (-ve) and each electron has a charge of $-e$, where $e=1.60\times10^{-19}\text{ C}$ is the \keypoint{elementary charge}.

Typically a body becomes electrically-charged by losing or gaining electrons.

\begin{itemize}
	\item[--] an object that gains electrons from elsewhere becomes negatively-charged
	
	\item[--] an object that loses electrons becomes positively-charged
\end{itemize}


The charge of any object must be an integer multiple of the elementary charge: $\tcbhighmath{Q=Ne}$ This is the same quantity as the charge on an electron or proton.

We say the charge of an object is \keypoint{quantised}. In the universe, there are a number of conserved quantities - This law of conservation, known as the \keypoint{conservation of electric charges} appears to be a very significant one. An unbalanced charge cannot be created nor be destroyed - ie you can cause a positron and electron to exist, as the total created charge is zero.


\example{A metal sphere carries a net charge of $+2.4 \times 10^{-9} \text{ C}$. Suggest whether the sphere has gained additional electrons or lost electrons, and also calculate how many electrons has the sphere gained or lost?}

\begin{soln} The sphere is positively-charged, so it has lost some of its electrons.

The number of lost electrons: $ N = \frac{Q}{e} = \frac{2.4\times10^{-9}}{1.60 \times10^{-19}} \RA N = 1.5 \times 10^{10}  $ \end{soln}



\subsection{Electric current}

A separation of charges gives rise to a force. Forces cause accelerations so it's inevitable that some carges move. The flow of electric charges gives rise to electric currents.

\begin{ilight}
	\centering \keypoint{Electric current} $I$ is defined as the amount of charge flow $Q$ per unit time: \begin{empheq}[box=\tcbhighmath]{equation*}{I=\frac{\Delta Q}{\Delta t}}\end{empheq}\index{current}
\end{ilight}


\cmt The unit of electric current: $[I] = \text{A}$ (ampere)

The Ampere is one of the S.I. base units\footnote{As of the 2019 redefinition of the SI base units, the ampere is defined by fixing the elementary charge e to be exactly $1.602176634 \times 10^{-19}$ C, which means an ampere is an electric current equivalent to $10^{19}$ elementary charges moving every $1.602176634$ seconds. Previously it was defined as the following: when two parallel infinitely long straight wires separated by a distance of 1 m carry two equal currents such that the force per metre between them is $2.0\times10^{-7}$N, then this current is of 1 A. - It was oddly circular that the notion of electric current was defined based on electric charge, and the unit of charge was defined based on the unit of current: 1 C = 1 A s.}

Though it causes some frustration - the direction of a current is defined as the direction of flow of positive charges. Of course, the current flowing in a conductor is usually due to motion of negatively-charged electrons so \keypoint{conventional current} is in opposite direction to the electron flow. 

In many cases we might use this equation to find charge flow from a current, we can use  $Q=It$ only if there is a constant current. For capacitors, AC circuits and in other non-constant circumstances we can use the area under an $I$-$t$ graph to give the total charge transferred.
It can sometimes be useful to find the average current and then proceed with $Q=\bar{I}t$. 

Note that current in a circuit can be measured with an \keypoint{ammeter}:
\begin{marginfigure}
\centering
\begin{tikzpicture}[]
    \draw (0,0) to[rmeter,t=A] (2,0);
  \end{tikzpicture}
  \caption{The symbol for an Ammeter}
\end{marginfigure}

\example{A current of 25 mA flows in a wire. How many electrons have passed a point in the wire in one hour?}

\begin{soln}
    
Charge flow: $Q = I t = 25\times10^{-3} \times 3600 \RA Q = 90 \text{ C}$

Number of electrons: $N = \frac{Q}{e} = \frac{90}{1.60 \times10^{-19}} \RA N \approx 5.6 \times 10^{20}$ \end{soln}

\begin{marginfigure}
	\vspace{-25pt}
	\begin{center}
		\begin{tikzpicture}[yscale=0.6]
		\draw[thick,blue] (0,1) -- (3,4);
		\draw[<->] (0,5) node[left]{$I$/A} -- (0,0) node[left]{0} -- (3.75,0) node[below]{$t$/s};
		\node[left] at (0,1) {2.0};
		\node[left] at (0,4) {8.0};
		\draw[dashed] (0,4) -- (3,4) -- (3,0) node[below] {3.0};
		\end{tikzpicture}
	\end{center}
	\vspace{-5pt}
\end{marginfigure}

\example{The current in a lamp is increased uniformly from 2.0 A to 8.0 A over 3.0 s. What is the charge flow?}

\begin{soln} Average current during this time is: $\bar{I} = \frac{2.0+8.0}{2} = 5.0 \text{ A} $

charge flow is then: $ Q = \bar{I} t = 5.0 \times 3.0 = 15 \text{ C} $

we can also use area under $I$-$t$ graph
\begin{equation*}
	Q = \frac{1}{2} \times (2.0 + 8.0) \times 3.0 = 15 \text{ C} 
\end{equation*}
\end{soln}

\subsection*{Microscopic view of electric currents}

The flow of electric current is essentially due to motion of charge carriers due to the electrostatic attraction and repulsion.\footnote{In most metallic conductors, charge carriers are those \emph{free electrons}. Charge carriers can also be thought about as positively-charged \emph{holes} (the absence of electrons) in some semi-conductors or \emph{ions} in chemical solutions.}
From an intuitive perspective, current would depend on number of charge carriers, and how fast they move. We'll see that this isn't far from correct. Let's take  a wire of cross section $A$, and see what follows in thinking about current.



\begin{figure}[ht]
	\centering
	\begin{tikzpicture}[scale=1.5]
	\draw[fill=orange!40, opacity=0.4] (5,-1) arc(-90:90:0.5 and 1) -- (2,1) arc(90:270:0.5 and 1) -- (5,-1);
	\draw[thick] (0,0) ellipse (0.5 and 1) node {$A$};
	\draw[dashed] (2,0) ellipse (0.5 and 1);
	\draw[thick] (0,-1) -- (5,-1) arc(-90:90:0.5 and 1) -- (0,1);
	\foreach \x/\y in {0.6/-.3, 0.9/0.5, 0.8/-0.7, 1.8/0.1, 2.6/0.6, 2.9/-0.5, 3.6/0.4, 4.1/-0.7, 4.4/0.1, 4.3/0.6}{
	\draw[->] (\x, \y) --++ (0.4,0);% node[above]{$v$};
	\shade[ball color = green] (\x, \y) circle (0.08);
	}
	\draw[thick, ->] (1.5,1.6) -- (3.5, 1.6) node[midway, above] {$I$};
	\draw[<->] (2,-1.6) -- (5,-1.6) node[above, midway]{$\Delta L = v \Delta t$};
	\draw (4.9, 0.75) --++ (0.7,0.6) node[above right]{$\Delta V$};
	\end{tikzpicture}
\end{figure}



We know current is transfer of charge so let's define $n$ as the density of charge carriers for the material - in other words the number of charge carriers per unit volume: $n = \frac{\text{Number of free charges}}{\text{Volume}}$ 

The total number $N$ of charge carriers in a volume $V$ can then be given by: $N = nV$

If each carrier has a charge of $q$ and moves with a speed $v$, then they contribute to a current:
\begin{equation*}
	I = \frac{\Delta Q}{\Delta t} = \frac{\Delta N q}{\Delta t} = \frac{n \Delta V q}{\Delta t} = \frac{n A \Delta L q}{\Delta t} = \frac{n A v \Delta t q}{\Delta t} \RA \end{equation*} \begin{empheq}[box=\tcbhighmath]{equation*}{I = nAvq}
\end{empheq}

In metallic conductors, free electrons act as charge carriers, the equation becomes: $I = nAve$ where $e$ is the charge on an electron.

As with many statistical physics phenomena, the charge carriers have a distribution of speeds as they move 
so it's more proper to say that $v$ is actually an average velocity, called the \keypoint{mean drift velocity}

The charge density $n$ is a property of the material in question - ie it's value depends on type of material, good conductors have large $n$, while semiconductors have smaller $n$. \footnote{We're past the point of considering the world in binary terms of \emph{conductor} and \emph{insulator}, though we still use those terms as shorthand for the extremes of range.}


\example{A length of silver wire of a diameter of $1.2 \text{ mm}$ carries a current of 2.0 A. The electron number density in silver is $5.9\times10^{28} \text{ m}^{-3}$. What is the mean drift velocity of the electrons?}

\begin{soln}\begin{equation*}
v = \frac{I}{nAe} = \frac{I}{n\cdot\frac{1}{4} \pi d^2 \cdot e} \RA \end{equation*}
\begin{equation*}\frac{2.0}{5.9\times10^{28} \times \frac{1}{4}\pi \times (1.2\times10^{-3})^2\times1.60\times10^{-19}} \approx 1.9\times10^{-4} \mps
\end{equation*}

\eqyskip This calculation shows that each free electron actually travel at very low speeds in a wire \end{soln}

\example{A metal conductor and a piece of semi-conductor of the same cross section are connected in series. When an electric current is driven through them, compare the mean drift speed of the charge carriers in the two materials.}

\begin{soln}
    
The same current and same cross section, so mean drift speed $ v = \frac{I}{nAe} \propto \frac{1}{n}$

The metal has larger $n$, so free electrons move at relatively lower speeds

The semi-conductor has smaller $n$, so charge carriers move at relatively higher speeds \end{soln}




\subsection{P.d. \& e.m.f.}\label{ch:potential-difference}

Potential difference and e.m.f.\footnote{e.m.f. stands for \emph{electromotive force}, which is a miserably misleading term due to historical reasons. It has nothing to do with a force, or motivation for that matter. It is better to simply use the abbreviation e.m.f.} are both defined as the energy transfer per unit charge, both terms are often lumped together as \emph{voltages} in less formal contexts.

\begin{ilight}
	The \keypoint{P.d. (potential difference)} between two points is the energy converted from 
electrical potential energy to some other form per coulomb of charge flowing from one point to the other.: \begin{empheq}[box=\tcbhighmath]{equation*}{V=\frac{W}{Q}}\end{empheq} \index{potential difference}
\end{ilight}

\begin{ilight}
	\keypoint{The e.m.f. (electromotive force)} of a power supply is the energy converted from some other form (e.g. chemical) to electrical potential energy per coulomb of charge flowing through the source: \begin{empheq}[box=\tcbhighmath]{equation*}{\mathcal{E}=\frac{W}{Q}} \end{empheq}\index{e.m.f.}
\end{ilight}


The unit of measurement for p.d./e.m.f.: $[V] = [\mathcal{E}] = \text{V}$ (volt), where $1 \text{ V} = 1 \text{ J C}^{-1}$

The p.d. can tell you the direction of current flow through a component, current always flows from a higher potential to a lower potential.

We measure the p.d. across a component can be measured with a \keypoint{voltmeter}. 
\begin{marginfigure}
\centering
\begin{tikzpicture}[]
    \draw (0,0) to[rmeter,t=V] (2,0);
  \end{tikzpicture}
  \caption{The symbol for a Voltmeter}
\end{marginfigure}


\example{How much electrical energy is transformed into thermal energy when a charge of 10 C flows through a heater whose p.d. across is 25 V?}

\begin{soln}\begin{equation*}
	V = \frac{W}{Q} \RA W = V Q = 25 \times 10 = 250 \text{ J} 
\end{equation*}
\end{soln}
\example{A fully-charged 12 V car battery can supply a total electrical energy of 0.90 MJ. The starter motor of the car requires an average current of 150 A for a period of 2.0s. The battery is not able to recharge due to a fault. How many times can the starter motor be used?}

\begin{soln} total charge that can be supplied: $Q_\text{total} = \frac{W}{\mathcal{E}} = \frac{0.90 \times 10^6}{12} = 75000 \text{ C}$

charge needed for each start: $Q = I t = 150 \times 2.0 = 300 \text{ C}$

number of starts possible: $N = \frac{Q_\text{total}}{Q} = \frac{75000}{300} = 250$ \end{soln}

\example{Electrons beams are accelerated between a heated cathode and an anode with a potential difference of 40 V. What is the speed of the electrons when they reach the anode?}

\begin{soln} electrical potential energy is transformed into kinetic energy of electrons:
\begin{equation*}
qV = \frac{1}{2}mv^2 \RA 
\end{equation*}
\begin{equation*}
v=\sqrt{\frac{2qV}{m}} = \sqrt{\frac{2 \times 1.60 \times 10^{-19} \times 40}{9.11 \times 10^{-31}}} \approx 3.75 \times 10^6 \mps 
\end{equation*}
\end{soln}



\subsection{Resistance}

\begin{ilight}
	The \keypoint{Resistance}\index{resistance} $R$ of an electrical component is defined as the potential difference across divided by the current flowing through it: \begin{empheq}[box=\tcbhighmath]{equation*}{R=\frac{V}{I}} \end{empheq}
\end{ilight}


The unit of resistance: $ [R] = \text{V A}^{-1} = \Omega $ (ohm) named after Georg Ohm\footnote{Ohm's law was probably the most important of the early quantitative descriptions of the physics of electricity. We consider it almost obvious today. When Ohm first published his work, this was not the case; critics reacted to his treatment of the subject with hostility. They called his work a "web of naked fancies" and the Minister of Education proclaimed that "a professor who preached such heresies was unworthy to teach science." The prevailing scientific philosophy in Germany at the time asserted that experiments need not be performed to develop an understanding of nature because nature is so well ordered, and that scientific truths may be deduced through reasoning alone...}.

When a p.d. of 1 V is applied and current flow is 1 A, then resistance is 1 $\Omega$ 

Note that:
\begin{itemize}
    \item Resistance of a component does not depend on p.d. applied or current through it.
 \item Resistance $R$ of a component mainly depends on the following factors:

\begin{compactitem}
	\item[--] $R$ is proportional to length $L$ of the component: $R \propto L$.
	
	\item[--] $R$ is inversely proportional to cross-sectional area $A$ : $R \propto \frac{1}{A}$.
	
	\item[--] $R$ depends on the material of the component.
\end{compactitem}
\end{itemize}
this can be written as: \begin{empheq}[box=\tcbhighmath]{equation*}{R=\rho\frac{L}{A}} \end{empheq}

The constant $\rho$ is called the \keypoint{resistivity}\index{resistivity} of the material, a typical value of resistivity for a conducting metal is approx. $\rho \sim 10^{-8} \text{ } \Omega \text{ m}$.


\example{A current of 0.10 A is driven through a heater when connected to a supply voltage of 220 V. (a) Find the resistance of the heater. (b) If the heater is made from a wire of resistivity $1.8\times10^{-5} \text{ }\Omega\text{ m}$ and a diameter of 0.24 mm, find the length of the wire.}

\begin{soln}\begin{equation*}
	R = \frac{V}{I} = \frac{220}{0.10} \RA R = 2200 \text{ }\Omega
\end{equation*}
\vspace*{-1.2em}\begin{equation*}
R = \frac{\rho L}{A} \RA L = \frac{RA}{\rho} = \frac{2200 \times \pi \times (0.12\times10^{-3})^2}{1.8 \times 10^{-5}} \RA L \approx 5.5 \text{ m} 
\end{equation*}
\end{soln}
\example{A uniform wire of resistance $R$ has a length of $L$ and a cross-sectional area of $A$. If the wire is stretched to twice the length while its volume remains constant, what is the new resistance $R'$ in terms of $R$?}

\begin{soln}volume: $V=LA = \text{constant}$, so $A \propto \frac{1}{L}$, hence doubling $L$ means $A$ would become halved

\vspace*{0.3em} resistance: $R = \frac{\rho L}{A} \propto \frac{L}{A} \RA \frac{R'}{R}  = \frac{L'}{L} \times \frac{A}{A'} = 2 \times 2 \RA R' = 4R $ \end{soln}

\subsection{Electrical power}\index{power}

Current flowing through any electrical component causes energy to be transferred from source to the compnent. The rate at which electrical energy is transferred - the electrical power - is computed as:
\begin{equation*}
P = \frac{\Delta W}{\Delta  t} \xlongequal{V=\frac{W}{Q}} \frac{\Delta Q \cdot V}{\Delta t} \xLongrightarrow{I = \frac{\Delta Q}{\Delta t}} \end{equation*}
\begin{empheq}[box=\tcbhighmath]{equation*}{P=IV}
\end{empheq}

Substituting $V=IR$, the formula for electrical power also takes two useful forms:

\begin{empheq}[box=\tcbhighmath]{equation*}{P=I^2 R} \quad \text{or} \quad {P=\frac{V^2}{R}}
\end{empheq}

From these we can see why the national grid, and it's step-up and step-down transformers is such a big part of GCSE - You can see that a very small change in the current flowing can have a large impact on the power dissipated in the overhead wires. 

The unit of electrical power allows us to roll out a classic physics joke  \footnote{
\begin{tikzpicture}
\calloutquote[author=Teacher,width=3cm,position={(-1,-0.2)},fill=red!30,rounded corners]{What is the unit of power?}
\end{tikzpicture}
\begin{tikzpicture}
\calloutquote[author=Student,width=3cm,position={(1,-0.2)},fill=blue!30,rounded corners]{Watt!}
\end{tikzpicture}
\begin{tikzpicture}
\calloutquote[author=Teacher,width=3cm,position={(-1,-0.2)},fill=red!30,rounded corners]{Yes, What is the unit of power?}
\end{tikzpicture}} as it's the $Watt$. If a p.d. of 1 V drives a current of 1 A, power produced is 1 W, i.e., $1 \text{ W} = 1 \text{ A} \cdot 1 \text{ V}$



%\example{An electric toaster has a power rating of 900 W and current of 4.0 A. What is the resistance in the toaster?}
%
%\solc\begin{equation*}
%	P = I^2 R \RA R = \frac{P}{I^2} = \frac{900}{4.0^2} \RA R \approx 56 \text{ }\Omega 
%\end{equation*}



\example{An electric toaster is labelled `220 V 900 W'. When it is operating normally, find (a) the current through it, (b) the resistance in the toaster.}

\begin{soln}current: $I = \frac{P}{V} = \frac{900}{220} \RA I \approx 4.09 \text{ A}$

\vspace*{0.3em} resistance: $R = \frac{V}{I} = \frac{220}{4.09}$, $\,\,$ or alternatively, $R = \frac{V^2}{P} = \frac{220^2}{900} \RA R \approx 53.8 \text{ }\Omega$ \end{soln}

\example{Two conducting cylinders $X$ and $Y$ are of the same diameter and the same material. The length of the cylinder $X$ is twice of that of the cylinder $Y$. If the two cylinders are connected in series to a power supply so that the current through them are equal, find the ratio of the electrical powers dissipated in the two cylinders $\frac{P_X}{P_Y}$.}

\begin{soln} for either cylinder, electrical power: $P = I^2 R = I^2 \cdot \frac{\rho L}{A} = I^2 \cdot \frac{4 \rho L}{\pi d^2}$

same current $I$, same diameter $d$, same resistivity $\rho$, so power: $P \propto L \RA	\frac{P_X}{P_Y} = \frac{L_X}{L_Y} = 2 $ \end{soln}


\section{Electrical components \& $I$-$V$ characteristics}

Electrical circuits are all about the interplay of Voltage and Current. As such the key to understanding how different components behave differently when a p.d. is applied is through an $I$-$V$ characteristics graph.


\subsection{Ohmic conductors}

\begin{marginfigure}
	\centering
	\begin{tikzpicture}[scale=0.5]
	\draw[->] (0,-3) -- (0,3) node[left]{$I$};
	\draw[->] (-4,0) -- (4,0) node[below]{$V$};
	\draw[thick, blue] (-3.6,-2.7) -- (3.6, 2.7);
	\end{tikzpicture}
	
\end{marginfigure}

The resistance of an ohmic conductor is constant for all currents, so current is directly proportional to p.d.\index{Ohm's law}. Metal wires are usually considered as perfect ohmic conductors as the resistance of metal wires is nearly constant at fixed temperature. In a V-I graph, the resistance of an ohmic conductor equals the inverse of the gradient.



\subsection{Filament lamps}

\begin{marginfigure}
	
	\centering
	\begin{tikzpicture}[scale=0.5]
	\draw[->] (0,-3) -- (0,3) node[left]{$I$};
	\draw[->] (-4,0) -- (4,0) node[below]{$V$};
	\draw[thick, blue] (0,0) -- (1, 1.2) [out=50.2, in=190] to (3.6,2.7);
	\draw[thick, blue] (0,0) -- (-1, -1.2) [out=230.2, in=10] to (-3.6,-2.7);
	\end{tikzpicture}
	
\end{marginfigure}

Lamp filaments are usually made from tungsten because it has a very even emission across the visible range when how. It's also capable of sustained high temperatures (a fillament is essentially a demonstration of black-body radiation, which is why some bulbd tend to look yellow - they're not hot enough to be white). For small currents, the filament doesn't get hot enough to emitt light and behaves like an ohmic component. Larger currents cause the filaments to be heated to high temperature and the resistance to increase.\footnote{This is because the vibration of the metal lattice increases as the temperature goes up, free electrons are more likely to be scattered, so the metal becomes less conductive and the current does not increase as much for the same p.d. increase.}
To find resistance of lamp filament, we read off values of $I$ and $V$, then the ratio $\frac{V}{I}$  gives us the resistance.

\subsection{Thermistors}

\begin{marginfigure}
		\centering
	\begin{circuitikz}[european resistors]
		\draw[thick] (0,0) to [thR,l_=$R_T$] (3,0);
	\end{circuitikz}
\end{marginfigure}

A \keypoint{Thermistor}\index{thermistor} has a resistance that varies with temperature in a way you might not expect. As temperature rises, the resistance of the thermistor becomes lower\footnote{In A-Levels, we consider \emph{NTC (negative temperature coefficient)} thermistors only. This means their resistance decreases as temperature rises. \emph{PTC (positive temperature coefficient)} thermistors also exist but let's pretend they don't.}

The behaviour of thermistors can be explained in terms of band theory of solids. The material is behaving just like a wire, but there's an additional dominant effect: To put it in simple words, we can say that electrons are bound to atoms and cannot move freely at low temperatures, but they can gain energy and break free from atoms if temperature goes up. There are more free electrons available to conduct electricity, so resistance decreases. At greater currents, current increases more for same p.d. increase. 

\begin{figure*}[ht]
	\centering
	\begin{minipage}{0.55\textwidth}
		\centering
		\begin{tikzpicture}[scale=0.96]
		\draw[->] (-1.6,0) -- (5,0) node[below] {$T$/\OC};
		\draw[->] (0,0) node[below]{0} -- (0,4) node[right] {$R_T$};
		\draw[blue,thick] (-1.5,4) [out=-60,in=170] to (4.8,0.8);
		\end{tikzpicture}
	\end{minipage}\hfil
	\begin{minipage}{0.35\textwidth}
		\centering
		\begin{tikzpicture}[scale=0.6]
		\draw[->] (0,-3.2) -- (0,3.2) node[left]{$I$};
		\draw[->] (-3.5,0) -- (3.5,0) node[below]{$V$};
		\draw[thick, blue] (0,0) -- (1.2, 0.75) [out=32, in=250] to (2.8,3);
		\draw[thick, blue] (0,0) -- (-1.2, -0.75) [out=212, in=70] to (-2.8,-3);
		\end{tikzpicture}
	\end{minipage}
\end{figure*}

To find resistance of thermistor, we also compute the ratio $\frac{V}{I}$

\subsection{Diodes}

\begin{marginfigure}
	\vspace*{-5pt}
	\centering
	\begin{circuitikz}[european resistors]
		\draw[thick] (0,0) to [stroke diode,l_=$D$] (3,0);
	\end{circuitikz}
\vspace*{-16pt}
\end{marginfigure}

A \keypoint{diode} only allows current to pass in one direction\index{diode}. \cmt When an \emph{ideal} diode is forward-biased, $R_D \to 0$.

\cmt When an \emph{ideal} diode is reversed-biased, $R_D \to \infty$.

\begin{figure}[ht]
	\centering
	\begin{tikzpicture}[xscale=1.4, yscale=2]
	\draw[->] (0,-1.5) -- (0,3) node[left]{$I$};
	\draw[->] (-3.5,0) -- (5.2,0) node[below]{$V$};
	\draw[thick, blue] (-3.2,0) -- (1,0) node[below, black]{$V_\text{th}$} [out=0, in=268] to (2.1,3);
	\draw[<-] (2.15, 2) --++ (0.6, 0.2) node[right, note]{once $V>V_\text{th}$, diode has\\nearly zero resistance};
	\draw[<-] (0.4,-0.1) --++ (0.2,-0.6) node[below right, note]{if $V<V_\text{th}$, diode does not conduct};
	\draw[<-] (-1.5,0.1) --++ (-0.3,0.7) node[above, note]{diode has very high resistance\\when p.d. is reversed-biased};
	\end{tikzpicture}
\end{figure}

A practical semiconductor diode become conductive once p.d. reaches a set value $V_\text{th}$ called the diode's \emph{threshold voltage} (about $0.6\sim0.7$ V for a silicon diode). When a p.d. applied is forward-biased and greater than $V_\text{th}$, the diode has very low resistance, so there is a sharp increase in current once $V>V_\text{th}$

Though diodes have very high resistance if reversed-biased, if the reverse voltage reaches the \emph{breakdown voltage}, diode would become conductive\footnote{Typical breakdown voltage for a diode range from a few volts to several hundred volts, depending on its desired application.}





	
\subsection{End-of-chapter questions}

\question{
	A technician reports that the charge of an ion in a solution is measured to be $-5.5 \times 10^{-19} \text{ C}$. State and explain whether this measurement is reliable.
}


\question{
	Calculate the current that gives a charge flow of 300 C in one minute.
}

\question{
	The manufacturer of a car battery claims that the battery can supply a steady current of 40 A for one hour when fully-charged. (a) What is the charge flow through a point in the circuit in this time? (b) If a car requires a current of 90 A to start, for how long could the battery be used before recharging?
}

\question{
	A beam of $2.0 \times 10^{10}$ electrons travel through a distance of 0.50 m and are stopped on an electrode in a time period of $5.0 \times 10^{-8} \text{ s}$. (a) Find the charge collected on the electrode. (b) Find the average current produced by the electron beam.
}

\question{
	State the SI base units of number density $n$ and electric charge $q$, and hence show that the equation $ I = nAqv$ has consistent units on both sides.
}

\question{
	The number density of free electrons in copper is about $8.0\times10^{28} \text{ m}^{-3}$. If the current in a copper wire of diameter 0.50 mm is 0.80 A, calculate the average drift speed of electrons in the wire.
}

\question{
	Two coppers wires $A$ and $B$ are joined together and carry a current. If wire $A$ has diameter $d$ and wire $B$ has diameter $2d$, calculate the ratio of the average drift velocity of free electrons in the wires $\frac{v_A}{v_B}$
}

\question{
	From the calculations, we find that the drift speed of electrons in metal wires is around a few mm per second, so it would take a very long time for the free electrons to travel a distance of a few metres in an electric circuit. But when we switch on an electric lamp, it lights up immediately. Give a reason for this observation.
}


\question{
	When a p.d. of 15 V is applied across a wire of resistance 4000 $\Omega$, how many electrons would have passed a point in the wire in one minute?
}

\question{
	A motor has a resistance of 40 $\Omega$ and carries a current of 2.0 A. (a) What is the p.d. across? (b) At what rate is electrical energy being dissipated in the motor?
}


\question{
	A cable of 50 m is made of 19 parallel strands of copper wires. Each strand has a diameter of 0.25 mm. Given that the resistivity of copper is $1.7\times10^{-8} \text{ }\Omega \text{ m}$, find the resistance of the cable.
}


\question{
	A pencil is used to draw a line of length 20 cm, width 1.5 mm, and thickness 0.10 $\mu$m. The resistance of the line is 30 k$\Omega$. What is the resistivity of the material in the pencil?
}

\question{
	One resistor $P$ has a resistance of 16 $\Omega$. A second resistor $Q$ is made of the same material, but has twice the length but half the diameter. What is the resistance of $Q$? 
}

\question{
	A coil is made by wounding a long insulated wire of diameter $d$ onto a cylindrical core of diameter $D$ where $D \gg d$. If the coil has $N$ turns, and the resistivity of the wire is $\rho$, show that the resistance in the coil is $\frac{4N\rho D}{d^2}$.
}


\newpage



\question{
	If we want to determine the resistivity of the material of a length of metal wire, what measurements do we need to take?
}

\question{
	What is the current in a 40 W light bulb when a supply of 200 V is applied?
}

\question{
	Fuses of current ratings 1 A, 5 A, 10 A, 20 A are available. Suggest and explain which fuse should be used for a 120 V, 900 W hairdryer.
}

\question{
	A 12 V battery supplies a continuous current of 5.0 A for 3.0 hours. (a) Find the charge supplied by the battery during this time. (b) Find the energy produced from the battery.
}

\question{
	An electric cooker is labelled `240 V, 900 W'. (a) What is the resistance of the cooker's filament at its operating temperature? (b) State and explain whether there is any change in the resistance if it is measured at room temperature.
}

\begin{marginfigure}
	\vspace*{-18pt}
	\centering
	\begin{circuitikz}[european resistors]
		\draw (2,1.2) to[R, l_=$R$] (0,1.2) to[battery, l_=$\mathcal{E}$] (0,-1.2) -- (2,-1.2);
		\draw (2,1.2) to[stroke diode] (2,-1.2); 
		\draw (2,1.2) -- (4,1.2) to[rmeter, t=V] (4,-1.2) -- (2,-1.2);
	\end{circuitikz}
	\vspace*{-16pt}
\end{marginfigure}

\question{
	The diagram shows an electric circuit incorporating a diode, connected in series with a 5.0 $\Omega$ resistor across a battery having an e.m.f. of 3.0 V. If the diode and the voltmeter are both ideal, suggest the reading on the voltmeter (a) for the circuit shown, (b) if the diode is reversed.
}



\begin{marginfigure}
	\vspace*{-8pt}
	\centering
	\begin{tikzpicture}
	\draw[very thin, gray!50, step=0.2] (0,0) grid (4,3);
	\draw[gray, step=1] (0,0) grid (4,3);
	\draw[thick, blue] (0,0) -- (1, 1.2) [out=50.2, in=190] to (4,2.8);
	\foreach \I in {0,1,2,3} \node[left] at (0,\I) {$\I$};
	\foreach \V in {0,2,4,6,8} \node[below] at (\V/2, 0) {$\V$};
	\node at (3.5, -0.7) {$V$/V};
	\node at (-0.7, 2.5) {$I$/A};
	\end{tikzpicture}
	\vspace*{-16pt}
\end{marginfigure}

\question{
	The diagram shows the $I$-$V$ characteristics of an electrical component $P$. (a) Suggest whether $P$ a filament lamp or a thermistor. (b) What is the resistance of $P$ for a current of 2.8 A? (c) Draw on the same graph the $I$-$V$ characteristics for a fixed resistor $Q$ of 4.0 $\Omega$. (d) Compare the power dissipated in $P$ and $Q$ for a p.d. of 4.0 V.
}








\include{current}
\include{electrostatics}
\chapter{Sub-atomic Particles}

\section{Atomic structure}

\subsection{Atomic model}\label{ch:atomic-model}


All matter is composed of tiny particles called \keypoint{atoms}, each different element has its own type of atom \& the diagram below illustrates a simple atomic model\index{atomic model}.
\begin{figure*}[h]
	\vspace*{5pt}
	\centering
	\includegraphics[height=180pt]{atomic-model.pdf}
	\vspace*{-20pt}
\end{figure*}

At the centre of each atom is the \keypoint{nucleus}\index{nucleus}. The radius of an atom $\sim 10^{-10} \sim 10^{-9} \text{ m}$ whereas the radius of a nucleus $\sim 10^{-15} \sim 10^{-14} \text{ m}$. Particles that make up the nucleus are called \keypoint{nucleons}\index{nucleon} and they come in two types, \keypoint{protons} and \keypoint{neutrons}\index{proton}\index{neutron}. Outside the nucleus are the \keypoint{electrons}\index{electron}, which move around the nucleus in a \emph{cloud}.

Protons, neutrons and electrons are the building blocks of all atoms and hence the building blocks for all matter

\cmt The properties of protons, neutrons and electrons are listed below:

\begin{compactenum}
	\item[--] charge of each proton and electron is the elementary charge unit: $e = 1.60 \times10^{-19} \text{ C}$
	
	\item[--] neutron has zero charge
	
	\item[--] a proton has very similar mass as a neutron
	
	\item[--] electron has much smaller mass than a nucleon (proton or neutron)
\end{compactenum}
\begin{figure*}
\begin{center}
	\begin{tabular}{|C{3.5cm}|C{1.6cm}|C{3.5cm}|C{2.8cm}|}
		\hline subatomic particle & charge & mass & location found \\ 
		\hline proton & $+e$ & $m_p = 1.67 \times 10^{-27}$ kg & nucleus \\ 
		\hline neutron & 0 & $m_n = 1.67 \times 10^{-27}$ kg & nucleus \\ 
		\hline electron & $-e$ & $m_e = 9.11 \times 10^{-31}$ kg & in outer atom \\ 
		\hline 
	\end{tabular}
\end{center}
\end{figure*}

\subsection{The $\alpha$-particle scattering experiment}

Hans Geiger, working in Rutherford's lab, did a series experiments in 1908 showing that alpha particles are "scattered" as they pass through thin layers of mica, and foils of gold and aluminium. In following year, joined by undergraduate Ernest Marsden, they did a series of experiments to untangle confusing results they observed. At the time of the experiment, the atom was thought to be analogous to a plum pudding (as proposed by J. J. Thomson), with the negatively-charged electrons (the plums) studded throughout a positive spherical matrix (the pudding). The scattering in this model was proposed to occur by many repeated collisions with the electrons and the alpha particles should only be deflected by small angles as they pass through.

The plum-pudding model predicts that deflection of $\alpha$-particles should be uniformly distributed around the central beam, but there is no reason that an $\alpha$-particle would undergo a large change in its direction of motion, like a cannonball hitting a piece of paper and rebounding backwards. The results of the scattering experiment meant that the plum-pudding model had to be overturned. In order to explain the experimental observations, Rutherford then put forward the new model which now bears his name.
\definecolor{mylightblue}{RGB}{220,228,255}
\definecolor{myblue}{RGB}{183,191,229}
\definecolor{mydarkblue}{RGB}{50,70,190}
\begin{marginfigure}
    \begin{tikzpicture}[scale=1]
  \coordinate (O)  at (0,0);
  \draw[mylightblue,fill] (O) circle (50pt);
  \draw[mydarkblue,thick] (O) circle (50pt) node[above right=34pt] {$+Ze$};
  \fill[radius=0.8pt,mydarkred]
    ( 0.20, 1.20)  circle node[above right=-1pt,scale=0.75] {$-e$}
    ( 0.10, 0.30)  circle node[above right=-1pt,scale=0.75] {$-e$}
    ( 0.68, 0.67)  circle node[above right=-1pt,scale=0.75] {$-e$}
    ( 1.05,-1.10)  circle node[above left =-1pt,scale=0.75] {$-e$}
    ( 0.75,-0.12)  circle node[      right= 2pt,scale=0.75] {$-e$}
    ( 0.21,-1.30)  circle node[above left =-1pt,scale=0.75] {$-e$}
    ( 0.18,-0.55)  circle node[above left =-1pt,scale=0.75] {$-e$}
    (-0.80, 0.60)  circle node[above left =-1pt,scale=0.75] {$-e$}
    (-0.50, 1.21)  circle node[above right=-1pt,scale=0.75] {$-e$}
    (-0.40, 0.05)  circle node[above left =-1pt,scale=0.75] {$-e$}
    (-0.58,-0.95)  circle node[above left =-1pt,scale=0.75] {$-e$}
    (-1.10,-0.20)  circle node[above left =-1pt,scale=0.75] {$-e$};
\end{tikzpicture}
\caption{The Plum pudding model of the atom.}
\end{marginfigure}

One of the most important experiments they carried out was the \keypoint{$\alpha$-particle scattering experiment}\index{$\alpha$-particle scattering experiment} in which \emph{Hans Geiger} and \emph{Ernest Marsden} studied atomic structure by firing $\alpha$-particle beam towards a thin gold foil at the University of Manchester.

\begin{figure}[!ht]
	\centering
	\includegraphics[height=180pt]{alpha-scattering}
	\caption{Rutherford's $\alpha$-particle scattering experiment}
\end{figure}

\begin{figure}[!ht]
	\centering
	\begin{tikzpicture}[scale=1.33]
		\shade[ball color=Gold] (0,0) circle (.2);
		\node[right, twoline] at (0.3,0) {gold\\nucleus}; 
		\draw[blue,->] (-4,0.03) -- (-0.4, 0.03) arc(-90:86:0.04) --++ (176:2.5);
		\draw[blue,->] (-4,1) -- (3, 1);
		\draw[blue,->] (-4,-1.2) -- (3, -1.2);
		\draw[blue,->] (-4,-0.5) -- (-1,-0.5) arc(90:65:4) --++ (-25:2.4);
	\end{tikzpicture}
	
	\caption{paths of $\alpha$-particles as they approach and pass by a gold nucleus}
\end{figure}

The main experimental results of the experiments are:

\begin{compactitem}
	\item[--] most $\alpha$-particles pass straight through the foil with almost no deflection
	
	\item[--] very few $\alpha$-particles (about 1 in $10^4$) are deflected by large angles.
\end{compactitem}

To explain the observations, Rutherford concluded the following about nature of atoms

\begin{compactitem}
	\item[--] most space in an atom is empty
	
	(so most $\alpha$-particles merely pass through empty space without deflection)
	
	\item[--] there is a tiny positively-charged core, called the nucleus, at centre of an atom
	
	(so only a small fraction of $\alpha$-particles would interact with the nucleus)
	
	\item[--] almost all mass of the atom is concentrated in the nucleus
	
	(so interaction is influential if nucleus happens to sit on the path of $\alpha$-particle beam)
\end{compactitem}

Based on these understandings, Rutherford proposed the atomic model introduced in \S\ref{ch:atomic-model}. This model is now known as \emph{Rutherford's model}, or \emph{planetary model of the atom}\footnote{\piste 
	 Rutherford originally suggested that electrons rotate around the nucleus in \emph{circular} orbits due to attractive electrostatic force, like the planets orbit around the sun. This is why Rutherford's model of atomic structure is also called \emph{the planetary model of the atom}.
	 
	 However, later studies showed that Rutherford's model had its own problems. Classical electromagnetic theory suggests that any charged particle in accelerated motion would radiate energy, so the orbiting electrons will gradually lose all its energy and fall into the nucleus. An atom can never be stable if the Rutherford's model is correct. The full description of the atomic structure requires the understanding of the \emph{quantum} behaviour of electrons, which are described by \emph{probability waves}. Electrons actually do not move in circular orbits but have a certain probability to be anywhere in space at any time, forming an \emph{electron cloud}.
 }




\subsection{Nuclide notation}

\keypoint{Nuclides} are specific types of atoms or nuclei. A nuclide is uniquely defined by its \keypoint{proton number} $Z$ and its \keypoint{nucleon number} $A$. A nuclide is usually denoted as $\tcbhighmath{_Z^A X}$, where $X$ is the chemical symbol for the element. 

\cmt By definition, $Z$ gives number of protons, $A$ gives number of nucleons, as a consequence, number of neutrons is given by $N = A-Z$. Proton number $Z$ can be further interpreted as \keypoint{charge number} of the particle \index{charge number}. Recall that each proton carries charge $+e$, and neutrons are chargeless so $Z$ further determines electrical charge of the nucleus: \begin{empheq}[box=\tcbhighmath]{equation*}{Q = +Ze}\end{empheq}.

This extension will be important when we deal with other particles later in this chapter.

The nucleon number $A$ can be interpreted as \keypoint{mass number} of the particle \index{mass number}, each proton and neutron is of approximately the same mass so the mass of nucleus is determined by total number of protons and neutrons: \begin{empheq}[box=\tcbhighmath]{equation*}{m = A\text{u}}\end{empheq}.

where $1 \text{ u} = 1.66\times10^{-27} \text{ kg}$ is the \keypoint{unified atomic mass unit}
\footnote[][-4cm]{Note that the unified atomic mass ($1.661\times10^{-27} \text{ kg}$), formally defined as one twelfth of the mass of an carbon-12 atom, is slightly less than the mass of a \emph{free} proton ($m_p=1.673\times10^{-27} \text{ kg}$ or that of a \emph{free} neutron ($m_n=1.675\times10^{-27} \text{ kg}$). This is because energy goes out when protons and neutrons combine to form a nucleus, and this would be associated with a decrease in the mass as suggested by \emph{Albert Einstein's mass-energy equivalence principle}. You can think of $\text{u}$ as the average mass of a proton or a neutron when a group of them squeeze together. If you want to find out the mass of a nucleus, you should use $\text{u}$. But if you are dealing with a single proton or a single neutron, you should use $m_p$ and $m_n$ instead.}

Using this notation, protons, neutrons and electrons can be represented by: \begin{empheq}[box=\tcbhighmath]{equation*}{^1_1\text{p}} \hspace{1cm}{^1_0\text{n}}\hspace{1cm}{^{\phantom{-}0}_{-1}\text{e}}\end{empheq}

\example{The nuclide uranium-238 can be denoted by $^{238}_{\phantom{2}92}\text{U}$. (a) State its proton number and its neutron number. (b) Find the charge and mass of the uranium-238 nucleus.}

\begin{soln} proton number: $Z=92$

neutron number: $N=A-Z=238-92=146$

nuclear charge: $Q=+Ze=+92\times1.60\times10^{-19}\approx1.47\times10^{-17}\text{ C}$

mass of nucleus: $m=A\text{u} = 238 \times 1.66\times10^{-27} \approx 3.95 \times 10^{-25}\text{ kg}$ \end{soln}


\example{Radius of a copper-64 nucleus is $4.8\times10^{-15} \text{ m}$. Calculate the mean nuclear density.}

\begin{soln}
\begin{equation*}
	\rho = \frac{m}{V} = \frac{A\text{u}}{\frac{4}{3}\pi r^3} = \frac{64 \times 1.66\times10^{-27}}{\frac{4}{3}\pi \times (4.8x10^{-15})^3} \RA \rho \approx 2.3 \times 10^{17} \text{ kg m}^{-3}
\end{equation*}

this is much greater than density of a copper block (about $8.9 \times  10^3 \text{ kg m}^{-3}$)

we can therefore infer that the nucleus indeed takes most mass of the atom \end{soln}


\subsection{Chemical properties of elements}

The chemical properties of an atom depend on the number of electrons, the proton number $Z$ gives us more than just the \keypoint{atomic number} of the element. In a neutral atom, there are same number of protons as electrons so atoms of same chemical properties, i.e., atoms of same element have same proton number. That said, atoms of same chemical properties can have different nuclei. \keypoint{Isotopes} of an element have same number of protons but a different number of neutrons\index{isotope}.


\example{State and explain whether carbon-12 ($^{12}_{\phantom{1}6}\text{C}$) and carbon-14  ($^{14}_{\phantom{1}6}\text{C}$) are different isotopes of the same element.}

\begin{soln} both carbon-12 and carbon-14 atoms have 6 protons

but carbon-12 has 6 neutrons, carbon-14 has 8 neutrons

same proton number but different neutron number, so different isotopes of same element \end{soln}




\subsection{Radioactivity}

\subsection{Radioactive decays}

If a nucleus is unstable then it may break down by releasing radiation from the nucleus, this process is called \keypoint{radioactive decay}\index{radioactive decay}. The three most common types of radioactive decays are $\alpha$-decay, $\beta$-decay and $\gamma$-decay. 

Radioactive decays are \emph{spontaneous} and decay events are independent of outside conditions, such as temperature, pressure, etc. Whether one nucleus is about to decay is also independent of the other nuclei around it. 

Radioactive decays depend only on the internal stability of the nucleus, they are truly \emph{random}. 

\subsection{Properties of $\alpha$-, $\beta$- \& $\gamma$-radiation}

\subsection*{Nature of \texorpdfstring{$\alpha$}{\textalpha}-radiation}

$\alpha$-particles are \emph{helium nuclei}, each made up of two protons and two neutrons\index{radioactive decay!$\alpha$-decay}

an $\alpha$-particle therefore carries a positive charge of 2$e$, and a mass of 4u

$\alpha$-particles emitted from a given source travel at almost the same speed

speed of $\alpha$-particles is usually less than one-tenth of the speed of light

\subsection*{Nature of \texorpdfstring{$\beta$}{\textbeta}-radiation}


$\beta$-radiation is a beam of high-speed \emph{electrons}\index{radioactive decay!$\beta$-decay}

$\beta$-particle is negatively charged and very light compared with a nucleus

$\beta$-particles travel at very high speeds close to the speed of light

\subsection*{Nature of \texorpdfstring{$\gamma$}{\textgamma}-radiation}

$\gamma$-radiation, or $\gamma$-photon, is a high-frequency electromagnetic wave\index{radioactive decay!$\gamma$-decay}

$\gamma$-radiation is electrically neutral and has no mass

like any other electromagnetic wave, $\gamma$-radiation travels at the speed of light in vacuum

\subsection*{Brief summary for nature of \texorpdfstring{$\alpha$}{\textalpha}-, \texorpdfstring{$\beta$}{\textbeta}- \& \texorpdfstring{$\gamma$}{\textgamma}-radiation}



\begin{center}
	{\renewcommand{\arraystretch}{1.2}
	\begin{tabular}{|C{2cm}|C{2.8cm}|C{2.8cm}|C{4.2cm}|}
	\hline  & $\alpha$-radiation & $\beta$-radiation & $\gamma$-radiation \\ 
	\hline nature & helium nucleus & electron & electromagnetic wave \\ 
	\hline notation & $^4_2\alpha$ & $^{\phantom{-}0}_{-1}\beta$ & $^0_0\gamma$ \\
	\hline mass & 4u & $\frac{1}{1800}\text{u} \approx 0$ & 0 \\[4pt]
	\hline charge & $+2e$ & $-e$ & 0 \\ 
	\hline speed & $0.1c$ & $0.99c$ & $c$ \\
	\hline
    \end{tabular}
    }
\end{center}

%\begin{question}
%	Some radioactive sources are found to contain helium gas. How do you explain that?
%\end{question}

%\begin{question}
%	$\alpha$-particles are used in the scattering experiment to investigate the structure of an atom. Is it appropriate to use a $\beta$-emitting source to do the same experiment? Explain your reasons.
%\end{question}


\example{If a single beam of $\alpha$-particles, $\beta$-particles and $\gamma$-radiation is sent into a uniform electric field as shown, sketch and explain the paths of the beam.}

\begin{figure}[ht]
	\centering
	\begin{tikzpicture}
	\draw[thick] (-3,-1) -- (3,-1) (-3,1) -- (3,1);
	\draw[thick] (0,-1) -- (0,-2) node[midway, right] {$-$};
	\draw[thick] (0,1) -- (0,2) node[midway, right] {$+$};
	\draw[thick, purple, ->] (-4,0) -- (4, 0) node[right] {$\gamma$};
	\draw[thick, Green, ->] (-4,0.05) -- (-3,0.05) parabola (-2.4,1) node[below right]{$\beta$};
	\draw[thick, blue, ->] (-4,-0.05) -- (-3,-0.05) parabola (3,-0.85) --++ (-14.93:1) node[right] {$\alpha$};
	\end{tikzpicture}
\end{figure}

\begin{soln} $\gamma$-radiation has zero charge, so it is not affected by electric field

$\alpha$-radiation is positively-charged, so it deflects towards the negative plate

$\beta$-radiation is negatively-charged, so it deflects towards the positive plate

$\beta$-particles are much lighter, so they have greater deflection \end{soln}


\subsection*{Penetration \& ionising power}

Different radiation have different abilities to pass through materials, $\alpha$-, $\beta$- and $\gamma$- radiation are able to knock out orbital electrons from an atom so radiation can cause the originally neutral atom to become charged. This process is called \keypoint{ionising}\index{ionising power}. 

Greater ionising ability would imply lower penetration power. 

$\gamma$-radiation is very penetrative but has very weak ionising power, $\gamma$-radiation is not charged, and hence it does not interact strongly with electrons. A $\gamma$-photon loses all its energy in single collision and gets absorbed i.e., one $\gamma$-photon can only ionise one atom so it has few collisions with electrons, so $\gamma$-radiation is not very ionising

$\alpha$- and $\beta$-radiation more ionising than $\gamma$-radiation because $\alpha$- and $\beta$-radiation are charged, making them interact with electrons more easily.

$\alpha$-particles are the most highly ionising, recall $\alpha$-particles are much more massive than electrons so one $\alpha$-particle can knock out many electrons before it loses all of its kinetic energy. Also $\alpha$-particles travel at low speeds, so frequent collision events take place in short range. $\alpha$-radiation therefore has strong ionising ability but weak penetration power.

The table below summarises penetration and ionising power of radioactive radiation:
\begin{figure*}

\begin{center}
	\begin{tabular}{|C{3cm}|C{3.2cm}|C{3.2cm}|C{3.2cm}|}
		\hline  & $\alpha$-radiation & $\beta$-radiation & $\gamma$-radiation \\ 
		\hline \multirow{2}{3.2cm}[-20pt]{penetration power} & low & fair & high \\
		 & stopped by thick cardboard or a few centimetres of air & stopped by metal plates of a few millimetres thick & stopped by thick lead or concrete of a few centimetres\\
		\hline ionising power & high & fair & low \\ 
		\hline
	\end{tabular} 
\end{center}
\end{figure*}


\subsection{Laws of conservation}

Some of the conserved quantities during any decay process are:
\begin{itemize}
    \item conservation of \emph{electric charge}
    \item conservation of \emph{charge number} ($Z$)\index{charge number}
    \item conservation of \emph{mass-energy}\index{mass-energy conservation}
    \item conservation of \emph{mass number} ($A$)\index{mass number}
    \item conservation of \emph{momentum} (review \S\ref{ch:momentum-conservation} if needed)
\end{itemize}

\subsection{Decay equations}

Using nuclide notation, radioactive decay processes can be described by \emph{decay equations} - the conserved quantities provide a simple guideline to write the decay equations 

\begin{compactenum}
	\item[--] sum of charge numbers ($Z$) on both sides of the equation should add up
	
	\item[--] sum of mass numbers ($A$) on both sides of the equation should add up
\end{compactenum}


The first reported nuclear reaction was credited to Ernest Rutherford. By firing $\alpha$-particles into pure nitrogen, Rutherford observed the ejection of hydrogen nuclei from the gas, which is now regarded as the discovery of protons. This reaction can be given by:
\begin{equation*}
_{\phantom{0}7}^{14} \text{N} + _2^4 \alpha \longrightarrow _{\phantom{0}8}^{17} \text{O} + _{1}^{1} \text{p} 
\end{equation*} 


One reaction that is widely used to power nuclear reactors is the induced fission of uranium-235. The reaction can be triggered by bombarding the uranium-235 nuclei with slow neutrons. The uranium nucleus would split up into two lighter nuclei and release a large amount of energy. One of such reactions can be described by:
\begin{equation*}
 _{\phantom{0}92}^{235} \text{U} + _0^1 \text{n} \longrightarrow _{\phantom{0}56}^{141} \text{Ba} + _{36}^{92} \text{Kr} + 3_0^1 \text{n} 
\end{equation*} 




\subsection*{Equation for \texorpdfstring{$\alpha$}{\textalpha}-decays}\index{radioactive decay!$\alpha$-decay}

If a nuclide $_Z^A X$ undergoes $\alpha$-decays, then we have:
\begin{empheq}[box=\tcbhighmath]{equation*}{_Z^A X \longrightarrow _{Z-2}^{A-4} Y + _2^4\alpha}
\end{empheq}

The original nucleus lost two protons and two neutrons during the decay, on both sides of the equation, there are $Z$ protons and $(A-Z)$ neutrons therefore $\alpha$-decay simply is an $\alpha$-particle escaping from the nucleus.



\example{The process in which a uranium-238 nucleus naturally decays into a thorium-234 nucleus through $\alpha$-emission can be written as:}
\begin{equation*}
	_{\phantom{0}92}^{238} \text{U} \longrightarrow _{\phantom{0}90}^{234} \text{Th} + _2^4\alpha 
\end{equation*}



\subsection*{Equation for \texorpdfstring{$\gamma$}{\textgamma}-decays}\index{radioactive decay!$\gamma$-decay}

$\gamma$-radiation has zero charge and zero mass, so decay equation is very straightforward:
\begin{empheq}[box=\tcbhighmath]{equation*}{_Z^A X \longrightarrow _{Z}^{A} X + _0^0\gamma}
\end{empheq}

since $\gamma$-radiation is pure energy, so there is no change of structure of the nucleus in any way



\subsection*{Problems with \texorpdfstring{$\beta$}{\textbeta}-decays}\index{radioactive decay!$\beta$-decay}

for $\beta$-decays, one might attempt to write: $
	_Z^A X \longrightarrow _{Z+1}^{\phantom{1+}A} Y + _{-1}^{\phantom{+}0}\beta$ 

on the left, nuclide $X$ has $Z$ protons and $(A-Z)$ neutrons

but on the right, nuclide $Y$ has $(Z+1)$ protons and $(A-Z-1)$ neutrons

one neutron has transformed into a proton through giving off an electron: $_0^1 \text{n} \longrightarrow _1^1 \text{p} + _{-1}^{\phantom{+}0}\text{e} \,\, \text{???}$

but how can neutrons and protons transform into one another?!

further experiments show that $\beta$-particles emitted for same source have \emph{a range of speeds}

this seems to indicate that energy released from a $\beta$-decay can be indefinite

both considerations imply something is wrong, our understanding of $\beta$-decay is not complete

\begin{compactenum}
	\item[--] protons and neutrons cannot be the most fundamental particles of nature
	
	there must exist some more fundamental constituent particles
	
	these are now known as \emph{quarks}, we will discuss them in the next section
	
	\item[--] there must be some other unseen particles released during $\beta$-decays
	
	these particles carry off a fraction of the energy released from the reaction
	
	for conservation laws to hold, they must be chargeless and very light (maybe massless)
	
	these ghostly particles are called \keypoint{neutrinos}, or more precisely, \keypoint{anti-neutrinos}\index{neutrino}
\end{compactenum}

\cmt we thereby can rewrite the $\beta$-decay equation for nuclide $_Z^A X$:
\begin{empheq}[box=\tcbhighmath]{equation*}
{_Z^A X \longrightarrow _{Z+1}^{\phantom{1+}A} Y + _{-1}^{\phantom{+}0}\beta + _0^0 \bar{\nu}} \quad \text{or} \quad {_0^1 \text{n} \longrightarrow _1^1 \text{p} + _{-1}^{\phantom{+}0}\beta + _0^0 \bar{\nu}} \label{eqn:beta-decay}
\end{empheq}

\cmt distribution of kinetic energy of $\beta$-particles can be then explained

only total energy of $\beta$-particle and anti-neutrino is constant

anti-neutrinos also carry some energy, so $\beta$-particles have a range of energies

\example{Radon $^{222}_{\phantom{2}86}\text{Rn}$ decays in a sequence of processes to form bismuth $^{214}_{\phantom{2}83}\text{Bi}$ by emitting $\alpha$-particles and $\beta$-particles. For the decay chain of each radon nucleus, how many $\alpha$-particles and $\beta$-particles are emitted?}

\begin{soln} let $x$ be number of $\alpha$-particles emitted and $y$ be number of $\beta$-particles emitted
\begin{equation*}
	^{222}_{\phantom{2}86}\text{Rn} \longrightarrow ^{214}_{\phantom{2}83}\text{Bi} + x \cdot ^4_2 \alpha + y \cdot _{-1}^{\phantom{+}0}\beta + y \cdot _0^0 \bar{\nu}
\end{equation*}

Mass number and charge number are conserved:\\ 
$\Bigg\{$

\vspace*{-1.58\baselineskip}\hspace*{18pt} $\begin{array}{l}
4x + 214 = 222\\[-5pt]
2x - y + 83 = 86
\end{array} \RA x = 2$, $y=1$ \end{soln}



\subsection{Fundamental particles}

In early 20th century, lots of new subatomic particles were discovered in \emph{cosmic rays} and \emph{particle accelerators}. Many of these particles did not fit into the model where proton, neutrons and electrons are the fundamental building blocks for the physical world. To incorporate these subatomic particles discovered at the time, physicist worked out the \keypoint{Standard Model of particle physics}.\footnote[][-3cm]{The theory of the Standard Model was developed in stages since the 1950s, through the work of many great scientists, with the current formulation being finalized in the 1970s.\index{Standard Model}
	\begin{compactitem}
		\item[--] \emph{Chen Ning Yang} and \emph{Robert Mills} developed the concept of \emph{gauge theory}, to provide an explanation for the interaction between elementary particles.
		
		\item[--] \emph{Sheldon Glashow} proposed the symmetry group that forms the basis of the accepted \emph{electroweak theory}, in which the electromagnetic interaction and weak interactions are unified into a single force.
		
		\item[--] \emph{Peter Higgs} and two other groups proposed the \emph{Higgs mechanism} that give rise to mass generation for elementary particles without violating gauge theory through a process called \emph{symmetry breaking}.
		
		\item[--] \emph{Steven Weinberg} and \emph{Abdus Salam} incorporated the Higgs mechanism into Glashow's theory, finalizing the unified electroweak theory.
		
		\item[--] \emph{David Gross}, \emph{Frank Wilczek}, \emph{David Politze}, and many others developed \emph{quantum chromodynamics}, the theory of the strong interaction, into its modern form.
		
		\item[--] $\cdots$
	\end{compactitem}

Being the most accurate accurate scientific theory known to human beings, the Standard Model is now regarded as one of the greatest triumphs of modern physics. The theory not only describes the particles known to scientists, but also predicted new particles, including the \emph{Higgs boson}. }


\subsection{Particles \& anti-particles}

For any particle $\text{p}$, there exists an associated \keypoint{anti-particle} $\bar{\text{p}}$\index{anti-particle}

\cmt an anti-particle has the same mass but opposite charge to its counterpart

\cmt when a particle meets its anti-particle, they \keypoint{annihilate} each other and produce two $\gamma$-photons

combined mass-energy of the pair is converted into electromagnetic energy
\footnote{The energy release from the annihilation of  particle pairs could be potentially used a energy source. The idea of using antimatter to power spaceships or weapons can be found in many science fiction stories.}



\example{Suggest the properties of the anti-particle of the electron.}

\begin{soln} anti-particle of electron has same electron mass: $m_e = 9.11\times10^{-31} \text{ kg}$

but it has a positive charge: $q = +e = + 1.60\times10^{-19} \text{ C}$

anti-particle of electron is usually called a \keypoint{positron} (positive electron) with the symbol $\text{e}^+$ \end{soln}


\subsection{Quarks \& hadrons}

One type of the elementary matter particle is the \keypoint{quark}\index{quark}

\cmt quarks come in six varieties, or six \keypoint{flavours}

they are \keypoint{up} (u), \keypoint{down} (d), \keypoint{strange} (s), \keypoint{chartm} (c), \keypoint{top} (t), \keypoint{bottom} (b)

\cmt each quark carries an electric charge, as given in the table below

\begin{center}
	{\renewcommand{\arraystretch}{1.35}
		\begin{tabular}{|C{3cm}|C{2cm}|C{2cm}|}
			\hline  & u & d \\
			 quarks & c & s \\
			  & t & b \\
			\hline charge & $+\frac{2}{3}e$ & $-\frac{1}{3}e$ \\[3pt]
			\hline
		\end{tabular}
	}
\end{center}

note that all quarks carry a \emph{fraction} of the elementary charge unit\footnote{Each horizontal line in the table is known as a \emph{generations} of quarks. As you can see, there are three generations of quarks.}

Quarks can combine to form composite particles called \keypoint{hadrons} \index{hadron}, there are two ways that several quarks can make up a hadron:

\begin{compactenum}
	\item[--] hadrons can be made up of three quarks (qqq), called \keypoint{baryons} \index{baryon}
	
	members of the baryon family include protons and neutrons
	
	a proton consists of two up quarks and one down quark (uud) \index{proton}
	
	a neutron consists of one up quark and two down quarks (udd) \index{neutron}
	
	\item[--] hadrons can also be made up of one quark and one anti-quark (q$\bar{\text{q}}$), called \keypoint{mesons} \index{meson}
	
	you don't need to memorise any example of meson
	
	you are only required to identify if a particle is a meson given the quark composition
\end{compactenum}

\cmt quarks and hadrons are affected by the \keypoint{strong nuclear force}\index{strong nuclear force}

strong nuclear force is a very short-ranged attractive force

\begin{compactenum}
	\item[--] it is responsible for hold quarks close together to from hadrons

	\item[--] it is also responsible for the attraction between hadrons
	
	for example, protons and neutrons are held together in a nucleus by strong force
	\footnote{For a small nucleus, strength of strong nuclear force is way greater than the electric repulsion between positively-charged protons, hence the force bind protons and neutrons together forming a stable nucleus. However, since the strong nuclear force has a very short range, the repulsive electrostatic force between the protons might dominate the attractive nuclear force for an over-sized nucleus. Such large nuclei become unstable and are likely to undergo \emph{radioactive decays}.}
\end{compactenum}


\cmt quarks only exist in hadrons, i.e., there are no single free quarks in nature\footnote{This is known as the \emph{confinement}, a consequence that is closely related to the fact that interaction between quarks are weak at high energies or smaller length scales but strong at low energies or large length scales, a feature of strong nuclear forces known as \emph{asymptotic freedom}.}


\example{By reference to the quark composition, explain the electric charge of protons and neutrons.}

\begin{soln} charge of proton (uud): $q_\text{p} = 2q_\text{u} + q_\text{d} = 2\times\left(+\frac{2}{3}e\right) + \left(-\frac{1}{3}e\right) = +e$

\eqyskip charge of neutron (udd): $q_\text{n} = q_\text{u} + 2q_\text{d} = \left(+\frac{2}{3}e\right) + 2\times\left(-\frac{1}{3}e\right) = 0$ \end{soln}

\example{A meson has an electric charge of $+e$ and is known to contain an up quark. Determine a possible flavour of the other quark.}

\begin{soln} charge of the other quark is: $(+e) - \left(+\frac{2}{3}e\right) = +\frac{1}{3}e$

this could be an anti-down ($\bar{\text{d}}$), an anti-strange ($\bar{\text{s}}$), or an anti-bottom quark ($\bar{\text{b}}$)\end{soln}



\subsection{Leptons}

Another type of the elementary matter particle is the \keypoint{lepton}\index{lepton} - they are not affected by strong nuclear forces. Members of the lepton family include electrons (e), neutrinos ($\nu$), muons ($\mu$), taons ($\tau$)\footnote{Just like the quarks, leptons also come in three \emph{generations}. Each generation of leptons consists of an electron-like particle and its associated neutrino.}

in A-Level exams, you only need to know electrons (e) and neutrinos ($\nu$)


\cmt for each particle, one can assign a \keypoint{lepton number} $L$

\begin{compactitem}
	\item[--] a lepton has lepton number $L=+1$
	
	\item[--] an anti-lepton has lepton number $L=-1$
	
	 \item[--]any other particle has lepton number $L=0$
\end{compactitem}

classification of sub-atomic particles within the A-Level syllabus is summarised below\footnote[][-2cm]{This is an over-simplified version of the Standard Model. I only included the particles that you need to know for the A-Level exams. There are a lot of missing pieces that are way beyond the scope of our course, so please don't take this mindmap too seriously.}

\begin{figure*}[ht]
	\centering
	\begin{tikzpicture}
%\draw[help lines, gray!30, step=0.5] (0,-5) grid (15,5);
%\foreach \x in {1,2,...,15} \node[below] at (\x,0) {\x};
%\foreach \y in {-5,-4,...,5} \node[left] at (0,\y) {\y};
	% root / fundamental particles
	\draw (2.5,0.4) -- (3.5,0.4) (4.5, -.6) -- (3.5,-0.6) -- (3.5, 2.5) -- (4.5, 2.5);
	\draw[rounded corners, fill=white] (0.1,-0.3) rectangle (2.5,1.1);
	\draw (1.3,0.4) node[twolinecap] {fundamental\\particles};
	% quarks
	\draw[rounded corners, fill=white] (4.3,2.2) rectangle (5.9, 2.8);
	\draw (5.1, 2.5) node {quarks};	
	\draw (5.0, 2.1) -- (4.6, 1.6);
	\draw (5.2, 2.1) -- (5.6, 1.6);
	\node at (4.6, 1.4) {u};
	\node at (5.6, 1.4) {d};
	\node at (4.6, 0.9) {{\scriptsize $+\frac{2}{3}e$}};
	\node at (5.6, 0.9) {{\scriptsize $-\frac{1}{3}e$}};
	% leptons
	\draw[rounded corners, fill=white] (4.3,-.4) rectangle (5.9, -1.0);
	\draw (5.1, -.7) node {leptons};
	\draw (5.0, -1.1) -- (4.6, -1.6) node[below]{e};
	\draw (5.2, -1.1) -- (5.6, -1.6) node[below]{$\nu$};
	\node at (4.6, -2.2) {{\scriptsize $-e$}};
	\node at (5.6, -2.2) {{\scriptsize $0$}};
	% hadrons
	\draw[dashed, ->] (5.9, 2.5) -- (8.1,2.5) node[midway, above]{make up};
	\draw[rounded corners, fill=white] (8.1,2.2) rectangle (9.9, 2.8);
	\draw (9, 2.5) node {hadrons};
	% baryons
	\draw (9.9, 2.5) -- (10.5, 2.5) -- (10.5, 3.5) -- (12, 3.5) node[midway, above]{(qqq)};
	\draw[rounded corners, fill=white] (12,3.2) rectangle (13.8, 3.8);
	\draw (12.9, 3.5) node {baryons};
	\draw (12.8, 3.1) -- (12.4, 2.6) node[below]{p};
	\draw (13.0, 3.1) -- (13.4, 2.6) node[below]{n};
	% mesons
	\draw (10.5, 2.5) -- (10.5, 1.5) -- (12, 1.5) node[midway, above]{(q$\bar{\text{q}}$)};
	\draw[rounded corners, fill=white] (12,1.2) rectangle (13.8, 1.8);
	\draw (12.9, 1.5) node {mesons};
	\end{tikzpicture}
	\caption{a coarse guide for classifying sub-atomic particles in A-Level physics}
\end{figure*}



\subsection{\texorpdfstring{$\beta$}{\textbeta}-decays}\index{radioactive decay!$\beta$-decay}

there exists positron, anti-particle of electron, so there are two types of $\beta$-particles: $\beta^-$ and $\beta^+$

hence two types of $\beta$-decays are possible: $\beta^-$-decay and $\beta^+$-decay



\subsection*{\texorpdfstring{$\beta^-$}{\textbeta\textsubscript{-}}-decay revisited}

a neutron changes into a proton during a $\beta^-$-decay: $\boxed{_0^1 \text{n} \longrightarrow _1^1 \text{p} + _{-1}^{\phantom{+}0}\beta + _0^0 \bar{\nu}} $

we have learned about the quark structures of protons (uud) and neutrons (udd)

therefore $\beta^-$-decay process at a more fundamental level is: \begin{empheq}[box=\tcbhighmath]{equation*}{\text{d} \longrightarrow \text{u} + \beta^- + \bar{\nu}} \end{empheq}


\subsection*{\texorpdfstring{$\beta^+$}{\textbeta\textsubscript{+}}-decay}


$\beta^+$-decay occurs when a proton changes into a neutron and emits a positron

$\beta^+$-decay process is described by the equation: \begin{empheq}[box=\tcbhighmath]{equation*}{_1^1 \text{p} \longrightarrow _0^1 \text{n} + _{+1}^{\phantom{+}0}\beta + _0^0 \nu} \end{empheq}

in terms of quarks, $\beta^+$-decay can be rewritten as: \begin{empheq}[box=\tcbhighmath]{equation*}{\text{u} \longrightarrow \text{d} + \beta^+ + \nu} \end{empheq}

\cmt nature of $\beta$-decay processes is the transformation of quark flavours

the interaction responsible for this transformation is the \keypoint{weak nuclear force}\index{weak nuclear force}

\cmt there is yet another law of conservation -- the \keypoint{conservation of lepton number}

we can use lepton number to predict whether a neutrino or an anti-neutrino is emitted

\example{Verify the conservation of (a) charge number, (b) mass number, and (c) lepton number for the $\beta^+$-decay process: $_1^1 \text{p} \longrightarrow _0^1 \text{n} + _{+1}^{\phantom{+}0}\beta + _0^0 \nu$}

\begin{soln} charge number conservation: $+1 = 0 + (+1) + 0$ \xskip \ding{51}

mass number conservation: $1 = 1 + 0 + 0$ \xskip \ding{51}

lepton number conservation: $0 = 0 + (-1) + (+1)$ \xskip \ding{51}

so this reaction does not violate any of these conservation law \end{soln}


\example{\emph{Electron capture} is a process where an electron in an atom's inner shell is drawn into the nucleus and combines with a proton. The result is to form a neutron and a neutrino. The process can be given by: $_1^1 \text{p} + _{-1}^{\phantom{+}0}\text{e} \longrightarrow _0^1 \text{n} +  + _0^0 \nu$. Show that the process satisfies the conservation of (a) charge number, (b) mass number, and (c) lepton number.}

\begin{soln} charge number conservation: $+1 + (-1) = 0 + 0$  \xskip \ding{51}

mass number conservation: $1 + 0 = 1 + 0$ \xskip \ding{51}

lepton number conservation: $0 + 1 = 0 + 1$  \xskip \ding{51}

so this reaction does not violate any of these conservation law \end{soln}

	
\section{End-of-chapter questions}

\subsection*{Atomic structure}

\question{
	List the number of sub-atomic particles in an argon-40 ($^{40}_{18}\text{Ar}$) atom.
}


 



\subsection*{Radioactive decays}
	
\question{
	A radioactive source emits a parallel beam of $\alpha$-, $\beta$- and $\gamma$-radiation. A detector sensitive to all forms of radiation has been calibrated for background radiation and is placed just 1 cm from the radioactive source. The detector registers 500 counts s$^{-1}$.  
		
	When the detector is 10 cm from the source, the radiation level drops to 200 counts s$^{-1}$.
	
	When a thin aluminium sheet is placed in front of the detector, the level falls to 50 counts s$^{-1}$.
	
	In what proportion is the source emitting $\alpha$-, $\beta$- and $\gamma$-particles?
}



\subsection*{Nuclear equations}

\question{
	An unstable isotope of phosphorus, $_{15}^{30}$P can be produced by bombarding aluminium, $_{13}^{27}$Al, with $\alpha$-particles. What is the by-product of this reaction?
}

\question{
	Carbon-14 ($^{14}_{\phantom{0}6}$ C) undergoes $\beta$-decay and forms an isotope of nitrogen (N). Write down the decay equation using the nuclide notations.
}

\question{
	Radon-222 ($^{222}_{\phantom{2}86}$Rn) undergoes a series of decays and forms lead-210 ($^{210}_{\phantom{2}82}$Pb). 
	How many $\alpha$-particles and $\beta$-particle are emitted during this process?
}

\subsection*{Fundamental particles}


\question{
	State the mass and electric charge of an anti-proton.
}

\question{
	$\pi^+$ is a particle with quark structure $u\bar{d}$. (a) State the family of this particle. (b) Show that it carries a charge of $+e$. (c) Suggest the charge of its anti-particle.
}



\question{
	$\Xi^0$ is a particle with quark structure $uss$. (a) State the family of $\Xi^0$. (b) Find its electric charge.
}

\question{
	Magnesium-23 ($^{23}_{12}$Mg) can decay into sodium-23 ($^{23}_{11}$Na). (a) Write down the decay equation using nuclide notation, and specify any other product particles produced. (b) Write an equation for this decay in terms of protons and neutrons. (c) Write an equation for this decay in terms of quark composition. (d) State the name of the force that is responsible for this decay.
}
	



\include{18-circular-motion}
\include{19-gravitation}
\include{20-oscillation}
\include{21-ideal-gases}
\include{22-thermal-physics}
\include{23-coulomb-law}
 \include{24-capacitance}
 \include{25-magnetic-fields}
\include{27-electromagnetic-induction}
 %\include{28-alternating-currents} not on course
 \include{29-quantum-physics}
 \include{30-nuclear-physics}
 \include{31-electronics}
% %\include{32-medical-imaging} not on course
% % \include{33-telecommunication} not on course
 \chapter{Jupyter}
\section{Introduction}
Jupyter/IPython Notebook Quick Start Guide
This document is a brief step-by-step tutorial on installing and running Jupyter (IPython) notebooks on local computer for new users who have no familiarity with python.

Briefly, if someone gave you a notebook to run and you don’t know what a notebook is, this document is for you.

Jupyter Notebook App (formerly IPython Notebook) is an application running inside the browser. This guide describes how to install and use Jupyter Notebook App as normal desktop application, without using any remote server.

For other use-cases, please refer to the Official Jupyter Documentation.

\section{What is the Jupyter Notebook?}

In this page briefly introduce the main components of the Jupyter Notebook environment. For a more complete overview see References.

\subsection{Notebook document}
Notebook documents (or “notebooks”, all lower case) are documents produced by the Jupyter Notebook App, which contain both computer code (e.g. python) and rich text elements (paragraph, equations, figures, links, etc…). Notebook documents are both human-readable documents containing the analysis description and the results (figures, tables, etc..) as well as executable documents which can be run to perform data analysis.
\subsection{Jupyter Notebook App}
The Jupyter Notebook App is a server-client application that allows editing and running notebook documents via a web browser. The Jupyter Notebook App can be executed on a local desktop requiring no internet access (as described in this document) or can be installed on a remote server and accessed through the internet.

In addition to displaying/editing/running notebook documents, the Jupyter Notebook App has a “Dashboard” (Notebook Dashboard), a “control panel” showing local files and allowing to open notebook documents or shutting down their kernels.

\subsection{kernel}
A notebook kernel is a “computational engine” that executes the code contained in a Notebook document. The ipython kernel, referenced in this guide, executes python code. Kernels for many other languages exist (official kernels).

When you open a Notebook document, the associated kernel is automatically launched. When the notebook is executed (either cell-by-cell or with menu Cell -> Run All), the kernel performs the computation and produces the results. Depending on the type of computations, the kernel may consume significant CPU and RAM. Note that the RAM is not released until the kernel is shut-down.

See also Close a notebook: kernel shut down.

References: from the official docs Opening Notebooks and Decoupled two-process model.

\subsection{Notebook Dashboard}
The Notebook Dashboard is the component which is shown first when you launch Jupyter Notebook App. The Notebook Dashboard is mainly used to open notebook documents, and to manage the running kernels (visualize and shutdown).

The Notebook Dashboard has other features similar to a file manager, namely navigating folders and renaming/deleting files.
Shut down the Jupyter Notebook App
Closing the browser (or the tab) will not close the Jupyter Notebook App. To completely shut it down you need to close the associated terminal.

In more detail, the Jupyter Notebook App is a server that appears in your browser at a default address (http://localhost:8888). Closing the browser will not shut down the server. You can reopen the previous address and the Jupyter Notebook App will be redisplayed.

You can run many copies of the Jupyter Notebook App and they will show up at a similar address (only the number after “:”, which is the port, will increment for each new copy). Since with a single Jupyter Notebook App you can already open many notebooks, we do not recommend running multiple copies of Jupyter Notebook App.

\subsection{Close a notebook: kernel shut down}
When a notebook is opened, its “computational engine” (called the kernel) is automatically started. Closing the notebook browser tab, will not shut down the kernel, instead the kernel will keep running until is explicitly shut down.

To shut down a kernel, go to the associated notebook and click on menu File -> Close and Halt. Alternatively, the Notebook Dashboard has a tab named Running that shows all the running notebooks (i.e. kernels) and allows shutting them down (by clicking on a Shutdown button).

\section{Notebook Basics}
\subsection{The Notebook dashboard}
When you first start the notebook server, your browser will open to the notebook dashboard. The dashboard serves as a home page for the notebook. Its main purpose is to display the notebooks and files in the current directory. For example, here is a screenshot of the dashboard page for the examples directory in the Jupyter repository:

\begin{figure}
    \centering
    \includegraphics[width=\linewidth]{dashboard_files_tab.png}
\end{figure}

The top of the notebook list displays clickable breadcrumbs of the current directory. By clicking on these breadcrumbs or on sub-directories in the notebook list, you can navigate your file system.\\

To create a new notebook, click on the "New" button at the top of the list and select a kernel from the dropdown (as seen below). Which kernels are listed depend on what's installed on the server. Some of the kernels in the screenshot below may not exist as an option to you.

\begin{marginfigure}
    \centering
    \includegraphics[]{dashboard_files_tab_new.png}
\end{marginfigure}

Notebooks and files can be uploaded to the current directory by dragging a notebook file onto the notebook list or by the "click here" text above the list.

The notebook list shows green "Running" text and a green notebook icon next to running notebooks (as seen below). Notebooks remain running until you explicitly shut them down; closing the notebook's page is not sufficient.

\begin{figure}
    \centering
    \includegraphics[width=0.5\linewidth]{dashboard_files_tab_run.png}
\end{figure}

To shutdown, delete, duplicate, or rename a notebook check the checkbox next to it and an array of controls will appear at the top of the notebook list (as seen below). You can also use the same operations on directories and files when applicable.

\begin{figure}
    \centering
    \includegraphics[width=0.5\linewidth]{dashboard_files_tab_btns.png}
\end{figure}

To see all of your running notebooks along with their directories, click on the "Running" tab:

\begin{figure}
    \centering
    \includegraphics[width=0.5\linewidth]{dashboard_running_tab.png}
\end{figure}

This view provides a convenient way to track notebooks that you start as you navigate the file system in a long running notebook server.

Overview of the Notebook UI
If you create a new notebook or open an existing one, you will be taken to the notebook user interface (UI). This UI allows you to run code and author notebook documents interactively. The notebook UI has the following main areas:

\titem Menu
\titem Toolbar
\titem Notebook area and cells
The notebook has an interactive tour of these elements that can be started in the "Help:User Interface Tour" menu item.

\subsection{Modal editor}
Starting with IPython 2.0, the Jupyter Notebook has a modal user interface. This means that the keyboard does different things depending on which mode the Notebook is in. There are two modes: edit mode and command mode.

\subsection{Edit mode}
Edit mode is indicated by a cell border and a prompt showing in the editor area:

\begin{figure}
    \centering
    \includegraphics[width=0.5\linewidth]{edit_mode.png}
\end{figure}

When a cell is in edit mode, you can type into the cell, like a normal text editor.

Enter edit mode by pressing `Enter` or using the mouse to click on a cell's editor area.
In this mode you can change the type of cell with the dropdown menu:  
\begin{figure}
    \centering
    \includegraphics[width=0.5\linewidth]{celltype.png}
\end{figure}

A \emph{markdown cell} allows you to type and format text using the markdown 


\subsection{Command mode}
Command mode is indicated by a grey cell border with a blue left margin:

\begin{figure}
    \centering
    \includegraphics[width=0.5\linewidth]{command_mode.png}
\end{figure}

When you are in command mode, you are able to edit the notebook as a whole, but not type into individual cells. Most importantly, in command mode, the keyboard is mapped to a set of shortcuts that let you perform notebook and cell actions efficiently. For example, if you are in command mode and you press c, you will copy the current cell - no modifier is needed.

Don't try to type into a cell in command mode; unexpected things will happen!
Enter command mode by pressing `Esc` or using the mouse to click *outside* a cell's editor area.

\subsection{Mouse navigation}
All navigation and actions in the Notebook are available using the mouse through the menubar and toolbar, which are both above the main Notebook area:

\begin{figure}
    \centering
    \includegraphics[width=0.5\linewidth]{menubar_toolbar.png}
\end{figure}

The first idea of mouse based navigation is that cells can be selected by clicking on them. The currently selected cell gets a grey or green border depending on whether the notebook is in edit or command mode. If you click inside a cell's editor area, you will enter edit mode. If you click on the prompt or output area of a cell you will enter command mode.\\

If you are running this notebook in a live session (not on https://nbviewer.jupyter.org) try selecting different cells and going between edit and command mode. Try typing into a cell.\\

The second idea of mouse based navigation is that cell actions usually apply to the currently selected cell. Thus if you want to run the code in a cell, you would select it and click the button in the toolbar or the "Cell:Run" menu item. Similarly, to copy a cell you would select it and click the button in the toolbar or the "Edit:Copy" menu item. With this simple pattern, you should be able to do most everything you need with the mouse.\\

Markdown cells have one other state that can be modified with the mouse. These cells can either be rendered or unrendered. When they are rendered, you will see a nice formatted representation of the cell's contents. When they are unrendered, you will see the raw text source of the cell. To render the selected cell with the mouse, click the button in the toolbar or the "Cell:Run" menu item. To unrender the selected cell, double click on the cell.\\

\section{Keyboard Navigation}
The modal user interface of the Jupyter Notebook has been optimized for efficient keyboard usage. This is made possible by having two different sets of keyboard shortcuts: one set that is active in edit mode and another in command mode.

The most important keyboard shortcuts are Enter, which enters edit mode, and Esc, which enters command mode.

In edit mode, most of the keyboard is dedicated to typing into the cell's editor. Thus, in edit mode there are relatively few shortcuts. In command mode, the entire keyboard is available for shortcuts, so there are many more. The Help->Keyboard Shortcuts dialog lists the available shortcuts.
\subsection{Shortcuts}
I recommend learning the command mode shortcuts in the following rough order:

Basic navigation: enter, shift-enter, up/k, down/j\\
Saving the notebook: s\\
Change Cell types: y, m, 1-6, t\\
Cell creation: a, b\\
Cell editing: x, c, v, d, z\\
Kernel operations: i, 0 (press twice)\\

\chapter{Conducting a practical}

Lets see how the previous chapter works in practice. To begin with we open a new jupyter file and name it sensibly:\\
\begin{figure}
    \centering
    \includegraphics[width=1\linewidth]{figures/SimplePrac.png}
\end{figure}
Then we start to add the detail\sidenote[][]{All practical work should contain: \begin{itemize} \item Title\\ \item Apparatus list and/or set-up diagram and commentary on why you have chosen the equipment. eg. Using a $\pm 0.001$mm micrometer would give you appropriate precision for measuring a thin wire\\ \item A method with appropriate safety precautions (this needn't be longer than "set up as in diagram" for some practicals!)\\ \item A results table\\ \item A concluding statement\\ \end{itemize}}: 
\begin{figure}
    \centering
    \includegraphics[width=1\linewidth]{figures/theoryprac.png}
\end{figure}
I've added an image by finding a good one online, saving it, and dragging and dropping it to the cell:\\
\begin{figure}[h!]
    \centering
    \includegraphics[width=1\linewidth]{figures/pracimage.png}
\end{figure}
Now we're into the code for your results etc. Some commented code explaining the function of each bit follows. A sensible student would use a basic version of this as a boilerplate (a partially complete template version) for all of your future write-ups.\\ 

\begin{figure}[h!]
    \centering
    \includegraphics[width=1\linewidth]{figures/praccode.png}
\end{figure}
\newpage
\subsection{Example code for your results}

This code is for the "g by freefall" experiment in which a ball is dropped, it's time of flight timed, and the acceleration due to gravity calculated.

\begin{fullwidth}
    

\begin{mintedbox}{python}
from matplotlib import pyplot as plt
from scipy import stats
from tabulate import tabulate

#Taking mean
# This is the number of different heights in this example
NumRepeats=3
#This is the resolution of the timer. 
Resolutiony = 0.05
Resolutionx = 0.01
#This is a list of list of the multiple readings of time at each height
Xmultraw=[[0.44,0.41,0.44],[0.49,0.47,0.50],[0.50,0.53,0.52],[0.57,0.56,0.54],[0.60,0.62,0.59],
[0.61,0.66,0.63],[0.67,0.67,0.67],[0.70,0.70,0.70]]
#This is the height measurements
yraw = [1.00,1.20,1.40,1.60,1.80,2.00,2.20,2.40]


# Looping over the number of heights. Calculating the average time and the uncertainty in those times and adding them to the 
# list xraw and uncertainty
xraw=[]
uncertaintyxraw=[]
for i in range (len(yraw)):
    sumx=0.0
    for j in range (NumRepeats):
        sumx+=Xmultraw[i][j]        
    HalfRange=(max(Xmultraw[i])-min(Xmultraw[i]))/2.0
    ave=sumx/NumRepeats
    if HalfRange < Resolutionx:
        HalfRange=Resolutionx
    uncertaintyxraw.append(HalfRange)
    xraw.append(ave)

# looping over the length of y and calculation 2h and t^2 so that the correct thing is plotted against each other
y=[]
x=[]
uncertaintyx=[]
uncertaintyy=[]
for i in range (len(xraw)):
    y.append(2*yraw[i])
    x.append(xraw[i]*xraw[i])
    uncertaintyx.append(2.0*(uncertaintyxraw[i]/xraw[i])*x[i])
    uncertaintyy.append(Resolutiony*2)

# merging the varaiables into one list of lists so tht tabulate can cope
merged_list = [(y[i],*Xmultraw[i],xraw[i],uncertaintyxraw[i],x[i],uncertaintyx[i]) for i in range (0,len(x))]
labels=['2 x h (m)','Time 1 (s)','Time 2 (s)','Time 3 (s)','Mean t (s)', 'unc t (s)', 'Mean t^2 (s)','unc t^2 (s^2)']
print(tabulate(merged_list, headers=labels,tablefmt='fancy_grid',floatfmt=".2f"))
\end{mintedbox}

Having run the code, it prints us a lovely table of results. Note that it displays the calculated means, some uncertainty and the calculated values of $t^2$. Note also that the table has headings with units for each column: 
\begin{figure}
    \centering
    \includegraphics[width=1\linewidth]{figures/resultstable.png}
\end{figure}
\end{fullwidth}
With the results taken and processed, it's time to do some analysis. In A-Level physics terms this means that we need to plot a graph, draw a line of best fit (a linear regression), and almost always - calculate a gradient. There are a few practicals where we draw non-linear relationships, or have no graph at all, but they are generally few and far between. 

For the results above, the perfect graph will look something like this\footnote{It's really important to realise - if you're reading this in Y12 - that youdon't need to learn all of this at once. We will learn this coding together}: 
\begin{figure}
    \centering
    \includegraphics[width=1\linewidth]{figures/graphper.png}
\end{figure}

\begin{fullwidth}
\begin{mintedbox}{python}
from scipy import stats
# Calling linear regression assigning the output to the variables on the left of the equation 
gradient, intercept,r,p,st_err = stats.linregress(x,y)
#looping over the length of x and applying y=mx+c using the m and c from the output above to find the regression line   
line=[]
for i in range (len(x)):
    line.append(gradient*x[i]+intercept)
# Plotting the data points and the best fit line
plt.scatter(x, y)
plt.errorbar(x,y,xerr=uncertaintyx,yerr=uncertaintyy,fmt='o',ecolor='cyan', color='black')
plt.plot(x,line)
plt.title('Best fit line using regression method')
plt.ylabel('2 x Height $2h$ (m)')
plt.xlabel('$Time^2$ $t^2$ (s)')

plt.show()
print ("Gradient=%.2f"%gradient,'ms-2')
print("Our y intercept should be zero it is equal to = ",intercept)
\end{mintedbox}
\end{fullwidth}
The output of this code is the above graph, but it will also print a gradient and intercept value: 
\begin{verbatim}
Gradient=9.18 ms-2
Our y intercept should be zero it is equal to =  0.29863797357735233  
\end{verbatim}

\section{Conclusions}
The final part of any practical write-up is the concluding statement. For this practical it would look like this:

\begin{fullwidth}
    \begin{mintedbox}{python}
print('Our calculated value of %.2f m/s^2 is close the value stated in our data booklet of 9.81 m/s^2' %gradient)
print('The line of best fit goes through all the error bars. This means we have verified the law and found a reasonable value of g.')
    \end{mintedbox}
\end{fullwidth}
We have done a few important things here:
\titem We have compared our value to a known value.
\titem We have quoted the source of information for the known value.
\titem We have concluded with a comment on the quality of our data and the implications of our experiment.

\newpage
\section{A complete example}
\includepdf[pages=-,nup=1x2]{fullprac.pdf}

\chapter{Markdown Guide}


You can use Markdown to format documentation you add to Markdown cells in your Jupyter notebook.

Here's how to format Markdown cells in Jupyter notebooks:

\subsection{Headings:}

Use the number sign (\#) followed by a blank space for notebook titles and section headings:
\begin{verbatim}    
# for titles
## for major headings
### for subheadings
#### for 4th level subheadings
\end{verbatim}

\subsection{Emphasis}
Use the following code to emphasize text:
\begin{verbatim}

Bold text: __string__ or **string**
Italic text: _string_ or *string*
\end{verbatim}

\section{Mathematical symbols}
Surround mathematical symbols with a dollar sign (\$), for example:
\begin{verbatim}
    $ mathematical symbols $
\end{verbatim}

\section{Monospace font}
Surround text with a grave accent (`) also called a back single quotation mark, for example:
\begin{verbatim}
    `string`
\end{verbatim}

You can use the monospace font for file paths, file names, message text that users see, or text that users enter.

\section{Line breaks}
Sometimes markdown doesn’t make line breaks when you want them. To force a linebreak, use the following code: \begin{verbatim}
    <br>
\end{verbatim}

\section{Indenting}
Use the greater than sign (>) followed by a space, for example:
\begin{verbatim}
> Text that will be indented when the Markdown is rendered. 
\end{verbatim}
Any subsequent text is indented until the next carriage return.

\section{Bullets}
To create a circular bullet point, use one of the following methods. Each bullet point must be on its own line.
A hyphen (-) followed by one or two spaces, for example: - Bulleted item
A space, a hyphen (-) and a space, for example: - Bulleted item
An asterisk (*) followed by one or two spaces, for example: * Bulleted item
To create a sub bullet, press Tab before entering the bullet point using one of the methods described above. For example:
- Main bullet point
     - Sub bullet point

\section{Numbered lists}
To create a numbered list, enter 1. followed by a space, for example:
\begin{verbatim}
1. Numbered item
1. Numbered item   
\end{verbatim}

For simplicity, you use 1. before each entry. The list will be numbered correctly when you run the cell.

To create a substep, press Tab before entering the numbered item, for example:
\begin{verbatim}
    1. Numbered item
     1. Substep
\end{verbatim}

\section{Colored note boxes}
Use one of the following <div> tags to display text in a colored box.
Restriction: Not all Markdown code displays correctly within <div> tags, so review your colored boxes carefully.\\
For example, to make a word bold, surround it with the HTML code for bold (<b>text</b> instead of the Markdown code.\\
The color of the box is determined by the alert type that you specify:
\subsection{Blue boxes (alert-info)}
\begin{verbatim}
   
<div class="alert alert-block alert-info">
<b>Tip:</b> Use blue boxes (alert-info) for tips and notes. 
If it’s a note, you don’t have to include the word “Note”.
</div> 
\end{verbatim}
\subsection{Yellow boxes (alert-warning)}
\begin{verbatim}
   
<div class="alert alert-block alert-warning">
<b>Example:</b> Use yellow boxes for examples that are not 
inside code cells, or use for mathematical formulas if needed.
</div> 
\end{verbatim}


\subsection{Green boxes (alert-success)}
\begin{verbatim}
  <div class="alert alert-block alert-success">
<b>Up to you:</b> Use green boxes sparingly, and only for some specific 
purpose that the other boxes can't cover. For example, if you have a lot 
of related content to link to, maybe you decide to use green boxes for 
related links from each section of a notebook.
</div>  
\end{verbatim}

\subsection{Red boxes (alert-danger)}
\begin{verbatim}
    
<div class="alert alert-block alert-danger">
<b>Just don't:</b> In general, avoid the red boxes. These should only be
used for actions that might cause data loss or another major issue.
</div>
\end{verbatim}

\section{Graphics}
You can attach image files directly to a notebook in Markdown cells by dragging and dropping it into the cell.
To add images to other cell types, use graphics that are hosted on the web with this code, substituting url/name with the full URL and name of the image:
\begin{verbatim}
<img src=“url/filename.gif” alt=“Alt text” title=“Title text” />
\end{verbatim}
Restriction: You cannot add captions to graphics.
\section{Geometric shapes}
Use \&\# followed by the decimal or hex reference number for the shape, for example:
\begin{verbatim}
 &#reference_number   
\end{verbatim}

For a list of reference numbers, see UTF-8 Geometric shapes.

\subsection{Horizontal lines}
On a new line, enter three asterisks:
\begin{verbatim}
    ***
\end{verbatim}

Internal links
To link to a section within your notebook, use the following code:
\begin{verbatim}
   [Section title](#section-title) 
\end{verbatim}

For the text inside the parentheses, replace any spaces and special characters with a hyphen. For example, if your section is called Analyzing customer purchasing habits, you'd enter:
\begin{verbatim}
  [Analyzing customer purchasing habits](#analyzing-customer-purchasing-habits)  
\end{verbatim}


Alternatively, you can add an ID above the section:
\begin{verbatim}
    <a id="section_ID"></a>
\end{verbatim}

Important: Each ID in the notebook must be unique.
To link to a section that has an ID, use the following code:
\begin{verbatim}
  [Section title](#section_ID)  
\end{verbatim}


Important: Test all internal links to ensure that they work.
External links
To link to an external site, use the following code:
\begin{verbatim}
    __[link text](http://url)__
\end{verbatim}

Surround the link with two underscores on each side


 \section{Practical Guide}

The following section was origionally written to cover the practical content of the AS and the A2 levels - though we don't split the content into AS and A2 any longer I have left the division as indicative of the Year 12 vs the Y13 content.

\includepdf[pages=3-]{prac.pdf}

%\chapter{Ticklists}
\section{Newtonian Physics}

Written examination: 2 hours 15 minutes
31.25\% of qualification
\subsection{Basic Physics}
Learners should be able to demonstrate and apply their knowledge and
understanding of:
\begin{itemize}
	\item[\Large{$\Square$}](a) the 6 essential base SI units (\sq kg, \sq m, \sq s, \sq A, \sq mol, \sq K)
	\item[\Large{$\Square$}](b) representing units in terms of the 6 base SI units and their prefixes
	\item[\Large{$\Square$}](c) checking equations for homogeneity using units
	\item[\Large{$\Square$}](d) the difference between scalar and vector quantities and to \sq give examples of each – displacement, velocity, acceleration, force, speed, time, density,
	pressure etc
	\item[\Large{$\Square$}](e) the addition and subtraction of coplanar vectors, and \sq perform mathematical calculations limited to \textbf{two} perpendicular vectors
	\item[\Large{$\Square$}](f) how to resolve a vector into two perpendicular components
	\item[\Large{$\Square$}](g) the concept of density and \sq how to use the equation \(\rho=\frac{m}{V}\) to calculate mass, density and volume
	\item[\Large{$\Square$}](h) what is meant by the turning effect of a force
	\item[\Large{$\Square$}](i) the use of the principle of moments
	\item[\Large{$\Square$}](j) the use of centre of gravity, for example in problems including stability:
	identify its position in a \sq cylinder, \sq sphere and \sq cuboid (beam) of uniform density
	\item[\Large{$\Square$}](k) when a body is in equilibrium the resultant force is zero and the net moment
	is zero, and \sq be able to perform simple calculations
	\subsection*{SPECIFIED PRACTICAL WORK}
	\item[\Large{$\Square$}] Measurement of the density of solids
	\item[\Large{$\Square$}] Determination of unknown masses by using the principle of moments
\end{itemize}
\subsection{Kinematics}
Learners should be able to demonstrate and apply their knowledge and
understanding of:
\begin{itemize}
	\item[\Large{$\Square$}](a) what is meant by \sq displacement, mean and instantaneous values of \sq speed, \sq velocity and \sq acceleration
	\item[\Large{$\Square$}](b) the representation of \sq displacement, \sq speed, \sq velocity and \sq acceleration by graphical methods
	\item[\Large{$\Square$}](c) the properties of \sq displacement-time graphs, \sq velocity-time graphs, and \sq interpret speed and displacement-time graphs for non-uniform acceleration
	\item[\Large{$\Square$}](d) how to derive and use equations which represent uniformly accelerated
	motion in a straight line
	\item[\Large{$\Square$}](e) how to describe the motion of bodies falling in a gravitational field \sq with and \sq without air resistance - terminal velocity
	\item[\Large{$\Square$}](f) the independence of vertical and horizontal motion of a body moving freely under gravity
	\item[\Large{$\Square$}](g) the explanation of the motion due to a uniform velocity in one direction and uniform acceleration in a perpendicular direction, and \sq perform simple
	calculations
	\subsection*{SPECIFIED PRACTICAL WORK}
	\item[\Large{$\Square$}] Measurement of \textit{g} by freefall
\end{itemize}
\subsection{Dynamics}
Learners should be able to demonstrate and apply their knowledge and
understanding of:
\begin{itemize}
	\item[\Large{$\Square$}](a) the concept of force and Newton's \(3^{\text{rd}}\) law of motion
	\item[\Large{$\Square$}](b) how free body diagrams can be used to represent forces on a particle or body
	\item[\Large{$\Square$}](c) the use of the relationship \(\Sigma F = ma\) in situations where mass is constant
	\item[\Large{$\Square$}](d) the idea that linear momentum is the product of mass and velocity
	\item[\Large{$\Square$}](e) the concept that force is the rate of change of momentum, applying this in situations where mass is constant
	\item[\Large{$\Square$}](f) the principle of conservation of momentum and use it to solve problems in one dimension involving \sq elastic collisions (where there is no loss of kinetic
	energy) and \sq inelastic collisions (where there is a loss of kinetic energy)
	\subsection*{SPECIFIED PRACTICAL WORK}
	\item[\Large{$\Square$}]Investigation of Newton’s \(2^{\text{nd}}\) law
\end{itemize}
\subsection{Energy Concepts}
Learners should be able to demonstrate and apply their knowledge and
understanding of:
\begin{itemize}
	\item[\Large{$\Square$}](a) the idea that work is the product of a force and distance moved in the
	direction of the force when the force is constant
	\item[\Large{$\Square$}](b) the calculation of the work done for constant forces, when the force is not along the line of motion (\(\text{work done} = Fx \cos \theta \))
	\item[\Large{$\Square$}](c) the principle of conservation of energy including knowledge of \sq gravitational potential energy (\(mg \Delta h\)) , \sq elastic potential energy (\(\frac{1}{2}kx^{2}\))  and \sq kinetic energy (\(\frac{1}{2}mv^{2}\))
	\item[\Large{$\Square$}](d) the work-energy relationship:\( Fx=\frac{1}{2}mv^{2}-\frac{1}{2}mv^{2}\)
	\item[\Large{$\Square$}](e) power being the rate of energy transfer
	\item[\Large{$\Square$}](f) dissipative forces for example, friction and drag cause energy to be
	transferred from a system and reduce the overall efficiency of the system
	\item[\Large{$\Square$}](g) the equation: $\text{efficiency} = \frac{\text{useful energy transfer}}{\text{total energy input}} \times 100\% $
	
\end{itemize}
\subsection{Circular Motion}
Learners should be able to demonstrate and apply their knowledge and
understanding of:
\begin{itemize}
	\item[\Large{$\Square$}](a) the terms \sq period of rotation, \sq frequency
	\item[\Large{$\Square$}](b) the definition of the unit radian as a measure of angle
	\item[\Large{$\Square$}](c) the use of the radian as a measure of angle
	\item[\Large{$\Square$}](d) the definition of angular velocity, $\omega$, for an object performing \sq circular motion and performing \sq simple harmonic motion
	\item[\Large{$\Square$}](e) the idea that the centripetal force is the resultant force acting on a body moving at constant speed in a circle
	\item[\Large{$\Square$}](f) the centripetal force and acceleration are directed towards the centre of the circular motion
	\item[\Large{$\Square$}](g) the use of the following equations relating to circular motion:
	\[\begin{aligned}
	\text{\sq} v=\omega r, &\text{\sq} a=\omega^{2} r, & \text{\sq} a=\frac{v^{2}}{r},& \text{\sq} F=\frac{mv^{2}}{r},&\text{\sq } F=m\omega^{2}r
	\end{aligned}
	\]
	
	
\end{itemize}
\subsection{Vibrations}
Learners should be able to demonstrate and apply their knowledge and
understanding of:
\begin{itemize}
	\item[\Large{$\Square$}](a) the definition of simple harmonic motion as a statement in words
	\item[\Large{$\Square$}](b) \(a=-\omega^{2} x\) as a mathematical defining equation of simple harmonic motion.
	\item[\Large{$\Square$}](c) the graphical representation of the variation of acceleration with displacement during simple harmonic motion
	\item[\Large{$\Square$}](d) \(x=A\cos(\omega t + \epsilon) \) as a solution to \(a=-\omega^{2} x\)
	\item[\Large{$\Square$}](e) the terms \sq frequency, \sq period, \sq amplitude and \sq phase
	\item[\Large{$\Square$}](f) period as \(\frac{1}{f}\) or \(\frac{2\pi}{\omega}\)
	\item[\Large{$\Square$}](g)  \(v=A\omega\sin(\omega t + \epsilon)\) for the velocity during simple harmonic motion
	\item[\Large{$\Square$}](h) the graphical representation of the changes in \sq displacement and \sq velocity with time during simple harmonic motion
	\item[\Large{$\Square$}](i) the equation \(T=2\pi \sqrt{\frac{m}{k}}\) for the period of a system having stiffness (force per unit extension) k and mass m
	\item[\Large{$\Square$}](j) the equation \(T=2\pi \sqrt{\frac{l}{g}}\) for the period of a simple pendulum
	\item[\Large{$\Square$}](k) the graphical representation of the interchange between kinetic energy and potential energy during undamped simple harmonic motion, and \sq perform simple calculations on energy changes
	\item[\Large{$\Square$}](l) free oscillations and the effect of damping in real systems
	\item[\Large{$\Square$}](m) practical examples of damped oscillations
	\item[\Large{$\Square$}](n) the importance of critical damping in appropriate cases such as vehicle
	suspensions
	\item[\Large{$\Square$}](o) forced oscillations and resonance, and to \sq describe practical examples
	\item[\Large{$\Square$}](p) the variation of the amplitude of a forced oscillation with driving frequency and \sq that increased damping broadens the resonance curve
	\item[\Large{$\Square$}](q) circumstances when resonance is useful for example, circuit tuning,
	microwave cooking and other circumstances in which it should be avoided for example, bridge design
	\subsection*{SPECIFIED PRACTICAL WORK}
	\item[\Large{$\Square$}] Measurement of g with a pendulum
	\item[\Large{$\Square$}] Investigation of the damping of a spring
\end{itemize}
\subsection{Kinetic Theory}
Learners should be able to demonstrate and apply their knowledge and
understanding of:
\begin{itemize}
	\item[\Large{$\Square$}](a) the equation of state for an ideal gas expressed as $pV = nRT$ where R is the 	molar gas constant and $pV = NkT$ where k is the Boltzmann constant
	\item[\Large{$\Square$}](b) the assumptions of the kinetic theory of gases which includes the random distribution of energy among the molecules
	\item[\Large{$\Square$}]	(c) the idea that molecular movement causes the pressure exerted by a gas, and \sq the use of \(p=\frac{1}{3}\rho \bar{c^{2}} = \frac{1}{3}\frac{N}{V} m\bar{c^{2}}\) where N is the number of molecules
	\item[\Large{$\Square$}]	(d) the definition of Avogadro constant $N_A$ and hence the mole
	\item[\Large{$\Square$}]	(e) the idea that the molar mass $M$ is related to the relative molecular mass $M_r$ by \(M/kg=\frac{M_{r}}{1000}\), and that \sq the number of moles $n$ is given by $	\frac{\text{total mass}}{\text{molar mass}}$
	\item[\Large{$\Square$}]	(f) how to combine \(pV=\frac{1}{3}Nm\bar{C^{2}} \) with \(pV = nRT\) and show that the total translational kinetic energy of a mole of a monatomic gas is given by \(\frac{3}{2}RT\) and \sq the mean kinetic energy of a molecule is \(\frac{3}{2}kT\) where \(k=\frac{R}{N_{A}}\) is the Boltzmann constant, and that T is proportional to the mean kinetic energy
\end{itemize}
\subsection{Thermal Physics}
Learners should be able to demonstrate and apply their knowledge and
understanding of:
\begin{itemize}
	\item[\Large{$\Square$}] (a) the idea that the internal energy of a system is the sum of the potential and kinetic energies of its molecules
	\item[\Large{$\Square$}] (b) absolute zero being the temperature of a system when it has minimum
	internal energy
	\item[\Large{$\Square$}]	(c) the internal energy of an ideal monatomic gas being wholly kinetic so it is given by \(U=\frac{3}{2}nRT\)
	\item[\Large{$\Square$}]	(d) the idea that heat enters or leaves a system through its boundary or container wall, according to whether the system's temperature is lower or higher than
	that of its surroundings, so heat is energy in transit and not contained within
	the system
	\item[\Large{$\Square$}]	(e) the idea that if no heat flows between systems in contact, then they are said to be in thermal equilibrium, and are at the same temperature
	\item[\Large{$\Square$}]	(f) the idea that energy can enter or leave a system by means of work, so work is also energy in transit
	\item[\Large{$\Square$}]	(g) the equation $W = p \Delta V $ can be used to calculate the work done by a gas under constant pressure
	\item[\Large{$\Square$}]	(h) the idea that even if $p$ changes, $W$ is given by the area under the $p – V$ graph
	\item[\Large{$\Square$}]	(i) the use of the first law of thermodynamics, in the form $\Delta U=Q-W $ and \sq know how to interpret negative values of $\Delta U $, $Q$, and $W$
	\item[\Large{$\Square$}]	(j) the idea that for a solid (or liquid), $W$ is usually negligible, so $Q = \Delta U$
	\item[\Large{$\Square$}]	(k) \(Q=mc\Delta \theta\) , for a solid or liquid, and this is the defining equation for specific	heat capacity, $ c $
	\subsection*{SPECIFIED PRACTICAL WORK}
	\item[\Large{$\Square$}] Estimation of absolute zero by use of the gas laws
	\item[\Large{$\Square$}] Measurement of the specific heat capacity for a solid
\end{itemize}
\section{Electricity and the Universe}
Written examination: 2 hours
31.25\% of qualification
\subsection{Conduction of Electricity}Learners should be able to demonstrate and apply their knowledge and
understanding of:
\begin{itemize}
	\item[\Large{$\Square$}]	(a) the fact that the unit of charge is the coulomb (C), and \sq an electron's charge, $e$,	is a very small fraction of a coulomb
	\item[\Large{$\Square$}]	(b) the fact that charge can flow through certain materials, called conductors
	\item[\Large{$\Square$}]	(c) electric current being the rate of flow of charge
	\item[\Large{$\Square$}]	(d) the use of the equation \(I=\frac{\Delta Q}{\Delta t}\)
	\item[\Large{$\Square$}]	(e) current being measured in ampères (A), where $A = C s^{-1}$
	\item[\Large{$\Square$}]	(f) the mechanism of conduction in metals as the drift of free electrons
	\item[\Large{$\Square$}]	(g) the derivation and use of the equation $I = nAve$ for free electrons
\end{itemize}
\subsection{Resistance}Learners should be able to demonstrate and apply their knowledge and
understanding of:
\begin{itemize}
	\item[\Large{$\Square$}] (a) the definition of potential difference
	\item[\Large{$\Square$}]	(b) the idea that potential difference is measured in volts (V) where $V = J C^{-1}$
	\item[\Large{$\Square$}]	(c) the characteristics of I – V graphs for \sq the filament of a lamp, and \sq a metal wire at constant temperature
	\item[\Large{$\Square$}]	(d) Ohm's law, the equation \(V = IR\) and the definition of resistance
	\item[\Large{$\Square$}]	(e) resistance being measured in ohms ($\Omega$), where $\Omega = V A^{-1}$
	\item[\Large{$\Square$}]	(f) the application of \(P=IV=I^{2}R=\frac{V^{2}}{R}\)
	\item[\Large{$\Square$}]	(g) collisions between free electrons and ions gives rise to electrical resistance, and \sq electrical resistance increases with temperature
	\item[\Large{$\Square$}]	(h) the application of \(R=\frac{\rho l}{A}\) , the equation for resistivity
	\item[\Large{$\Square$}]	(i) the idea that the resistance of metals varies almost linearly with temperature	over a wide range
	\item[\Large{$\Square$}]	(j) the idea that ordinarily, collisions between free electrons and ions in metals increase the random vibration energy of the ions, so the temperature of the
	metal increases
	\item[\Large{$\Square$}]	(k) what is meant by \sq superconductivity, and \sq superconducting transition temperature
	\item[\Large{$\Square$}]	(l) the fact that most metals show superconductivity, and have transition temperatures a few degrees above absolute zero ($–273 C$)
	\item[\Large{$\Square$}]	(m) certain materials (high temperature superconductors) having transition temperatures above the boiling point of nitrogen ($–196 C$)
	\item[\Large{$\Square$}]	(n) some uses of superconductors for example, MRI scanners and particle
	accelerators
	\subsection*{SPECIFIED PRACTICAL WORK}
	\item[\Large{$\Square$}]Investigation of the I-V characteristics of the filament of a lamp and a metal wire at constant temperature
	\item[\Large{$\Square$}] Determination of resistivity of a metal
	\item[\Large{$\Square$}]Investigation of the variation of resistance with temperature for a metal wire
\end{itemize}
\subsection{D.C. Circuits}Learners should be able to demonstrate and apply their knowledge and
understanding of:
\begin{itemize}
	\item[\Large{$\Square$}](a) the idea that the current from a source is equal to the sum of the currents in	the separate branches of a parallel circuit, and that \sq this is a consequence of
	conservation of charge
	\item[\Large{$\Square$}]	(b) the sum of the potential differences across components in a series circuit is	equal to the potential difference across the supply, and that \sq this is a	consequence of conservation of energy
	\item[\Large{$\Square$}]	(c) potential differences across components in parallel are equal
	\item[\Large{$\Square$}]	(d) the application of equations for the combined resistance of resistors in \sq series	and \sq parallel
	\item[\Large{$\Square$}]	(e) the use of a potential divider in circuits (including circuits which contain \sq LDRs and \sq thermistors)
	\item[\Large{$\Square$}]	(f) what is meant by the emf of a source
	\item[\Large{$\Square$}]	(g) the unit of emf is the volt (V), which is the same as that of potential difference
	\item[\Large{$\Square$}]	(h) the idea that sources have internal resistance and to use the equation \(V=E-Ir\)
	\item[\Large{$\Square$}]	(i) how to calculate current and potential difference in a circuit containing one cell or cells in series
	\subsection*{SPECIFIED PRACTICAL WORK}
	\item[\Large{$\Square$}] Determination of the internal resistance of a cell
\end{itemize}
\subsection{Capacitance}Learners should be able to demonstrate and apply their knowledge and
understanding of:
\begin{itemize}
	\item[\Large{$\Square$}] (a) the idea that a simple parallel plate capacitor consists of a pair of equal parallel metal plates separated by a vacuum or air
	\item[\Large{$\Square$}]		(b) a capacitor storing energy by transferring charge from one plate to the other,	so that the plates carry equal but opposite charges (the net charge being zero)
	\item[\Large{$\Square$}]		(c) the definition of capacitance as \(C=\frac{Q}{V}\)
	\item[\Large{$\Square$}]		(d) the use of \(C=\frac{\epsilon_{0} A}{d}\) for a parallel plate capacitor, with no dielectric
	\item[\Large{$\Square$}]		(e) the idea that a dielectric increases the capacitance of a vacuum-spaced
	capacitor
	\item[\Large{$\Square$}]		(f) the $E$ field within a parallel plate capacitor being uniform and the use of the
	equation \(E=\frac{V}{d}\)
	\item[\Large{$\Square$}]	(g) the equation \(U=\frac{1}{2}QV \) for the energy stored in a capacitor
	\item[\Large{$\Square$}]		(h) the equations for capacitors in \sq series and in \sq parallel
	\item[\Large{$\Square$}]		(i) the process by which a capacitor charges and discharges through a resistor
	\item[\Large{$\Square$}]		(j) the equations: \sq \(Q=Q_{0} \left( 1-e^{-\frac{t}{RC}}\right) \) and \sq \(Q=Q_{0}e^{-\frac{t}{RC}} \) where RC is the time
	constant
	\subsection*{SPECIFIED PRACTICAL WORK}
	\item[\Large{$\Square$}] Investigation of the charging and discharging of a capacitor to determine the	time constant
	\item[\Large{$\Square$}]Investigation of the energy stored in a capacitor
\end{itemize}
\subsection{Solids Under Stress}Learners should be able to demonstrate and apply their knowledge and
understanding of:
\begin{itemize}
	\item[\Large{$\Square$}] (a) Hooke’s law and use \(F=kx\) where the spring constant k is the force per unit
	extension
	\item[\Large{$\Square$}]	(b) the ideas that for materials the tensile stress, \sq \(\sigma=\frac{F}{A}\) and the tensile strain, \sq \(\epsilon= \frac{\Delta l}{l}  \)
	and the Young modulus, \sq \(E=\frac{\sigma}{\epsilon}  \) when Hooke’s law applies
	\item[\Large{$\Square$}]		(c) the work done in deforming a solid being equal to the area under a force-extension graph, which is \(\frac{1}{2} Fx\) if Hooke’s law is obeyed
	\item[\Large{$\Square$}]		(d) the classification of solids as \sq crystalline, \sq amorphous (to include glasses and ceramics) and \sq polymeric
	\item[\Large{$\Square$}]	(e) the features of a force-extension (or stress-strain) graph for a metal such as	copper, to include
	\begin{itemize}
		\item[\Large{$\Square$}] elastic and \sq plastic strain
		\item[\Large{$\Square$}] the effects of dislocations, and the strengthening of metals by introducing barriers to dislocation movement, such as foreign atoms, other dislocations, and more grain boundaries
		\item[\Large{$\Square$}] necking and ductile fracture
	\end{itemize}
	\item[\Large{$\Square$}](f) the features of a force-extension (or stress-strain) graph for a brittle material such as glass, to include:
	\begin{itemize}
		\item[\Large{$\Square$}] elastic strain and obeying Hooke’s law up to fracture
		\item[\Large{$\Square$}] brittle fracture by crack propagation, the effect of surface imperfections on breaking stress, and \sq how breaking stress can be increased by reducing surface imperfections (as in thin fibres) or by \sq putting surface under compression (as in toughened glass or pre-stressed	concrete)
	\end{itemize}
	\item[\Large{$\Square$}](g) the features of a force-extension (or stress-strain) graph for rubber, to include
	\begin{itemize}
		\item Hooke’s law only approximately obeyed, low Young modulus and the \sq extension due to straightening of chain molecules against thermal opposition
		\item hysteresis
	\end{itemize}	
	\subsection*{SPECIFIED PRACTICAL WORK}
	\item[\Large{$\Square$}] Determination of Young modulus of a metal in the form of a wire
	\item[\Large{$\Square$}] Investigation of the force-extension relationship for rubber
\end{itemize}
\subsection{Electrostatic and Gravitational Fields of Force}Learners should be able to demonstrate and apply their knowledge and
understanding of:
\begin{itemize}
	\item[\Large{$\Square$}] (a) the features of electric and gravitational fields as specified in figure~\ref{fieldtab}
	\begin{figure}
		\centering
		\label{fieldtab}
		\begin{tabular}{p{0.4\textwidth} p{0.4\textwidth}}
			\toprule
			\textbf{Electric Fields}                                                 & \textbf{Gravitational fields}    \\ \midrule
			\begin{Large}$\Square$\end{Large} Electric field strength, $E$, is the force per unit charge on a small positive test charge placed at the point     & \begin{Large}$\Square$\end{Large} Gravitational field strength, g, is the force per
			unit mass on a small test mass placed at the
			point                                        \\ \midrule
			\begin{Large}$\Square$\end{Large} Inverse square law for the force between two electric charges in the form: \[F=\frac{1}{4 \pi \epsilon_{0}}\frac{Q_{1}Q_{2}}{r^2}\]
			(Coulomb's law) & \begin{Large}$\Square$\end{Large} Inverse square law for the force between two masses in the form: \[F=G\frac{M_{1}M_{2}}{r^2}\]
			(Newton's law of gravitation) \\ \midrule
			\begin{Large}$\Square$\end{Large} $F$ can be attractive or repulsive       & \begin{Large}$\Square$\end{Large} $F$ is attractive only                                                                                                                      \\ \midrule
			\begin{Large}$\Square$\end{Large} \[E=\frac{1}{4 \pi \epsilon_{0}}\frac{Q}{r^2}\]  for the field strength due to a point charge in free space or air     &\begin{Large}$\Square$\end{Large} \[g=\frac{GM}{r^2}\] for the field strength due to a point mass                                                                             \\ \midrule
			\begin{Large}$\Square$\end{Large} Potential at a point due to a point charge in terms of the work done in bringing a unit positive charge from infinity to that point & \begin{Large}$\Square$\end{Large} Potential at a point due to a point charge in terms of the work done in bringing a unit positive charge from infinity to that point \\ \midrule
			\begin{Large}$\Square$\end{Large} \[V_{E}=\frac{1}{4 \pi \epsilon_{0}}\frac{Q}{r}\] and \[PE=\frac{1}{4 \pi \epsilon_{0}}\frac{Q_{1}Q_{2}}{r}\] & \begin{Large}$\Square$\end{Large} \[V_{g}=-\frac{GM}{r}\] and \[PE=-\frac{GM_{1}M_{2}}{r}\] \\ \midrule
			\begin{Large}$\Square$\end{Large} Change in potential energy of
			a point charge moving in any electric field
			$= q \Delta V_{E}$  & \begin{Large}$\Square$\end{Large} Change in potential energy of a point mass moving in any gravitational field $= m \Delta V_{g}$\\ \midrule 
			\begin{Large}$\Square$\end{Large} Field strength at a point is
			given by
			$E = - \text{slope of the }V_{E}-r$ graph at that point & \begin{Large}$\Square$\end{Large} Field strength at a point is given by $g = - \text{ slope of the }V_{g}-r$ graph at that point \\ \bottomrule
		\end{tabular}
		\caption{Fields Table}
	\end{figure}
	\item[\Large{$\Square$}]		(b) the idea that the gravitational field outside spherical bodies such as the Earth is essentially the same as if the whole mass were concentrated at the centre
	\item[\Large{$\Square$}]		(c) field lines (or lines of force) giving the direction of the field at a point, thus, for	a positive point charge, the field lines are radially outward
	\item[\Large{$\Square$}]		(d) equipotential surfaces joining points of equal potential and are therefore spherical for a point charge
	\item[\Large{$\Square$}]		(e) how to calculate the net potential and resultant field strength for a number of point charges or point masses
	\item[\Large{$\Square$}]		(f) the equation $\Delta U_{P} = mg\Delta h$ for distances over which the variation of g is	negligible
\end{itemize}
\subsection{Using Radiation to Investigate Stars}Learners should be able to demonstrate and apply their knowledge and
understanding of:
\begin{itemize}
	\item[\Large{$\Square$}] (a) the idea that the stellar spectrum consists of a continuous emission spectrum, from the dense gas of the surface of the star, and a line absorption spectrum arising from the passage of the emitted electromagnetic radiation through the tenuous atmosphere of the star
	\item[\Large{$\Square$}]		(b) the idea that bodies which absorb all incident radiation are known as black bodies and that \sq stars are very good approximations to black bodies
	\item[\Large{$\Square$}]		(c) the shape of the black body spectrum and that the peak wavelength is inversely proportional to the absolute temperature (defined by:
	\(T (K) = \theta (C) + 273.15)\)
	\item[\Large{$\Square$}]		(d) Wien's displacement law, \sq Stefan's law and \sq the inverse square law to investigate the properties of stars – \sq luminosity, \sq size, \sq temperature and
	\sq distance [N.B. stellar brightness in magnitudes will not be required]
	\item[\Large{$\Square$}]		(e) the meaning of multiwavelength astronomy and that by studying a region of	space at different wavelengths (different photon energies) the different processes which took place there can be revealed
\end{itemize}
\subsection{Obits and the Wider Universe}Learners should be able to demonstrate and apply their knowledge and
understanding of:
\begin{itemize}
	\item[\Large{$\Square$}] (a) Kepler's three laws of planetary motion: \sq 1\sq  2\sq  3 
	\item[\Large{$\Square$}]		(b) Newton's law of gravitation \( F=G\frac{M_{1}M_{2}}{r^2} \) in simple examples, including the	motion of planets and satellites
	\item[\Large{$\Square$}]		(c) how to derive Kepler's 3rd law, for the case of a circular orbit from Newton's law of gravity and the formula for centripetal acceleration
	\item[\Large{$\Square$}]		(d) how to use data on orbital motion, such as period or orbital speed, to calculate the mass of the central object
	\item[\Large{$\Square$}]		(e) how the orbital speeds of objects in spiral galaxies implies the existence of dark matter
	\item[\Large{$\Square$}]		(f) how the recently discovered Higgs boson may be related to dark matter
	\item[\Large{$\Square$}]		(g) how to determine the position of the centre of mass of two spherically	symmetric objects, given their masses and separation, and \sq calculate their mutual orbital period in the case of circular orbits
	\item[\Large{$\Square$}]		(h) the Doppler relationship in the form \( \frac{\Delta \lambda}{\lambda}=\frac{v}{c} \)
	\item[\Large{$\Square$}]		(i) how to determine a star's radial velocity (i.e. the component of its velocity along the line joining it and an observer on the Earth) from data about the Doppler shift of spectral lines
	\item[\Large{$\Square$}]		(j) the use of data on the variation of the radial velocities of the bodies in a	double system (for example, a star and orbiting exo-planet) and their orbital period to determine the masses of the bodies for the case of a circular orbit edge on as viewed from the Earth
	\item[\Large{$\Square$}]		(k) how the Hubble constant ($H_0$) relates galactic radial velocity ($v$) to distance ($D$) and it is defined by $v= H_{0}D$
	\item[\Large{$\Square$}]		(l) why \(\frac{1}{H_0}\) approximates the age of the universe
	\item[\Large{$\Square$}] (m) how the equation \( \rho_{c} = \frac{3H_{0}^{2}}{8\pi G} \) for the critical density of a 'flat' universe can be derived very simply using conservation of energy
\end{itemize}
\section{Light, Nuclei and Options}
Written examination: 2 hours 15 minutes
37.5\% of qualification
\subsection{The Nature of Waves}Learners should be able to demonstrate and apply their knowledge and
understanding of:
\begin{itemize}
	\item[\Large{$\Square$}] (a) the idea that a progressive wave transfers energy without any transfer of
	matter
	\item[\Large{$\Square$}]		(b) the difference between transverse and longitudinal waves
	\item[\Large{$\Square$}]		(c) the term polarisation
	\item[\Large{$\Square$}]		(d) the terms ``in phase'' and ``in antiphase''
	\item[\Large{$\Square$}]		(e) the terms \sq displacement, \sq amplitude, \sq wavelength, \sq frequency, \sq period and \sq velocity of a wave
	\item[\Large{$\Square$}]		(f) graphs of \sq displacement against time, and \sq displacement against position for transverse waves only
	\item[\Large{$\Square$}]		(g) the equation $c = f\lambda$
	\item[\Large{$\Square$}]		(h) the idea that all points on wavefronts oscillate in phase, and that \sq wave propagation directions (rays) are at right angles to wavefronts
	\subsection*{SPECIFIED PRACTICAL WORK}
	\item[\Large{$\Square$}] Measurement of the intensity variations for polarisation
\end{itemize}
\subsection{Wave Properties}Learners should be able to demonstrate and apply their knowledge and
understanding of:
\begin{itemize}
	\item[\Large{$\Square$}] (a) diffraction occuring when waves encounter slits or obstacles
	\item[\Large{$\Square$}]	(b) the idea that there is little diffraction when $\lambda$ is much smaller than the dimensions of the obstacle or slit
	\item[\Large{$\Square$}]	(c) the idea that if $\lambda$ is equal to or greater than the width of a slit, waves spread	as roughly semicircular wavefronts, but if $\lambda$ is less than the slit width the main	beam spreads through less than 180C
	\item[\Large{$\Square$}]	(d) how two source interference occurs
	\item[\Large{$\Square$}]	(e) the historical importance of Young’s experiment
	\item[\Large{$\Square$}]	(f) the principle of superposition, giving appropriate sketch graphs
	\item[\Large{$\Square$}]	(g) the path difference rules for constructive and destructive interference between waves from in phase sources
	\item[\Large{$\Square$}]	(h) the use of \( \lambda= \frac{a \Delta y}{D} \)
	\item[\Large{$\Square$}]	(i) the derivation and use of $d sin \theta = n \lambda$ for a diffraction grating
	\item[\Large{$\Square$}]	(j) the idea that for a diffraction grating a very small $d$ makes beams (“orders”) much further apart than in Young’s experiment, and that the large number of
	slits makes the bright beams much sharper
	\item[\Large{$\Square$}]	(k) the idea that coherent sources are monochromatic with wavefronts
	continuous across the width of the beam and, (when comparing more than one source) with a constant phase relationship
	\item[\Large{$\Square$}]	(l) examples of \sq coherent and \sq incoherent sources
	\item[\Large{$\Square$}]	(m) the idea that for two source interference to be observed, the sources must have a zero or constant phase difference and have oscillations in the same
	direction
	\item[\Large{$\Square$}]	(n) the differences between stationary and progressive waves
	\item[\Large{$\Square$}]	(o) the idea that a stationary wave can be regarded as a superposition of two	progressive waves of equal amplitude and frequency, travelling in opposite directions, and that \sq the internodal distance is $ \frac{\lambda}{2}$
	\subsection*{SPECIFIED PRACTICAL WORK}
	\item[\Large{$\Square$}] Determination of wavelength using Young’s double slits
	\item[\Large{$\Square$}] Determination of wavelength using a diffraction grating
	\item[\Large{$\Square$}] Determination of the speed of sound using stationary waves
\end{itemize}
\subsection{Refraction of Light}Learners should be able to demonstrate and apply their knowledge and
understanding of:
\begin{itemize}
	\item[\Large{$\Square$}] (a) the refractive index, $n$, of a medium being defined as \( \frac{c}{v} \), in which \(v\) is the speed of light in the medium and \(c\) is the speed of light in a vacuum
	\item[\Large{$\Square$}]	(b) the use of the equations: \sq \( n_{1}v_{1}=n_{2}v_{2} \) and \sq \( n_{1} \sin \theta_{1} = n_{2}\sin \theta_{2}  \) (regarded as	Snell’s law)
	\item[\Large{$\Square$}]	(c) how Snell's law relates to the wave model of light propagation and for diagrams of plane waves approaching a plane boundary obliquely, and being refracted
	\item[\Large{$\Square$}]	(d) the conditions for total internal reflection
	\item[\Large{$\Square$}]	(e) the derivation and use of the equation for the critical angle \( n_{1}\sin \theta_{c} = n_{2} \)
	\item[\Large{$\Square$}]	(f) how to apply the concept of total internal reflection to multimode optical fibres
	\item[\Large{$\Square$}]	(g) the problem of multimode dispersion with optical fibres in terms of limiting the	rate of data transfer and transmission distance
	\item[\Large{$\Square$}]	(h) how the introduction of monomode optical fibres has allowed for much greater transmission rates and distances
	\subsection*{SPECIFIED PRACTICAL WORK}
	\item[\Large{$\Square$}] Measurement of the refractive index of a material
\end{itemize}
\subsection{Photons}Learners should be able to demonstrate and apply their knowledge and
understanding of:
\begin{itemize}
	\item[\Large{$\Square$}] (a) the fact that light can be shown to consist of discrete packets (photons) of energy
	\item[\Large{$\Square$}]	(b) how the photoelectric effect can be demonstrated
	\item[\Large{$\Square$}]	(c) how a vacuum photocell can be used to measure the maximum kinetic
	energy, $E_{k\text{ max}}$, of emitted electrons in $eV$ and hence in \(J\)
	\item[\Large{$\Square$}]	(d) the graph of  $E_{k\text{ max}}$ against frequency of illuminating radiation
	\item[\Large{$\Square$}]	(e) how a photon picture of light leads to Einstein's equation,
	$E_{k\text{ max}}=hf-\phi$, and \sq how this equation correlates with the graph of  $E_{k\text{ max}}$ against frequency
	\item[\Large{$\Square$}]	(f) the fact that the visible spectrum runs approximately from 700 nm (red end) to 400 nm (violet end) and \sq the orders of magnitude of the wavelengths of the other named regions of the electromagnetic spectrum
	\item[\Large{$\Square$}]	(g) typical photon energies for these radiations
	\item[\Large{$\Square$}]	(h) how to produce line emission and line absorption spectra from atoms
	\item[\Large{$\Square$}]	(i) the appearance of such spectra as seen in a diffraction grating
	\item[\Large{$\Square$}]	(j) simple atomic energy level diagrams, together with the photon hypothesis, line emission and line absorption spectra
	\item[\Large{$\Square$}]	(k) how to determine ionisation energies from an energy level diagram
	\item[\Large{$\Square$}]	(l) the demonstration of electron diffraction and that particles have a wave-like	aspect
	\item[\Large{$\Square$}]	(m) the use of the relationship \( p=\frac{h}{\lambda} \) for both particles of matter and photons
	\item[\Large{$\Square$}]	(n) the calculation of radiation pressure on a surface absorbing or reflecting	photons
	\subsection*{SPECIFIED PRACTICAL WORK}
	\item[\Large{$\Square$}] Determination of h using LEDs
\end{itemize}
\subsection{Lasers}Learners should be able to demonstrate and apply their knowledge and
understanding of:
\begin{itemize}
	\item[\Large{$\Square$}] (a) the process of stimulated emission and \sq how this process leads to light emission that is coherent
	\item[\Large{$\Square$}]	(b) the idea that a population inversion (N2 > N1) is necessary for a laser to operate
	\item[\Large{$\Square$}]	(c) the idea that a population inversion is not (usually) possible with a 2-level energy system
	\item[\Large{$\Square$}]	(d) how a population inversion is attained in 3 and 4-level energy systems
	\item[\Large{$\Square$}]	(e) the process of pumping and its purpose
	\item[\Large{$\Square$}]	(f) the structure of a typical laser i.e. an amplifying medium between two mirrors, one of which partially transmits light
	\item[\Large{$\Square$}]	(g) the advantages and uses of a semiconductor laser i.e. small, cheap, far more	efficient than other types of laser, and it is used for CDs, DVDs, telecommunication etc
\end{itemize}
\subsection{Nuclear Decay}Learners should be able to demonstrate and apply their knowledge and
understanding of:
\begin{itemize}
	\item[\Large{$\Square$}] (a) the spontaneous nature of nuclear decay; the nature of $\alpha$, $\beta$ and $\gamma$ radiation, and	equations to represent the nuclear transformations using the $_{z}^{x}A $ notation
	\item[\Large{$\Square$}]	(b) different methods used to distinguish between $\alpha$, $\beta$ and $\gamma$ radiation and the connections between the nature, penetration and range for ionising particles
	\item[\Large{$\Square$}]	(c) how to make allowance for background radiation in experimental measurements
	\item[\Large{$\Square$}]	(d) the concept of the half-life, $T_{1/2}$
	\item[\Large{$\Square$}]	(e) the definition of the activity, \(A\), and the Becquerel
	\item[\Large{$\Square$}]	(f) the decay constant, \( \lambda \), and the equation \( A= \lambda N \)
	\item[\Large{$\Square$}]	(g) the exponential law of decay in graphical and algebraic form, 
	\begin{align*}
	\text{\sq }N=N_{0}e^{- \lambda t}  & &\text{and} & &\text{\sq }A=A_{0}e^{-\lambda t}& &\text{or \ldots}\\
	\text{\sq}N=\frac{N_{0}}{2^{x}} & &\text{and}& &\text{\sq }A=\frac{A_{0}}{2^{x}}
	\end{align*}
	where x is the number of half-lives elapsed - not necessarily an integer
	\item[\Large{$\Square$}]	(h) the derivation and use of \( \lambda=\frac{\ln 2}{T_{1/2}}  \)
	\subsection*{SPECIFIED PRACTICAL WORK}
	\item[\Large{$\Square$}] Investigation of radioactive decay - a dice analogy
	\item[\Large{$\Square$}] Investigation of the variation of intensity of gamma radiation with distance
\end{itemize}
\subsection{Particles and Nuclear Structure}Learners should be able to demonstrate and apply their knowledge and
understanding of:
\begin{itemize}
	\item[\Large{$\Square$}](a) the significance of the results of the Rutherford alpha particle scattering experiment
	\item[\Large{$\Square$}](b) how to approximate the maximum size of the Coulomb repulsion force
	between an alpha particle and a gold atom / nucleus for both the plum pudding model and the Rutherford model
	\item[\Large{$\Square$}](c) the idea that matter is composed of quarks and leptons and that there are three generations of quarks and leptons, although no questions will be set involving second or third generations 
	\begin{figure}
		\resizebox{\linewidth}{!}{%
			
			\begin{tabular}{c|c|c||c|c|}
				\cline{2-5}
				&\multicolumn{2}{c||}{Leptons}& \multicolumn{2}{c|}{Quarks}\\
				\hline
				\multicolumn{1}{|c|}{Particle (symbol)}& \sq electron ($e^-$)&\sq electron neutrino ($v_e$)&\sq up(u)& \sq down(d)\\
				\hline
				\hline
				\multicolumn{1}{|c|}{charge (e)}& -1& 0 & $+\frac{2}{3}$&$-\frac{1}{3}$\\
				\hline
			\end{tabular}
			
		} \caption{Table of Quarks}
	\end{figure}
	
	\item[\Large{$\Square$}](d) the idea that antiparticles exist for the particles given in figure 2, \sq that the properties of an antiparticle are identical to those of its corresponding particle apart from having opposite charge, and that particles and \sq antiparticles annihilate
	\item[\Large{$\Square$}](e) symbols for a \sq positron and for antiparticles of \sq quarks and \sq hadrons
	\item[\Large{$\Square$}](f) the idea that quarks and antiquarks are never observed in isolation, but are bound into composite particles called hadrons, or three types of baryon (combinations of 3 quarks), or antibaryons (combinations of 3 antiquarks) or mesons (quark-antiquark pairs)
	\item[\Large{$\Square$}](g) the quark compositions of the \sq neutron and \sq proton \item[\Large{$\Square$}](h) how to use data in the table in figure 2 to suggest the quark make-up of less well known first generation baryons and of charged pions
	\item[\Large{$\Square$}](i) the properties of the four forces or interactions experienced by particles as summarized in figure 3:
	\begin{figure}
		\resizebox{\linewidth}{!}{
			\begin{tabular}{|c|c|c|p{5cm}|}
				\hline
				Interaction &Experienced by & Range & Comments\\
				\hline
				\sq Gravitational & all particles & infinite&
				very weak – negligible
				except in the context of
				large objects such as
				planets and stars
				\\
				\hline
				\sq Weak& all particles&
				very short
				range&
				only significant in cases
				where the electromagnetic
				and strong interactions do
				not operate\\
				\hline
				\sq Electromagnetic&
				all charged
				particles&
				infinite&
				also experienced by neutral
				hadrons because they are
				composed of quarks\\
				\hline
				\sq Strong &quarks& short range&
				experienced by quarks and
				particles composed of
				quarks\\
				\hline
		\end{tabular}}
		\caption{The fundamental forces}
	\end{figure}
	
	\item[\Large{$\Square$}](j) how to apply conservation of \sq charge, \sq lepton number and \sq baryon number (or quark number) to given simple reactions
	\item[\Large{$\Square$}](k) the idea that neutrino involvement and quark flavour changes are exclusive to weak interactions
	
\end{itemize}
\subsection{Nuclear Energy}Learners should be able to demonstrate and apply their knowledge and
understanding of:
\begin{itemize}
	\item[\Large{$\Square$}] (a) the association between mass and energy and that \(E = mc^{2}\)
	\item[\Large{$\Square$}]	(b) the binding energy for a nucleus and hence the binding energy per nucleon, making use, where necessary, of the unified atomic mass unit (u)
	\item[\Large{$\Square$}]	(c) how to calculate binding energy and binding energy per nucleon from given masses of nuclei
	\item[\Large{$\Square$}]	(d) the conservation of mass / energy to particle interactions - for example: fission, fusion
	\item[\Large{$\Square$}]	(e) the relevance of binding energy per nucleon to nuclear \sq fission and \sq fusion making reference when appropriate to the binding energy per nucleon versus nucleon number curve
\end{itemize}
\subsection{Magnetic Fields}Learners should be able to demonstrate and apply their knowledge and
understanding of:
\begin{itemize}
	\item[\Large{$\Square$}] (a) how to determine the direction of the force on a current carrying conductor in a magnetic field
	\item[\Large{$\Square$}]	(b) how to calculate the magnetic field, B, by considering the force on a current carrying conductor in a magnetic field i.e. understand how to use \( F=BIl \sin \theta \)
	\item[\Large{$\Square$}]	(c) how to use \(F=Bqv \sin \theta \) for a moving charge in a magnetic field
	\item[\Large{$\Square$}]	(d) the processes involved in the production of a Hall voltage and \sq understand that \( V_{H} \propto B\) for constant \(I\)
	\item[\Large{$\Square$}]	(e) the shapes of the magnetic fields due to a current in \sq a long straight wire and \sq a	long solenoid
	\item[\Large{$\Square$}]	(f) the equations:
	\begin{align*}
	\text{\sq } B= \frac{\mu_{0} I}{2 \pi a} && \text{and} && \text{\sq } B= \mu_{0} nI
	\end{align*} for the field strengths due to a long straight wire and in a long solenoid
	\item[\Large{$\Square$}]	(g) the fact that adding an iron core increases the field strength in a solenoid
	\item[\Large{$\Square$}]	(h) the idea that current carrying conductors exert a force on each other and \sq to predict the directions of the forces
	\item[\Large{$\Square$}]	(i) quantitatively, how ion beams of charged particles, are deflected in uniform \sq electric and \sq magnetic fields
	\item[\Large{$\Square$}]	(j) the motion of charged particles in magnetic and electric fields in linear accelerators, \sq cyclotrons and \sq synchrotrons
	\subsection*{SPECIFIED PRACTICAL WORK}
	\item[\Large{$\Square$}] Investigation of the force on a current in a magnetic field
	\item[\Large{$\Square$}] Investigation of magnetic flux density using a Hall probe
\end{itemize}
\subsection{Electromagnetic Induction}Learners should be able to demonstrate and apply their knowledge and
understanding of:
\begin{itemize}
	\item[\Large{$\Square$}] (a) the definition of magnetic flux as \( \phi = AB \cos \theta \) and
	\(\text{Flux linkage} = N \phi \)
	\item[\Large{$\Square$}]	(b) the laws of Faraday and Lenz
	\item[\Large{$\Square$}]	(c) how to apply the laws of Faraday and Lenz (i.e. emf = - rate of change of flux	linkage)
	\item[\Large{$\Square$}]	(d) the idea that an emf is induced in a linear conductor moving at right angles to	a uniform magnetic field
	\item[\Large{$\Square$}]	(e) qualitatively, how the instantaneous emf induced in a coil rotating at right angles to a magnetic field is related to the position of the \sq coil, flux \sq density, \sq coil area and \sq angular velocity
\end{itemize} 
\section*{Option A: Alternating Currents}
\subsection*{Overview}
This topic begins by using Faraday’s law to derive an expression for the induced emf
in a coil rotating in a B-field and then develops the idea of the rms value of an
alternating current or voltage. The phase relationships between current and voltage
for capacitors and inductors in a.c. circuits are studied. The terms phase angle and
impedance are discussed and the phenomenon of resonance in an LCR circuit is
introduced.

\subsection*{Learners should be able to demonstrate and apply their knowledge and
	understanding of:}
\begin{itemize}
	\item[\Large{$\Square$}](a) using Faraday's law, the principle of electromagnetic induction applied to a
	rotating coil in a magnetic field
	
	\item[\Large{$\Square$}](b) the idea that the flux linkage of a rotating flat coil in a uniform magnetic
	B-field is \(BAN \cos \omega t\) because the angle between the coil normal and the field
	can be expressed as \(\theta=\omega t\)
	\item[\Large{$\Square$}](c) the equation \(V=-\omega BAN \sin (\omega t) \) for the induced emf in a rotating flat coil in a uniform B-field
	\item[\Large{$\Square$}](d) the terms frequency, period, peak value and rms value when applied to
	alternating potential differences and currents
	\item[\Large{$\Square$}](e) the idea that the rms value is related to the energy dissipated per cycle, and use the relationships \(I= \frac{I_0}{\sqrt{2}}\) and \(V= \frac{V_0}{\sqrt{2}}\)
	\item[\Large{$\Square$}](f) the idea that the mean power dissipated in a resistor is given by \(P=VI=I^{2}R\) where V and I are the rms values
	\item[\Large{$\Square$}](g) the use of an oscilloscope (CRO or PC based via USB or sound card) to
	measure
	\begin{itemize}
		\item[\Large{$\Square$}] a.c. and d.c. voltages and currents
		\item[\Large{$\Square$}] frequencies
	\end{itemize}
	\item[\Large{$\Square$}](h) the \(90 degrees\) phase lag of current behind potential difference for an inductor in a
	sinusoidal a.c. circuit
	\item[\Large{$\Square$}](i) the idea that \(X_{L}=\frac{V_{rms}}{I_{rms}}\)
	is called the reactance, \(X_{L}\) , of the inductor, and to
	use the equation \(X_{L}=\omega L \)
	\item[\Large{$\Square$}](j) the 90 degree  phase lead of current ahead of potential difference for a capacitor in a sinusoidal a.c. circuit, and to use the equation \( X_{C}= \frac{V_{rms}}{I_{rms}} \), where \( X_{C}= \frac{1}{\omega C} \)
	\item[\Large{$\Square$}](k) the idea that the mean power dissipation in an inductor or a capacitor is zero
	\item[\Large{$\Square$}](l) how to add potential differences across series RC, RL and RCL combinations
	using phasors
	\item[\Large{$\Square$}](m) how to calculate phase angle and impedance, \(Z\), (defined as \( Z= \frac{V_{rms}}{I_{rms}} \) for such circuits)
	\item[\Large{$\Square$}](n) how to derive an expression for the resonance frequency of an RCL series
	circuit
	\item[\Large{$\Square$}](o) the idea that the Q factor of a RCL circuit is the ratio \( \frac{V_{L}}{V_{R}}(=\frac{V_{C}}{V_{R}}) \) at resonance
	\item[\Large{$\Square$}](p) the idea that the sharpness of the resonance curve is determined by the
	Q factor of the circuit
\end{itemize}
\newpage
\section*{Option B: Medical Physics}

\subsection*{Overview}
This topic begins by studying the nature and properties of X-rays and the uses of Xrays
in imaging soft tissue. Techniques of radiography are studied together with the
use of a rotating beam X-ray computed tomography scanner. The generation and
detection of ultrasound, its use for diagnosis and the study of blood flow are
introduced. The principles of magnetic resonance are discussed together with the
use of MRI in obtaining diagnostic information. The uses of radionuclides as tracers
is covered, together with the use of the gamma camera and positron emission
tomography scanning.

\subsection*{Learners should be able to demonstrate and apply their knowledge and
	understanding of:
}
\begin{itemize}
	\item[\Large{$\Square$}](a) the nature and properties of X-rays
	\item[\Large{$\Square$}](b) the production of X-ray spectra including methods of controlling the beam
	intensity and photon energy
	\item[\Large{$\Square$}](c) the use of high energy X-rays in the treatment of patients (therapy) and low
	energy X-rays in diagnosis
	\item[\Large{$\Square$}](d) the equation \( I=I_{0}e^{(-\mu x)} \)for the attenuation of X-rays
	\item[\Large{$\Square$}](e) the use of X-rays in imaging soft tissue, and fluoroscopy to produce real time X-rays using image intensifiers
	\item[\Large{$\Square$}](f) techniques of radiography including using digital image receptors
	\item[\Large{$\Square$}](g) the use of a rotating beam X-ray computed tomography (CT) scanner
	\item[\Large{$\Square$}](h) the generation and detection of ultrasound using piezoelectric transducers
	\item[\Large{$\Square$}](i) scanning with ultrasound for diagnosis including A-scans and
	B-scans incorporating examples and applications
	\item[\Large{$\Square$}](j) the significance of acoustic impedance, defined by \(Z= c \rho \)for the reflection and transmission of sound waves at tissue boundaries, including the need for
	a coupling medium
	\item[\Large{$\Square$}](k) the use of the Doppler equation \( \frac{\Delta f}{f_{0}}=\frac{2v}{c}\cos \theta \)to study blood flow using an ultrasound probe
	\item[\Large{$\Square$}](l) the principles of magnetic resonance with reference to precession nuclei,
	resonance and relaxation time, and to apply the equation \(f=42.6 \times 10^{6}B \) for the Lamour frequency
	\item[\Large{$\Square$}](m) the use of MRI in obtaining diagnostic information about internal structures
	\item[\Large{$\Square$}](n) the advantages and disadvantages of ultrasound imaging, X-ray imaging and
	MRI in examining internal structures
	\item[\Large{$\Square$}](o) the effects of \(\alpha\), \(\beta\), and \(\gamma\) radiation on living matter
	\item[\Large{$\Square$}](p) the Gray (Gy) as the unit of absorbed dose and the Sievert (Sv) as the unit of equivalent dose and effective dose. Define absorbed dose as energy per kilogram
	\item[\Large{$\Square$}](q) the use of the equations 
	\begin{itemize}
		\item[\Large{$\Square$}]\(\text{equivalent dose} = \text{absorbed dose} \times \text{(radiation) weighting factor}\) expressed as \(H=DW_{R}\) and \ldots
		\item[\Large{$\Square$}]\(\text{effective dose} = \text{equivalent dose} \times \text{(tissue) weighting factor}\) expressed as \(E=HW_{T}\)
	\end{itemize}
	\item[\Large{$\Square$}](r) the uses of radionuclides as tracers to image body parts with particular
	reference to technetium-99m (Tc-99m)
	\item[\Large{$\Square$}](s) the use of the gamma camera including the principles of the collimator,
	scintillation counter and photomultiplier / CCD
	\item[\Large{$\Square$}](t) positron emission tomography (PET) scanning and its use in detecting
	tumours
\end{itemize}
\newpage
\section*{Option C: The Physics of Sports}
\subsection*{Overview}
This topic studies the use of the centre of gravity in explaining how stability and
toppling is achieved in various sporting contexts. The concept of moment of inertia is
introduced together with the principle of conservation of angular momentum and their
application to different sporting contexts is studied. Projectile motion and the
Bernouilli equation and their application to sporting events is also studied.
\subsection*{Learners should be able to demonstrate and apply their knowledge and
	understanding of:}
\begin{itemize}
	\item[\Large{$\Square$}] (a) how to use the centre of gravity to explain how stability and toppling is
	achieved in various sporting contexts
	\item[\Large{$\Square$}] (b) how to use the principle of moments to determine forces within
	\begin{itemize}
		\item[\Large{$\Square$}] Various muscle systems in the human body and
		\item[\Large{$\Square$}] Other sporting contexts, for example, sailing
	\end{itemize}
	\item[\Large{$\Square$}] (c) how to use Newton’s 2nd law in the form \(Ft=mv-mu\) in various sporting
	contexts
	\item[\Large{$\Square$}] (d) the coefficient of restitution as \(e = \frac{\text{Relative speed after collision}}{\text{Relative speed before collision}}\) and also use it in the form\( e= \sqrt{\frac{h}{H}}\) where h is the bounce height and H is the drop height
	\item[\Large{$\Square$}] (e) what is meant by the moment of inertia of a body
	\item[\Large{$\Square$}] (f) how to use equations to determine the moment of inertia, \(I\), for example
	\begin{itemize}
		
		\item[\Large{$\Square$}]a solid sphere \( I= \frac{2}{5}mr^{2} \)
		\item[\Large{$\Square$}]a thin spherical shell \( I=\frac{2}{3}mr^{2} \) where m is the mass and r is the radius
		
	\end{itemize}\item[\Large{$\Square$}] (g) the idea that angular acceleration,\(\alpha \), is defined as the rate of change of angular velocity, \(\omega\), and how to use the equation \( \alpha = \frac{\omega_{2} - \omega_{1}}{t} \)
	\item[\Large{$\Square$}] (h) the idea that torque, \(\tau\), is given as \(\tau=I\alpha \)
	\item[\Large{$\Square$}] (i) angular momentum, \(J\), is given as \( J=I\omega \) where \(\omega\) is the angular velocity
	\item[\Large{$\Square$}] (j) the principle of conservation of angular momentum and use it to solve
	problems in sporting contexts
	\item[\Large{$\Square$}] (k) how to use the equation for the rotational kinetic energy, \(\text{Rotational}_{KE}=\frac{1}{2}I\omega^{2}\)
	\item[\Large{$\Square$}] (l) how to use the principle of conservation of energy including the use of linear and rotational kinetic energy as well as gravitational and elastic potential
	energy in various sporting contexts
	\item[\Large{$\Square$}] (m) how to use projectile motion theory in sporting contexts
	\item[\Large{$\Square$}] (n) how to use Bernoulli’s equation \(p=p_{0} - \frac{1}{2}\rho v^{2}\) in sporting contexts
	\item[\Large{$\Square$}] (o) how to determine the magnitude of the drag force using \( F_{D}=\frac{1}{2}\rho v^{2}AC_{D} \) where \(C_{D}\) is the drag coefficient
\end{itemize}
\newpage
\section*{Option D: Energy and the Environment}
\subsection*{Overview}
In this topic, learners will consider different factors which affect the rate at which the
temperature of the Earth rises. Common sources of renewable and non-renewable
energy are discussed and their development as sources of energy, both in the UK
and internationally are compared. Learners study the effect of insulation on thermal
energy loss and perform quantitative calculations on comparative uses of energy
transfer.
\subsection*{Learners should be able to demonstrate and apply their knowledge and
	understanding of:}
\begin{itemize}
	\item[\Large{$\Square$}](a) how the following affect the rate at which the temperature of the Earth rises
	including:
	\begin{itemize}
		\item[\Large{$\Square$}](i) the need for thermal equilibrium: that is the balance between energy
		inflow from the Sun and energy re-radiated from the Earth in the
		context of global energy demand and the effect of \(CO_{2}\) levels in the
		atmosphere
		\item[\Large{$\Square$}](ii) the origin and means of transmission of solar energy and the form of
		the Sun’s power spectrum including the idea that wavelengths are converted into the near infrared in the atmosphere
		\item[\Large{$\Square$}](iii) the use of Wien’s law \((\lambda_{max} T = \text{constant})\) and Stefan-Boltzman \(T^{4}\) law in the context of solar power
		\item[\Large{$\Square$}](iv) use of the density equation and Archimedes’ principle to explain why
		rising sea levels are due to melting ice caps and that the melting of ice on land increases sea levels but melting icebergs do not 
		
	\end{itemize}
	\item[\Large{$\Square$}](b) the common sources of renewable and non-renewable energy and be able to
	compare their development and use both in the UK and internationally
	\begin{itemize}
		\item[\Large{$\Square$}](i) solar power:
		\begin{itemize}
			\item[\Large{$\Square$}] the idea that the main branch of the proton-proton chain is the
			main energy production mechanism in the Sun
			\item[\Large{$\Square$}] the intensity of power from the Sun \(I=\frac{P}{A} \)and the inverse square
			law for a point source
			\item[\Large{$\Square$}] how to perform energy conversions using photovoltaic cells
			(including efficiency calculations)
			
		\end{itemize}
		\item[\Large{$\Square$}](ii) wind power:
		\begin{itemize}
			\item[\Large{$\Square$}] the power available from a flowing fluid \(P=\frac{1}{2}A \rho v^{3} \)
			\item[\Large{$\Square$}] the factors affecting the efficiency of wind turbines
			
		\end{itemize}
		\item[\Large{$\Square$}](iii) tidal barrages, hydroelectric power and pumped storage:
		\begin{itemize}
			\item[\Large{$\Square$}] the principles of energy conversion (\(E_p\) to \(E_k\)) in tidal barrage,
			hydroelectric and pumped storage schemes and be able to carry out energy and power calculations related to these schemes and compare with the energy produced from wind
			
		\end{itemize}
		\item[\Large{$\Square$}](iv) nuclear fission and fusion:
		\begin{itemize}
			\item[\Large{$\Square$}] the principles underlying breeding and enrichment in nuclear
			fission applications
			\item[\Large{$\Square$}] the difficulties in producing sustained fusion power
			- fusion triple product
			
		\end{itemize}
	\end{itemize}
	\item[\Large{$\Square$}](c) the principles of fuel cell operation and the benefits of fuel cells particularly regarding greenhouse gas emissions
	\item[\Large{$\Square$}](d) the thermal conduction equation in the form \( \frac{\Delta Q}{\Delta t}= - A K \frac{\Delta \theta}{\Delta x} \)
	\item[\Large{$\Square$}](e) the effect of insulation on thermal energy loss and be able to calculate the heat loss for parallel surfaces using the rate of energy transfer \( =UA\Delta \theta\) including cases
	where different materials are in contact
	
\end{itemize}
\setcounter{tocdepth}{10}
% \listoftodos
%\phantomsection
%\cleardoublepage
%\begin{fullwidth} 

%\addcontentsline{toc}{section}{\indexname}
%\newgeometry{left=6cm,right=6cm,bottom=1cm}
\printindex
%\end{fullwidth}
\end{document}
